\documentclass[11pt]{article}

\usepackage[T1]{fontenc}
\usepackage[utf8]{inputenc}
\usepackage[english]{babel}
\usepackage[margin=1in]{geometry}
\usepackage{lmodern}
\usepackage{microtype}
\usepackage{mathtools}
\usepackage{amsmath,amssymb,amsthm}
\usepackage{array}
\usepackage{booktabs}
\usepackage{tabularx}
\usepackage{float}
\usepackage{graphicx}
\usepackage{pdflscape}
\usepackage{caption}
\usepackage{enumitem}
\usepackage{fancyhdr}
\usepackage[numbers]{natbib}
\usepackage{doi}
\usepackage{xurl}
\usepackage{hyperref}
\usepackage{cleveref}
\usepackage{orcidlink}
\usepackage{titling}



\newcommand{\axisAccess}{Access}
\newcommand{\axisClosure}{Closure}
\newcommand{\axisStability}{Stability}
\newcommand{\axisTransport}{Transport}
\newcommand{\axisStatus}{Status}
\newcommand{\axisRecords}{Records}
\newcommand{\axisSymmetry}{Symmetry}
\newcommand{\axisInterfaces}{Interfaces}
\newcommand{\axisSelfReference}{Self-reference}

\newcommand{\chapAccess}{Access}
\newcommand{\chapStability}{Stability}
\newcommand{\chapTransport}{Transport}
\newcommand{\chapSymmetry}{Symmetry}
\newcommand{\chapStatusRecordsCoherence}{Status, Records, Coherence}
\newcommand{\chapSelfReferenceLimits}{Self-Reference and Limits}
\newcommand{\chapClosureRealityBoundary}{Closure--Reality Boundary}

\newcommand{\FIV}{Foundations~IV}
\newcommand{\FoundI}{Foundations~I}
\newcommand{\FoundII}{Foundations~II}
\newcommand{\FoundIII}{Foundations~III}
\newcommand{\NumLaws}{52}
\newcommand{\NumMechanizeNow}{46}
\newcommand{\NumAnchor}{2}
\newcommand{\NumPartial}{2}
\newcommand{\NumObligation}{2}
\newcommand{\NumAxes}{8}
\newcommand{\NumChapters}{6}

\newcommand{\currQ}{Q}                 \newcommand{\predM}{M}                 \newcommand{\splitObs}{\Delta}         \newcommand{\quotMap}{\pi}             \newcommand{\accessQ}{q_I}             \newcommand{\obsSig}{\sigma_I}         \newcommand{\closureForm}{\mathsf{cl}} \newcommand{\BirdInt}{\mathsf{BirdInt}}

\newcommand{\Kdom}{K^{D}}              \newcommand{\Kloss}{K^{L}}             \newcommand{\Schur}{\mathsf{Schur}}    \newcommand{\ClSB}{\mathrm{Cl}_{\mathrm{SB}}}  \newcommand{\etaSB}{\eta_{\mathrm{SB}}} \newcommand{\XiSB}{\Xi_{\mathrm{SB}}}   

\newcommand{\Jinv}{J}                  \newcommand{\Fixinv}{\mathrm{Fix}(\Jinv)}  \newcommand{\trAX}{\operatorname{tr}}  

\newcommand{\NormalForm}[1]{\textbf{Normal form #1}}
\newcommand{\SetupBlock}{\textit{Setup.}}
\newcommand{\ProofSpine}{\textit{Proof spine.}}
\newcommand{\CaseEnum}{\textit{Cases.}}
\newcommand{\BirdIntStatus}{\textit{Interaction status.}}
\newcommand{\Nonclaims}{\textit{Nonclaims.}}

\newcommand{\statTheorem}{Lean theorem}
\newcommand{\statPartial}{partial theorem (axiom-applied)}
\newcommand{\statObligation}{obligation (axiom-applied, opacity-guarded)}
\newcommand{\statAnchor}{anchor (substrate-delegated)}
\newcommand{\statSubstratePending}{substrate-pending (axiom-applied)}

\newcommand{\PaperAR}{the adequacy-residuals paper}
\newcommand{\PaperSDTC}{the self-dual trace confinement paper}
\newcommand{\PaperAOR}{the audited-operational-realisability paper}
\newcommand{\PaperNeedleKiller}{the needle-killer paper}
\newcommand{\PaperNS}{the Navier--Stokes paper}
\newcommand{\PaperRH}{the Riemann Hypothesis paper}
\newcommand{\PaperPvNP}{the P versus NP paper}
\newcommand{\PaperCSL}{the hiddenness paper}

\newcommand{\GuardPvNP}{saturated-SAT source predicate}
\newcommand{\GuardCSL}{CSL diagnosis source predicate}
\newcommand{\GuardDuality}{duality-fixity substrate-anchor predicate}
\newcommand{\GuardSSB}{future branch-selection formation source predicate}

\newcommand{\thmDescentRepair}{Descent--Repair Normal Form}
\newcommand{\thmHolonomyMemory}{Holonomy-Memory Repair Normal Form}
\newcommand{\thmLocalGlobalObstruction}{Local--Global Obstruction Normal Form}
\newcommand{\thmRecombinationComparison}{Recombination Comparison}
\newcommand{\thmNoNeedles}{No-Needles Stability}
\newcommand{\thmSufficiencyClosure}{Sufficiency Closure}
\newcommand{\thmStrictExtension}{Strict Extension}
\newcommand{\thmNoUnstatusedResidual}{No Unstatused Residual}
\newcommand{\thmQuotientality}{Quotientality}
\newcommand{\thmAdequacy}{Adequacy / No-Overread}
\newcommand{\thmNoFreeDistinction}{No-Free-Distinction}
\newcommand{\thmHiddenness}{Hiddenness Normal Form}
\newcommand{\thmCSLFormedClosure}{CSL Formed-Closure Law}
\newcommand{\thmLayerMultiplicity}{Layer Multiplicity}
\newcommand{\thmDualityFixity}{Duality Fixity}
\newcommand{\thmSelectionObstruction}{Selection Obstruction}
\newcommand{\thmSpontaneousSSB}{Spontaneous Symmetry Breaking}
\newcommand{\thmAnomalySymmetry}{Anomaly / Symmetry Obstruction}
\newcommand{\thmDualityEquivalence}{Duality Equivalence}
\newcommand{\thmFactRecordFormation}{Fact / Record Formation}
\newcommand{\thmHorizonSufficiency}{Horizon Sufficiency}
\newcommand{\thmClosureRealityBoundary}{Closure--Reality Boundary}
\newcommand{\thmClosureCompletenessBoundary}{Closure Completeness Boundary}
 
\newtheorem{theorem}{Theorem}[section]
\newtheorem{lemma}[theorem]{Lemma}
\newtheorem{proposition}[theorem]{Proposition}
\newtheorem{corollary}[theorem]{Corollary}
\theoremstyle{definition}
\newtheorem{definition}[theorem]{Definition}
\newtheorem{remark}[theorem]{Remark}
\floatstyle{ruled}
\newfloat{algorithm}{tbp}{loa}
\floatname{algorithm}{Algorithm}

\urlstyle{same}
\captionsetup{
  font=small,
  labelfont=bf,
  labelsep=period
}
\captionsetup[table]{skip=6pt,position=top}
\captionsetup[figure]{skip=6pt}
\captionsetup[algorithm]{skip=6pt,position=top}
\crefname{algorithm}{algorithm}{algorithms}
\Crefname{algorithm}{Algorithm}{Algorithms}
\crefname{theorem}{theorem}{theorems}
\Crefname{theorem}{Theorem}{Theorems}
\crefname{lemma}{lemma}{lemmas}
\Crefname{lemma}{Lemma}{Lemmas}
\crefname{proposition}{proposition}{propositions}
\Crefname{proposition}{Proposition}{Propositions}
\crefname{corollary}{corollary}{corollaries}
\Crefname{corollary}{Corollary}{Corollaries}
\crefname{definition}{definition}{definitions}
\Crefname{definition}{Definition}{Definitions}
\crefname{remark}{remark}{remarks}
\Crefname{remark}{Remark}{Remarks}
\setlist[enumerate]{leftmargin=*, itemsep=2pt, topsep=4pt}
\setlength{\droptitle}{-3.2em}
\pretitle{\begin{center}\LARGE\bfseries}
\posttitle{\par\end{center}\vspace{0.5em}}
\preauthor{\begin{center}\normalsize}
\postauthor{\par\end{center}\vspace{-0.4em}}
\predate{\begin{center}\small}
\postdate{\par\end{center}\vspace{0.6em}}
\setlength{\emergencystretch}{2em}
\clubpenalty=10000
\widowpenalty=10000
\displaywidowpenalty=10000
\interfootnotelinepenalty=10000
\allowdisplaybreaks[2]
\fancypagestyle{firstpage}{\fancyhf{}
  \fancyfoot[C]{\scriptsize Automorph Inc., Wilmington, DE, USA \quad \texttt{ioannis@automorph.io} \quad \textcopyright\ 2026 \quad CC-BY 4.0 \quad \doi{10.5281/zenodo.20713187}}
  \renewcommand{\headrulewidth}{0pt}
  \renewcommand{\footrulewidth}{0pt}
}

\title{Six Birds Foundations IV:\\
A Catalog of Layer-Agnostic Structural Laws}
\author{Ioannis Tsiokos\,\orcidlink{0009-0009-7659-5964}}
\date{\today}
\hypersetup{
  pdftitle={Six Birds Foundations IV: A Catalog of Layer-Agnostic Structural Laws},
  pdfauthor={Ioannis Tsiokos},
  hidelinks
}

\begin{document}

\sloppy
\hbadness=10000

\maketitle
\thispagestyle{firstpage}

\begin{abstract}
Foundations IV is a catalog of layer-agnostic structural laws built on
the Foundations I/II/III apparatus. Its organizing question is what
must be true, before choosing a substrate---whether analysis,
arithmetic, physics, biology, chemistry, economics and finance,
linguistics, cognitive and social science, or computation---for a
layer to have lawful access, stable records, transport, symmetry,
closure, and limits.  The contribution is
a uniform normal-form template applied across 52 laws; each row states
its setup, theorem, proof spine, case structure, interaction status,
and nonclaims, so recurring obstructions such as failed descent, hidden
predictive structure, local-global failure, and record nonclosure can
be compared in one register.  The formalization status is explicit: the
catalog contains 46 direct Lean theorems, 2 substrate-delegated anchors
(No-Needles Stability and Closure--Reality Boundary), 2 partial
wrappers (Hiddenness and Duality Fixity), and 2 obligation wrappers
(CSL Formed-Closure and Spontaneous Symmetry Breaking); the four
trust-base axioms used by the wrappers are disclosed through an
opacity-guard pattern, each guarded by an opaque source predicate that
prevents the axiom from being instantiated without the corresponding
substrate or future-source certificate.  Foundations IV is not a
foundational reset of Foundations I/II/III, not a unification claim
about the eight axes, and not an exhaustive catalog.
\end{abstract}

\medskip
\noindent\textbf{Keywords.} quotient packaging; normal-form structural law;
catalog discipline; opacity-guard pattern; substrate-anchor disclosure;
layer-agnostic statement; Six Birds theory; emergence calculus.

\section{Introduction}\label{sec:intro}

Let \(q:X\to Q\) be a quotient map, let \(F:X\to Y\) be a map one
wants to transport through that quotient, and let \(r:Y\to R\) be the
target identification under which the transported output is to be
read.  The elementary question is whether the composite \(rF\)
depends only on the \(q\)-class of the input.  Equivalently, does
there exist a map \(\overline F:Q\to R\) with
\[
  \overline F q = rF?
\]
If \(x\) and \(x'\) have the same image in \(Q\) but \(rF(x)\ne
rF(x')\), then no such \(\overline F\) can exist.  Conversely, in
the usual quotient-map situation, the absence of such pairs is
exactly the condition for descent.  There are two canonical repairs.
One may refine the source quotient so that the offending pair is no
longer identified, or one may coarsen the target by quotienting \(R\)
by the smallest equivalence relation generated by the identifications
\[
  rF(x)\sim rF(x')
  \qquad\text{whenever}\qquad
  q(x)=q(x').
\]
The latter repair forgets precisely the distinction that the source
quotient cannot support.  This is a familiar mathematical pattern:
an obstruction is not merely a failure, but a diagnostic object whose
minimal repair is determined by the same data that produced the
failure.

The example is intentionally elementary.  It uses no special
operator theory, no hidden regularity hypothesis, and no category
larger than the quotient construction itself.  Nevertheless, the
same shape appears in more elaborate settings.  A local family may
agree on overlaps but fail to come from a global object.  A symmetric
construction may be unable to choose a non-fixed branch without
adding symmetry-breaking data.  A formally closed record may exist
without being physically instantiated.  A finite observation may
determine all declared future probes, or fail precisely because some
future probe separates states the present quotient cannot see.  In
each case there is a setup, an obstruction, a repair or boundary, and
a statement of what the theorem does not claim.  The purpose of this
paper is to put those repeated shapes into a common register.

\FIV\ is a catalog of such patterns.  The word ``catalog'' is meant
literally: the paper does not introduce a new primitive foundation,
but organizes recurring structural obstructions that appear across
the Six Birds corpus into named normal-form laws.  A normal-form law,
in this paper, is a theorem presented in a fixed template: setup,
theorem, proof spine, case enumeration, interaction status, and
nonclaims.  The inherited apparatus is used as a black box.
\FoundI\ supplies the closure-formation language, meaning the
operations by which a structure is completed under declared closure
rules \citep{Tsiokos2026EmergenceCalculus}.  \FoundII\ supplies the
admissibility and quotient-packaging discipline, where quotient
packaging is just the standard act of retaining only equivalence
classes of a structure \citep{Tsiokos2026FoundationsII}.
\FoundIII\ supplies the finite audited interaction judgment used to
record when a proposed interaction is legitimately visible to the
apparatus \citep{Tsiokos2026FoundationsIII}.  \FIV\ takes these
ingredients as already established and uses them to state, name, and
formalize the recurring structural laws that had previously appeared
only as fragmentary substructures in individual papers.

This catalog posture is useful only if the entries are comparable.
For that reason the paper does not present the laws as isolated
lemmas with unrelated notation.  Each entry is phrased in the same
normal-form template, so that a reader can ask the same questions
from one axis to the next: what are the objects, what is the
obstruction, what is the repair or classification, what is the proof
spine, and what is explicitly outside the claim?  The template is
not meant to flatten the mathematics.  Rather, it records enough
common structure to let a quotient-theoretic result in the Transport
chapter, an access result about observables, and a closure boundary
result at the end of the paper be compared without pretending that
they are the same theorem.

The opening descent example is the prototype for the style of the
paper.  Later we call the set of offending pairs a split-pair
obstruction, denoted \(\splitObs\): it is the set of pairs that are
identified by the source quotient but separated after transport and
target identification.  Similar quotient language recurs throughout
the catalog.  A current quotient \(\currQ\) is the equivalence
relation the system presently imposes, while a predictive quotient
\(\predM\) is a refinement predicted by the apparatus but not yet
available at the current observational level.  These terms are not
additional structure imposed on the reader; they are names for
ordinary quotient comparisons that recur often enough to warrant
uniform notation.

The same caveat applies to the word ``law.''  A law in this
paper is not a physical law and not a universal principle detached
from hypotheses.  It is a named structural theorem, stated at the
level of the data in its setup.  Some laws are ordinary quotient
factorization statements.  Some are finite classification
statements.  Some are boundary statements saying that a desirable
implication cannot hold without strengthening the data.  The uniform
language is there to make the hypotheses and the nonclaims visible,
not to erase their differences.

The catalog contains \NumLaws\ rows arranged along eight
primitive-necessity axes---access, closure or stability, transport,
status, records, symmetry, interfaces, and self-reference---grouped
for readability into six drafting chapters and a separate
Closure--Reality Boundary anchor; the per-status coverage counts are
recorded in \Cref{tab:lean-coverage-by-axis}.

The chapter grouping is likewise practical.  The Access chapter
studies what an instrument can distinguish; its lead theorem says
that a lawful observation signature determines a canonical access
quotient.  The Stability chapter studies when declared future
behavior is already controlled by present data, and when residuals
must be accounted for.  The Transport chapter is the easiest entry
point, because its laws are closest to standard descent, gluing, and
quotient comparison.  The Symmetry chapter treats obstructions
coming from involutions, branch selection, and anomaly-style failures
of descent or gluing.  The Status/Records/Coherence chapter collects
laws in which records, facts, persistence, explanation, and
counterfactuals are themselves part of the mathematical object.  The
Self-Reference/Limits chapter records horizon, information,
singular-limit, locality, topology, and completeness boundaries.  A
final Closure--Reality Boundary anchor separates formal closure from
bare physical realisability and records the corrected reflective
closure statement.

The catalog includes substrate-delegated anchors and axiom-applied
wrappers whose Lean-disclosure discipline is recorded in
\Cref{app:formalization}; per-theorem footnotes in the body cite the
Lean realization at the appropriate fidelity.

The substrate-instance map is visible from the start.  No-Needles
Stability is the layer-agnostic form instantiated by
\PaperNeedleKiller\ \citep{Tsiokos2026NeedleKiller} together with
\PaperNS\ \citep{Tsiokos2026NavierStokes}.  Hiddenness Normal Form
points to \PaperPvNP\ \citep{Tsiokos2026PvNP}.  CSL Formed-Closure
points to \PaperCSL\ \citep{Tsiokos2026CSL}.  Duality Fixity points
to \PaperSDTC\ \citep{Tsiokos2026SDTC} together with \PaperRH\
\citep{Tsiokos2026RH}.  The Closure--Reality Boundary anchor points
to \PaperAOR\ \citep{Tsiokos2026AOR}.  In each case the statement in
this paper is the layer-agnostic version: it quantifies over the
kind of closure, quotient, record, or symmetry data at issue, while
the substrate paper specializes that pattern to a particular
carrier.  The anchor row is not a re-proof of the substrate theorem;
its contribution is to isolate the reusable structural form that the
substrate theorem instantiates.  This is also why \FIV\ is a fourth
foundations installment rather than a foundational reset or a
meta-theory paper: it neither duplicates the substrate proofs nor
hides them.

The same posture explains the balance between abstraction and
example.  The laws are stated abstractly because their role is to
identify reusable structure, but the introduction and the Transport
chapter begin with concrete quotient situations to make the
abstractions inspectable.  A reader who knows the universal property
of a quotient already knows the core move behind many entries: show
that a proposed construction is constant on the relevant fibers, or
identify the precise pair, record, route, or probe that prevents it
from being constant.  The rest of the paper repeats that move in
progressively less elementary settings.

We suggest reading the introduction, then \Cref{sec:apparatus}, and
then \Cref{sec:transport}, whose descent--repair, holonomy-memory,
local--global, and recombination laws form the most familiar
quotient-theoretic entry point into the catalog.  After that, the
axes can be read independently.  Readers already familiar with the
foundations papers may skim \Cref{sec:apparatus} for notation and
use the axis summaries.  Readers coming from a substrate paper should start from the
corresponding anchor row: No-Needles for \PaperNeedleKiller\
\citep{Tsiokos2026NeedleKiller} and \PaperNS\
\citep{Tsiokos2026NavierStokes}, Hiddenness Normal Form for
\PaperPvNP\ \citep{Tsiokos2026PvNP}, CSL Formed-Closure for
\PaperCSL\ \citep{Tsiokos2026CSL}, Duality Fixity for \PaperSDTC\
\citep{Tsiokos2026SDTC} and \PaperRH\ \citep{Tsiokos2026RH}, and
Closure--Reality Boundary for \PaperAOR\ \citep{Tsiokos2026AOR}.

Each catalog entry in the body is stated as a formal theorem, or
when appropriate a definition, lemma, or proposition, followed by a
full proof.  A footnote attached to each theorem records the Lean
realization, citing the declaration's namespaced identifier and the
appropriate wording for its coverage status.  The mechanization
bookkeeping---per-row coverage, the trust-base axiom register, the
opacity-guard pattern, and the narrowing disclosures for the four
axiom-applied wrappers---lives in \Cref{app:formalization}.  The
body of the paper is self-contained as a mathematical text; the
footnotes and the formalization appendix give the parallel account
of what is verified in Lean and at what fidelity.

The catalog is not exhaustive.  Several structural candidates remain
as sketches because the recognition source carrying the relevant
structure has not yet been identified; \Cref{sec:sketches} lists
those items explicitly.  What is closed here is the catalog scope:
\Cref{sec:apparatus} fixes the normal-form template by which the
\NumLaws\ laws are stated, compared, and disclosed.
 \section{Apparatus and conventions}\label{sec:apparatus}

\FIV\ is written so that the catalog can be read without first
working through the earlier foundations papers.  This section records
the imported apparatus in ordinary mathematical language, fixes the
quotient notation used throughout the paper, and states the uniform
normal-form template used by the catalog.  The imported foundations
results are used as background, not reproved here.

\FoundI\ supplies closure formation
\citep{Tsiokos2026EmergenceCalculus}; \FoundII\ supplies admissibility
and quotient packaging \citep{Tsiokos2026FoundationsII}; and
\FoundIII\ supplies the finite audited interaction judgment
\citep{Tsiokos2026FoundationsIII}.  The following definitions are the
forms of that apparatus used in this paper.

\begin{definition}[Closure formation]\label{def:closure-formation}
A \emph{closure formation} is the passage from initial data to the
object obtained by closing that data under a declared family of closure
rules.  We write \(\closureForm\) for the formed-closure operation.  A
\emph{formed closure} is the completed object relative to the declared
rules.  A \emph{closure-of-record} is a formed closure together with an
audit record of the closure data used to form it.
\end{definition}

Thus completing a substructure under an operation, taking the closure
of a relation, or closing a set of rules under composition are familiar
examples.  The \FoundI\ vocabulary does not introduce a new species of
closure; it records which rules have been declared and which completed
object is being used.

\begin{definition}[Current and predictive quotients]\label{def:quotient-packaging}
Quotient packaging is the standard quotient construction: an
equivalence relation is imposed on a carrier and only equivalence
classes are retained.  The quotient map is denoted \(\quotMap\) when no
more specific symbol is needed.  The \emph{current quotient}
\(\currQ\) is the quotient presently visible to the apparatus.  The
\emph{predictive quotient} \(\predM\) is a finer quotient predicted by
the apparatus but not necessarily visible at the current level.  Thus
\(\predM\) refines \(\currQ\): a \(\currQ\)-class may contain several
\(\predM\)-classes.
\end{definition}

The gap between \(\currQ\) and \(\predM\) is the basic way the paper
records a distinction that is predicted but not currently observable.
The same comparison appears in the Access, Stability, Transport, and
Symmetry chapters.

\begin{definition}[BirdInt judgment]\label{def:birdint}
A \emph{\(\BirdInt\) judgment} is the \FoundIII\ finite audited
interaction judgment.  In this paper it records that the interaction
named in a catalog law is legitimately observable by the declared
finite audit data.  It is not an informal plausibility assertion; it is
the inherited finite-witness condition for observable interaction.
\end{definition}

\begin{figure}[t]
\centering
\small
\setlength{\tabcolsep}{4pt}
\renewcommand{\arraystretch}{1.12}
\begin{tabularx}{\linewidth}{@{}X X X@{}}
\fbox{\begin{minipage}[t]{0.88\linewidth}
\textbf{Access}\\
Atlas axis: access\\
\([T]\) F10--F12, F24\\
\([P]\) F13a \quad \([O]\) F13b
\end{minipage}}
&
\fbox{\begin{minipage}[t]{0.88\linewidth}
\textbf{Stability}\\
Atlas axis: closure/stability\\
\([T]\) F7--F9\\
\([A]\) F6
\end{minipage}}
&
\fbox{\begin{minipage}[t]{0.88\linewidth}
\textbf{Transport}\\
Atlas axis: transport\\
\([T]\) F2--F5\\
source of split-pair discipline
\end{minipage}}
\\[0.8em]
\fbox{\begin{minipage}[t]{0.88\linewidth}
\textbf{Symmetry}\\
Atlas axis: symmetry\\
\([T]\) F15a, F40--F41\\
\([P]\) F14 \quad \([O]\) F15b
\end{minipage}}
&
\fbox{\begin{minipage}[t]{0.88\linewidth}
\textbf{Status / Records / Coherence}\\
Atlas axes: status, records, interfaces\\
\([T]\) F17--F23, F25--F31
\end{minipage}}
&
\fbox{\begin{minipage}[t]{0.88\linewidth}
\textbf{Self-reference / Limits}\\
Atlas axis: self-reference\\
\([T]\) F33--F39, F42--F52\\
F52 closes catalog coherence
\end{minipage}}
\end{tabularx}

\vspace{0.75em}

\fbox{\begin{minipage}{0.94\linewidth}
\textbf{Cross-axis dependency hints.}
F2 supplies split-pair/descent discipline reused by F4, F5, F15a,
and the F33+ limit rows.
F10 supplies access-quotient discipline reused across Access and
Status/Records/Coherence.
F6 is the needle-killer / Navier--Stokes substrate-delegated anchor;
F53 is the audited-operational-realisability Closure--Reality
Boundary anchor outside the six drafting chapters.
F13a and F14 are partial wrappers; F13b and F15b are
obligation-guarded wrappers.
\end{minipage}}

\vspace{0.75em}

\fbox{\begin{minipage}{0.58\linewidth}
\textbf{Separate anchor component:}
\([A]\) F53, Closure--Reality Boundary.
\end{minipage}}
\qquad
\begin{minipage}{0.31\linewidth}
\textbf{Legend.}\\
\([T]\) direct Lean theorem\\
\([A]\) substrate-delegated anchor\\
\([P]\) partial theorem wrapper\\
\([O]\) obligation-guarded wrapper
\end{minipage}

\caption{Axis topology for \FIV.  The six drafting chapters group the
eight atlas axes, with the Closure--Reality Boundary shown as a
separate anchor component; the annotations record only the
cross-axis dependencies used later in the body chapters, not the full
52-row dependency graph.}
\label{fig:axis-topology}
\end{figure}
 
\begin{definition}[Split-pair obstruction]\label{def:split-pair}
Let \(q:X\to Q\) be a source quotient, \(F:X\to Y\) a map, and
\(r:Y\to R\) a target identification.  The \emph{split-pair
obstruction} is
\[
  \splitObs(q,F,r)
  =
  \{(x,x')\in X\times X
    \mid q(x)=q(x')\ \text{and}\ rF(x)\ne rF(x')\}.
\]
It is the set of pairs identified by the source quotient but separated
after transport to the target quotient.
\end{definition}

The descent example from \Cref{sec:intro} is exactly
\Cref{def:split-pair}: the map \(rF\) descends through \(q\) precisely
when \(\splitObs(q,F,r)=\emptyset\).  Equivalently, there is a map
\(\overline F:Q\to R\) with
\[
  \overline F\circ q=r\circ F .
\]
When the obstruction is nonempty, the source-side repair refines
\(q\), while the target-side repair coarsens \(R\) by the smallest
equivalence relation generated by the identifications
\(rF(x)\sim rF(x')\) whenever \(q(x)=q(x')\).

\begin{definition}[Observation signature; access quotient]\label{def:access-quotient}
An \emph{observation signature} is a declared map
\(\obsSig:H\to \Sigma\) from a carrier \(H\) to observation symbols.
It determines the access kernel
\[
  h\sim_I h'
  \quad\Longleftrightarrow\quad
  \obsSig(h)=\obsSig(h').
\]
The \emph{access quotient} \(\accessQ\) is the quotient of \(H\) by
this kernel.  It is the quotient carrying exactly the distinctions made
by the declared observation signature.
\end{definition}

Any observable constant on the access kernel factors uniquely through
\(\accessQ\).  The Access chapter turns this universal property into
the Quotientality theorem; here it fixes the notation used by later
laws.

\begin{figure}[t]
\centering
\fbox{\begin{minipage}{0.88\linewidth}
\centering
\(\SetupBlock\)
\(\longrightarrow\)
\(\NormalForm{theorem}\)
\(\longrightarrow\)
\(\ProofSpine\)
\(\longrightarrow\)
\(\CaseEnum\)
\(\longrightarrow\)
\(\BirdIntStatus\)
\(\longrightarrow\)
\(\Nonclaims\)
\end{minipage}}
\caption{The six-part normal-form template used for each catalog law.}
\label{fig:normal-form-template}
\end{figure}

\begin{definition}[Normal-form template]\label{def:normal-form}
Each catalog law is presented by a six-part normal-form template:
\emph{setup}, \emph{theorem}, \emph{proof spine}, \emph{case
enumeration}, \emph{interaction status}, and \emph{nonclaims}.  The
setup names the data on which the law operates.  The theorem states the
layer-agnostic claim.  The proof spine gives the structural argument.
The case enumeration lists the mutually exclusive outcomes when the
statement branches.  The interaction status records the relevant
\(\BirdInt\) legitimacy condition.  The nonclaims field states what the
law does not assert.  \Cref{fig:normal-form-template} gives the
schematic form.
\end{definition}

The template is editorial rather than categorical.  It does not assert
that the laws are equivalent or that the axes are canonical.  Its role
is to put hypotheses, proof shape, case split, audit condition, and
scope fence in the same order for every row.

The catalog contains four trust-base axioms used by the partial and
obligation wrappers in \Cref{sec:access,sec:symmetry}.  Each is paired
with an opaque source predicate, or a definitional disjunction of such
predicates, whose lack of public constructors prevents the axiom from
deriving \(\bot\) on its own.  \Cref{app:formalization} records the
pattern in detail, the per-axiom guards, and the registered axiom names.

The Lean correspondence for the catalog rows is recorded in
\Cref{app:formalization}.  Body chapters state the mathematical law and
its proof spine; that appendix records the per-row mapping, coverage
discipline, and trust-base register.
 \section{Axis I — Access}\label{sec:access}

The Access axis is about what an instrument can distinguish.  An
observation signature \(\obsSig\), in the sense fixed in
\Cref{sec:apparatus}, is a declared map from states to observation
symbols; its access quotient \(\accessQ\) identifies exactly those
states that no declared observation separates.  The laws in this
chapter govern how such signatures determine quotients, when those
quotients are adequate for accepted claims, when new distinctions
require real apparatus data rather than a free relabeling, and when
structure remains hidden behind the current quotient \(\currQ\) even
though a predictive quotient \(\predM\) would separate it.

The lead theorem is Quotientality: a lawful observation signature
determines a canonical access quotient, and observables constant on
the signature factor through it.  Adequacy / No-Overread then says
that accepted claims must factor through the available access, so the
instrument cannot legitimately read beyond its quotient.
No-Free-Distinction bridges these two laws: a new distinction must
enter through memory, bridge data, residual budget, or a change of
scope, rather than through an unaccounted functorial gate, meaning a
monotone map on quotients.  The chapter closes with Layer
Multiplicity, which explains why nontrivial limited access
generically forces additional layers.

Between those direct rows are two axiom-applied wrappers.  Hiddenness
Normal Form compares current and predictive observation: hiddenness
is the situation in which states agree at \(\currQ\) but differ at
\(\predM\), while exposure is the refinement of \(\currQ\) that
removes that gap.  Its substrate source is the PvNP paper
\citep{Tsiokos2026PvNP}.  The CSL Formed-Closure Law is the
corresponding diagnosis-grade closure statement from the CSL paper
\citep{Tsiokos2026CSL}.  The guard details and per-row formal record
are deferred to \Cref{app:formalization}, especially
\Cref{tab:tagged-axiom-register}.

Readers can start this axis immediately after \Cref{sec:apparatus}.
The first three subsections are ordinary quotient and factorization
statements, with Quotientality as the access-theoretic entry point.  A
reader who wants to stay entirely within direct quotient arguments
may read Quotientality, Adequacy / No-Overread,
No-Free-Distinction, and Layer Multiplicity first, and return to the
guarded wrappers only when the source-paper status matters.

\subsection{Quotientality (Layer-Sameness via Quotient)}\label{sec:access:quotientality}

\begin{definition}[Observation signature]\label{def:obs-signature}
Let \(H\) be a closure carrier.  An observation signature on \(H\)
is a declared map
\[
  \obsSig:H\to\Sigma
\]
from states to observation symbols.  Its access kernel is the
equivalence relation
\[
  h\sim_I h'
  \quad\Longleftrightarrow\quad
  \obsSig(h)=\obsSig(h'),
\]
and its access quotient is the quotient map
\[
  \accessQ:H\to Q_I:=H/{\sim_I}.
\]
For a calibration family of observation signatures, the same
construction is applied to the joint access kernel.
\end{definition}

\begin{theorem}[Quotientality]\label{thm:F10}
Let \(\obsSig:H\to\Sigma\) be an observation signature with access
quotient \(\accessQ:H\to Q_I\).  Every observable \(o:H\to R\) that
is constant on access classes factors uniquely through
\(\accessQ\):
\[
  o=\overline o\circ \accessQ
\]
for a unique \(\overline o:Q_I\to R\).  Moreover, \(Q_I\) is the
coarsest quotient carrying the declared observations: if
\(\pi:H\to Q'\) is any quotient whose fibers are observation-invisible,
then \(\accessQ\) factors through \(\pi\).
\end{theorem}
\footnote{Verified in Lean as
\texttt{Access.Quotientality.quotientality}; see
\Cref{app:formalization}.}
\begin{proof}
Define \(\overline o([h]_I)=o(h)\).  This is well defined because
\([h]_I=[h']_I\) means \(h\sim_I h'\), and the hypothesis on \(o\)
gives \(o(h)=o(h')\).  The equation
\(o=\overline o\circ\accessQ\) follows from the definition on
representatives.  If \(\overline o'\) is another such map, then for
every access class \([h]_I\) we have
\[
  \overline o'([h]_I)
  =
  \overline o'(\accessQ(h))
  =
  o(h)
  =
  \overline o([h]_I),
\]
so \(\overline o'=\overline o\).

For coarsestness, suppose \(\pi:H\to Q'\) is a quotient whose fibers
are observation-invisible: \(\pi(h)=\pi(h')\) implies
\(\obsSig(h)=\obsSig(h')\).  Define \(u:Q'\to Q_I\) on reachable
classes by \(u(\pi(h))=\accessQ(h)\).  The same invisibility
hypothesis makes this assignment independent of the representative.
Hence \(\accessQ=u\circ\pi\).  Thus any quotient already carrying the
declared observations refines the canonical access quotient.
\end{proof}

The interaction status is inherited from the declaration of
\(\obsSig\).  Observation signatures are declared interactions, and
the access quotient is \(\BirdInt\)-legitimate precisely because it
records only distinctions made by those declared observations.  The
theorem does not produce the coarsest quotient possible for all
observers; it produces the coarsest quotient carrying the given
signature.  It also does not create new observable distinctions and
does not instantiate any particular carrier theory.

\subsection*{Layer instantiations}\label{layer-inst:F10}

\begin{description}
\item[\emph{Myhill--Nerode minimal {DFA} (formal language theory).}]
For a regular language \(L\subseteq\Sigma^*\), the Nerode
right-congruence \(\equiv_L\), defined by
\(w\equiv_L w'\Leftrightarrow\forall z.\,(wz\in L\Leftrightarrow w'z\in L)\),
is the coarsest right-congruence on \(\Sigma^*\) under which the
indicator \(\mathbf 1_L\) is constant on classes: every
\(L\)-respecting right-congruence refines \(\equiv_L\), and the
quotient \(\Sigma^*/\!\!\equiv_L\) is the state set of the minimal
DFA recognising \(L\)
\citep{Nerode1958LinearAutomatonTransformations}.

\item[\emph{Kalman observability canonical decomposition (control
engineering).}]
For a linear time-invariant system \(\dot x=Ax+Bu,\ y=Cx\), the
unobservable subspace
\(N=\bigcap_{k\ge 0}\ker(CA^k)\) is the kernel of the observability
matrix.  Two state vectors agree under \(\mathbb R^n/N\) iff they
produce identical output trajectories for every admissible input,
and the Kalman canonical form factors every input--output
observable uniquely through \(\mathbb R^n/N\); any quotient with
output-trajectory-preserving fibres must refine through this
observable quotient
\citep{Kalman1963MathematicalDescription}.

\item[\emph{Bernstein 3{NF} synthesis (database engineering).}]
Given a relation schema with attributes \(U\) and a
functional-dependency set \(F\), the 3NF decomposition produced by
Bernstein
synthesis from a chosen minimal cover of \(F\) is the canonical
dependency-preserving 3NF decomposition: every dependency-preserving
observable (query, FD-checker) factors through the projections onto
the component schemas, and any further coalescing of components
either loses a declared FD or introduces a 3NF violation
\citep{Bernstein1976Synthesizing3NF}.

\item[\emph{Trubetzkoy phonemic principle (structural linguistics).}]
For a language \(L\), the phoneme inventory \(\mathcal P_L\) is the
coarsest equivalence on phonetic realisations preserving every
morpheme-distinguishing minimal-pair contrast.  Two phones collapse
to the same phoneme iff they never produce a minimal pair, and
every morpheme-discriminating phonological observable factors
through the phoneme map; any coarser partition collapses a
morpheme distinction \(L\) tracks
\citep{Trubetzkoy1939GrundzugePhonologie}.

\item[\emph{Categorical perception of voice onset time (cognitive
psychophysics).}]
For a synthesised voice-onset-time continuum spanning \(/{\!b}/\) to
\(/{\!p}/\), human listeners exhibit \emph{categorical perception}:
within-category {ABX} discrimination collapses to chance while
between-category discrimination is sharp at the boundary near
\(+25\)~ms.  The phonetic-category partition is the coarsest
equivalence on the acoustic continuum preserving the declared
behavioural discriminability probes (ABX and identification
tasks), and each such probe factors uniquely through the
phonetic-category map
\citep{LibermanHarrisHoffmanGriffith1957Discrimination}.
\end{description}

\subsection{Adequacy / No-Overread}\label{sec:access:adequacy}

Fix the same closure carrier \(H\) and declared observable family as
in the previous subsection, and let \(\accessQ:H\to Q_I\) be the
access quotient constructed from that family.  A claim is accepted at
the current layer only after it passes the finite audited status
discipline inherited from \FoundIII; the no-overread suppression
there says that an accepted claim may not depend on content that the
instrument has suppressed without an explicit bridge
\citep{Tsiokos2026FoundationsIII}.

\begin{theorem}[Adequacy / No-Overread]\label{thm:F11}
Let \(F:H\to W\) be a current-layer claim over a formed access
package.  If \(F\) is accepted as an exact current claim, then it is
constant on the fibers of \(\accessQ\).  Equivalently, there is a map
\(\overline F:Q_I\to W\) such that
\[
  F=\overline F\circ\accessQ .
\]
\end{theorem}
\footnote{Verified in Lean as
\texttt{Access.Adequacy.adequacy\_no\_overread}; see
\Cref{app:formalization}.}
\begin{proof}
The accepted-claim audit permits a current exact claim to use only
the distinctions available to the current access package.  Hence if
\(\accessQ(h)=\accessQ(h')\), the states \(h\) and \(h'\) are
indistinguishable by the declared observations, and no accepted
current claim may separate them.  Therefore \(F(h)=F(h')\).

By Quotientality, every map constant on access classes factors
uniquely through \(\accessQ\).  Applying \Cref{thm:F10} to \(F\)
gives a unique \(\overline F:Q_I\to W\) with
\(F=\overline F\circ\accessQ\).  Thus an accepted current claim has
no overread defect.
\end{proof}

The interaction status is \(\BirdInt\)-compatible because the
accepted-claim audit is itself a finite audited interaction
judgment.  The theorem does not forbid recording distinctions beyond
current access; such distinctions may belong to the predictive
quotient \(\predM\), to memory, or to a later access extension.  It
also does not determine the audit data, bridge strengths, thresholds,
or residual budgets required for acceptance.

\subsection*{Layer instantiations}\label{layer-inst:F11}

\begin{description}
\item[\emph{Rao--Blackwell theorem (statistical decision theory).}]
For a sufficient statistic \(T:X\to\mathcal T\) on a parametric
family \(\{P_\theta\}\), the Rao--Blackwell construction
\(\delta_{\mathrm{RB}}(x)=\mathbb E[\delta(X)\mid T(x)]\) takes any
decision rule \(\delta\) and produces a \(\sigma(T)\)-measurable
rule whose risk under a convex loss is at most that of \(\delta\)
almost surely.  Hence any admissible decision rule factors as
\(\delta=\bar\delta\circ T\) a.s.; a rule that uses within-fibre
detail of \(X\) beyond \(T(X)\) is the overread failure
\citep{Blackwell1947ConditionalExpectation}.

\item[\emph{Post-processing invariance of differential privacy
(cryptography / privacy).}]
For an \((\varepsilon,\delta)\)-differentially-private mechanism
\(M:\mathcal D\to\mathcal R\), the post-processing theorem of
Dwork--McSherry--Nissim--Smith states that for any (possibly
randomised) \(g:\mathcal R\to\mathcal W\), the composition
\(g\circ M\) is itself \((\varepsilon,\delta)\)-DP.  Releases
issued under the declared mechanism contract are therefore exactly
the post-processings of \(M\); a published quantity whose
neighboring-database distinguishability exceeds the
\(e^{\varepsilon}\!\cdot\!\!\ (\cdot)+\delta\) bound is the
overread failure --- evidence of access outside \(M\)'s
randomisation
\citep{DworkMcSherryNissimSmith2006CalibratingNoise}.

\item[\emph{Registered Reports / preregistered analyses (scientific
methodology).}]
In the Registered Reports format the access quotient is the
\emph{Stage~1 preregistered protocol together with the locked
dataset}: confirmatory claims must be computable from those two
inputs alone, and any two execution histories sharing them must
produce the same confirmatory verdict.  HARKing, p-hacking, and
outcome switching are overread failures: a published claim
presented as confirmatory but computed from post-hoc analytic
choices fails to factor through the declared protocol-plus-data
quotient
\citep{Chambers2013RegisteredReports}.

\item[\emph{Federal Rules of Evidence: admissibility (legal
procedure).}]
Under the U.S. Federal Rules of Evidence (Rules 401--403 on
relevance, Rule 802 on hearsay), a finding of fact must factor
through the admitted-evidence quotient of the case record: two
factually distinct records with identical admitted-evidence content
must yield the same verdict.  A verdict resting on inadmissible
evidence is reversible error unless harmless, instantiating the
F11 overread failure
\citep{FRE2024FederalRulesEvidence}.

\item[\emph{Halpern--Moses common knowledge (distributed-systems
epistemic discipline).}]
For a process \(i\) in a distributed system with local-state map
\(r_i\) and induced indistinguishability \(\sim_i\), the Kripke
semantics for \(K_i\varphi\) says that the truth of
\(K_i\varphi\) at a global state factors through \(r_i\): two
global states with the same local state for \(i\) make
\(K_i\varphi\) hold or fail together.  Accepted knowledge claims
by process \(i\) are therefore exactly those determined by its
view; a claim that depends on global state beyond \(r_i\) is the
overread failure
\citep{HalpernMoses1990KnowledgeCommonKnowledge}.
\end{description}

\subsection{No-Free-Distinction}\label{sec:access:no-free-distinction}

Keep the formed access quotient \(\accessQ:H\to Q_I\) from
Quotientality, together with the no-overread discipline from
Adequacy.  A proposed new current distinction is a readout
\(d:H\to D\).  Its split-pair obstruction at fixed access is
\[
  \splitObs_{\mathrm{dist}}(\accessQ,d)
  =
  \{(h,h')\mid \accessQ(h)=\accessQ(h')\ \text{and}\ d(h)\ne d(h')\}.
\]

\begin{theorem}[No-Free-Distinction]\label{thm:F12}
A functorial gate using only the existing access quotient cannot
produce a strict new current distinction.  If a readout
\(d:H\to D\) is current-visible without additional apparatus data,
then
\[
  \splitObs_{\mathrm{dist}}(\accessQ,d)=\varnothing,
  \qquad\text{equivalently}\qquad
  d=\overline d\circ\accessQ
\]
for some \(\overline d:Q_I\to D\).  If the obstruction is nonempty,
the distinction requires declared memory data, bridge data, residual
budget, scope change, or an enlarged calibration family.
\end{theorem}
\footnote{Verified in Lean as
\texttt{Access.NoFreeDistinction.no\_free\_distinction}.}
\begin{proof}
A gate that uses only the existing quotient can depend only on the
class \(\accessQ(h)\).  Therefore any readout it produces is
constant on the fibers of \(\accessQ\), and the displayed
split-pair obstruction is empty.  By \Cref{thm:F10}, this is
equivalent to the existence of a factorization
\(d=\overline d\circ\accessQ\).

Conversely, if \(\splitObs_{\mathrm{dist}}(\accessQ,d)\) is
nonempty, then \(d\) separates two histories that the current access
quotient identifies.  Treating \(d\) as a current exact distinction
would therefore overread the available access, contradicting
\Cref{thm:F11}.  The only lawful way to use \(d\) is to add declared
apparatus data: a memory record, a bridge, a budgeted residual, a
scope restriction, or a stronger calibration family.
\end{proof}

The theorem is not anti-refinement.  Calibration families may enrich
the apparatus, and strict extensions may add real distinctions.  The
claim is only that the enrichment must be declared and statused, not
smuggled in as a free consequence of the old access quotient.  Nor
does this theorem classify every possible source of new data; later
axes treat memory, bridges, residual budgets, scope changes, and
predictive refinements in their own normal forms.

\subsection*{Layer instantiations}\label{layer-inst:F12}

\begin{description}
\item[\emph{Landauer's principle (thermodynamics of computation).}]
Establishing a new reliable logical bit in a physical register
(e.g., a reset-to-zero operation that erases an unknown prior bit)
costs at least \(k_B T \ln 2\) of heat dissipated to the
environment.  The new logical distinction is not free: the
thermodynamic apparatus must export entropy at the Landauer bound,
and a protocol claiming a lower cost would compress phase-space
volume without compensating heat export, contradicting the second
law
\citep{Landauer1961IrreversibilityHeatGeneration}.

\item[\emph{Wolpert no-free-lunch theorem (statistical learning).}]
For supervised learning on a finite input space, averaged uniformly
over all possible target functions, every learning algorithm
achieves the same expected off-training-set error.  Off-sample
generalization is not a distinction obtainable from training data
alone; it requires the declared apparatus extension of an
inductive bias --- a hypothesis class restriction, a prior, or a
regulariser --- without which Wolpert's averaging identity collapses
every learner to chance
\citep{Wolpert1996NoFreeLunch}.

\item[\emph{Wootters--Zurek no-cloning theorem (quantum mechanics).}]
No unitary on \(\mathcal H_S\otimes\mathcal H_A\) sends
\(|\psi\rangle|0\rangle\) to \(|\psi\rangle|\psi\rangle\) for every
unknown \(|\psi\rangle\); the linearity of unitary evolution
forbids perfect copying of an unknown quantum state.  Acquiring a
distinguishing measurement of \(|\psi\rangle\) requires either
declaring a measurement basis (with collapse) or accepting the
fidelity bound of an approximate cloner --- the F12 apparatus
cost in either case
\citep{WoottersZurek1982NoCloning}.

\item[\emph{Goldwasser--Micali semantic security (cryptography).}]
For an IND-CPA secure probabilistic public-key encryption scheme,
no probabilistic polynomial-time adversary holding only the public
key and a challenge ciphertext can distinguish encryptions of two
adversarially chosen messages with non-negligible advantage.
Producing such a distinction requires apparatus extension --- the
secret key, a decryption oracle, or super-polynomial computation
breaking the underlying trapdoor (quadratic residuosity, factoring)
\citep{GoldwasserMicali1984ProbabilisticEncryption}.

\item[\emph{Myerson revenue equivalence (mechanism design).}]
In a symmetric independent-private-values auction, Myerson's
envelope-formula characterisation shows that an allocation rule
alone does not distinguish bidder types: interim expected payments
are fixed by the allocation up to a constant, and any
incentive-compatible mechanism extracting private values must pay
the bidders an information rent through the declared payment
schedule.  An allocation-only mechanism with no transfers collapses
to type-pooling
\citep{Myerson1981OptimalAuctionDesign}.
\end{description}

\subsection{Hiddenness Normal Form}\label{sec:access:hiddenness}

\begin{definition}[Hiddenness interface]\label{def:hiddenness-exposure}
A \emph{hiddenness interface} \(I\) consists of a closure carrier
\(\mathcal H_I\), a current quotient
\(\currQ_I:\mathcal H_I\to Q_I\), a predictive quotient
\(\predM_I:\mathcal H_I\to M_I\) refining \(\currQ_I\), and the
following predicates:
\[
  \operatorname{determiningStateTyped}(I),\quad
  \operatorname{minimalRelevantUpstream}(I),\quad
  \operatorname{lawfulCurrentExposure}(I),
\]
together with
\[
  \operatorname{QuotientExposed}(I),\quad
  \operatorname{PredictiveCollapse}(I),\quad
  \operatorname{PredictiveSurplus}(I),\quad
  \operatorname{HiddenFromCurrent}(I).
\]
The hidden-from-current predicate records that one current fiber
contains at least two predictive classes.  The quotient-exposed
predicate records that the lawful exposure quotient removes exactly
those predictive splits.  The hiddenness normal form is the
five-conjunct property
\begin{align*}
  \operatorname{HiddennessNormalForm}(I)
  :={}&
  \operatorname{determiningStateTyped}(I)
  \wedge
  \operatorname{minimalRelevantUpstream}(I)\\
  &{}\wedge
  \bigl(\operatorname{lawfulCurrentExposure}(I)
    \leftrightarrow
    \operatorname{QuotientExposed}(I)\bigr)\\
  &{}\wedge
  \bigl(\operatorname{PredictiveCollapse}(I)
    \leftrightarrow
    \operatorname{lawfulCurrentExposure}(I)\bigr)\\
  &{}\wedge
  \bigl(\operatorname{PredictiveSurplus}(I)
    \leftrightarrow
    \operatorname{HiddenFromCurrent}(I)\bigr).
\end{align*}
\end{definition}

\begin{theorem}[Hiddenness Normal Form]\label{thm:F13a}
Let \(I\) be a hiddenness interface and let
\[
  \mathsf{src}:
  \operatorname{Paper7SaturatedSATInstance}(I)
\]
be the opaque saturated-SAT source certificate.  Then
\begin{align*}
  &\operatorname{determiningStateTyped}(I),\\
  &\operatorname{minimalRelevantUpstream}(I),\\
  &\operatorname{lawfulCurrentExposure}(I)
    \iff \operatorname{QuotientExposed}(I),\\
  &\operatorname{PredictiveCollapse}(I)
    \iff \operatorname{lawfulCurrentExposure}(I),\\
  &\operatorname{PredictiveSurplus}(I)
    \iff \operatorname{HiddenFromCurrent}(I).
\end{align*}
That is, \(I\) satisfies
\(\operatorname{HiddennessNormalForm}(I)\).
\end{theorem}
\footnote{The Lean theorem at
\texttt{Access.HiddennessNormalForm.hiddenness\_normal\_form}
proves the layer-agnostic statement under a substrate-pending axiom
\texttt{hiddenness\_pending} guarded by an opaque source predicate;
the substrate proof appears in \citep{Tsiokos2026PvNP} and will
discharge the axiom when its Lean ships.  See
\Cref{app:formalization} for the narrowing disclosure.}
\begin{proof}
Apply the same-paper substrate-pending axiom
\(\mathrm{hiddenness\_pending}\), specialized to the hiddenness
interface \(I\) under the saturated-SAT source certificate \(\mathsf{src}\).
The axiom delivers the conjunction
\[
  \operatorname{HiddennessNormalForm}(I),
\]
that is, the five displayed conjuncts in the statement:
\begin{align*}
  &\operatorname{determiningStateTyped}(I),\qquad
  \operatorname{minimalRelevantUpstream}(I),\\
  &\operatorname{lawfulCurrentExposure}(I)
    \iff \operatorname{QuotientExposed}(I),\\
  &\operatorname{PredictiveCollapse}(I)
    \iff \operatorname{lawfulCurrentExposure}(I),\\
  &\operatorname{PredictiveSurplus}(I)
    \iff \operatorname{HiddenFromCurrent}(I).
\end{align*}
The substrate-side proof of the axiom lives in the saturated-SAT
analysis of \citet{Tsiokos2026PvNP}; within this paper the
wrapper records the layer-agnostic conjunction as the hiddenness
normal form.  The opacity-guard pattern of \Cref{sec:apparatus}
ensures that no instantiation of \(\mathrm{hiddenness\_pending}\) can
derive \(\mathrm{False}\) from the axiom alone: the guard
\(\operatorname{Paper7SaturatedSATInstance}\) has no public
constructor, so an actual instance must come from the saturated-SAT
substrate Lean.
\end{proof}

This theorem does not re-prove the saturated-SAT closure result of
the PvNP paper \citep{Tsiokos2026PvNP}, and it does not assert any
conclusion about a specific computational complexity class.  It
records the layer-agnostic access law: hiddenness means that a
predictive distinction exists upstream of current access, and
exposure is the minimal quotient refinement that removes that gap.
The classical quantum-foundations instance of a non-contextual
hidden-variable obstruction is the Kochen--Specker theorem
\citep{KochenSpecker1967Hidden}; that is a substrate-specific
result, not a special case of the layer-agnostic law above.

\subsection*{Layer instantiations}\label{layer-inst:F13a}

\begin{description}
\item[\emph{Hidden Markov model identifiability (statistical
latent-variable analysis).}]
For an HMM specified by transition matrix \(A\), emission matrix
\(B\), and initial distribution \(\pi\), the Allman--Matias--Rhodes
tensor-rank theorem establishes parameter identifiability up to
label-permutation when \(B\) has full column rank and \(A\) is
non-singular: distinct hidden Markov dynamics then produce distinct
joint emission distributions on observation triplets.  Predictive
surplus --- distinct future observation distributions inside a
current-observation fiber --- is the named witness of the hidden
state distinguishing upstream, with lawful exposure via Baum--Welch
EM or tensor-decomposition spectral methods
\citep{AllmanMatiasRhodes2009Identifiability}.

\item[\emph{Pearl back-door criterion and \(do\)-calculus (causal
inference).}]
In a causal DAG with observed variables \(V\) and unobserved
confounders \(U\), an unblocked back-door path from \(X\) to \(Y\)
through a common ancestor in \(U\) is the canonical witness of a
hidden upstream determining state.  The observational/interventional
gap \(P(Y\mid X)\ne P(Y\mid \mathrm{do}(X))\) is then the named
predictive surplus inside a \(P(V)\)-fiber, and \(do\)-calculus
characterises when causal effects remain identifiable via a
back-door adjustment set, an instrumental variable, or a front-door
mediator
\citep{Pearl1995CausalDiagrams}.

\item[\emph{Le Verrier's prediction of Neptune (celestial
mechanics).}]
In 1846, Le Verrier and (independently) Adams analysed the residual
perturbations of Uranus's orbit after accounting for all
gravitational influences then catalogued.  The unexplained
predictive surplus in Uranus's observed position was interpreted as
the gravitational signature of an undiscovered upstream determining
state, and Le Verrier computed an ecliptic-longitude prediction
which Galle confirmed telescopically on 1846 September 23 within
\(1^\circ\) of Le Verrier's position: the lawful exposure of the
hidden planet Neptune
\citep{LeVerrier1846RecherchesUranus}.

\item[\emph{{\AA}str{\"o}m belief-MDP reduction for {POMDPs}
(reinforcement learning).}]
For a partially observable Markov decision process, two histories
sharing the current observation can yield different distributions
over future observation--reward sequences exactly when their belief
states (posteriors over the hidden Markov state) differ.
{\AA}str{\"o}m's 1965 reduction theorem realises the lawful exposure:
the POMDP is equivalent to a fully observed MDP on the
continuous-belief simplex \(\Delta(S)\), with the belief state as
the minimal upstream-quotient refinement that resolves the
observation-fiber surplus
\citep{Astrom1965OptimalControlIncomplete}.

\item[\emph{Fisher's epistatic variance component (quantitative
genetics).}]
Under the Cockerham--Kempthorne orthogonal decomposition,
phenotypic variance partitions as
\(\operatorname{Var}(Y)=V_A+V_D+V_I+V_E\) with
\(V_I\) the epistatic interaction component.  A non-zero \(V_I\) is
the named witness of an upstream multi-locus determining state:
the phenotypic variability within a single-locus marginal fiber
that additive single-locus regression cannot explain corresponds
under the declared decomposition to hidden gene--gene interaction
structure
\citep{Fisher1918CorrelationRelatives}.
\end{description}

\subsection{CSL Formed-Closure Law}\label{sec:access:csl-formed-closure}

\begin{definition}[Diagnosis source data]\label{def:diagnosis-source}
Diagnosis source data for a closure carrier consists of a certified
source predicate saying that the carrier has the declared
measurement context, record update discipline, and stability
conditions required for a formed predictive closure.  In the CSL
case, the source predicate is supplied by the diagnosis-grade
apparatus of \PaperCSL\ \citep{Tsiokos2026CSL}.
\end{definition}

\begin{theorem}[CSL Formed-Closure Law]\label{thm:F13b}
Let \(\mathsf{FPC}\) be the class of formed predictive closures
\(\mathcal C\), each carrying a hiddenness interface with current
quotient \(\currQ\) and predictive quotient \(\predM\), together with
four declared side-condition predicates
\begin{align*}
  &\operatorname{genuinelyFormed}(\mathcal C),\qquad
  \operatorname{saturatedCurrentizers}(\mathcal C),\\
  &\operatorname{determiningStateDeclared}(\mathcal C),\qquad
  \operatorname{sourceReadoutSeparated}(\mathcal C).
\end{align*}
Write
\begin{align*}
  \operatorname{FormedClosureReady}(\mathcal C)
  :={}&
  \operatorname{genuinelyFormed}(\mathcal C)
  \wedge\operatorname{saturatedCurrentizers}(\mathcal C)\\
  &{}\wedge\operatorname{determiningStateDeclared}(\mathcal C)
  \wedge\operatorname{sourceReadoutSeparated}(\mathcal C).
\end{align*}
Let \(\operatorname{CSLSrc}:\mathsf{FPC}\to\mathrm{Prop}\) be the
opaque CSL diagnosis-source predicate.
Then for every \(\mathcal C\in\mathsf{FPC}\),
\begin{align*}
  \operatorname{CSLSrc}(\mathcal C)
  \Longrightarrow{}&
    \operatorname{FormedClosureReady}(\mathcal C)\\
  \Longrightarrow{}&
    \bigl(
      \operatorname{PredictiveSurplus}(\mathcal C)
      \iff
      \operatorname{HiddenFromCurrent}(\mathcal C)
    \bigr).
\end{align*}
Here \(\operatorname{PredictiveSurplus}\) and
\(\operatorname{HiddenFromCurrent}\) are the predictive-surplus and
hidden-from-current predicates of \Cref{def:hiddenness-exposure}.  In
a formed predictive closure with the CSL diagnosis source and the
four declared side conditions, predictive surplus across the
\(\predM/\currQ\) gap survives exactly when the determining state is
hidden from current observation.
\end{theorem}
\footnote{The Lean theorem
\texttt{Access.CSLFormedClosure.csl\_formed\_closure}
proves the layer-agnostic statement under an obligation axiom
\texttt{csl\_formed\_closure} guarded by an opaque source predicate;
no instantiation of the axiom can derive \texttt{False} from the
axiom alone.  The diagnosis context appears in
\citep{Tsiokos2026CSL}.  See \Cref{app:formalization} for the guard
register.}
\begin{proof}
Fix a formed predictive closure \(\mathcal C\in\mathsf{FPC}\) and
suppose \(\operatorname{CSLSrc}(\mathcal C)\).  Apply the same-paper
obligation axiom \(\mathrm{csl\_formed\_closure}\), specialized to
\(\mathcal C\) under this opaque CSL diagnosis source.  The axiom
delivers
\begin{align*}
  \operatorname{CSLFormedClosureLaw}(\mathcal C)
  ={}&
  \operatorname{FormedClosureReady}(\mathcal C)\\
  &\Longrightarrow
  \bigl(
    \operatorname{PredictiveSurplus}(\mathcal C)
    \iff
    \operatorname{HiddenFromCurrent}(\mathcal C)
  \bigr).
\end{align*}
Assume further that the four formed-closure side conditions hold,
i.e.\ \(\operatorname{FormedClosureReady}(\mathcal C)\).  The axiom
then gives the displayed iff:
\[
  \operatorname{PredictiveSurplus}(\mathcal C)
  \iff
  \operatorname{HiddenFromCurrent}(\mathcal C).
\]
Unpacking: predictive surplus is the existence of a
\(\predM\)-distinct pair inside a \(\currQ\)-fiber;
hidden-from-current is the existence of two upstream-distinct states
identified by \(\currQ\) and separated only at \(\predM\).  The axiom
records that, in a genuinely formed, currentizer-saturated,
determining-state-declared, source-readout-separated closure, these
two conditions are the same: predictive surplus survives precisely on
the hidden-determining-state fibers.

The substrate-side derivation of the axiom lives in
\citep{Tsiokos2026CSL}; within this paper the wrapper records the
layer-agnostic equivalence as the CSL formed-closure law.  The
opacity-guard pattern of \Cref{sec:apparatus} ensures that no
instantiation of \(\mathrm{csl\_formed\_closure}\) can derive
\(\mathrm{False}\) from the axiom alone: the guard
\(\operatorname{CSLSrc}\) (Lean
\texttt{Paper8CSLDiagnosisSource}) has no public constructor, so an
actual instance must come from the CSL substrate Lean.
\end{proof}

This subsection does not instantiate the CSL diagnosis source, does
not assert a measurement-theoretic contextuality theorem at the
\FIV\ layer, and does not close the source obligation
unconditionally.  It records the guarded source needed to treat CSL
formed-closure as a diagnosis-grade entry in the access catalog.

\subsection*{Layer instantiations}\label{layer-inst:F13b}

\begin{description}
\item[\emph{Bayesian clinical diagnostic testing (medicine).}]
For a formed clinical assay with declared sensitivity \(\mathrm{Se}\)
and specificity \(\mathrm{Sp}\), Bayes' rule maps a declared prior
\(P(\text{disease})\) to a posterior \(P(\text{disease}\mid r)\) for
each result \(r\).  The post-test probability gap between
positive- and negative-result populations within an otherwise
identical declared cohort is the named predictive surplus
witnessing the patient's true disease status as the hidden upstream
determining state, with lawful exposure under the
evidence-based-medicine discipline (validated test, calibrated
cutoff, declared prevalence)
\citep{Sackett2000EvidenceBasedMedicine}.

\item[\emph{Shewhart statistical process control (industrial
quality engineering).}]
In a formed manufacturing closure with declared in-control mean
\(\mu\), standard deviation \(\sigma\), and \(3\sigma\) control
limits calibrated from a baseline period, an out-of-control signal
(point outside the control limits or a non-random run pattern) is
the named predictive surplus that triggers root-cause investigation
into an assignable cause --- tool wear, raw-material change,
operator effect, environmental drift --- under FormedClosureReady
conditions (declared sampling rule, baseline independence and
identical distribution).  The control-chart signal is a witness, not
a proof: false alarms exist under the in-control distribution at the
declared rate
\citep{Shewhart1931EconomicControl}.

\item[\emph{Model checking against a temporal-logic specification
(formal verification).}]
For a finite-state model \(M\) and a linear-temporal-logic
specification \(\varphi\), an explicit-state or symbolic
model-checker produces a counterexample trace \(\tau\) when
\(M\not\models\varphi\).  The trace is the named predictive surplus
witnessing a hidden upstream defect in the implementation that the
declared specification was meant to forbid, with lawful exposure via
the automated checker (SPIN, NuSMV, TLA+) under FormedClosureReady
conditions (closed system, finite state, fair scheduler, sound LTL
semantics)
\citep{ClarkeEmersonSistla1986AutomaticVerification}.

\item[\emph{Exoplanet detection by radial velocity and transit
(observational astronomy).}]
For a formed exoplanet-detection closure with declared instrument
(spectrograph + Doppler pipeline; or photometric transit pipeline)
and declared sensitivity (Doppler precision in m\,s\(^{-1}\);
photometric precision in ppm), a periodic radial-velocity
oscillation or photometric dimming inconsistent with stellar
activity is the named predictive surplus witnessing a hidden
orbiting body.  Mayor and Queloz realised this lawful exposure with
the 1995 detection of 51 Pegasi b, the first confirmed exoplanet
around a Sun-like star
\citep{MayorQueloz199551Pegasi}.

\item[\emph{{LHC} Higgs-boson discovery as resonance detection
(high-energy physics).}]
For a formed LHC search closure with declared integrated luminosity,
detector acceptance \(\times\) efficiency, and background-only
prediction (Standard-Model processes excluding the hypothesised
Higgs), an observed event yield exceeding the background-only model
by \(\ge 5\sigma\) in a signal region is the named predictive
surplus witnessing a hidden upstream resonance / particle predicted
by the broader theory.  The ATLAS 2012 observation of a 125 GeV
resonance in declared diphoton and four-lepton mass windows realised
the lawful exposure, with follow-up spin-parity and branching-ratio
measurements identifying the Standard-Model Higgs boson
\citep{ATLAS2012HiggsObservation}.
\end{description}

\subsection{Layer Multiplicity}\label{sec:access:layer-multiplicity}

Let \(H\) be a closure carrier with a formed current access quotient
\(q:H\to Q\).  A calibration or role signature is another declared
map \(s:H\to S\): it may record a future probe, an interaction role,
a lawhood signature, a branch readout, or a calibrated instrument
response.  The signature is already carried by the old layer exactly
when it descends through \(q\), equivalently when it is constant on
the \(q\)-fibers.

\begin{theorem}[Layer Multiplicity]\label{thm:F24}
Let \(H\) be a history space, let \(q:H\to Q\) be a surjective access
quotient, and let \(s:H\to S\) be a declared role readout.  Define
the role obstruction and joint role-layer quotient by
\[
  \mathcal O_s
  :=
  \{(h,h')\in H^2:q(h)=q(h')\ \text{and}\ s(h)\ne s(h')\},
  \qquad
  q_s(h):=(q(h),s(h)).
\]
Then
\begin{align*}
  s\ \text{descends through}\ q
  &\iff \mathcal O_s=\emptyset,\\
  \operatorname{RoleSplit}(q,s)
  &\iff \mathcal O_s\ne\emptyset\\
  &\iff q_s\ \text{is a strict refinement of}\ q.
\end{align*}
When \(\operatorname{RoleSplit}(q,s)\) holds, exactly one resolution
case is selected by the layer-formation discipline:
\begin{align*}
  &\operatorname{MemoryLayer}
  \mathbin{\dot\vee}
  \operatorname{HiddenUpstreamRole}
  \mathbin{\dot\vee}
  \operatorname{BridgeMediatedRole}\\
  &\mathbin{\dot\vee}
  \operatorname{BudgetedRole}
  \mathbin{\dot\vee}
  \operatorname{ScopedRole}
  \mathbin{\dot\vee}
  \operatorname{CoarsenedRole}\\
  &\mathbin{\dot\vee}
  \operatorname{OutsideRoleScope}
  \mathbin{\dot\vee}
  \operatorname{BlockedNonClosure}.
\end{align*}
\end{theorem}
\footnote{Verified in Lean as
\texttt{Access.LayerMultiplicity.layer\_multiplicity}; see
\Cref{app:formalization}.}
\begin{proof}
Assume \(q:H\to Q\) is surjective and \(s:H\to S\) is declared.  If
\(s\) descends through \(q\), then \(s=\bar s\circ q\) for some
\(\bar s:Q\to S\), so
\[
  q(h)=q(h')\Longrightarrow s(h)=s(h')
\]
and \(\mathcal O_s=\emptyset\).  Conversely, if
\(\mathcal O_s=\emptyset\), define
\[
  \bar s(q(h)):=s(h).
\]
The empty obstruction makes \(\bar s\) independent of the
representative \(h\), and surjectivity of \(q\) supplies all quotient
classes.  Thus \(s\) descends through \(q\).

If \(\mathcal O_s\ne\emptyset\), choose
\((h,h')\in\mathcal O_s\).  Then \(q(h)=q(h')\) but
\[
  q_s(h)=(q(h),s(h))\ne(q(h'),s(h'))=q_s(h'),
\]
so \(q_s\) strictly refines \(q\).  Conversely, a strict refinement of
\(q\) by \(q_s\) means that some \(q\)-fiber is split by the second
coordinate \(s\), giving an element of \(\mathcal O_s\).

Finally, the layer-formation discipline decides the resolution state
of a role split.  Its cases are mutually exclusive by construction:
memory, hidden upstream role, bridge mediation, budget, scope,
coarsening, outside-scope, and blocked nonclosure are distinct status
families.  They are exhaustive because every declared split must be
formed, hidden, mediated, budgeted, scoped, coarsened, excluded, or
blocked.
\end{proof}

The interaction status is \(\BirdInt\)-compatible because every
signature that contributes to the layer decision is declared: as a
calibration, future probe, role readout, bridge, budgeted residual,
or scoped restriction.  The discipline is that generic pressure
toward new layers is theorem-grade, but a specific new layer still
requires admissibility, closure formation, and status.
Layer Multiplicity does not classify the number of layers in a
carrier, does not prescribe how layers must be presented, and does
not assert that multiplicity is finite in every setting.

\subsection*{Layer instantiations}\label{layer-inst:F24}

\begin{description}
\item[\emph{Operating-system privilege protection (computer
security).}]
In a multi-user operating system, the privilege quotient classifies
operations by the trust level of the executing principal.  When a
process needs an operation outside its declared privilege, the
layer-formation discipline names a resolution: a user-to-kernel
system-call trap, a capability-bridge invocation, an
inter-process-communication transfer, a sandbox or namespace
restriction, or
outright denial --- each instantiating a specific F24
layer-formation category (memory layer, bridge, scope, blocked
non-closure)
\citep{SaltzerSchroeder1975ProtectionInformation}.

\item[\emph{{ANSI} {SQL} transaction isolation levels (database
concurrency).}]
For concurrent database transactions, the isolation-level quotient
classifies the visibility of in-flight operations.  When a
transaction's read or write role would split that quotient ---
producing a dirty read, a non-repeatable read, or a phantom ---
the declared SQL resolution discipline forces the system into one
of the four ANSI isolation levels (read-uncommitted, read-committed,
repeatable-read, serializable), each pricing different role-split
obstructions as budgeted or scoped
\citep{BerensonBernsteinGrayMeltonONeilONeil1995SQLIsolation}.

\item[\emph{Grammatical case-marking (structural linguistics).}]
In a case-marking language the morphological word-form quotient
classifies nouns by inflection.  When a syntactic role (subject,
object, oblique, possessor) splits that morphological quotient, the
language's case system resolves the obstruction by assigning a
named case --- nominative, accusative, ergative, absolutive,
dative, genitive, locative, instrumental --- from Blake's
typological inventory.  Different languages instantiate different
subsets of the inventory, but each role-split is closed by a named
case-formation category
\citep{Blake2001Case}.

\item[\emph{{OSI} seven-layer networking reference model (computer
networking).}]
In a networking stack the protocol-layer quotient classifies
communication functions.  When a new communication role splits the
current layer quotient, the ISO/IEC 7498-1 model resolves the
obstruction by placing the function at one of seven named layers
(physical, data-link, network, transport, session, presentation,
application), each licensing specific function types and obstructing
others
\citep{ISOIEC7498OSI}.

\item[\emph{Polymorphism taxonomy (programming-language theory).}]
In a typed language the type quotient classifies expressions by
their monomorphic type.  When a function's required role is
generic across multiple types --- a role-split obstruction in the
monomorphic quotient --- the Cardelli--Wegner taxonomy resolves
the obstruction by assigning the function to a named polymorphism
category: parametric (universal quantification), ad-hoc (type
classes / overloading), or inclusion (subtyping).  Each category
licenses a distinct layer-formation mechanism for the generic role
\citep{CardelliWegner1985UnderstandingTypes}.
\end{description}

\subsection*{Access at a glance}\label{sec:access:summary}
\begin{table}[t]
  \centering
  \footnotesize
  \begin{tabularx}{\textwidth}{@{}
      >{\raggedright\arraybackslash}p{0.24\textwidth}
      >{\raggedright\arraybackslash}X
      >{\raggedright\arraybackslash}p{0.18\textwidth}
      >{\raggedright\arraybackslash}p{0.12\textwidth}@{}}
    \toprule
    Paper-prose name & One-line claim & Status & Substrate cite \\
    \midrule
    Quotientality theorem &
    Lawful observation signatures determine canonical access
    quotients. &
    Theorem &
    \textemdash \\
    Adequacy / No-Overread theorem &
    Accepted claims factor through available access. &
    Theorem &
    \textemdash \\
    No-Free-Distinction theorem &
    Refinements require declared apparatus data, not free gates. &
    Theorem &
    \textemdash \\
    Hiddenness Theorem &
    Exposure increments are bounded by the current/predictive gap. &
    Partial wrapper &
    \PaperPvNP \\
    CSL Formed-Closure &
    Opaque CSL diagnosis source \(\Rightarrow\) formed-closure law. &
    Obligation wrapper &
    \PaperCSL \\
    Layer-Multiplicity Theorem &
    Limited access forces multiple coexisting layers. &
    Theorem &
    \textemdash \\
    \bottomrule
  \end{tabularx}
  \caption{Access axis at a glance: the six laws of \Cref{sec:access}.}
  \label{tab:axis-access-summary}
\end{table}
  \section{Axis II — Stability}\label{sec:stability}

The Stability axis asks when declared future behavior is already
controlled by present data, and when a residual must be recorded
instead of ignored.  In the language of \Cref{sec:apparatus}, the
question is whether the current quotient carries the probes the
package has declared, or whether closure-stability fails under a
structural perturbation.  A perturbation may be a new future probe,
a layer-dissolving move at a membrane, a strict extension of a
quotient, or an unpriced residual.  The axis records the normal
forms that decide which of these outcomes is legitimate.

The chapter has four laws.  No-Needles Stability is the lead theorem:
a formed stability package cannot admit an unpriced
layer-dissolving intrusion.  Sufficiency Closure then gives the
factorization criterion for declared futures: the current quotient is
sufficient exactly when every declared future probe descends through
it.  Strict Extension records when genuine layer growth has occurred,
namely when the old quotient is retained and a new quotient splits at
least one old fiber.  No Unstatused Residual closes the axis by
requiring every native residual of a formed stability package to
carry a typed status.

No-Needles Stability is anchored to the substrate results in
\PaperNeedleKiller\ \citep{Tsiokos2026NeedleKiller} and the
Navier--Stokes application in \PaperNS\
\citep{Tsiokos2026NavierStokes}.  Here a membrane means the
layer-dissolving boundary of a closure: the place where a distinction
between ``inside'' and ``outside'' the layer can cease to be
meaningful.  Strict Extension also has descriptive substrate context,
but it is a direct law here: the Schur-calculus intuition comes from
\PaperAR\ \citep{Tsiokos2026XiCompanion}, where residuals compose
along refinement chains.  The Stability axis can be read
independently of the Access axis; its only prerequisite is the
notation of \Cref{sec:apparatus}.

\subsection{No-Needles Stability}\label{sec:stability:no-needles}

The setup is a formed stability package: a closure carrier equipped
with the inherited closure apparatus, quotient discipline, audit
ledger, and a layer-dissolving membrane.  The membrane is the
boundary across which a native layer distinction is allowed to
dissolve only when the package has recorded the appropriate
apparatus footprint.  A needle is a structural intrusion across that
boundary with no such footprint: no declared audit data, no residual
budget, no transport price, and no typed status for the defect it
would create.

\begin{theorem}[No-Needles Stability]\label{thm:F6}
Let \(\mathcal P\) be a formed stability package with a
layer-dissolving membrane.  Suppose \(\mathcal P\) declares:
\begin{itemize}
  \item a currency block \(C\) on its exterior carrier of size \(e\);
  \item a layer-bridge block \(L\) of size \(y\times e\) and a
  dissolving-direction block \(D\) of size \(z\times e\), with block
  currency pieces \(K_{LL}\), \(K_{LD}\), \(K_{DL}\), and \(K_{DD}\);
  \item the symmetry conditions \(K_{DL}=K_{LD}^{\!\top}\) and
  \(K_{LL}^{\dagger}=(K_{LL}^{\dagger})^{\!\top}\), together with the
  Moore--Penrose projection identity
  \[
    K_{LL}\,K_{LL}^{\dagger}\,K_{LD}=K_{LD};
  \]
  \item currency and residual budgets
  \[
    K_{LL}\preceq \Theta_Y,
    \qquad
    \Xi(C,L,D,K_{LL}^{\dagger})\preceq \Omega_Z,
  \]
  where \(\Xi\) is the adequacy residual;
  \item a transport selector \(S\) of size \(d\times z\) satisfying
  the endpoint budget
  \[
    S\bigl(\widehat E\,\Theta_Y\,\widehat E^{\!\top}
      +\Omega_Z\bigr)S^{\!\top}\preceq\Theta_D,
  \]
  where \(\widehat E\) is the optimal native explanation determined
  by \((C,L,D,K_{LL}^{\dagger})\).
\end{itemize}
Then the transported dissolving block satisfies the Loewner endpoint
\[
  S\,K_{DD}\,S^{\!\top}\preceq\Theta_D .
\]
Equivalently, every layer-dissolving move accepted by \(\mathcal P\)
is bounded by its declared transport budget; no unpriced
layer-dissolving needle survives the formed package.
\end{theorem}
\footnote{Anchored to the substrate Lean at
\texttt{Stability.NoNeedles.\allowbreak{}no\_needles\_stability};
the substrate proof appears in
\citep{Tsiokos2026NeedleKiller, Tsiokos2026NavierStokes}.  See
\Cref{app:formalization} for the anchor-disclosure discipline.}
\begin{proof}
Assume the symmetry, projection, currency-budget, residual-budget,
and transport-budget hypotheses displayed in the theorem.  The
Moore--Penrose identity
\[
  K_{LL}\,K_{LL}^{\dagger}\,K_{LD}=K_{LD}
\]
gives the layer-dissolving decomposition inherited from the
substrate master theorem:
\[
  K_{DD}
  =
  \widehat E\,K_{LL}\,\widehat E^{\!\top}
  +\Xi(C,L,D,K_{LL}^{\dagger}),
  \qquad
  \widehat E:=K_{DL}\,K_{LL}^{\dagger}.
\]
Here \(\widehat E\) is the optimal native explanation and \(\Xi\) is
the adequacy residual.  Applying Loewner monotonicity under the
declared bounds \(K_{LL}\preceq\Theta_Y\) and
\(\Xi(C,L,D,K_{LL}^{\dagger})\preceq\Omega_Z\) yields
\[
  K_{DD}
  \preceq
  \widehat E\,\Theta_Y\,\widehat E^{\!\top}+\Omega_Z .
\]
Conjugation by \(S\) preserves the Loewner order, so
\[
  S\,K_{DD}\,S^{\!\top}
  \preceq
  S\bigl(\widehat E\,\Theta_Y\,\widehat E^{\!\top}
    +\Omega_Z\bigr)S^{\!\top}.
\]
The declared transport-side endpoint budget gives
\[
  S\bigl(\widehat E\,\Theta_Y\,\widehat E^{\!\top}
    +\Omega_Z\bigr)S^{\!\top}
  \preceq \Theta_D .
\]
Transitivity of \(\preceq\) therefore gives the endpoint
\[
  S\,K_{DD}\,S^{\!\top}\preceq\Theta_D .
\]

The carrier-specific realization of the decomposition, the
Moore--Penrose identity, and the optimal native explanation is the
layer-dissolving master theorem of \citet{Tsiokos2026NeedleKiller};
the Navier--Stokes instantiation appears in
\citet{Tsiokos2026NavierStokes}.  This anchor row records
the layer-agnostic Loewner endpoint instantiated by those substrate
carriers, while \Cref{app:formalization} records the anchor-disclosure
discipline.
\end{proof}

The theorem does not re-prove the Needle Killer or Navier--Stokes
substrate instances, does not classify every possible
layer-dissolving move, and does not assert that the substrate carrier
is unique.  It records the common stability normal form: at a formed
membrane, layer-dissolving behavior must be priced, statused, or
rejected.

\subsection*{Layer instantiations}\label{layer-inst:F6}

\begin{description}
\item[\emph{Schur-complement positivity criterion (linear algebra).}]
A symmetric block matrix
\(\bigl[\begin{smallmatrix}A&B\\B^{\!\top}&C\end{smallmatrix}\bigr]\)
is positive semidefinite iff \(A\succeq0\) and
\(C-B^{\!\top}A^{+}B\succeq0\).  This is the canonical block-matrix
instantiation of \Cref{thm:F6}: the Moore--Penrose pseudoinverse
\(A^{+}\) plays the role of the structural inverse, and the
transported endpoint is the Schur complement of the layered
partition \citep{Zhang2005SchurApplications}.

\item[\emph{Lindblad--Uhlmann monotonicity of quantum relative entropy
(quantum information).}]
For any completely positive trace-preserving (CPTP) quantum channel
\(T\) and density operators \(\rho,\sigma\), the relative entropy
satisfies \(S(T(\rho)\,\|\,T(\sigma))\le S(\rho\,\|\,\sigma)\).  The
decomposition is the Stinespring dilation
\(T(\rho)=\operatorname{Tr}_E[V\rho V^{*}]\); the monotonicity
combines joint convexity of relative entropy (Lindblad's lemma) with
isometric embedding and partial trace.  The preorder is scalar
\(\le\) on entropy functionals, not the Loewner block ordering
\citep{Lindblad1975CompletelyPositiveMaps}.

\item[\emph{Csiszár \(f\)-divergence data-processing inequality
(information theory).}]
For any convex \(f:(0,\infty)\to\mathbb R\) with \(f(1)=0\), the
\(f\)-divergence
\(D_f(P\,\|\,Q)=\mathbb E_Q\!\bigl[f(\mathrm dP/\mathrm dQ)\bigr]\) is
monotone under every Markov kernel \(K\):
\(D_f(K\!\cdot\!P\,\|\,K\!\cdot\!Q)\le D_f(P\,\|\,Q)\).  The
decomposition is the Csiszár chain rule for divergences; the
monotonicity is the joint convexity of \((p,q)\mapsto qf(p/q)\) plus
Jensen on the kernel mixture.  The preorder is scalar \(\le\) on a
divergence value rather than the Loewner cone
\citep{LieseVajda2006Divergences}.

\item[\emph{Strassen monotone-coupling stochastic dominance
(probability theory).}]
By Strassen's theorem, \(X\le_{\mathrm{st}}Y\) iff there is a
coupling \(\pi\) of \(\mu_X\) and \(\mu_Y\) supported on
\(\{(x,y):x\le y\}\); applying a stochastically monotone Markov
kernel \(K\) to both coordinates preserves the order, so
\(K\!\cdot\!\mu_X\le_{\mathrm{st}}K\!\cdot\!\mu_Y\).  The
decomposition is the coupling existence (a measure-theoretic object
on a product space); the preorder is stochastic dominance defined by
tail-probability inequalities, with no positive-semidefinite cone
involved \citep{Strassen1965ExistenceProbabilityMeasures}.

\item[\emph{Lyapunov's direct method (dynamical systems).}]
For a dissipative dynamical system \(\dot x=f(x)\) with a
positive-definite Lyapunov functional \(V\) and a non-negative
dissipation
residual \(W\), the Lie-derivative inequality
\(\dot V=\nabla V\!\cdot\!f(x)\le-W(x)\le 0\) along trajectories
propagates \(V(x(t))\le V(x_0)\) for all \(t\ge 0\).  The
decomposition is the pair \((V,W)\) chosen against the vector field's
structure; transitivity composes through the flow semigroup
\(\varphi_t\), giving the F6 endpoint along an infinite-dimensional
flow rather than via a matrix Schur complement
\citep{Lyapunov1892StabilityMotion}.
\end{description}

\subsection{Sufficiency Closure}\label{sec:stability:sufficiency-closure}

Fix a closure package with history carrier \(H\), current quotient
map \(q:H\to\currQ\), and a declared family \(\mathcal F\) of future
probes.  The family is the apparatus's announcement of which future
observables will be tested at this interface; it is not a claim of
sufficiency for every possible future question.

\begin{theorem}[Sufficiency Closure]\label{thm:F7}
The current quotient \(\currQ\) is sufficient for the declared future
family \(\mathcal F\) if and only if every declared future probe
\(f\in\mathcal F\) descends through \(q\).  Equivalently, for every
\(f:H\to F_f\) in \(\mathcal F\) there is a map
\(\overline f:\currQ\to F_f\) such that
\[
  f=\overline f\circ q .
\]
\end{theorem}
\footnote{Verified in Lean as
\texttt{Stability.SufficiencyClosure.\allowbreak{}sufficiency\_closure};
see \Cref{app:formalization}.}
\begin{proof}
Assume first that every declared future probe descends through
\(q\).  Then each future value \(f(h)\) is determined by the current
class \(q(h)\).  Hence two histories with the same current quotient
have the same value under every declared future probe, so no
declared future distinction remains hidden inside a current fiber.
This is precisely sufficiency for the declared family.

Conversely, suppose \(\currQ\) is sufficient for the declared future
family.  For a probe \(f:H\to F_f\), define a candidate lower map on
a current class \(q(h)\) by \(\overline f(q(h))=f(h)\).  Sufficiency
is exactly the assertion that this assignment is independent of the
chosen representative: if \(q(h)=q(h')\), then the declared future
record gives \(f(h)=f(h')\).  Therefore \(\overline f\) is
well-defined and satisfies \(f=\overline f\circ q\).  Since \(f\)
was arbitrary, every declared future probe descends through the
current quotient.
\end{proof}

The interaction status is \(\BirdInt\)-compatible because the future
probes are declared interactions and the sufficiency claim is audited
relative to that declaration.  The theorem does not enumerate the
declared future probes; that choice belongs to the apparatus.  It
also does not assert sufficiency for undeclared probes, which lie
outside the current audit scope.

\subsection*{Layer instantiations}\label{layer-inst:F7}

\begin{description}
\item[\emph{Information bottleneck (information theory / machine
learning).}]
For input \(X\) and a declared prediction target \(Y\), the
Tishby--Pereira--Bialek information bottleneck seeks an encoder
\(T=q(X)\) minimising the compression cost \(I(X;T)\) under the
sufficiency constraint \(I(T;Y)=I(X;Y)\).  The constraint is exactly
the F7 factorization for the declared family
\(\mathcal F=\{Y\text{-predictor}\}\): the Bayes predictor
\(P(Y\mid X)\) factors as \(P(Y\mid T)\circ q\).  Over-compressing the
encoder breaks the factorization, and the obstruction is the gap
\(I(X;Y)-I(T;Y)\)
\citep{TishbyPereiraBialek1999InformationBottleneck}.

\item[\emph{Predictive state representations (reinforcement
learning).}]
A PSR summarizes the history of a partially observable dynamical
system by the vector \(p(h)\) of probabilities of a declared family of
core future tests.  The Singh--James--Rudary sufficiency theorem
states that every declared future test factors linearly through
\(p(h)\): \(\Pr(t\mid h)=m_t\cdot p(h)\).  The summary is sufficient
for the declared test family while remaining hidden-state-free, and
enlarging the declared family beyond the linear span of the core
forces a finer summary
\citep{SinghJamesRudary2004PredictiveStateRepresentations}.

\item[\emph{Query-complete materialized views (databases).}]
For a base schema \(R\) and a declared workload \(\mathcal Q\) of
conjunctive queries, a set of materialized views \(V\) is
\emph{query-complete} for \(\mathcal Q\) iff every \(Q\in\mathcal Q\)
admits a rewriting \(Q'\) over the view symbols with
\(Q(I)=Q'(V(I))\) for every instance \(I\).  The rewriting plays the
F7 lower map; the view materialization plays the summary.  The
sufficiency is workload-relative: a query needing an attribute
projected away by \(V\) does not factor through the views
\citep{Halevy2001AnsweringQueriesUsingViewsSurvey}.

\item[\emph{Quantum sufficient subalgebras (quantum information).}]
For a parametric family of density operators
\(\{\rho_\theta:\theta\in\Theta\}\) and a CPTP quantum channel \(T\),
Petz's theorem says \(T\) is sufficient for the family iff there is a
recovery channel \(R\) (the Petz recovery map) with
\(\rho_\theta=R(T(\rho_\theta))\) for every \(\theta\).  Every declared
measurement statistic \(\operatorname{Tr}(\rho_\theta M_i)\) then
factors through \(T\), and the obstruction is strict contraction of
relative entropy by \(T\), which prevents recovery
\citep{Petz1986SufficientSubalgebras}.

\item[\emph{Noninterference and information-flow security (security /
programming languages).}]
A program operating over states with high-security (secret) and
low-security (public) components satisfies noninterference iff every
declared low-observer factors through the low-projection
\(q:H\to H{\restriction}L\).  Side-channel attacks --- timing, cache,
fault --- appear as family-extension bridges: the declared
low-observer family is enlarged to include an observer for which the
low-projection is no longer sufficient
\citep{GoguenMeseguer1982SecurityPolicies,SabelfeldMyers2003LanguageBasedInformationFlow}.
\end{description}

\subsection{Strict Extension}\label{sec:stability:strict-extension}

Let \(H\) be the carrier under discussion, let
\[
  q_{\mathrm{old}}:H\to \currQ_{\mathrm{old}}
\]
be the old quotient, and let
\[
  q_{\mathrm{new}}:H\to \currQ_{\mathrm{new}}
\]
be a candidate new quotient.  A new presentation is not
automatically a new layer: it may be a relabeling, a coarsening, or
an incomparable rewrite.  The Schur-calculus intuition from
\PaperAR\ \citep{Tsiokos2026XiCompanion} is that residuals compose
along refinement chains, so genuine layer growth preserves the old
quotient while accounting for the new split.

\begin{lemma}[Strict Extension]\label{lem:F8}
The extension from \(q_{\mathrm{old}}\) to \(q_{\mathrm{new}}\) is
strict if and only if there is a retention map
\[
  \rho:\currQ_{\mathrm{new}}\to \currQ_{\mathrm{old}},
  \qquad
  q_{\mathrm{old}}=\rho\circ q_{\mathrm{new}},
\]
and there are histories \(h,h'\in H\) such that
\[
  q_{\mathrm{old}}(h)=q_{\mathrm{old}}(h')
  \quad\text{but}\quad
  q_{\mathrm{new}}(h)\ne q_{\mathrm{new}}(h').
\]
\end{lemma}
\footnote{Verified in Lean as
\texttt{Stability.StrictExtension.\allowbreak{}strict\_extension};
see \Cref{app:formalization}.}
\begin{proof}
Assume the extension is strict.  By the retention part of strictness,
the old quotient is still readable from the new one; hence there is
a map \(\rho:\currQ_{\mathrm{new}}\to\currQ_{\mathrm{old}}\) with
\(q_{\mathrm{old}}=\rho\circ q_{\mathrm{new}}\).  By the genuine
growth part of strictness, the new quotient is not merely a relabeling
of the old quotient.  Therefore some old fiber is split, so there
are \(h,h'\in H\) with \(q_{\mathrm{old}}(h)=q_{\mathrm{old}}(h')\)
but \(q_{\mathrm{new}}(h)\ne q_{\mathrm{new}}(h')\).

Conversely, suppose such a retention map and split old fiber are
given.  The equation \(q_{\mathrm{old}}=\rho\circ q_{\mathrm{new}}\)
says that every new class lies inside an old class, so the candidate
retains the old quotient rather than replacing it.  The displayed
pair \(h,h'\) shows that at least one old class has been split by
the new quotient.  Thus the refinement is proper and retentive,
which is exactly strict extension.
\end{proof}

The interaction status is \(\BirdInt\)-compatible when the new split
is declared by the apparatus as a lawful probe, memory, role, or
extension record.  This lemma does not re-derive the Schur calculus,
and it does not classify every kind of layer growth.  It only gives
the quotient test for genuine retentive growth.

\subsection*{Layer instantiations}\label{layer-inst:F8}

\begin{description}
\item[\emph{Ring species (evolutionary biology).}]
For a ring-distributed population complex (the \emph{Ensatina}
salamander complex around California's Central Valley, the
\emph{Larus} gull ring), the local interbreeding equivalence
\(q_{\mathrm{old}}\) identifies adjacent populations that interbreed
in the wild --- yielding a single ring-wide class.  The reproductive
equivalence \(q_{\mathrm{new}}\) at the ring closure refines this:
the terminal-pair populations are reproductively isolated.  The
retention map is the inclusion of new classes into the single old
class, and the witness split is the closure pair, reproductively
isolated yet connected through the chain of locally interbreeding
intermediates \citep{IrwinBenschPrice2001SpeciationInARing}.

\item[\emph{Hematopoietic CMP \(\to\) MEP/GMP bifurcation (stem-cell
biology).}]
The common myeloid progenitor (CMP) population, prospectively
isolated by Lin\(^{-}\) Sca-1\(^{-}\) c-Kit\(^{+}\) CD34\(^{+}\)
Fc\(\gamma\)RII/III\(^{\mathrm{lo}}\) surface markers, splits into
megakaryocyte--erythroid (MEP) and granulocyte--macrophage (GMP)
progenitors distinguished by CD34 and Fc\(\gamma\)RII/III.  The CMP
gating is retained by both daughters, and the witness split is two
clonally related siblings of one CMP division that commit to MEP and
GMP respectively \citep{AkashiTraverMiyamotoWeissman2000CMPClonogenic}.

\item[\emph{Refinement mappings in formal verification (computer
science).}]
For a high-level specification \(S_2\) and a candidate implementation
\(S_1\), the Abadi--Lamport refinement-mapping theorem states that
\(S_1\) implements \(S_2\) if and only if, possibly after adjoining
history and prophecy variables, there exists a state-level map
\(f:\mathrm{States}(S_1)\to\mathrm{States}(S_2)\) sending each \(S_1\)
step to a legal \(S_2\) step.  The map \(f\) is the F8 retention map;
the witness split is two implementation behaviors with identical
\(S_2\)-observable trace but distinct internal scheduling or
auxiliary-variable values
\citep{AbadiLamport1991RefinementMappings}.

\item[\emph{Parisi replica-symmetry-breaking hierarchy (statistical
mechanics).}]
In Parisi's mean-field solution of the Sherrington--Kirkpatrick spin
glass, the one-step replica-symmetry-breaking ansatz partitions
overlaps between pure states into two values \(\{q_0,q_1\}\); full RSB
refines this to a continuous overlap function \(q(x)\) for
\(x\in[0,1]\).  The 1-RSB cluster structure is the retention; the
ultrametric witness split is a triple of pure states with all three
pairwise overlaps equal under 1-RSB but with
\(q_{\alpha\beta}\ne q_{\alpha\gamma}\) under full RSB
\citep{Parisi1983OrderParameterSpinGlasses}.

\item[\emph{Verner's Law refining Grimm's Law (historical
linguistics).}]
Grimm's Law sends PIE \({*}p,{*}t,{*}k,{*}k^w\) uniformly to
Proto-Germanic \({*}f,{*}\theta,{*}h,{*}h^w\); Verner's Law refines
this output by voicing the
medial fricative when the preceding PIE syllable was unaccented.  The
grammatischer-Wechsel alternation within Germanic strong-verb
paradigms is the canonical witness: cognates that PIE accent
distinguished by voicing remain merged under Grimm alone but split
under Grimm + Verner, even though the conditioning accent is itself
lost in attested Germanic
\citep{Verner1877AusnahmeLautverschiebung}.
\end{description}

\subsection{No Unstatused Residual}\label{sec:stability:no-unstatused-residual}

A formed stability package has a native residual registry: the
defects, gaps, budget failures, split pairs, critical pairs, or
endpoint failures that its own claims make relevant.  Each native
residual is not merely a number or a warning flag.  It comes with a
record saying which judgment family is responsible for it and which
status classifier applies.

\begin{theorem}[No Unstatused Residual]\label{thm:F9}
Let \(R\) be a residual type, let \(\mathcal P\) be a formed
stability package over \(R\), and let
\[
  \operatorname{st}_{\mathcal P}:R\longrightarrow\mathcal S
\]
be the declared typed-status classifier supplied by the package.
Assume the package discipline makes \(\operatorname{st}_{\mathcal P}\)
single-valued on native residual records.  Then
\begin{align*}
  &\forall r\in R,\quad
    \operatorname{st}_{\mathcal P}(r)\in
    \mathcal S_{\mathrm{typed}},\\
  &\forall r\in R,\quad
    \exists!\,s\in\mathcal S_{\mathrm{typed}},\
      s=\operatorname{st}_{\mathcal P}(r),\\
  &\forall r\in R,\quad
    \neg\operatorname{UnstatusedResidual}_{\mathcal P}(r).
\end{align*}
\end{theorem}
\footnote{Verified in Lean as
\texttt{Stability.NoUnstatusedResidual.\allowbreak{}no\_unstatused\_residual};
see \Cref{app:formalization}.}
\begin{proof}
Assume the formed stability package \(\mathcal P\), the residual type
\(R\), and the classifier \(\operatorname{st}_{\mathcal P}\).  The
F-II residual-coverage condition for a formed package says that every
native residual has a declared status record:
\[
  \forall r\in R,\ \exists s\in\mathcal S_{\mathrm{typed}},\
  s=\operatorname{st}_{\mathcal P}(r).
  \tag{1}
\]
The package's declared-fiber discipline says that the classifier is
single-valued:
\[
  \operatorname{st}_{\mathcal P}(r)=s
  \ \text{and}\
  \operatorname{st}_{\mathcal P}(r)=s'
  \quad\Longrightarrow\quad
  s=s'.
  \tag{2}
\]
Together, (1) and (2) give the unique typed status for every
\(r\in R\).

If \(\operatorname{UnstatusedResidual}_{\mathcal P}(r)\) held for
some \(r\), then \(r\) would have no typed status in
\(\mathcal S_{\mathrm{typed}}\).  That contradicts (1).  Hence no
native residual of \(\mathcal P\) is unstatused.
\end{proof}

The interaction status is \(\BirdInt\)-compatible because the
residual record is part of the audited interaction data: the package
must say whether the residual is visible, hidden with a record,
outside scope, blocked, priced, or accepted.  The theorem does not
enumerate or redefine the inherited status families; that belongs to
the admissibility apparatus.  It also does not extend to unformed
packages.

\subsection*{Layer instantiations}\label{layer-inst:F9}

\begin{description}
\item[\emph{{GHS} chemical hazard classification (regulatory
registry).}]
The UN Globally Harmonized System of Classification and Labelling of
Chemicals requires that every recognized hazard endpoint of a
substance be classified into exactly one category within its hazard
class (Cat.~1/1A/1B/1C, 2, 3, 4 depending on the class).  Criteria
thresholds are disjoint, the ``most severe category applies'' rule
breaks ties, and EU CLP Article~4(1) and OSHA HazCom transpose the
discipline into enforceable national law
\citep{UN2025GHSRev11}.

\item[\emph{Progress + preservation for a sound type system
(programming-language semantics).}]
The Wright--Felleisen syntactic type-soundness theorem states that
every well-typed reachable configuration of a source term retains
its typing judgment (preservation) and is either a value or admits
a reduction step (progress).  The progress class
\(\{\text{value},\,\text{can-step}\}\) together with the retained
typing judgment is the F9 typed status.  The orphan residual is a
\emph{stuck term}: a closed non-value with no reduction rule and no
type, forbidden by progress together with preservation
\citep{WrightFelleisen1994SyntacticTypeSoundness}.

\item[\emph{Genetic code via ribosomal translation (molecular
biology).}]
The ribosome decodes each mRNA codon into a unique amino acid (the
61 sense codons) or a termination signal (the 3 stop codons),
following the standard genetic code table; codon--anticodon pairing
at the A-site enforces single-valuedness.  Stalled ribosomes at
non-stop codons lacking cognate charged tRNAs constitute the
translation-failure residual; tmRNA / trans-translation in bacteria
and Pelota--Hbs1 ribosome-associated quality control in eukaryotes
rescue the residual by reassigning a non-stop status and tagging the
peptide for degradation
\citep{Crick1968OriginGeneticCode}.

\item[\emph{Wick--Wightman--Wigner superselection (foundational
quantum physics).}]
The WWW superselection rules require that every physically
admissible pure state lie in a unique joint eigenspace of the
superselection observables --- electric charge \(Q\), baryon number
\(B\), lepton number \(L\), and univalence \((-1)^{2J}\).
Coherent superpositions across distinct sectors are excluded from
the physical state space because the relative phase is unobservable
(the superselection observables commute with every observable), so
the would-be superposition reduces to a sector-mixture rather than a
pure state
\citep{WickWightmanWigner1952IntrinsicParity}.

\item[\emph{Gale--Shapley deferred-acceptance matching
(mechanism design).}]
In a two-sided matching market with complete strict preferences and
equal market sides, the Gale--Shapley deferred-acceptance algorithm
terminates with a stable matching \(\mu\) that is a function: every
agent is assigned exactly one partner, and truth-telling is dominant
on the proposing side.  An orphan residual --- an agent with no
\(\mu(x)\) output --- cannot occur, because deferred-acceptance
termination is the F9 coverage discipline and the function
structure of \(\mu\) is the F9 single-value discipline
\citep{GaleShapley1962CollegeAdmissions}.
\end{description}

\subsection*{Stability at a glance}\label{sec:stability:summary}
\begin{table}[t]
  \centering
  \footnotesize
  \begin{tabularx}{\textwidth}{@{}
      >{\raggedright\arraybackslash}p{0.24\textwidth}
      >{\raggedright\arraybackslash}X
      >{\raggedright\arraybackslash}p{0.18\textwidth}
      >{\raggedright\arraybackslash}p{0.12\textwidth}@{}}
    \toprule
    Paper-prose name & One-line claim & Status & Substrate cite \\
    \midrule
    No-Needles Theorem &
    Formed stability packages cannot admit unpriced
    layer-dissolving needles. &
    Anchor &
    \PaperNeedleKiller, \PaperNS \\
    Sufficiency Closure Theorem &
    \(\currQ\) is sufficient for declared futures iff each future
    probe descends through \(\currQ\). &
    Theorem &
    \textemdash \\
    Strict Extension Lemma &
    An extension is strict iff it is a refinement with a split old
    fiber. &
    Theorem &
    \PaperAR \\
    No-Unstatused-Residual Theorem &
    Every native residual of a formed package carries a typed status. &
    Theorem &
    \textemdash \\
    \bottomrule
  \end{tabularx}
  \caption{Stability axis at a glance: the four laws of \Cref{sec:stability}.}
  \label{tab:axis-stability-summary}
\end{table}
  \section{Axis III — Transport}\label{sec:transport}

The Transport axis is about moving structure across quotients.  Its
basic question is the familiar one: when a map, comparison, or family
of local data is defined before quotienting, does it still make sense
after the quotient has forgotten some distinctions?  This is the
easiest entry point into the catalog because the laws are closest to
standard descent, gluing, and quotient comparison.  The closure
apparatus supplies the bookkeeping, but the mathematical shape is the
ordinary universal-property shape: a construction is legitimate only
when it respects the equivalence relation through which it is being
read.

The chapter has four direct laws.  Descent--Repair Normal Form is
the lead theorem: it makes formal the scenario from the
introduction, where a
map fails to descend through a quotient and the obstruction is exactly
the split-pair set \(\splitObs\).  Holonomy-Memory Repair Normal Form
then treats the case where two routes look identical at the current
quotient \(\currQ\) but leave different residue for a predictive
quotient \(\predM\).  Local--Global Obstruction Normal Form records
the gluing situation: locally compatible data globalizes exactly when
it lies in the image of the declared restriction map.  Recombination
Comparison treats branch-merge transport, where branchwise equality
need not determine merged behavior unless the recombination quotient
factors through the branchwise quotient.

All four laws in this chapter are direct quotient-theoretic laws:
there are no substrate anchors, partial wrappers, or obligation
wrappers here.  This makes the Transport axis self-contained after
\Cref{sec:apparatus}.  It also makes the axis a toolbox for the rest
of the paper.  The split-pair discipline introduced in the
Descent--Repair Normal Form reappears in the local--global and
recombination laws, and later in the selection, horizon, and
self-reference limits axes.

Readers coming cold to the framework should read
\Cref{sec:apparatus} and then this chapter before choosing another
axis.  The examples here use only quotients, factorization, local
compatibility, and comparison maps; later chapters reuse the same
patterns with additional status, access, symmetry, or self-reference
apparatus.

\subsection{Descent--Repair Normal Form}\label{sec:transport:descent-repair}

The first transport law is the quotient descent criterion already
introduced in \Cref{sec:intro}.  We use the split-pair obstruction of
\Cref{def:split-pair}.

\begin{definition}[Canonical target-repair quotient]\label{def:target-repair}
Let \(q:X\to Q\), \(F:X\to Y\), and \(r:Y\to R\).  The canonical
target-repair relation \(\equiv_F\) on \(R\) is the smallest
equivalence relation generated by the identifications
\[
  r(Fx)\equiv_F r(Fx')
  \qquad\text{whenever}\qquad
  q(x)=q(x').
\]
The \emph{canonical target-repair quotient} is
\[
  R_F:=R/{\equiv_F},
\]
with quotient map \(c_F:R\to R_F\).
\end{definition}

\begin{theorem}[Descent-Repair Normal Form]\label{thm:F2}
Let \(q:X\to Q\) be a quotient map, let \(F:X\to Y\) be a map, and
let \(r:Y\to R\) be a target readout.  Then \(r\circ F\) descends
through \(q\) if and only if
\[
  \splitObs(q,F,r)=\emptyset .
\]
If the obstruction is nonempty, the canonical repairs are: source
refinement, which separates the offending source fibers, and target
coarsening by the quotient \(c_F:R\to R_F\) of
\Cref{def:target-repair}.  The target repair is minimal among target
coarsenings that make \(r\circ F\) descend.
\end{theorem}
\footnote{Verified in Lean as
\texttt{Transport.DescentRepair.\allowbreak{}descent\_repair\_normal\_form};
see \Cref{app:formalization}.}

\begin{proof}
Suppose first that \(r\circ F\) descends through \(q\).  Then there is
a map \(\overline F:Q\to R\) such that
\(\overline F\circ q=r\circ F\).  If \(q(x)=q(x')\), then applying
\(\overline F\) gives
\[
  r(Fx)=\overline F(qx)=\overline F(qx')=r(Fx').
\]
Thus no pair \((x,x')\) with \(q(x)=q(x')\) can be separated by
\(rF\), so \(\splitObs(q,F,r)=\emptyset\).

Conversely, assume \(\splitObs(q,F,r)=\emptyset\).  Then whenever
\(q(x)=q(x')\), the pair \((x,x')\) is not a split pair, so
\(r(Fx)=r(Fx')\).  Therefore \(r\circ F\) is constant on the fibers
of \(q\).  By the quotient universal property, there is a unique map
\(\overline F:Q\to R\) satisfying
\(\overline F\circ q=r\circ F\).

The source repair is obtained by replacing \(q\) with any refinement
that no longer identifies the split pairs; the least such repair is
the quotient whose fibers refine \(q\)-fibers by the target values
\(r(Fx)\).  For the target repair, quotient \(R\) by the generated
relation \(\equiv_F\).  If \(q(x)=q(x')\), then
\(r(Fx)\equiv_F r(Fx')\), hence
\[
  c_F(r(Fx))=c_F(r(Fx')).
\]
Thus \(c_F\circ r\circ F\) is constant on \(q\)-fibers and descends
through \(q\).

It remains to check minimality.  Let \(c:R\to C\) be any target
coarsening such that \(c\circ r\circ F\) descends through \(q\).  Then
for every generator \(r(Fx)\equiv_F r(Fx')\), the equality
\(q(x)=q(x')\) implies
\[
  c(r(Fx))=c(r(Fx')).
\]
Since equality under \(c\) is an equivalence relation on \(R\), it
contains the equivalence relation generated by those identifications.
Therefore \(c\) is constant on \(\equiv_F\)-classes and factors
uniquely through \(c_F:R\to R_F\).  Hence \(R_F\) is the finest
coarsening that repairs descent.
\end{proof}

\subsection*{Layer instantiations}\label{layer-inst:F2}

\begin{description}
\item[\emph{Coequalizer in a category (mathematics; target
coarsening).}]
The coequalizer of a parallel pair \(u, v : A \rightrightarrows B\)
is the universal arrow \(c : B \to C\) with \(c \circ u = c \circ
v\). Given a candidate target readout \(r : B \to R\), the
canonical minimal target coarsening that restores descent through
the \(\langle u, v\rangle\)-generated equivalence on \(B\) is
\(R \to \operatorname{coeq}(r \circ u, r \circ v)\): the
coequalizer of the parallel pair pushed forward into \(R\).
\citep{MacLane1998Categories}.

\item[\emph{Gauge-invariant observables (physics; source
refinement).}]
In a gauge theory the gauge group \(G\) is given a priori by the
theory's symmetry structure, and the orbit quotient
\(q : X \to X/G\) on field configurations is the discipline-fixed
source quotient. A candidate observable \(r : X \to \mathbb{R}\)
descends through \(q\) iff it is gauge-invariant. F2's canonical
source refinement is the joint quotient
\(\langle q, r\rangle : X \to (X/G) \times \mathbb{R}\); declared
gauge-fixing data is one concrete implementation. F2's target
coarsening yields the gauge-invariant observable algebra
\citep{PeskinSchroeder1995QFT}.

\item[\emph{Hash collisions in dictionary lookup (computer
science; source refinement).}]
A hash function \(q : K \to B\) is a discipline-given source
quotient on a key universe. A downstream lookup operation
\(r : K \to R\) descends through \(q\) iff no two collision-merged
keys carry distinct intended values. When they do, F2's canonical
source refinement is the joint quotient
\(\langle q, r\rangle : K \to B \times R\); FKS two-level perfect
hashing implements this joint refinement by introducing a
secondary hash that separates the offending colliders
\citep{CormenLeisersonRivestStein2009CLRS}.

\item[\emph{Lehmann--Scheff\'e minimal sufficient statistic
(statistics; source refinement).}]
A researcher picks a candidate statistic \(T : X \to \mathcal{T}\)
and asks whether the likelihood ratio descends through \(T\). When
\(T\) is not sufficient, the F2 split pairs are sample pairs with
the same \(T\) but distinct likelihood ratios. F2's canonical
source refinement is the joint quotient
\(\langle T, L(\cdot;-)/L(\cdot;x_0)\rangle\); when \(T\) is coarser
than the minimal sufficient statistic \(T^{*}\), this joint
quotient equals \(T^{*}\) (Lehmann--Scheff\'e / Bahadur uniqueness)
\citep{LehmannScheffe1950Completeness}.

\item[\emph{Genetic code, codon-to-amino-acid quotient (biology;
source refinement).}]
The standard genetic code is a discipline-given quotient
\(q : \mathrm{Codons} \to \mathrm{AminoAcids}\) (61 sense codons
to 20 amino acids). A codon-level cellular readout such as
translation efficiency \(\tau : \mathrm{Codons} \to \mathbb{R}\)
does not in general descend through \(q\): synonymous codons can
differ in translation rate due to tRNA-pool composition (the
\emph{codon usage bias} literature). F2's canonical source
refinement is the joint quotient
\(\langle q, \tau\rangle : \mathrm{Codons} \to \mathrm{AminoAcids}
\times \mathbb{R}\); codon usage tables and codon-adaptation
indices are operational approximations \citep{Crick1968GeneticCode}.
\end{description}

\subsection{Holonomy-Memory Repair Normal Form}\label{sec:transport:holonomy-memory}

The second transport law concerns route dependence: two routes can be
indistinguishable at the current quotient while still leaving a
predictive difference.

\begin{definition}[Route residue]\label{def:route-residue}
Let \(H\) be a carrier, let \(\gamma,\eta:H\to Z\) be two declared
routes, let \(\currQ:Z\to Q\) be the current quotient, and let
\(\predM:Z\to M\) be a predictive quotient refining the current
readout.  The \emph{route residue} of \(\gamma\) and \(\eta\) at
\(h\in H\) is present when
\[
  \currQ(\gamma h)=\currQ(\eta h)
  \quad\text{but}\quad
  \predM(\gamma h)\ne\predM(\eta h).
\]
\end{definition}

\begin{theorem}[Holonomy-Memory Repair Normal Form]\label{thm:F3}
For two declared routes \(\gamma,\eta:H\to Z\), current quotient
\(\currQ\), and predictive quotient \(\predM\), memory is forced
exactly when there is route residue.  Equivalently, if two route
outputs agree at \(\currQ\) but differ at \(\predM\), then no
representation through the current quotient alone is sufficient for
the declared future probes.  Recording the route residue gives the
canonical memory repair.
\end{theorem}
\footnote{Verified in Lean as
\texttt{Transport.HolonomyMemoryRepair.\allowbreak{}holonomy\_memory\_repair\_normal\_form};
see \Cref{app:formalization}.}

\begin{proof}
There are three cases for a point \(h\in H\).  If
\(\currQ(\gamma h)\ne\currQ(\eta h)\), the difference is already
visible at the current quotient and no hidden memory claim is needed.
If \(\predM(\gamma h)=\predM(\eta h)\), then the routes remain
identified even after passing to the predictive quotient, so the
residue is predictively flat.

The remaining case is exactly the route-residue case of
\Cref{def:route-residue}: the two route outputs agree at \(\currQ\)
but differ at \(\predM\).  Any construction that factors only through
\(\currQ\) assigns the same value to \(\gamma h\) and \(\eta h\), so it
cannot recover the declared future distinction seen by \(\predM\).
Thus the current quotient is not sufficient for the route comparison.

The repair is to enlarge the recorded state by the residue witness,
or equivalently to promote the route comparison to the predictive
quotient on the relevant fiber.  Once this memory witness is recorded,
the two cases that were equal at \(\currQ\) and separated at \(\predM\)
are no longer represented by the same memory state.  The repair is
canonical relative to the declared future probes because any quotient
that makes those probes sufficient must separate every such
\(\predM\)-split inside a current fiber, hence must factor through the
quotient obtained by recording precisely these residues.
\end{proof}

\subsection*{Layer instantiations}\label{layer-inst:F3}

\begin{description}
\item[\emph{Holonomy of a connection on a fiber bundle
(mathematics).}]
Given a connection \(\nabla\) on a fiber bundle \(E \to M\) and two
piecewise-smooth paths \(\gamma, \eta\) in \(M\) with common
endpoints, parallel-transport produces lifts in the total space
\(Z = E\) that project to the same base point under \(\pi : E \to
M\) (the current quotient) but generally end at distinct fiber
elements (the predictive readout). The canonical memory witness
is the holonomy
\(P_\eta^{-1} \circ P_\gamma \in \operatorname{Hol}(\nabla, x)\),
or globally the holonomy representation
\(\operatorname{Hol}(\nabla, \cdot) : \pi_1(M, x) \to G\)
\citep{KobayashiNomizu1963Foundations}.

\item[\emph{Berry phase in adiabatic quantum mechanics (physics).}]
A quantum system carried adiabatically around a closed loop in
Hamiltonian parameter space returns to the same instantaneous
eigenstate but acquires an additional geometric phase.  Two such
loops sharing endpoints agree at the current quotient (eigenstate
band index) but differ in the enriched
\((\text{band}, \text{phase})\) readout whenever they enclose
distinct flux of the Berry curvature.  The Berry connection
\(\mathbf A = i\langle \psi \mid \nabla \psi\rangle\) is the
canonical memory witness; its holonomy is the Berry phase
\citep{Berry1984QuantalPhase}.

\item[\emph{Event sourcing in distributed systems (computer
science).}]
In an event-sourced system the authoritative state is an
append-only log of domain events; the current materialized snapshot
is a fold over that log.  Two distinct event histories can produce
identical snapshots while remaining distinguishable under
retroactive policy replay, audit, or future business rules — the
enriched \((\text{snapshot}, \text{retro\,results})\) readout
differs even when the snapshot agrees.  The event log itself is the
canonical memory witness retained precisely because the snapshot
is lossy \citep{Kleppmann2017DDIA}.

\item[\emph{Epigenetic memory of cell fate (biology).}]
Two somatic cells derived from a common zygote share an essentially
identical genome (the current quotient) but exhibit divergent
transcriptional responses to the same regulatory cues.  The
enriched
\((\text{genome}, \text{chromatin}, \text{predicted response})\)
readout differs because chromatin marks (DNA methylation, histone
modifications, bivalent domains) encode the developmental lineage.
The chromatin state, propagated through mitosis, is the canonical
memory witness \citep{Reik2007Epigenetic}.

\item[\emph{Preisach hysteresis and return-point memory
(engineering).}]
The current magnetization \(M(t)\) of a hysteretic material is
reachable by many distinct input-field histories \(H(t)\); two
histories arriving at the same \((H, M)\) state predict divergent
responses to the next field reversal.  The Preisach model's
canonical memory witness is the \emph{return-point staircase}: the
stack of dominant past extrema of the input history that the
Preisach plane records.  Mayergoyz's wiping-out theorem
establishes this staircase as the minimal residue distinguishing
otherwise-equivalent presents \citep{Mayergoyz2003Hysteresis}.
\end{description}

\subsection{Local--Global Obstruction Normal Form}\label{sec:transport:local-global}

The third transport law is the local-to-global image criterion.  It
separates overlap tests from genuine global existence.

\begin{definition}[Local-global compatibility]\label{def:local-global-compat}
Let \(I\) index a declared finite patch system, with local object sets
\((Q_i)_{i\in I}\).  Let \(\mathcal C\subseteq \prod_i Q_i\) be the
subset of local families satisfying all declared overlap comparisons.
Let \(Q\) be the global object set and
\[
  R:Q\to \prod_i Q_i
\]
the declared restriction map.  A local family \(a\in\prod_i Q_i\) is
\emph{compatible} if \(a\in\mathcal C\), and it \emph{globalizes} if
\(a\in\operatorname{im}(R)\).
\end{definition}

\begin{theorem}[Local-Global Obstruction Normal Form]\label{thm:F4}
Let \(\mathcal G\) be a set of global states, let \(\mathcal F\) be a
set of local-family data, let
\(\operatorname{Comp}:\mathcal F\to\mathrm{Prop}\) be the
overlap-compatibility predicate, and let
\[
  R:\mathcal G\longrightarrow\mathcal F
\]
be a declared restriction map such that
\(\operatorname{Comp}(R(g))\) for every \(g\in\mathcal G\).  For
\(a\in\mathcal F\), write
\[
  \operatorname{Globalizes}_R(a)
  \quad:\Longleftrightarrow\quad
  \exists g\in\mathcal G,\ R(g)=a.
\]
Then
\begin{align*}
  \operatorname{Comp}(a)
  &\Longrightarrow
  \bigl(\operatorname{Globalizes}_R(a)
  \iff a\in R(\mathcal G)\bigr),\\
  \operatorname{ExactlyOneStatus}(a)
  &\iff
  \operatorname{OverlapInconsistent}(a)
  \mathbin{\dot\vee}
  \operatorname{CompatibleNonGlobal}(a)\\
  &\hspace{3.5em}\mathbin{\dot\vee}
  \operatorname{GloballyAmbiguous}(a)
  \mathbin{\dot\vee}
  \operatorname{GloballyExact}(a),
\end{align*}
where the four status predicates are determined by
\(\operatorname{Comp}(a)\), membership in \(R(\mathcal G)\), and the
cardinality of the fiber \(R^{-1}(a)\).
\end{theorem}
\footnote{Verified in Lean as
\texttt{Transport.LocalGlobalObstruction.\allowbreak{}local\_global\_obstruction\_normal\_form};
see \Cref{app:formalization}.}

\begin{proof}
Assume the compatibility predicate \(\operatorname{Comp}\), the
restriction map \(R\), and the hypothesis
\(\operatorname{Comp}(R(g))\) for every \(g\in\mathcal G\).  Fix
\(a\in\mathcal F\) and suppose \(\operatorname{Comp}(a)\).  By
definition, \(a\) globalizes exactly when there exists a global state
restricting to it:
\[
  \operatorname{Globalizes}_R(a)
  \iff
  \exists g\in\mathcal G,\ R(g)=a
  \iff
  a\in R(\mathcal G).
\]
This proves the local-to-global image criterion.

The status enumeration is a case split on compatibility, image
membership, and the fiber \(R^{-1}(a)\).  If
\(\neg\operatorname{Comp}(a)\), then \(a\) is overlap-inconsistent.
If \(\operatorname{Comp}(a)\) and \(a\notin R(\mathcal G)\), then
\[
  R^{-1}(a)=\emptyset,
\]
so \(a\) is compatible non-global.  If
\(\operatorname{Comp}(a)\), \(a\in R(\mathcal G)\), and
\(|R^{-1}(a)|>1\), then \(a\) is globally ambiguous.  If
\(\operatorname{Comp}(a)\), \(a\in R(\mathcal G)\), and
\(|R^{-1}(a)|=1\), then \(a\) is globally exact.  These four cases
are mutually exclusive and exhaustive.
\end{proof}

\subsection*{Layer instantiations}\label{layer-inst:F4}

\begin{description}
\item[\emph{\v{C}ech cohomology and sheaf gluing (mathematics).}]
For an abelian sheaf \(\mathcal F\) on a space \(X\) with finite
open cover \(\{U_i\}\), local sections agree pairwise on overlaps
iff the 0-cocycle condition holds; the obstruction to lifting a
compatible local family to a global section is the \v{C}ech class
in \(\check H^1\).  The four F4 statuses correspond to non-cocycle
data, nontrivial \(\check H^1\) class (canonical example:
\(\mathcal O(1)\) on \(\mathbb P^1\)), nontrivial kernel of the
restriction, and acyclic Leray covers
\citep{Hartshorne1977AlgebraicGeometry}.

\item[\emph{Topological defects in ordered media (physics).}]
In a system with order-parameter space \(M = G/H\), patchwise
order-parameter sections must respect the group law of \(\pi_n(M)\)
on defect linking charges across patch overlaps.  Forbidden total
charges (compatible local but no global texture), non-abelian
disclination entanglements (multiple distinct globals with
identical local tabulation), and defect-free uniform textures
illustrate the four-status partition in condensed-matter physics
\citep{Mermin1979TopologicalDefects}.

\item[\emph{Schema mapping in data integration (computer science).}]
In data integration, source instances over local schemas are
glued to a mediated global schema through declared
tuple-generating and equality-generating dependencies (tgds, egds).
Local-overlap compatibility is the egd-and-tgd satisfaction check
at the chase; the four statuses are egd violation,
non-terminating chase, multiple universal solutions resolved by
certain-answer semantics, and unique chase-output solution under
weakly acyclic dependencies \citep{Fagin2005DataExchange}.

\item[\emph{Genome assembly from overlapping reads (biology).}]
Sequencing reads are local sequence data on intervals of an
unknown genome; overlap-layout-consensus and de Bruijn graph
assemblers test pairwise overlaps for sequence agreement and
attempt to stitch compatible reads into a single contig.  The
four statuses appear as overlap conflicts, fragmented contigs
without a single traversal, repeat-induced chimeric-vs-unique
ambiguity, and unique reference assembly
\citep{PevznerTangWaterman2001Eulerian}.

\item[\emph{Pose-graph SLAM with loop closures (engineering).}]
A robot trajectory and map are estimated from local sensor
observations whose loop closures impose pairwise pose constraints.
The four statuses appear as chi-squared loop-closure rejection
(perceptual aliasing), nonzero cocycle drift around a cycle (no
global \(SE(3)\) embedding), gauge ambiguity (global rigid
transform, monocular scale), and converged g2o\,/\,GTSAM solution
with positive-definite information matrix
\citep{Cadena2016SLAMSurvey}.
\end{description}

\subsection{Recombination Comparison}\label{sec:transport:recombination}

The final transport law compares branchwise equality with equality
after a declared branch merge.

\begin{definition}[Bridge composition]\label{def:bridge-composition}
Let \(B\) be the carrier of branch data, let \(K:B\to Q_K\) be the
branchwise quotient, and let \(R:B\to Q_R\) be the recombination
quotient.  The \emph{bridge composition} is the use of \(K\) as the
source quotient and \(R\) as the target readout in the split-pair
criterion of \Cref{def:split-pair}.  Its obstruction is
\[
  \{(b,b')\in B\times B:
      K(b)=K(b')\ \text{and}\ R(b)\ne R(b')\}.
\]
\end{definition}

\begin{theorem}[Recombination Comparison]\label{thm:F5}
Let \(K:B\to Q_K\) be a branchwise quotient and
\(R:B\to Q_R\) a recombination quotient on the same branch carrier.
Then branchwise equality determines recombination equality if and only
if \(R\) factors through \(K\), that is, if and only if there exists
\(\rho:Q_K\to Q_R\) such that
\[
  R=\rho\circ K .
\]
Equivalently, the branch-merge split-pair obstruction of
\Cref{def:bridge-composition} is empty.
\end{theorem}
\footnote{Verified in Lean as
\texttt{Transport.RecombinationComparison.\allowbreak{}recombination\_comparison};
see \Cref{app:formalization}.}

\begin{proof}
If \(R=\rho\circ K\), then any two branch records with
\(K(b)=K(b')\) satisfy
\[
  R(b)=\rho(K(b))=\rho(K(b'))=R(b').
\]
Thus the bridge-composition obstruction is empty: branchwise equality
determines equality after recombination.

Conversely, assume the obstruction is empty.  Then \(R\) is constant
on the fibers of \(K\): whenever \(K(b)=K(b')\), emptiness of the
obstruction gives \(R(b)=R(b')\).  By the quotient universal property,
there is a unique map \(\rho:Q_K\to Q_R\) with \(R=\rho\circ K\).
This proves the factorization criterion.

This is exactly \Cref{thm:F2} applied with \(K\) as source quotient,
the identity branch transport as the pre-quotient map, and \(R\) as
target readout.  A recombination failure is therefore not a one-sided
strict-refinement slogan; it is a two-sided quotient comparison
witnessed by a split pair at the merged level.
\end{proof}

\subsection*{Layer instantiations}\label{layer-inst:F5}

\begin{description}
\item[\emph{Meiotic recombination and genetic linkage (biology).}]
In meiosis the chromatid-of-origin assignment per allele
(branchwise quotient) and the gamete genotype after recombination
(target readout) coincide iff there is no crossing-over at the
relevant loci.  The bridge-composition obstruction is the
\emph{recombinant gamete}: two meiotic configurations with the
same parental-locus assignment but distinct gamete genotypes
\citep{Morgan1911Coupling}.

\item[\emph{Church--Rosser confluence (computer science / logic).}]
For a term-rewriting system, the source term being reduced
(branchwise quotient \(K\)) determines the normal form
(recombination target \(R\)) independently of the reduction
strategy iff the system is confluent.  Bridge obstructions are
\emph{non-confluence witnesses}: two reductions from the same
source term producing distinct normal forms
\citep{ChurchRosser1936Conversion}.

\item[\emph{Path integral with gauge equivalence (physics).}]
For a quantum-mechanical path integral with gauge-orbit
quotient \(K\) on field-configuration space, the exponentiated
action \(R = \exp(iS[\text{path}])\) descends through \(K\) iff
the classical action is gauge-invariant.  Bridge obstructions are
gauge-related paths with different action values, signalling a
\emph{gauge-invariance failure} of the action
\citep{FeynmanHibbs1965PathIntegrals}.

\item[\emph{Tversky--Kahneman framing effect (cognitive psychology
/ behavioural economics).}]
For a participant presented with a choice problem under various
framings, the objective lottery underlying the problem (branchwise
quotient \(K\), with frame forgotten) and the participant's modal
choice (recombination target \(R\)) coincide --- i.e., \(R\)
descends through \(K\) --- iff the participant exhibits
frame-invariance.  The Tversky--Kahneman Asian disease problem
demonstrates that participants systematically violate
frame-invariance: presenting the same objective lottery under
gain-frame vs.\ loss-frame phrasings yields different modal
choices, producing the bridge-composition obstruction --- a
presentation pair sharing the underlying lottery but differing in
choice
\citep{TverskyKahneman1981Framing}.
\end{description}

\subsection*{Transport at a glance}\label{sec:transport:summary}
\begin{table}[t]
  \centering
  \footnotesize
  \begin{tabularx}{\textwidth}{@{}
      >{\raggedright\arraybackslash}p{0.24\textwidth}
      >{\raggedright\arraybackslash}X
      >{\raggedright\arraybackslash}p{0.18\textwidth}
      >{\raggedright\arraybackslash}p{0.12\textwidth}@{}}
    \toprule
    Paper-prose name & One-line claim & Status & Substrate cite \\
    \midrule
    Descent--Repair Theorem &
    \(\splitObs\) is empty iff descent holds; canonical target
    repair coarsens \(R\) by the generated equivalence. &
    Theorem &
    \textemdash \\
    Holonomy-Memory Theorem &
    Routes equal at \(\currQ\) may differ at \(\predM\); route
    residue is the recorded memory witness. &
    Theorem &
    \textemdash \\
    Local--Global Obstruction Theorem &
    Locally compatible data globalizes iff it lies in the image of
    the declared restriction map. &
    Theorem &
    \textemdash \\
    Recombination Comparison &
    Branchwise equality determines merged equality iff the
    recombination quotient factors through the branchwise quotient. &
    Theorem &
    \textemdash \\
    \bottomrule
  \end{tabularx}
  \caption{Transport axis at a glance: the four laws of \Cref{sec:transport}.}
  \label{tab:axis-transport-summary}
\end{table}
  \section{Axis IV — Symmetry}\label{sec:symmetry}

The Symmetry axis collects obstructions caused by involutions,
branch selection, and anomaly-style failures.  In this chapter a
symmetry is not merely a convenient change of notation.  It is part
of the apparatus that controls which branches, traces, and quotient
classes may be distinguished without adding new breaking data.  The
basic question is therefore familiar: when a construction is
invariant under the declared symmetry, what can it select, transport,
or record?

The five laws proceed from fixed data to failure modes and then to
equivalence.  Duality Fixity, or Self-Dual Trace Confinement, states
the layer-agnostic self-dual trace confinement law, with substrate
provenance in \PaperSDTC\ \citep{Tsiokos2026SDTC} and
\PaperRH\ \citep{Tsiokos2026RH}.  Selection Obstruction / Repair
Normal Form is a direct theorem: exact symmetry cannot select a
non-fixed branch without declared breaking data.  Spontaneous
Symmetry Breaking / Branch-Selection Formation is an axiom-applied
future-source wrapper.  Anomaly / Symmetry Obstruction and Duality
Equivalence are direct theorems about failed symmetry descent and
role-preserving equivalence of formed packages.

The Duality Fixity source condition has a small shape difference from
the other guarded rows: either the SDTC apparatus source or the RH
trace-specialization source is enough to supply the trace-confinement
bridge.  The direct selection-obstruction theorem reuses the
Descent--Repair split-pair discipline: a forbidden selection is a
split-pair obstruction at the level of symmetry orbits.  The later
anomaly and duality-equivalence subsections keep the same posture:
the apparatus must say when symmetry descends, when it fails, and
when two symmetry presentations preserve the same formed roles.

\subsection{Duality Fixity (Self-Dual Trace Confinement)}\label{sec:symmetry:duality-fixity}

\begin{definition}[Involutive ledger; anti-invariant readout]\label{def:involutive-ledger}
An involutive ledger is a carrier \(X\) equipped with an involution
\(\Jinv:X\to X\), so \(\Jinv^2=\mathrm{id}\).  Its fixed locus is
\[
  \Fixinv=\{x\in X:\Jinv(x)=x\}.
\]
An anti-invariant readout is a quantity whose sign changes under
\(\Jinv\).  It is separating when zero trace against the
anti-invariant readout forces the corresponding visible
anti-invariant mass to vanish.
\end{definition}

\begin{theorem}[Duality Fixity / Self-Dual Trace Confinement]\label{thm:F14}
Let a formed closure carry an involutive ledger with separating
anti-invariant readout and trace functional \(\trAX\).  Assume a
source instance supplying the self-dual trace-confinement bridge.
If a completed domination ladder satisfies
\[
  A_X \preceq B_n
  \qquad\text{and}\qquad
  \trAX(B_n)\to 0,
\]
then \(A_X=0\).  Equivalently, visible mass native to the
self-dual trace layer is confined to \(\Fixinv\).
\end{theorem}
\footnote{The Lean theorem
\texttt{Symmetry.DualityFixity.duality\_fixity} proves the
layer-agnostic statement under a substrate-pending axiom
\texttt{duality\_fixity\_pending} guarded by a definitional
disjunction of two opaque source predicates (one for the SDTC
apparatus, one for the RH trace specialization); the substrate
proofs appear in \citep{Tsiokos2026SDTC, Tsiokos2026RH} and either
suffices to discharge the axiom.  See \Cref{app:formalization} for
the narrowing disclosure.}
\begin{proof}
Let \(I\) be the duality-fixity interface determined by the formed
closure, involutive ledger, trace functional, domination ladder, and
fixed-locus readout in the statement.  Apply the same-paper
substrate-pending axiom \(\mathrm{duality\_fixity\_pending}\),
specialized to \(I\) under the disjunctive source hypothesis
\begin{align*}
  \mathsf{src}:
  \operatorname{DualityFixitySubstrateAnchor}(I)
  &=
  \operatorname{Paper2SDTCApparatus}(I)\\
  &\quad{}\lor
  \operatorname{Paper6RHTraceSpecialization}(I).
\end{align*}
Either disjunct of \(\mathsf{src}\) suffices: the axiom is specified
to accept the disjunction itself, not a preferred source branch, and
it treats the SDTC-apparatus and RH-trace-specialization paths
symmetrically.  Given \(\mathsf{src}\), the axiom delivers
\begin{align*}
  \operatorname{DualityFixityNormalForm}(I)
  \equiv{}&
  \operatorname{SDTCDominationSqueezeHypotheses}(I)\\
  &{}\Longrightarrow
  \operatorname{FixedLocusConfinement}(I).
\end{align*}
Unpacked for the displayed domination ladder, this says that under
SDTC-style domination-squeeze hypotheses, visible mass confines to
\(\Fixinv\) and the anti-invariant readout vanishes.

The substrate-side derivation of the axiom lives in the SDTC
apparatus theorem \citep{Tsiokos2026SDTC} and the Riemann Hypothesis
trace specialization \citep{Tsiokos2026RH}; either substrate Lean
discharges the guard.  The opacity-guard pattern of
\Cref{sec:apparatus} ensures that no instantiation of
\(\mathrm{duality\_fixity\_pending}\) can derive \(\mathrm{False}\)
from the axiom alone: both disjuncts of the guard are opaque
predicates with no public constructor.
\end{proof}

This subsection does not re-prove the SDTC apparatus, does not prove
the RH specialization, and does not instantiate the law to a specific
arithmetic carrier.  It records the layer-agnostic trace-confinement
normal form and the source condition under which the wrapper becomes
a discharged theorem.

\subsection*{Layer instantiations}\label{layer-inst:F14}

\begin{description}
\item[\emph{Symmetric/skew decomposition under transpose involution
(linear algebra).}]
For any \(A\in\mathrm{Mat}_n(\mathbb R)\), the transpose
\(J(A)=A^{\!\top}\) is an involution; the unique decomposition
\(A=\tfrac12(A+A^{\!\top})+\tfrac12(A-A^{\!\top})\) separates
\(J\)-symmetric and \(J\)-skew parts.  The Frobenius norm
\(\lVert A-A^{\!\top}\rVert_F\) is the F14 anti-invariant readout,
and vanishing of this norm along a Loewner-dominating sequence
controlled by a self-adjoint trace budget forces \(A\) into the
fixed locus \(\mathrm{Sym}_n(\mathbb R)\)
\citep{HornJohnson2013MatrixAnalysis}.

\item[\emph{Bender--Boettcher {PT}-symmetric Hamiltonians
(foundational physics).}]
For a spin-\(0\) quantum system, the composite \(\mathcal{PT}\)
operator is antilinear with \((\mathcal{PT})^2=\mathrm{id}\), an
involution on the Hilbert space.  Bender and Boettcher showed that
\(\mathcal{PT}\)-symmetric non-Hermitian Hamiltonians (such as
\(H=p^2+x^2(ix)^{\varepsilon}\), including \(H=p^2+ix^3\) at
\(\varepsilon=1\)) have entirely real spectra precisely when the
\(\mathcal{PT}\) symmetry is unbroken: each physical eigenstate is
a \(\mathcal{PT}\)-fixed eigenstate, and vanishing of the
\(\mathcal{PT}\)-skew spectral component (the imaginary part of an
eigenvalue) along a domination ladder of trial states confines the
spectrum to the unbroken fixed locus
\citep{BenderBoettcher1998PTSymmetric}.

\item[\emph{Hermitian-symmetric spectrum of a real signal (signal
processing).}]
For a discrete-time signal \(x[n]\in\mathbb R\) over \(\mathbb Z\),
the discrete-time Fourier transform satisfies
\(X(e^{-i\omega})=\overline{X(e^{i\omega})}\) under the
complex-conjugation involution
\(J(X)(\omega)=\overline{X(-\omega)}\).  The
\(J\)-anti-invariant spectral component is the F14 anti-invariant
readout, and vanishing of its integrated power forces the spectrum
into the Hermitian-symmetric fixed locus
\citep{OppenheimSchafer2010DiscreteTimeSP}.

\item[\emph{Galois conjugation and the fixed-field theorem
(algebraic number theory).}]
For a Galois extension \(L/K\) with cyclic order-2 Galois group
\(\langle\sigma\rangle\), the fundamental theorem of Galois theory
identifies the base field \(K\) with the fixed locus \(L^{\sigma}\)
of the involution \(\sigma\).  The anti-invariant component
\(z-\sigma(z)\) is the F14 anti-invariant readout, and vanishing
of any nontrivial norm respecting the Galois structure forces
\(z\in K\) --- confinement to the fixed field
\citep{Galois1846Memoire}.

\item[\emph{Smith theory: \(\mathbb Z/2\)-cohomology and
fixed-point sets (algebraic topology).}]
P.~A. Smith's theorem relates the
\(\mathbb Z/2\)-cohomology of a topological space \(X\) under a
continuous involution \(J:X\to X\) to its fixed-point set
\(X^J\).  Under finite-type hypotheses,
\(\dim H^*(X^J;\mathbb Z/2)\le\dim H^*(X;\mathbb Z/2)\), with
equality when \(X\) is totally non-homologous to zero.  The
\(\mathbb Z/2\)-anti-invariant cohomology classes are the F14
anti-invariant readout, and vanishing of the
Smith--Floyd spectral-sequence differentials confines all
equivariant cohomology to the fixed-point set
\citep{Smith1938TransformationsFinitePeriod}.
\end{description}

\subsection{Selection Obstruction / Repair Normal Form}\label{sec:symmetry:selection-obstruction}

\begin{definition}[Selection obstruction]\label{def:selection-obstruction}
Let a declared symmetry group \(G\) act on a carrier \(X\), and let
\[
  p:X\to X/G
\]
be the orbit quotient.  For a selectedness predicate
\(\chi:X\to\{0,1\}\), the selection obstruction is the split-pair set
\[
  \splitObs(p,\chi)
  =
  \{(x,x')\in X\times X:
      p(x)=p(x')\ \text{and}\ \chi(x)\ne\chi(x')\}.
\]
\end{definition}

\begin{theorem}[Selection Obstruction / Repair Normal Form]\label{thm:F15a}
Let \(G\) act on a set \(X\) and on an orbit set \(O\) by
\(\rho_X:G\times X\to X\) and \(\rho_O:G\times O\to O\).  Let
\(\pi:X\to O\) be a surjective orbit map, let
\(s:X\to\{0,1\}\) be a selectedness predicate, let
\(\sigma:O\to X\) be a section, and let \(b:X\to B\) be a breaker
readout.  Define
\[
  \mathcal O_s
  :=
  \{(x,x')\in X^2:\pi(x)=\pi(x')\ \text{and}\ s(x)\ne s(x')\}.
\]
Then
\begin{align*}
  &\operatorname{EquivariantSelector}(\rho_O,\rho_X,\sigma)
    \wedge \operatorname{TrivialOrbitAction}(\rho_O)\\
  &\hspace{4em}\Longrightarrow
    \forall o\in O,\ \forall g\in G,\
      \rho_X(g,\sigma(o))=\sigma(o),\\
  s\ \text{descends through}\ \pi
  &\iff
  \mathcal O_s=\emptyset
  \iff
  \exists\,\bar s:O\to\{0,1\}\ \text{with}\ \bar s\circ\pi=s,\\
  \operatorname{HasSelectionSplit}(\pi,s)
  &\iff
  \mathcal O_s\ne\emptyset.
\end{align*}
Thus an equivariant selector can select only common fixed points
unless a declared breaker \(b\), memory, condition, or quotient
refinement supplies additional data.
\end{theorem}
\footnote{Verified in Lean as
\texttt{Symmetry.SelectionObstruction.selection\_obstruction}; see
\Cref{app:formalization}.}
\begin{proof}
Assume the actions \(\rho_X,\rho_O\), the surjective orbit map
\(\pi\), the selectedness predicate \(s\), and the section
\(\sigma\).  If \(\sigma\) is equivariant and the \(G\)-action on
\(O\) is trivial, then for every \(g\in G\) and \(o\in O\),
\[
  \pi(\rho_X(g,\sigma(o)))
  =
  \rho_O(g,\pi(\sigma(o)))
  =
  \rho_O(g,o)
  =
  o.
\]
Equivariance of the selector gives
\(\rho_X(g,\sigma(o))=\sigma(\rho_O(g,o))=\sigma(o)\), so
\(\sigma(o)\) lies in the common fixed locus of the \(G\)-action on
\(X\).

For selectedness, if \(s=\bar s\circ\pi\), then
\(\pi(x)=\pi(x')\) implies \(s(x)=s(x')\), hence
\(\mathcal O_s=\emptyset\).  Conversely, if
\(\mathcal O_s=\emptyset\), define
\[
  \bar s(\pi(x)):=s(x).
\]
The empty obstruction makes the definition independent of the
representative \(x\), and surjectivity of \(\pi\) supplies
representatives for all \(o\in O\).  There are two cases.  If
\(\mathcal O_s=\emptyset\), selectedness is a quotient-level
predicate.  If \(\mathcal O_s\ne\emptyset\), a split pair
\((x,x')\) witnesses a non-descending selection, so a repair must add
declared breaker data \(b\), memory, a condition, or a refinement of
\(\pi\).
\end{proof}

The interaction status is \(\BirdInt\)-compatible when the group
action, orbit quotient, selection interface, and any breaker record
are declared interaction data.  The theorem does not classify which
breaking data are allowed and does not instantiate the obstruction to
a particular symmetry group.

\subsection*{Layer instantiations}\label{layer-inst:F15a}

\begin{description}
\item[\emph{Vitali sets and the axiom of choice (mathematical
logic).}]
For the free \(\mathbb Q\)-action on \(\mathbb R\) by translation,
no equivariant section of the orbit projection
\(\mathbb R\to\mathbb R/\mathbb Q\) can exist (there are no fixed
points).  The axiom of choice supplies a non-equivariant Vitali
set of representatives, which is necessarily non-measurable: the
F15a declared breaker is the axiom of choice itself, licensing a
non-equivariant selection where no equivariant one is available
\citep{Vitali1905ProblemaMisura}.

\item[\emph{Faddeev--Popov gauge fixing and the Gribov--Singer
obstruction (gauge theory).}]
Singer's theorem states that the gauge orbit space
\(\mathcal A/\mathcal G\) of Yang--Mills theory admits no
continuous global section.  Faddeev--Popov gauge fixing provides
the F15a declared breaker through a local gauge condition
\(F(A)=0\) together with the Faddeev--Popov determinant
\(\det(\delta F/\delta\omega)\) and ghost fields, licensing local
path-integral computations; the global Gribov--Singer obstruction
witnesses the non-existence of a clean equivariant selector
\citep{FaddeevPopov1967YangMills}.

\item[\emph{Hairy ball theorem / Poincar{\'e}--Hopf (differential
topology).}]
On the 2-sphere \(S^2\), no continuous nonvanishing tangent vector
field exists.  The Poincar{\'e}--Hopf theorem identifies the
obstruction with the Euler characteristic
\(\chi(S^2)=2\neq 0\): no \(\mathrm{SO}(3)\)-equivariant
nonvanishing section of the tangent bundle \(TS^2\) is possible.
Selecting a nonvanishing direction field requires a declared
breaker --- a non-equivariant atlas, a preferred polar axis, or a
stochastic perturbation
\citep{Poincare1885CourbesEquationsDifferentielles}.

\item[\emph{Schelling focal points in coordination games (game
theory).}]
In a pure-coordination game with multiple symmetric equilibria, no
symmetric pure strategy can pin down a unique equilibrium.  Schelling
identified \emph{focal points} --- salient landmarks, precedents,
or conventions distinguishing one equilibrium from its symmetric
partners --- as the F15a declared breaker; mixed-strategy
randomisation provides an equivariant but non-selective
alternative
\citep{Schelling1960StrategyConflict}.

\item[\emph{Buridan's ass and symmetric tie-breaking (decision
theory).}]
Under exact symmetry between two equally optimal choices, no
equivariant deterministic decision rule selects one option.  The
declared breaker is a tie-breaking rule --- random coin flip, a
lexicographic preference axis, or a memory of past choices --- that
supplies the non-symmetric apparatus needed for a definite
selection
\citep{Rescher1959ChoiceWithoutPreference}.
\end{description}

\subsection{Spontaneous Symmetry Breaking / Branch-Selection Formation}\label{sec:symmetry:spontaneous-ssb}

The setup is a formed closure under an exact declared symmetry, with
the orbit quotient and selection interface already in place as in the
previous subsection.  A spontaneous symmetry-breaking claim asks for
more than a set-theoretic section of the orbit map.  It asks for a
formed branch-selection record: a typed witness that a breaker source
has been generated, stabilized, and carried by the layer.

\begin{theorem}[Spontaneous Symmetry Breaking / Branch-Selection Formation]\label{thm:F15b}
Let \(\mathcal C\) be a closure space with a \(G\)-action, let
\(\Omega:\mathcal C\to\mathrm{Prop}\) be the ground-state predicate,
and let \(\mathcal B=\mathcal C/G\) be the branch quotient.  Suppose
a future branch-selection formation source is given for
\((\mathcal C,G,\Omega)\).  Then there is a formation operator
\[
  F:\mathcal C\longrightarrow\mathcal B
\]
and data \(c\in\mathcal C\), \(g\in G\) such that
\[
  \Omega(c)
  \quad\text{and}\quad
  g\cdot F(c)\ne F(c).
\]
Equivalently,
\[
  \exists F:\mathcal C\to\mathcal B,\,
  \exists c\in\mathcal C,\,
  \exists g\in G,\,
  \Omega(c)\wedge g\cdot F(c)\ne F(c).
\]
\end{theorem}
\footnote{The Lean theorem
\texttt{Symmetry.SpontaneousSSB.spontaneous\_ssb\_formation}
proves the layer-agnostic statement under an obligation axiom
\texttt{spontaneous\_ssb\_formation} guarded by an opaque source
predicate; no instantiation of the axiom can derive \texttt{False}
from the axiom alone.  The parallel diagnosis context appears in
\citep{Tsiokos2026CSL}.  See \Cref{app:formalization} for the guard
register.}
\begin{proof}
Assume the closure space \(\mathcal C\), the \(G\)-action, the
ground-state predicate \(\Omega\), and the future source instance.
Apply the same-paper obligation axiom to this source instance.  The
axiom asserts exactly the existence of branch-selection formation
data:
\[
  \exists F:\mathcal C\to\mathcal B,\,
  \exists c\in\mathcal C,\,
  \exists g\in G,\,
  \Omega(c)\wedge g\cdot F(c)\ne F(c).
\]
Choose \(F,c,g\) from this witness.  The inequality says that the
selected branch is not \(G\)-fixed, while \(\Omega(c)\) says it is
reachable from a ground state.  The formation operator \(F\) is the
typed branch-selection record, so the asymmetry is carried by formed
selection data rather than overread from the symmetric quotient.
\end{proof}

This subsection does not instantiate the future branch-selection
source, does not discharge the source condition unconditionally, and
does not classify spontaneous symmetry breaking at a substrate level.
It records the layer-agnostic formation claim that becomes direct
once the future source supplies the required branch-selection
certificate.  The canonical particle-physics instances of spontaneous
symmetry breaking are surveyed in Goldstone, Salam, and Weinberg
\citep{Goldstone1962BrokenSymmetries}, and the broader emergent-layer
reading of multi-level breaking appears in
Anderson's \emph{More is Different} \citep{Anderson1972MoreIsDifferent}.

\subsection*{Layer instantiations}\label{layer-inst:F15b}

\begin{description}
\item[\emph{Bose--Einstein condensation (atomic physics).}]
When a dilute gas of bosons is cooled below its critical
temperature, a macroscopic fraction of atoms condenses into a
single one-particle quantum state.  The condensate's overall
\(U(1)\) phase is arbitrary --- a continuous symmetry of the
underlying dynamics --- yet each realisation selects a definite
phase, breaking the symmetry.  The 1995 Anderson--Ensher--
Matthews--Wieman--Cornell experiment on rubidium-87 confirmed
Einstein's 1925 prediction; the macroscopic phase selection is the
F15b formation operator
\citep{AndersonEnsherMatthewsWiemanCornell1995BEC}.

\item[\emph{Rayleigh--B{\'e}nard convection (fluid dynamics).}]
A horizontal fluid layer heated from below transitions, at a
critical Rayleigh number (\(\mathrm{Ra}_c\approx 1708\) for the
standard rigid no-slip boundary case and lower for stress-free
boundaries), from a translation- and rotation-invariant
conduction state to a patterned state of convection rolls or
hexagonal cells.  Each experimental realisation fixes a particular
roll axis or hexagonal sublattice, selecting a branch in the
continuous symmetry orbit
\citep{Chandrasekhar1961HydrodynamicStability}.

\item[\emph{Soai chiral autocatalysis (chemistry).}]
In the Soai reaction the product catalyses its own formation with
the same chirality; starting from a near-racemic mixture, the
autocatalytic amplification drives the system to nearly
enantiopure product, spontaneously selecting one mirror-image
branch.  The underlying kinetics are chirally symmetric, but the
autocatalytic formation operator picks a single enantiomer per
realisation
\citep{SoaiShibataMoriokaChoji1995AsymmetricAutocatalysis}.

\item[\emph{Biological homochirality of life (origin-of-life
biology).}]
Terrestrial biology employs L-amino acids and D-sugars almost
exclusively, despite a mirror-symmetric prebiotic chemistry.
Candidate F15b formation-mechanism classes include
Frank-style autocatalytic amplification, photochemical bias from
circularly polarised light, and evolutionary locking after an
early enantiomeric bias; each names a declared branch-selection
source by which the biosphere settled into one chirality branch
\citep{Bonner1991OriginAmplificationBiomolecularChirality}.

\item[\emph{Crystallisation from supercooled liquid (condensed
matter).}]
A homogeneous, isotropic liquid below its freezing point is
thermodynamically metastable; nucleation of a crystalline phase
breaks continuous translations and \(\mathrm{SO}(3)\) rotations to
a discrete crystallographic subgroup, with each nucleation event
selecting a particular lattice orientation and translation phase
\citep{Turnbull1950FormationCrystalNuclei}.
\end{description}

\subsection{Anomaly / Symmetry Obstruction}\label{sec:symmetry:anomaly}

Let a formed closure have carrier \(H\), layer quotient
\(q:H\to Q\), declared readouts and audit records, and a would-be
symmetry gate \(g\) acting on the raw carrier.  The gate is a genuine
layer symmetry only when it acts functorially on the quotient and
preserves the declared structure of the layer.

\begin{theorem}[Anomaly / Symmetry Obstruction]\label{thm:F40}
Let \(G\) act on a history space \(H\) by
\(\alpha:G\times H\to H\), let \(q:H\to Q\) be a surjective quotient,
let \(\bar\alpha:G\times Q\to Q\) be a candidate descended action,
let \(D:G\times H\to\{0,1\}\) be a domain predicate, and let
\(r:H\to V\) be a readout with value action
\(\rho:G\times V\to V\).  Define
\[
  \mathcal O_{\alpha}
  :=
  \{(g,h,h'):\ q(h)=q(h')\ \text{and}\
    q(\alpha(g,h))\ne q(\alpha(g,h'))\},
\]
and define \(\mathcal O_D\) analogously by replacing
\(q(\alpha(g,h))\) with \(D(g,h)\).  Then
\begin{align*}
  \operatorname{SymmetryDescends}(q,\alpha)
  &\iff \mathcal O_{\alpha}=\emptyset,\\
  \operatorname{DomainDescends}(q,D)
  &\iff \mathcal O_D=\emptyset,\\
  \operatorname{GroupLawPreserved}(\bar\alpha)
  &\iff
  \forall g,g'\in G,\forall \bar h\in Q,\,
    \bar\alpha(gg',\bar h)=
    \bar\alpha(g,\bar\alpha(g',\bar h)),\\
  \operatorname{ReadoutPreserved}(r,\alpha,\rho)
  &\iff
  \forall g\in G,\forall h\in H,\,
    r(\alpha(g,h))=\rho(g,r(h)).
\end{align*}
An anomaly is precisely the failure of at least one of these four
conditions.
\end{theorem}
\footnote{Verified in Lean as
\texttt{Symmetry.Anomaly.anomaly\_symmetry\_obstruction}.}
\begin{proof}
Assume the action \(\alpha\), quotient \(q\), candidate descended
action \(\bar\alpha\), domain predicate \(D\), readout \(r\), and
value action \(\rho\).  The action descends through \(q\) exactly when
the formula
\[
  \bar\alpha(g,q(h)):=q(\alpha(g,h))
\]
is independent of the representative \(h\).  That independence is
equivalent to \(\mathcal O_{\alpha}=\emptyset\).  The domain predicate
descends by the identical argument applied to \(D(g,h)\), giving
\(\mathcal O_D=\emptyset\).

The descended action is a group action exactly when it respects
multiplication:
\[
  \bar\alpha(gg',\bar h)
  =
  \bar\alpha(g,\bar\alpha(g',\bar h))
  \quad
  \forall g,g'\in G,\ \bar h\in Q.
\]
The readout is \(G\)-equivariant exactly when
\[
  r(\alpha(g,h))=\rho(g,r(h))
  \quad
  \forall g\in G,\ h\in H.
\]
If all four displayed conditions hold, the symmetry is a formed-layer
symmetry.  If any one fails, the failed descent, domain, group-law, or
readout condition is the anomaly residue.
\end{proof}

The interaction status is \(\BirdInt\)-compatible when the raw action,
quotient action, overlap comparisons, and audit transformations are
declared interaction data.  The theorem does not classify anomalies
into carrier-specific forms and does not assert anything about
non-symmetric gates.  The standard quantum-field-theory treatment of
anomalies, including the Wess--Zumino consistency condition and the
descent-equation classification, is collected in
\citet{Bertlmann2000Anomalies}.

\subsection*{Layer instantiations}\label{layer-inst:F40}

\begin{description}
\item[\emph{{IEEE} 754 floating-point non-associativity (numerical
computing).}]
The would-be commutative ring symmetry of real arithmetic fails to
descend to the IEEE 754 floating-point representation:
\((a+b)+c\) generically differs from \(a+(b+c)\) due to rounding
at each operation.  The F40 reading is a failure of
GroupLawPreserved on the floating-point readout; numerical
libraries declare explicit accumulation orders (Kahan or
pairwise summation) as the apparatus extension that mitigates the
anomaly
\citep{IEEE7542019FloatingPoint}.

\item[\emph{Gimbal lock in Euler-angle parametrisations
(3D graphics).}]
The Lie group \(\mathrm{SO}(3)\) is a smooth 3-manifold, but its
Euler-angle parametrisation has coordinate singularities --- gimbal
lock --- where two rotation axes align.  At those configurations
DomainDescends fails: the parametrisation domain breaks down and
the orientation readout loses a degree of freedom.  Quaternion or
rotation-matrix representations are the declared apparatus repair
\citep{Kuipers1999QuaternionsRotationSequences}.

\item[\emph{Lattice discretisation and continuum
translational symmetry (computational physics).}]
A continuum PDE such as Navier--Stokes has continuous translational
symmetry, with momentum conservation following from Noether's
theorem.  A finite-difference discretisation on a fixed spatial
lattice respects only a discrete translation subgroup: both
SymmetryDescends and DomainDescends fail for the continuous group,
and the downstream momentum readout fails to be preserved.
Symplectic and energy-momentum-preserving discretisations
(variational integrators) are the apparatus repair
\citep{HairerLubichWanner2006GeometricNumerical}.

\item[\emph{Triangular arbitrage in currency exchange (finance).}]
In a frictionless ideal market, exchange rates would compose
multiplicatively to a triangular product equal to 1
(e.g.\ \(\mathrm{USD/EUR}\cdot\mathrm{EUR/JPY}\cdot
\mathrm{JPY/USD}=1\)).  Bid--ask spreads, transaction fees, and
latency break GroupLawPreserved on the round-trip composition,
producing triangular-arbitrage opportunities when the product
exceeds or falls below unity.  Modelling friction explicitly or
competitive-arbitrage pressure are the declared apparatus repairs
\citep{Aliber1973InterestRateParity}.

\item[\emph{Galilean vs special-relativistic velocity addition
(foundational physics).}]
The would-be Galilean velocity-addition group \(v_{AB}=v_{AC}+v_{CB}\)
fails to be closed on the physical velocity domain \(|v|<c\):
DomainDescends fails when Galilean composition of sublight
velocities can exit the declared domain, and the velocity readout
fails to be preserved.  Einstein's relativistic addition
\(v_{AB}=(v_{AC}+v_{CB})/(1+v_{AC}\cdot v_{CB}/c^2)\) --- equivalently
rapidity additivity --- is the declared apparatus repair
\citep{Einstein1905Elektrodynamik}.
\end{description}

\subsection{Duality Equivalence}\label{sec:symmetry:duality-equivalence}

Let \(\mathcal L_A\) and \(\mathcal L_B\) be formed closures with
declared symmetry presentations.  Each presentation has an object
quotient, transport data, declared observables, role registry,
residual registry, and status discipline.  A proposed duality is not
just a dictionary between symbols.  It is a pair of translations
between the two formed packages, with the claim that the translations
preserve the roles being compared.

\begin{theorem}[Duality Equivalence]\label{thm:F41}
Let \(\mathcal C_A,\mathcal C_B\) be formed closures with quotients
\(q_A:\mathcal C_A\to Q_A\) and \(q_B:\mathcal C_B\to Q_B\).  Let
\(F:\mathcal C_A\to\mathcal C_B\) and
\(G:\mathcal C_B\to\mathcal C_A\) descend to
\(\bar F:Q_A\to Q_B\) and \(\bar G:Q_B\to Q_A\).  Let
\(r_A:R_A\to S_A\), \(r_B:R_B\to S_B\) be role readouts, let
\(\tau:R_A\to R_B\) be role translation, let
\(\rho:S_A\to S_B\) be the role-value map, and let
\(\phi:\mathcal R_A\to\mathcal R_B\) and
\(\kappa:\mathcal K_A\to\mathcal K_B\) translate routes and claims.
Then a role-preserving duality equivalence holds if and only if
\begin{align*}
  \bar G\circ\bar F&=\mathrm{id}_{Q_A}
  \quad\text{and}\quad
  \bar F\circ\bar G=\mathrm{id}_{Q_B},\\
  r_B\circ\tau&=\rho\circ r_A,\\
  \bar F\circ R_A(\gamma)
  &=R_B(\phi(\gamma))\circ\bar F
  \quad\forall \gamma\in\mathcal R_A,\\
  \operatorname{status}_B(\kappa(k))
  &=\operatorname{status}_A(k)
  \quad\forall k\in\mathcal K_A.
\end{align*}
\end{theorem}
\footnote{Verified in Lean as
\texttt{Symmetry.DualityEquivalence.duality\_equivalence}.}
\begin{proof}
Assume the quotients \(q_A,q_B\), translations \(F,G\), descended
maps \(\bar F,\bar G\), role readouts \(r_A,r_B\), role translation
\(\tau\), route map \(\phi\), and claim translation \(\kappa\).  If
the presentations are role-preservingly dual, then translating to
\(B\) and back to \(A\) changes no quotient-level object, and
similarly in the other direction:
\[
  \bar G\bar F=\mathrm{id}_{Q_A},
  \qquad
  \bar F\bar G=\mathrm{id}_{Q_B}.
\]
Role preservation gives \(r_B\circ\tau=\rho\circ r_A\).  Route
preservation says that translating after following a route is the
same as following the translated route after translating:
\[
  \bar F(R_A(\gamma)(x))
  =
  R_B(\phi(\gamma))(\bar F(x)).
\]
Claim preservation gives equality of translated statuses.

Conversely, suppose the four displayed conditions hold.  The first
line makes \(Q_A\) and \(Q_B\) equivalent quotient presentations.
The second line transports role observables, the third transports
routes, and the fourth transports claim statuses.  These are exactly
the declared role data of the formed packages, so the two symmetry
presentations determine equivalent formed roles.
\end{proof}

The interaction status is \(\BirdInt\)-compatible when the
translations, role maps, observable families, and status-preservation
checks are declared interaction data.  The theorem does not classify
symmetry presentations and does not assert a unique role presentation
beyond observable equivalence.

\subsection*{Layer instantiations}\label{layer-inst:F41}

\begin{description}
\item[\emph{Pontryagin duality (harmonic analysis).}]
For a locally compact abelian group \(G\), the Pontryagin dual
\(\widehat G\) of continuous characters \(G\to S^1\) satisfies the
Pontryagin duality theorem
\(G\cong\widehat{\widehat G}\), with the Fourier transform
\(L^2(G)\to L^2(\widehat G)\) a unitary isomorphism.  The
category-level translations \(F=(-)^\wedge\) and \(G=(-)^\wedge\)
are mutually inverse on quotients, preserve role readouts
(characters and Haar-measure pairings), commute with route action
(convolution \(\leftrightarrow\) pointwise multiplication), and
preserve declared status (group structure, topology)
\citep{Pontryagin1966TopologicalGroups}.

\item[\emph{Stone duality (logic and topology).}]
Stone's representation theorem gives a contravariant equivalence
between Boolean algebras and Stone spaces (compact totally
disconnected Hausdorff spaces): \(\mathrm{Spec}\) sends a Boolean
algebra to its ultrafilter space, \(\mathrm{Clop}\) sends a Stone
space to its Boolean algebra of clopen subsets, and
\(A\cong\mathrm{Clop}(\mathrm{Spec}(A))\),
\(X\cong\mathrm{Spec}(\mathrm{Clop}(X))\).  Role readouts (Boolean
operations \(\leftrightarrow\) topological operations on clopens),
routes (homomorphisms \(\leftrightarrow\) continuous maps), and
status (the Boolean-algebra / Stone-space structures) are
preserved on both sides
\citep{Stone1936RepresentationsBooleanAlgebras}.

\item[\emph{Gelfand duality (operator algebras and topology).}]
The Gelfand--Naimark theorem gives a contravariant equivalence
between commutative unital \(C^*\)-algebras and compact Hausdorff
spaces: \(\mathrm{Spec}\) returns the space of characters and
\(C(\cdot,\mathbb C)\) returns the algebra of continuous functions,
with the Gelfand transform an isometric \(*\)-isomorphism
\(A\cong C(\mathrm{Spec}(A))\) and homeomorphism
\(X\cong\mathrm{Spec}(C(X))\).  Role readouts (\(*\)-algebra
operations \(\leftrightarrow\) continuous-function operations),
routes (\(*\)-homomorphisms \(\leftrightarrow\) continuous maps),
and status (commutative-\(C^*\) versus compact-Hausdorff) are
preserved
\citep{GelfandNaimark1943NormedRings}.

\item[\emph{Curry--Howard correspondence (logic and type theory).}]
For intuitionistic natural deduction and the simply typed
\(\lambda\)-calculus, Curry--Howard gives an exact syntactic
equivalence between propositions and types, proofs and
\(\lambda\)-terms: introduction/elimination rules correspond to
abstraction/application, and cut-elimination corresponds to
\(\beta\)-normalisation.  Roles (formula \(\leftrightarrow\) type,
proof \(\leftrightarrow\) program), routes (proof rules
\(\leftrightarrow\) typing rules), and status (provability
\(\leftrightarrow\) inhabitation, normalisation) are preserved by
mutually inverse translations between the two formal systems
\citep{Howard1980FormulaeAsTypes}.

\item[\emph{{AdS}/{CFT} holographic correspondence (string
theory; conjectural).}]
Maldacena's AdS/CFT correspondence conjectures that quantum gravity
on a \((d+1)\)-dimensional anti-de Sitter spacetime is dual to a
conformal field theory on its \(d\)-dimensional asymptotic
boundary, with the canonical example type IIB string theory on
\(\mathrm{AdS}_5\times S^5\) and \(\mathcal N=4\) super-Yang--Mills
on the boundary.  Within the accepted holographic dictionary, bulk
gravitational operators map to boundary CFT operators, partition
functions match, and AdS isometries map to CFT conformal
symmetries --- the proposed F41 role/route/status preservation
across the duality
\citep{Maldacena1998LargeNSuperconformal}.
\end{description}

\subsection*{Symmetry at a glance}\label{sec:symmetry:summary}
\begin{table}[t]
  \centering
  \footnotesize
  \begin{tabularx}{\textwidth}{@{}
      >{\raggedright\arraybackslash}p{0.24\textwidth}
      >{\raggedright\arraybackslash}X
      >{\raggedright\arraybackslash}p{0.18\textwidth}
      >{\raggedright\arraybackslash}p{0.12\textwidth}@{}}
    \toprule
    Paper-prose name & One-line claim & Status & Substrate cite \\
    \midrule
    Duality-Fixity Theorem &
    Completed domination plus trace vanishing forces visible mass
    into \(\Fixinv\). &
    Partial wrapper &
    \PaperSDTC, \PaperRH \\
    Selection-Obstruction Theorem &
    Exact symmetry cannot select a non-fixed branch without declared
    breaking data. &
    Theorem &
    \textemdash \\
    Spontaneous SSB &
    Opaque future branch-selection source implies a formed
    branch-selection record. &
    Obligation wrapper &
    \PaperCSL\ (parallel) \\
    Anomaly Theorem &
    Anomaly occurs when a declared symmetry gate fails descent,
    gluing, or audit preservation. &
    Theorem &
    \textemdash \\
    Duality-Equivalence Theorem &
    Two symmetry presentations are equivalent iff they agree on
    declared observables. &
    Theorem &
    \textemdash \\
    \bottomrule
  \end{tabularx}
  \caption{Symmetry axis at a glance: the five laws of \Cref{sec:symmetry}.}
  \label{tab:axis-symmetry-summary}
\end{table}
  \section{Axis V — Status, Records, Coherence}\label{sec:status-records-coherence}

The Status, Records, and Coherence axis treats records as part of the
mathematical object rather than as commentary about it.  A fact is not
merely a sentence someone writes down after a construction succeeds; a
formed fact has a claim record, a status family, and a place in the
coherence ledger that records how it persists, transports, composes,
and supports explanation.  This chapter collects the laws governing
that bookkeeping when the bookkeeping itself is structural.

The fourteen laws proceed in document order from shared objects to
counterfactual use.  Objectivity as Common Quotient identifies when
multiple presentations share an object.  Lawhood Descent tests when a
law factors through the available quotient, and Object Persistence
tests when the same object survives transport.  Fact / Record
Formation is the lead theorem: it makes formal what it means for a
fact or record to form inside the apparatus.  Reflexive Nonclosure,
Interface Mediation, Probability as Stable Unresolved Fiber, and
Causality as Stable Intervention Transport then describe how records
behave under self-audit, interfaces, unresolved fibers, and
interventions.

The remaining laws move from records to coherence.  Moduli /
Contingency separates stable parameters from contingent choices.
Conservation as Orbit Descent describes conservation as a quotient
law on orbits.  Composability and Presentation Invariance say when
formed records survive composition and change of presentation.
Explanation as Compression with Descent treats explanation as a
compressed account that still factors through the declared quotient,
and Counterfactual Legitimacy states when counterfactual claims are
lawful rather than overreads.

All fourteen rows in this chapter are direct theorem rows.  After
\Cref{sec:apparatus}, the chapter is self-contained: it uses the same
quotient, descent, and interaction-status vocabulary fixed there, but
does not rely on substrate wrappers or obligations.  Because this is
the densest axis, the at-a-glance table at the end groups the laws by
sub-axis: status, records, and coherence.

\subsection{Objectivity as Common Quotient}\label{sec:status-records-coherence:objectivity}

\begin{theorem}[Objectivity as Common Quotient]\label{thm:F17}
Let \(H\) be a carrier with a family of access quotients
\(q_i:H\to Q_i\) indexed by \(i\in I\), a public quotient
\(p:H\to P\), and a target readout \(r:H\to R\).  Assume each
\(q_i\) is surjective, \(p\) is surjective, and the family satisfies
the common-quotient universal property
\(\operatorname{CommonQuotientUniversal}(q,p)\): \(p\) factors
through every \(q_i\), and any quotient through which all \(q_i\)
factor also factors through \(p\).  Then
\begin{align*}
  &\operatorname{CommonAcrossAccesses}(q,r)
    \iff \operatorname{PublicDeterminesObject}(p,r),\\
  &\operatorname{CommonAcrossAccesses}(q,r)\\
  &\qquad{}\iff
    \forall i\in I,\,
      \operatorname{AccessObjectivityDefectEmpty}(q_i,r),\\
  &\operatorname{PublicDeterminesObject}(p,r)\\
  &\qquad{}\iff
    \operatorname{ObstructionEmpty}
      \bigl(\operatorname{PublicToObjectDefect}(p,r)\bigr),
\end{align*}
together with
\begin{align*}
  &\operatorname{CommonQuotientVisible}(q,p),\\
  &\operatorname{RetainsPublicQuotient}
      \bigl(p,\operatorname{objectiveExtension}(p,r)\bigr),\\
  &\operatorname{CarriesObject}
      \bigl(\operatorname{objectiveExtension}(p,r),r\bigr),
\end{align*}
and the status partition
\[
  \operatorname{ExactlyOneObjectivityStatus}(p,r).
\]
\end{theorem}
\footnote{\raggedright Verified in Lean as
\texttt{StatusRecordsCoherence.\allowbreak{}ObjectivityAsCommonQuotient.\allowbreak{}objectivity\_as\_common\_quotient};
see \Cref{app:formalization}.}
\begin{proof}
Assume the surjectivity of each \(q_i\), the surjectivity of \(p\),
and \(\operatorname{CommonQuotientUniversal}(q,p)\).  We prove the
seven conjuncts in turn.

\emph{(1) Common across accesses iff public determines object.}
The universal property supplies comparison maps
\[
  p=\pi_i\circ q_i\quad(i\in I),
  \qquad
  s=\lambda\circ p
  \quad\text{whenever all }q_i\text{ factor through }s .
\]
Thus \(p\) is the public quotient generated by the family
\((q_i)_{i\in I}\).  A readout \(r\) is common across the access
quotients exactly when it is constant on the common public fibers:
\[
  \forall h,h'\in H,\quad
  p(h)=p(h')\Longrightarrow r(h)=r(h').
\]
Since \(p\) is surjective, this is equivalent to the existence of
\(\bar r:P\to R\) with
\[
  r=\bar r\circ p,
\]
which is precisely \(\operatorname{PublicDeterminesObject}(p,r)\).

\emph{(2) Common across accesses iff every access defect is empty.}
For \(i\in I\), the access objectivity defect is the set
\[
  \Delta_i(q_i,r)
  =
  \{(h,h')\in H\times H:
      q_i(h)=q_i(h')\ \text{and}\ r(h)\neq r(h')\}.
\]
The equality \(\Delta_i(q_i,r)=\emptyset\) is equivalent to
\[
  \forall h,h'\in H,\quad
  q_i(h)=q_i(h')\Longrightarrow r(h)=r(h').
\]
Requiring this for every \(i\in I\) is exactly
\(\operatorname{CommonAcrossAccesses}(q,r)\), giving the second
equivalence.

\emph{(3) Public determines object iff the public-to-object
obstruction is empty.}
The public-to-object defect is
\[
  \Delta_P(p,r)
  =
  \{(h,h')\in H\times H:
      p(h)=p(h')\ \text{and}\ r(h)\neq r(h')\}.
\]
Its emptiness is the standard quotient-descent condition for \(r\)
through \(p\).  By surjectivity of \(p\), this is equivalent to a
well-defined \(\bar r:P\to R\) satisfying \(r=\bar r\circ p\), hence
to \(\operatorname{PublicDeterminesObject}(p,r)\).

\emph{(4) Common quotient visible.}
The same universal property exhibits \(p\) as a declared common
quotient of the family \((q_i)_{i\in I}\):
\[
  \operatorname{CommonQuotientUniversal}(q,p)
  \Longrightarrow
  \operatorname{CommonQuotientVisible}(q,p).
\]
Visibility here is the statement that the common public quotient is
one of the quotients carried by the declared structure, not a quotient
introduced after the fact.

\emph{(5) The objective extension retains the public quotient.}
Let
\[
  \operatorname{objectiveExtension}(p,r)
  :=
  \langle p,r\rangle:H\longrightarrow P\times R,
  \qquad
  h\longmapsto (p(h),r(h)).
\]
Projection to the first coordinate recovers the public quotient:
\[
  \operatorname{pr}_P\circ\langle p,r\rangle=p.
\]
Therefore the objective extension retains \(p\).

\emph{(6) The objective extension carries the object.}
Projection to the second coordinate recovers the object readout:
\[
  \operatorname{pr}_R\circ\langle p,r\rangle=r.
\]
Thus the same extension carries the object \(r\).

\emph{(7) Exactly one objectivity status.}
The four status predicates are determined by whether the access
defects \(\Delta_i(q_i,r)\) and the public defect \(\Delta_P(p,r)\)
are empty.  The cases are
\[
\begin{array}{c|c|c}
\text{status} & \forall i,\ \Delta_i(q_i,r)=\emptyset
  & \Delta_P(p,r)=\emptyset\\
\hline
\text{objective} & \text{yes} & \text{yes}\\
\text{access-defective} & \text{no} & \text{yes}\\
\text{public-defective} & \text{yes} & \text{no}\\
\text{multi-defective} & \text{no} & \text{no}
\end{array}
\]
Every configuration of the two Boolean defect tests falls into one
row, and no configuration falls into two rows.  Hence
\(\operatorname{ExactlyOneObjectivityStatus}(p,r)\).
\end{proof}

The theorem does not determine which presentations are declarable and
does not enumerate the kinds of objects that may be shared.  The
closest physical instantiation is quantum Darwinism, in which
objectivity emerges as redundant agreement of the environment's
records across many access channels \citep{Zurek2009QuantumDarwinism};
the layer-agnostic statement above is independent of that
substrate.

\subsection*{Layer instantiations}\label{layer-inst:F17}

\begin{description}
\item[\emph{Fixed-effect meta-analysis (statistics).}]
Under a declared fixed-common-effect model, a fixed-effect
meta-analysis combines independent study-level effect-size
estimates by inverse-variance weighting to produce a pooled
consensus estimate that the study-level estimates factor through
when the model holds.  Cochran's \(Q\) and \(I^2\) measure the
F17 access-objectivity defect: heterogeneity is the witness that
the declared common-effect model is violated
\citep{HedgesOlkin1985MetaAnalysis}.

\item[\emph{State-machine replication via Paxos (distributed
systems).}]
In a Paxos-style replicated state machine, each replica observes
its own log of accepted operations while the consensus protocol
commits a single sequence that all non-faulty replicas converge
on.  The per-replica log views are F17 access quotients, the
committed sequence is the public quotient, and the safety property
of consensus is the common-across-accesses condition; split-brain
configurations are the F17 access-objectivity defect
\citep{Lamport1998PartTimeParliament}.

\item[\emph{The {WGS-84} geodetic reference frame (geodesy).}]
National and regional geodetic datums (NAD-83, ETRS-89, JGD2011)
each provide a local access quotient on geodetic position, and
each maps to WGS-84 via published transformations.  WGS-84 is the
F17 public quotient through which the local datums factor, and
agreement under WGS-84 reduction is the common-across-accesses
property
\citep{NGA2014WGS84}.

\item[\emph{The {SI} as public measurement quotient (metrology).}]
National metrology institutes (NIST, NPL, PTB, etc.) maintain
their own primary standards realising the seven SI base units
through the seven defining constants
\((c,h,e,k,N_A,K_{\mathrm{cd}},\Delta\nu_{\mathrm{Cs}})\).  The SI
is the public quotient through which all national standards
factor, and SI-traceable measurement agreement across laboratories
under the CIPM Mutual Recognition Arrangement is the F17
common-across-accesses property; unresolved Key-Comparison
discrepancies are the access-objectivity defect
\citep{BIPM2019SIBrochure}.

\item[\emph{Thermodynamic ensemble equivalence (statistical
mechanics).}]
In the thermodynamic limit and for short-range systems away from
phase transitions, the microcanonical, canonical, and grand-canonical
ensembles yield equivalent equations of state for macroscopic
intensive quantities.  Each ensemble is an access quotient on
phase space, and the thermodynamic equation of state is the F17
public quotient through which all three factor.  Ensemble
inequivalence at long-range systems / phase-transition
pathologies is the access-objectivity defect
\citep{Ruelle1969StatisticalMechanics}.
\end{description}

\subsection{Lawhood Descent}\label{sec:status-records-coherence:lawhood-descent}

\begin{theorem}[Lawhood Descent]\label{thm:F18}
Let a formed closure have available access quotient
\(\accessQ:H\to Q\).  A declared candidate law \(P\) is lawful at the
layer if and only if it factors through \(\accessQ\); for a predicate,
this means \(P(h)=P(h')\) whenever \(\accessQ(h)=\accessQ(h')\).
\end{theorem}
\footnote{\raggedright Verified in Lean as
\texttt{StatusRecordsCoherence.LawhoodDescent.\allowbreak{}lawhood\_descent};
see \Cref{app:formalization}.}
\begin{proof}
If \(P=\bar P\circ\accessQ\), then \(P\) is constant on access
classes by substitution.  Conversely, if \(P\) is constant on access
classes, define \(\bar P([h])=P(h)\).  Constancy makes this
well-defined, and the defining equation gives
\(P=\bar P\circ\accessQ\).  Relations and transition rules are handled
by applying the same representative-independence argument in each
argument.  Thus lawhood is precisely quotient descent.
\end{proof}

The theorem does not enumerate lawful candidates and does not
construct the access quotient; that is the role of Quotientality.

\subsection*{Layer instantiations}\label{layer-inst:F18}

\begin{description}
\item[\emph{Wilson loops in gauge theory (physics).}]
In Yang--Mills theory, a candidate physical observable is a
lawful observable at the gauge-quotient iff it is
gauge-invariant.  Wilson loops --- traces of the holonomy around
closed loops --- are the canonical gauge-invariant
observables, descending through the gauge-orbit projection
\(\mathcal A\to\mathcal A/\mathcal G\)
\citep{Wilson1974ConfinementQuarks}.

\item[\emph{Parnas information hiding (software engineering).}]
Parnas's principle of information hiding declares that a module
exposes a public interface and hides its implementation behind
that interface.  A client predicate is a lawful use of the
module iff it depends only on the declared public operations;
predicates that read hidden state are overreads.  The interface
is the F18 access quotient through which lawful client claims
descend
\citep{Parnas1972InformationHiding}.

\item[\emph{Bisimulation invariance in process calculi
(concurrency theory).}]
In Milner--Park bisimulation theory, a Hennessy--Milner
modal-logic formula is a lawful process property iff it is
bisimulation-invariant.  The bisimulation quotient on the
labelled transition system is the F18 access quotient through
which lawful process-level laws descend
\citep{Milner1989CommunicationConcurrency}.

\item[\emph{Congruences and quotient algebras (universal
algebra).}]
A predicate or operation on an algebra \(A\) descends to a
well-defined predicate on the quotient \(A/{\sim}\) iff
\(\sim\) is a congruence (i.e., \(\sim\) is compatible with all
operations and the predicate is \(\sim\)-invariant); laws over
quotient algebras correspond exactly to congruence-respecting
predicates on the original algebra
\citep{Gratzer1979UniversalAlgebra}.

\item[\emph{Stone--{\v C}ech compactification and continuous
extensions (topology).}]
For a Tychonoff space \(X\), the Stone--{\v C}ech
compactification \(\beta X\) is characterised by the universal
property that every continuous bounded real-valued function on
\(X\) extends uniquely to a continuous function on \(\beta X\).
A predicate declared on a dense subset is lawful on the closure
iff it admits a continuous extension --- the F18 descent
criterion for topological laws
\citep{Cech1937BicompactSpaces}.
\end{description}

\subsection{Object Persistence}\label{sec:status-records-coherence:object-persistence}

\begin{theorem}[Object Persistence]\label{thm:F19}
Let \(H_i,H_j\) be carriers with access quotients
\(q_i:H_i\to Q_i\) and \(q_j:H_j\to Q_j\), and let
\(\gamma:H_i\to H_j\) be a declared transport.  Assume \(q_i\) and
\(q_j\) are surjective.  Let
\(\operatorname{viable}:Q_j\to\mathrm{Prop}\) be a viability
predicate, \(C:H_i\times H_j\to\mathrm{Prop}\) a relation predicate,
and let
\[
  \bar\gamma:Q_i\to Q_j,\qquad
  \bar\eta:Q_j\to Q_k,\qquad
  \operatorname{compBar}:Q_i\to Q_k
\]
be arbitrary holonomy data (no compatibility between
\(\operatorname{compBar}\) and \(\bar\eta\circ\bar\gamma\) is
assumed).  Then
\begin{align*}
  &\operatorname{Persists}(q_i,\gamma,q_j)
    \iff
    \operatorname{PersistenceObstructionEmpty}(q_i,\gamma,q_j),\\
  &\operatorname{SplitObstructionEmpty}(q_i,\gamma,q_j)
    \iff
    \operatorname{Persists}(q_i,\gamma,q_j),\\
  &\operatorname{ViablePersistence}(\gamma,q_j,\operatorname{viable})
    \iff
    \operatorname{DissolutionObstructionEmpty}
      (\gamma,q_j,\operatorname{viable}),\\
  &\operatorname{RelationValuedPersistence}(q_i,q_j,C)
    \iff
    \operatorname{RelationPersistenceObstructionEmpty}(q_i,q_j,C),\\
  &\operatorname{IdentityHolonomyCoherent}
      (\bar\gamma,\bar\eta,\operatorname{compBar})\\
  &\qquad{}\iff
    \operatorname{IdentityHolonomyDefectEmpty}
      (\bar\gamma,\bar\eta,\operatorname{compBar}),
\end{align*}
together with the comparison-status partition
\[
  \operatorname{ExactlyOnePersistenceComparisonStatus}
    (q_i,\gamma,q_j),
\]
and the retention and carrying conditions
\begin{align*}
  &\operatorname{RetainsOldIdentity}
      \bigl(q_i,\operatorname{persistenceExtension}(q_i,\gamma,q_j)\bigr),\\
  &\operatorname{CarriesLaterIdentity}
      \bigl(\operatorname{persistenceExtension}(q_i,\gamma,q_j),
        \gamma,q_j\bigr).
\end{align*}
\end{theorem}
\footnote{\raggedright Verified in Lean as
\texttt{StatusRecordsCoherence.ObjectPersistence.\allowbreak{}object\_persistence};
see \Cref{app:formalization}.}
\begin{proof}
Assume the surjectivity of \(q_i\) and \(q_j\), the predicates
\(\operatorname{viable}\) and \(C\), and the holonomy data
\(\bar\gamma,\bar\eta,\operatorname{compBar}\).  We prove the eight
conjuncts in turn.

\emph{(1) Persists iff the persistence obstruction is empty.}
The persistence obstruction is
\[
  \Delta_{\mathrm{pers}}(q_i,\gamma,q_j)
  =
  \{(h,h')\in H_i\times H_i:
      q_i(h)=q_i(h')\ \text{and}\
      q_j(\gamma h)\ne q_j(\gamma h')\}.
\]
Its emptiness is exactly the quotient-descent condition for
\(q_j\circ\gamma:H_i\to Q_j\) through \(q_i\).  Since \(q_i\) is
surjective, this is equivalent to the existence of a descended map
\(\bar\gamma:Q_i\to Q_j\) satisfying
\[
  \bar\gamma\circ q_i=q_j\circ\gamma,
\]
which is \(\operatorname{Persists}(q_i,\gamma,q_j)\).

\emph{(2) Split-obstruction empty iff persists.}
The split obstruction is the same pair set written in the comparison
vocabulary:
\[
  \Delta_{\mathrm{split}}(q_i,\gamma,q_j)
  =
  \Delta_{\mathrm{pers}}(q_i,\gamma,q_j).
\]
Thus \(\Delta_{\mathrm{split}}=\emptyset\) iff the persistence
obstruction is empty, and by (1) this is equivalent to persistence.

\emph{(3) Exactly-one persistence comparison status.}
The comparison statuses are decided by two split-pair defects on the
declared continuation \(\gamma\): the persistence/split obstruction
\(\Delta_{\mathrm{split}}\) of case~(1), and the merge obstruction
\[
  \Delta_{\mathrm{merge}}(q_i,\gamma,q_j)
  =
  \{(h,h')\in H_i\times H_i:
      q_j(\gamma h)=q_j(\gamma h')\ \text{and}\ q_i(h)\ne q_i(h')\}.
\]
The four statuses partition the configurations by the joint
\(\Delta_{\mathrm{split}}\)\textendash{}\(\Delta_{\mathrm{merge}}\)
emptiness pattern:
\[
\begin{array}{c|c|c}
\text{status} & \Delta_{\mathrm{split}}=\emptyset
  & \Delta_{\mathrm{merge}}=\emptyset\\
\hline
\text{exact endurance} & \text{yes} & \text{yes}\\
\text{persistent with merge} & \text{yes} & \text{no}\\
\text{identity branching} & \text{no} & \text{yes}\\
\text{identity rewrite} & \text{no} & \text{no}
\end{array}
\]
Every configuration of the two Boolean defect tests falls into
exactly one row.  Existence is by case enumeration; exclusivity is by
definitional disjointness of the four status predicates.  The
dissolution and relation-valued obstructions of cases~(4) and~(5) are
separate typed gates and do not enter this partition.

\emph{(4) Viable persistence iff the dissolution obstruction is
empty.}
The dissolution obstruction is
\[
  \Delta_{\mathrm{diss}}(\gamma,q_j,\operatorname{viable})
  =
  \{h\in H_i:
      \neg\,\operatorname{viable}(q_j(\gamma h))\}.
\]
Viable persistence says that every transported \(q_j\)-class lies in
the viable region:
\[
  \forall h\in H_i,\quad
  \operatorname{viable}(q_j(\gamma h)).
\]
This is precisely \(\Delta_{\mathrm{diss}}=\emptyset\).

\emph{(5) Relation-valued persistence iff its obstruction is empty.}
The relation-persistence obstruction is
\[
  \Delta_C(q_i,q_j,C)
  =
  \{(h,h')\in H_i\times H_i:
    q_i(h)=q_i(h')\ \text{and}\
    C(h,\gamma h)\not\iff C(h,\gamma h')\}.
\]
Emptiness of \(\Delta_C\) says that the truth of the relation is
independent of the chosen representative of a \(q_i\)-class.  Hence
\(C\) descends to a well-defined relation on
\(Q_i\times Q_j\), which is exactly relation-valued persistence.

\emph{(6) Identity holonomy coherent iff the holonomy defect is
empty.}
The holonomy defect is
\[
  \Delta_{\mathrm{hol}}(\bar\gamma,\bar\eta,\operatorname{compBar})
  =
  \{a\in Q_i:
      \operatorname{compBar}(a)\ne\bar\eta(\bar\gamma(a))\}.
\]
Thus \(\Delta_{\mathrm{hol}}=\emptyset\) iff
\[
  \forall a\in Q_i,\quad
  \operatorname{compBar}(a)=\bar\eta(\bar\gamma(a)),
\]
which is the pointwise form of
\(\operatorname{compBar}=\bar\eta\circ\bar\gamma\), i.e. identity
holonomy coherence.

\emph{(7) The persistence extension retains the old identity.}
Define the persistence extension as the joint quotient
\[
  \operatorname{persistenceExtension}(q_i,\gamma,q_j)
  :=
  \langle q_i,q_j\circ\gamma\rangle:H_i\to Q_i\times Q_j.
\]
Projection to the first coordinate recovers the old identity:
\[
  \operatorname{pr}_{Q_i}\circ
  \langle q_i,q_j\circ\gamma\rangle=q_i.
\]
Therefore the persistence extension retains \(q_i\).

\emph{(8) The persistence extension carries the later identity.}
Projection to the second coordinate recovers the transported later
identity:
\[
  \operatorname{pr}_{Q_j}\circ
  \langle q_i,q_j\circ\gamma\rangle=q_j\circ\gamma.
\]
Thus the same extension carries \(q_j\) along the declared transport
\(\gamma\).
\end{proof}

The theorem does not classify persistent objects or say which
transports should be declared.

\subsection*{Layer instantiations}\label{layer-inst:F19}

\begin{description}
\item[\emph{Mendelian inheritance (genetics).}]
Through meiotic segregation and fertilisation, parental allele
identity at a locus survives into the offspring's diploid
genotype.  The F19 reading: allele identity persists when no
mutation or non-disjunction breaks the descent; mutation and
non-disjunction events are the persistence obstruction
\citep{Mendel1866VersuchePflanzenHybriden}.

\item[\emph{Comoving galaxy identity in FLRW cosmology
(cosmology).}]
In a Friedmann--Lema{\^\i}tre--Robertson--Walker universe, galaxies
moving with the Hubble flow keep their comoving coordinates as the
scale factor evolves.  Comoving-coordinate identity persists under
the cosmic-time transport; peculiar velocities and gravitational
mergers are the persistence obstruction
\citep{Friedmann1922KrummungRaumes}.

\item[\emph{Round-trip identity under {JSON} serialisation
(computer science).}]
For a data value \(D\), the round-trip
\(D\equiv\mathrm{deserialise}(\mathrm{serialise}(D))\) is the F19
persistence condition for the declared canonical-form identity
quotient.  Locale-dependent number formatting, time-zone-aware
datetimes, hash-randomised dictionary ordering, and floating-point
precision loss are common persistence obstructions
\citep{Crockford2006JSON}.

\item[\emph{Pure-state identity under unitary evolution (quantum
mechanics).}]
For a closed quantum system, time evolution is the unitary
\(U(t)=\exp(-iHt/\hbar)\); the projective ray
\([|\psi\rangle]\) (state modulo global phase) evolves through the
descended projective evolution.  Pure-state identity persists
under closed unitary evolution; decoherence from system--
environment coupling is the persistence obstruction
\citep{vonNeumann1932MathematischeGrundlagen}.

\item[\emph{Corporate personhood across statutory merger
(corporate law).}]
Under the Model Business Corporation Act \S\,11.07, a statutory
merger transfers the corporate personhood of constituent
corporations to a single surviving (or new) corporation, which
inherits the rights, properties, and obligations of the
constituents.  The legal-identity transport descends through the
corporate-personhood quotient under declared statutory procedures;
defective merger filings or undocumented dissolution without
successor are the persistence obstruction
\citep{MBCA2016Section1107}.
\end{description}

\subsection{Fact / Record Formation}\label{sec:status-records-coherence:fact-record}

\begin{definition}[Record formation]\label{def:record-formation}
In a formed closure \(\closureForm\), a fact is a claim candidate
that is closure-of-record and whose evidence descends through the
available access.  A record is the audited witness, carried by the
\(\BirdInt\) ledger, that the claim and evidence have formed at the
declared layer.
\end{definition}

\begin{theorem}[Fact / Record Formation]\label{thm:F20}
Let \(H_0,H_1,H_s,H_t\) be history spaces with event readout
\(\varepsilon:H_0\to E\), continuation \(\gamma:H_0\to H_1\),
record readouts
\[
  \rho_1:H_1\to R_1,\qquad
  \rho_s:H_s\to R_s,\qquad
  \rho_t:H_t\to R_t,
\]
transport \(\theta:H_s\to H_t\), descended transport
\(\bar\theta:R_s\to R_t\), fact-value maps
\(d_s:R_s\to E_s\), \(d_t:R_t\to E_t\), update
\(u:E_s\to E_t\), and a packaging map \(C:H\to H\) on an auxiliary
carrier \(H\).  Define the record signature
\[
  s:=\rho_1\circ\gamma:H_0\to R_1,
\]
and assume that \(s\) and \(\rho_s\) are surjective.  Define
\begin{align*}
  \mathcal O_{\mathrm{ev}}
  &:=
  \{(h,h')\in H_0^2:
      \rho_1(\gamma h)=\rho_1(\gamma h')
      \text{ and } \varepsilon(h)\ne\varepsilon(h')\},\\
  \mathcal O_{\mathrm{red}}
  &:=
  \{(h,h')\in H_0^2:
      \varepsilon(h)=\varepsilon(h')
      \text{ and } \rho_1(\gamma h)\ne\rho_1(\gamma h')\},\\
  \mathcal O_{\mathrm{rec}}
  &:=
  \{(x,y)\in H_s^2:
      \rho_s(x)=\rho_s(y)
      \text{ and } \rho_t(\theta x)\ne\rho_t(\theta y)\},\\
  \mathcal O_{\mathrm{val}}
  &:=
  \{r\in R_s:
      d_t(\bar\theta r)\ne u(d_s r)\}.
\end{align*}
Then
\begin{align*}
  \operatorname{EventRecorded}(\varepsilon,\gamma,\rho_1)
  &\iff \mathcal O_{\mathrm{ev}}=\emptyset,\\
  \operatorname{RecordStable}(\rho_s,\theta,\rho_t)
  &\iff \mathcal O_{\mathrm{rec}}=\emptyset,\\
  \operatorname{FactValueStable}(\bar\theta,d_s,d_t,u)
  &\iff \mathcal O_{\mathrm{val}}=\emptyset,\\
  \operatorname{StableFact}
  &\iff
  \mathcal O_{\mathrm{ev}}=\emptyset
  \wedge
  \mathcal O_{\mathrm{rec}}=\emptyset
  \wedge
  \mathcal O_{\mathrm{val}}=\emptyset.
\end{align*}
together with the record-comparison status partition
\[
  \operatorname{ExactlyOneRecordComparisonStatus}
    (\varepsilon,\gamma,\rho_1),
\]
the retention and carrying conditions for the record-event extension
\(\langle\rho_1\circ\gamma,\varepsilon\rangle:H_0\to R_1\times E\),
\begin{align*}
  &\operatorname{RetainsRecord}
      \bigl(\rho_1\circ\gamma,
        \langle\rho_1\circ\gamma,\varepsilon\rangle\bigr),\\
  &\operatorname{CarriesEvent}
      \bigl(\langle\rho_1\circ\gamma,\varepsilon\rangle,
        \varepsilon\bigr),
\end{align*}
and the idempotent-packaging fixed-point property
\[
  \operatorname{IdempotentPackaging}(C)
  \Longrightarrow
  \forall h\in H,\ C(C(h))=C(h).
\]
The record-comparison status takes one of four values
(\emph{exact}, \emph{redundant}, \emph{ambiguous}, or
\emph{incomparable}) determined by the joint emptiness pattern of
\(\mathcal O_{\mathrm{ev}}\) and \(\mathcal O_{\mathrm{red}}\).
\end{theorem}
\footnote{\raggedright Verified in Lean as
\texttt{StatusRecordsCoherence.FactRecordFormation.\allowbreak{}fact\_record\_formation};
see \Cref{app:formalization}.}
\begin{proof}
Assume the displayed maps, the surjectivity of \(s=\rho_1\circ\gamma\),
and the surjectivity of \(\rho_s\).  We prove the eight conjuncts in
turn.

\emph{(1) Event recorded iff event-record obstruction empty.}
If \(\operatorname{EventRecorded}(\varepsilon,\gamma,\rho_1)\) holds,
then the event readout is constant on record-signature fibers.  Hence
for all \(h,h'\in H_0\),
\[
  \rho_1(\gamma h)=\rho_1(\gamma h')
  \quad\Longrightarrow\quad
  \varepsilon(h)=\varepsilon(h'),
\]
so \(\mathcal O_{\mathrm{ev}}=\emptyset\).  Conversely, if
\(\mathcal O_{\mathrm{ev}}=\emptyset\), define the descended event
readout on the record signature \(s\) by
\[
  \bar\varepsilon(s(h)):=\varepsilon(h).
\]
The emptiness of \(\mathcal O_{\mathrm{ev}}\) makes this definition
independent of the representative \(h\), and surjectivity supplies
representatives for all record signatures.  This is precisely the
event-recording factorization.

\emph{(2) Record stable iff record-stability obstruction empty.}
The record-stability equivalence is the same quotient argument for
\(\rho_s\).  The condition \(\mathcal O_{\mathrm{rec}}=\emptyset\) is
exactly the independence needed to define
\[
  \bar\theta(\rho_s x):=\rho_t(\theta x),
  \qquad\text{so that}\qquad
  \bar\theta\circ\rho_s=\rho_t\circ\theta.
\]

\emph{(3) Fact-value stable iff the fact-value defect is empty.}
The fact-value defect at \(r\in R_s\) is the failure of the
commutative square
\[
  d_t\circ\bar\theta=u\circ d_s.
\]
Equivalently,
\[
  \mathcal O_{\mathrm{val}}=\emptyset
  \quad\Longleftrightarrow\quad
  \forall r\in R_s,\ d_t(\bar\theta r)=u(d_s r),
\]
which is precisely \(\operatorname{FactValueStable}(\bar\theta,d_s,d_t,u)\).

\emph{(4) Stable fact iff the three obstructions are empty.}
By definition, a stable fact requires event recording, record
stability, and fact-value stability.  Combining (1)--(3) gives
\[
  \operatorname{StableFact}
  \iff
  \mathcal O_{\mathrm{ev}}=\emptyset
  \wedge
  \mathcal O_{\mathrm{rec}}=\emptyset
  \wedge
  \mathcal O_{\mathrm{val}}=\emptyset.
\]

\emph{(5) Exactly-one record-comparison status.}
The four statuses are decided by the joint
\((\mathcal O_{\mathrm{ev}},\mathcal O_{\mathrm{red}})\) emptiness
pattern:
\[
\begin{array}{c|c|c}
\text{status} & \mathcal O_{\mathrm{ev}}=\emptyset
  & \mathcal O_{\mathrm{red}}=\emptyset\\
\hline
\text{exact record} & \text{yes} & \text{yes}\\
\text{redundant record} & \text{yes} & \text{no}\\
\text{ambiguous record} & \text{no} & \text{yes}\\
\text{incomparable record} & \text{no} & \text{no}
\end{array}
\]
Every configuration of the two Boolean defect tests falls into one
row, and no configuration falls into two rows.  Existence is by case
enumeration; exclusivity is by definitional disjointness of the four
status predicates.

\emph{(6) Retains record.}
Define the canonical record-event extension
\[
  \langle\rho_1\circ\gamma,\varepsilon\rangle:
  H_0\longrightarrow R_1\times E,
  \qquad
  h\longmapsto(\rho_1(\gamma h),\varepsilon(h)).
\]
Projection to the first coordinate recovers the record signature:
\[
  \operatorname{pr}_{R_1}
  \circ\langle\rho_1\circ\gamma,\varepsilon\rangle
  =
  \rho_1\circ\gamma.
\]

\emph{(7) Carries event.}
Projection to the second coordinate recovers the event readout:
\[
  \operatorname{pr}_{E}
  \circ\langle\rho_1\circ\gamma,\varepsilon\rangle
  =
  \varepsilon.
\]

\emph{(8) Idempotent packaging fixes every \(C(h)\).}
Suppose \(\operatorname{IdempotentPackaging}(C)\), i.e.
\[
  \forall h\in H,\quad C(C(h))=C(h).
\]
For each \(h\in H\), set \(h':=C(h)\).  Then
\[
  C(h')=C(C(h))=C(h)=h',
\]
so \(C(h)\) is a fixed point of \(C\).
\end{proof}

The theorem does not determine which candidates count as evidence and
does not identify the fact with the record.

\subsection*{Layer instantiations}\label{layer-inst:F20}

\begin{description}
\item[\emph{{FAIR} scientific-data principles (research
methodology).}]
The FAIR Guiding Principles --- Findable, Accessible,
Interoperable, Reusable --- declare that scientific data form
stable facts iff generation events are recorded as provenance
metadata, records persist under storage / migration transport via
persistent identifiers and access protocols, and derived analyses
preserve fact-value consistency through declared processing
pipelines.  Each obstruction (unrecorded provenance, broken DOI,
undocumented pipeline change) is an F20-style stable-fact failure
\citep{WilkinsonDumontier2016FAIRPrinciples}.

\item[\emph{Blockchain ledger and confirmation depth (distributed
computer science).}]
In a blockchain, each transaction is recorded by cryptographic
signature, included in a block, and chained to prior blocks via
hash pointers; ``stable fact'' is properly indexed by
confirmation depth.  Signature-verification failure is the
event-recording obstruction, hash tampering or chain
reorganisations and forks are the record-stability obstruction,
and a double-spend constitutes the fact-value obstruction
\citep{Nakamoto2008Bitcoin}.

\item[\emph{Forensic chain of custody (criminal procedure).}]
Physical evidence is admissible only when its chain of custody
is fully documented from collection through analysis: each
transfer between custodians is the recorded event, an unbroken
custodial log is the stable record, and laboratory results
matching the seized item are the consistent fact-value.  Gaps,
mishandling, or tampering are F20 stable-fact failures
\citep{Rice2008ElectronicEvidence}.

\item[\emph{Pacioli double-entry bookkeeping (accounting).}]
Pacioli's double-entry method requires every transaction to be
posted as two equal opposing entries, preserving the accounting
identity \(\text{Assets}=\text{Liabilities}+\text{Equity}\).  An
unrecorded transaction is the event-recording obstruction, a
posting error is the record-stability obstruction, and a
trial-balance discrepancy is the fact-value obstruction
\citep{Pacioli1494Summa}.

\item[\emph{Paleontological taphonomy and the fossil record
(natural science).}]
Fossil-record formation is the natural-history analogue of F20:
preserved burial events, geological persistence, and consistent
inferential analyses yield robust paleontological facts.  Pre-burial
scavenging and decay are the event-recording obstruction;
metamorphism, erosion, and exhumation damage are the
record-stability obstruction; conflicting species-identification
or dating analyses are the fact-value obstruction.  The fossil
record is therefore selective rather than complete, with stable
facts indexed by preservation-and-inferential quality
\citep{KidwellHolland2002QualityFossilRecord}.
\end{description}

\subsection{Reflexive Nonclosure}\label{sec:status-records-coherence:reflexive-nonclosure}

\begin{theorem}[Reflexive Nonclosure]\label{thm:F21}
Let \(\mathcal C\) be a space of claims, let
\(\pi\) be a same-layer audit policy on \(\mathcal C\), and let
\(\mu\) be a meta-layer audit relation from \(\mathcal C\) to a
space \(\mathcal M\) of meta-claims.  For every
\(c\in\mathcal C\),
\begin{align*}
  &\operatorname{SelfCertificationBlocked}_{\pi}(c),\\
  &\operatorname{AcceptedScopedLowerAudit}_{\pi}(c)
    \Longrightarrow
    \neg\operatorname{CompleteSelfSoundness}_{\pi}(c),\\
  &\operatorname{HasMetaLayerCertificate}_{\mu}(c)
    \Longrightarrow
    \exists m\in\mathcal M,\,
      \operatorname{Audits}_{\mu}(m,c)
      \wedge
      \neg\operatorname{SelfSound}_{\pi}(c),\\
  &\operatorname{CompleteSelfSoundness}_{\pi}(c)
    \Longrightarrow
    \operatorname{OutsideScope}_{\pi}(c).
\end{align*}
\end{theorem}
\footnote{\raggedright Verified in Lean as
\texttt{StatusRecordsCoherence.ReflexiveNonclosure.\allowbreak{}reflexive\_nonclosure};
see \Cref{app:formalization}.}
\begin{proof}
Assume the same-layer policy \(\pi\) and the meta-layer audit
\(\mu\), and fix \(c\in\mathcal C\).  If \(c\) certified its own
complete soundness at the same layer, then the certification record
would have to be accepted by the policy whose soundness it certifies:
\[
  \pi(c)\ \text{accepts}\ \operatorname{Sound}(\pi,c)
  \quad\text{using only}\quad \pi(c).
\]
This is a circular same-layer dependency, so
\(\operatorname{SelfCertificationBlocked}_{\pi}(c)\) holds.

If \(\operatorname{AcceptedScopedLowerAudit}_{\pi}(c)\) holds, the
audit is scoped below complete self-soundness.  Thus it can certify
the scoped lower claim but not the full assertion
\(\operatorname{CompleteSelfSoundness}_{\pi}(c)\).

There are two cases for a legitimate reflexive audit.  In the
same-layer case, the previous paragraph blocks complete
self-soundness.  In the meta-layer case, a certificate supplies
\(m\in\mathcal M\) with
\[
  \operatorname{Audits}_{\mu}(m,c)
  \quad\text{and}\quad
  \neg\operatorname{SelfSound}_{\pi}(c),
\]
so the audit has moved to a strict meta-layer.  Therefore any claim
of complete same-layer self-soundness lies outside the scope of
\(\pi\).
\end{proof}

This does not forbid scoped self-audit; it forbids unrestricted
same-layer self-closure.  The classical proof-theoretic instance of
the same obstruction is G\"odel's second incompleteness theorem
\citep{Godel1931Unentscheidbare}; the modal-logic reformulation in
terms of a provability operator is developed in
\citet{Boolos1993LogicProvability}.

\subsection*{Layer instantiations}\label{layer-inst:F21}

\begin{description}
\item[\emph{{ISO}/{IEC} 17025 laboratory accreditation
(industrial metrology).}]
Under ISO/IEC 17025, a testing or calibration laboratory
cannot self-certify its measurement competence: an
accredited laboratory must be assessed by a recognised
accreditation body (an ILAC member such as A2LA, UKAS, or
DAkkS), which acts as the meta-layer.  Internal
quality-management records are a scoped lower audit,
sufficient
for internal corrective action but not for legal or
commercial recognition of competence; the meta-layer
accreditation certificate is the legitimate reflexive
audit witness \citep{ISOIEC17025_2017}.

\item[\emph{Sarbanes-Oxley independent-audit mandate
(corporate-finance regulation).}]
The Sarbanes-Oxley Act of 2002 forbids public-company
financial statements from being audited by the issuer's
own accounting staff.  Management's internal-controls
assessment under \S404(a) is the scoped lower audit; the
independent CPA firm's attestation under \S404(b) is the
meta-layer certificate, and the auditor is itself
overseen by the Public Company Accounting Oversight
Board, producing a tower of meta-audits.  Complete
self-attestation by the issuer is outside the scope of
the regime \citep{SarbanesOxley2002Act}.

\item[\emph{X.509 PKI external root-of-trust
(cryptographic protocol design).}]
In the X.509 public-key infrastructure, an endpoint
certificate cannot self-certify its trustworthiness to a
relying party: validation requires a chain of signatures
terminating at a root certificate authority in the
client's externally maintained trust store, which acts
as the meta-layer.  The leaf's own self-signature and
local syntactic integrity checks are the scoped lower
audit; self-signed leaves are rejected by default by
mainstream TLS clients.  Complete in-certificate
self-trust is outside the scope of the protocol
\citep{Cooper2008RFC5280}.

\item[\emph{Rice's theorem on undecidability of semantic
properties (computability theory).}]
Rice's theorem rules out any same-layer universal
semantic decider for non-trivial behavioural properties
of partial recursive functions; specialised to
self-application, no program can act as a complete
same-layer
semantic self-certifier for a non-trivial behavioural
property.  Bounded-resource testing on a finite input
set within a step bound is a scoped lower audit, accepted
but insufficient for the full semantic claim.  Legitimate
reflexive certification of a particular instance must
invoke a strictly stronger meta-formalism (proof-assistant
formalisation against a specification, abstract
interpretation sound for the target property, or
model-checking of a decidable abstraction), which
certifies particular claims rather than serving as a
same-layer universal decider \citep{Rice1953Classes}.

\item[\emph{M\"unchhausen trilemma in epistemology
(analytic epistemology).}]
Hans Albert's M\"unchhausen trilemma argues that any
attempted self-grounding of a knowledge claim faces
three exhaustive options, none of which constitutes
complete same-layer self-certification: infinite
regress, logical circularity, or arbitrary termination
in unargued axioms.  A locally consistent dialectical
chain of reasons is the scoped lower audit; a legitimate
epistemic certificate must invoke a meta-layer
commitment (foundationalist intuition, coherentist
network, externalist reliability, or transcendental
argument), each of which exits the same-layer chain of
reasons \citep{Albert1985TreatiseCriticalReason}.
\end{description}

\subsection{Interface Mediation}\label{sec:status-records-coherence:interface-mediation}

\begin{theorem}[Interface Mediation]\label{thm:F22}
Let \(q:X\to Q\) be a source quotient, let
\(\beta:Y\to B\) be a boundary readout, let
\(U:X\times Y\to X'\) be an interface update, and let
\(q':X'\to Q'\) be the target quotient.  Assume \(q\), \(\beta\),
and \(q'\) are surjective.  Define
\[
  \mathcal O_U
  :=
  \{((x,y),(x',y')) :
    q(x)=q(x'),\ \beta(y)=\beta(y'),\
    q'(U(x,y))\ne q'(U(x',y'))\}.
\]
Then
\begin{align*}
  \operatorname{InterfaceMediated}(q,\beta,U,q')
  &\iff \mathcal O_U=\emptyset\\
  &\iff
  \exists\,\bar U:Q\times B\to Q'\ \text{with}\
    \bar U(q(x),\beta(y))=q'(U(x,y))\ \forall x,y,\\
  \mathcal O_U=\emptyset
  &\Longrightarrow
  \operatorname{BoundaryLeakageAbsent}\wedge
  \operatorname{InternalInsufficiencyAbsent},\\
  \operatorname{ActiveMediatedInfluence}(\bar U)
  &\iff \neg\operatorname{ExternallyInert}(\bar U).
\end{align*}
\end{theorem}
\footnote{\raggedright Verified in Lean as
\texttt{StatusRecordsCoherence.InterfaceMediation.\allowbreak{}interface\_mediation};
see \Cref{app:formalization}.}
\begin{proof}
Assume the source quotient \(q\), boundary readout \(\beta\),
interface update \(U\), and target quotient \(q'\).  If
\(\operatorname{InterfaceMediated}(q,\beta,U,q')\) holds, then
changing representatives inside the same \(q\)-class and
\(\beta\)-class cannot change the target \(q'\)-class; hence
\(\mathcal O_U=\emptyset\).

Conversely, suppose \(\mathcal O_U=\emptyset\).  Define
\[
  \bar U(q(x),\beta(y)):=q'(U(x,y)).
\]
If \(q(x)=q(x')\) and \(\beta(y)=\beta(y')\), the emptiness of
\(\mathcal O_U\) gives
\[
  q'(U(x,y))=q'(U(x',y')),
\]
so \(\bar U\) is independent of representatives.  Thus the
interface update descends through \(Q\times B\).  The absence of
\(\mathcal O_U\) also rules out the two possible mediation failures:
boundary-equivalent inputs cannot leak distinct target classes, and
source-equivalent internal states cannot create an undeclared target
difference.  Finally, \(\bar U\) has active mediated influence
exactly when it is not constant in the external argument, which is
the negation of external inertness.
\end{proof}

The theorem does not determine the interface's internal structure and
does not apply to undeclared interfaces.

\subsection*{Layer instantiations}\label{layer-inst:F22}

\begin{description}
\item[\emph{Abstract-data-type encapsulation (software
engineering).}]
In an abstract data type, the carrier is the space of
concrete representations (e.g., a stack realised as a
linked list, a dynamic array, or a deque); the source
quotient collapses representations satisfying the same
representation invariant.  The boundary is the public
method invocation; the update is the implementation's
response; the target quotient reads the post-state at
the public-interface level.  Encapsulation is the
requirement that two representations in the same source
class respond identically (modulo the target quotient)
to the same public-method call; private-state exposure
through the public interface is the boundary leakage
failure mode \citep{LiskovZilles1974ADT}.

\item[\emph{Black-box security reduction
(theoretical cryptography).}]
A standard security proof shows that any adversary
breaking a scheme can be turned into an algorithm
solving an underlying hard problem; in the black-box
reduction discipline, the reduction accesses the
adversary only through oracle queries.  The carrier is
the adversary's internal state, the source quotient
collapses adversaries with the same query/response
distribution, the boundary is the oracle transcript,
and the update is the adversary's next step.  The
reduction's success thereby depends only on the
boundary I/O profile, not on the adversary's internal
code.  The random-oracle methodology, in which the hash
function is itself modelled as a uniformly random
oracle, is one influential setting within this
discipline \citep{BellareRogaway1993RandomOracles}.

\item[\emph{Pharmacophore / structure-activity relationship
(medicinal chemistry).}]
Within the declared pharmacophore model of a target
receptor, the source quotient collapses candidate ligands
to the spatial / electronic arrangement of features
required for binding (hydrogen-bond donors and acceptors,
hydrophobic regions, charged groups); the boundary is the
receptor binding-site geometry; the update is binding
and downstream signalling; the target quotient is the
pharmacological response class.  Two molecules sharing
the same pharmacophore are predicted, within the model,
to produce the same response class at the same receptor;
deviations attributable to off-pharmacophore substituents
(allosteric effects, distinct ADMET profiles) are
diagnosed as boundary-leakage corrections outside the
declared interface \citep{Wermuth1998IUPACGlossary}.

\item[\emph{Vienna Convention on Diplomatic Relations
(international law).}]
The 1961 Vienna Convention formalises the diplomatic
interface between sovereign states.  The carrier is the
sending state's internal political and administrative
state; the source quotient collapses internal states
producing the same officially endorsed position; the
boundary is the recognised diplomatic channel (note
verbale, official meeting, protected correspondence)
delivered via accredited diplomatic agents; the update
is the receiving state's policy response through its
foreign ministry; the target quotient is the receiving
state's diplomatic posture.  Article 41 explicitly
forbids diplomats from interfering in the receiving
state's internal affairs, codifying the prohibition on
source-internal leakage past the diplomatic boundary
\citep{ViennaConventionDiplomatic1961}.

\item[\emph{Thermodynamic state functions and boundary
protocol (classical thermodynamics).}]
For a thermodynamic system following a declared
boundary protocol (e.g., quasi-static, adiabatic, or
isothermal exchange under specified constraints on
volume, pressure, and particle exchange), the carrier
is the microstate space; the source quotient is the
macroscopic state (energy, volume, particle number);
the boundary is the protocol-fixed heat and work
exchange together with its constraint structure; the
update is microscopic dynamical evolution; the target
quotient reads the post-process macrostate.  Within the
declared protocol, the macrostate change depends only
on the initial macrostate and the protocol-fixed
boundary data, not on the microstate trajectory; this
is the existence of state functions \citep{Callen1985Thermodynamics}.
\end{description}

\subsection{Probability as Stable Unresolved Fiber}\label{sec:status-records-coherence:probability}

\begin{theorem}[Probability as Stable Unresolved Fiber]\label{thm:F23}
Let \(q:H\to Q\) be a quotient of histories, let
\(\lambda:Q\to L\) be a ledger map, let
\(S\subseteq L\times H\) be a support relation, let
\(m:L\to (H\to\mathrm{Prop})\to W\) be a weight assignment, and let
\(r:H\to Z\) be a readout.  For \(a\in Q\) and \(v\in Z\), set
\[
  \operatorname{Pr}_{q,r}(a,v)
  :=
  m(\lambda(a))\bigl(\{h\in H:q(h)=a\ \text{and}\ r(h)=v\}\bigr).
\]
Then
\begin{align*}
  r\ \text{descends through}\ q
  &\iff
  \forall h,h'\in H,\ q(h)=q(h')\Rightarrow r(h)=r(h'),\\
  \operatorname{Unresolved}(q,r)
  &\iff
  \exists h,h'\in H,\ q(h)=q(h')\ \text{and}\ r(h)\ne r(h'),\\
  \operatorname{StableProbability}(q,r,\lambda,m)
  &\iff
  \operatorname{Unresolved}(q,r)
  \wedge
  \operatorname{Stable}(\operatorname{Pr}_{q,r}).
\end{align*}
\end{theorem}
\footnote{\raggedright Verified in Lean as
\texttt{StatusRecordsCoherence.\allowbreak{}ProbabilityAsStable\allowbreak{}UnresolvedFiber.\allowbreak{}probability\_as\_stable\_\allowbreak{}unresolved\_fiber};
see \Cref{app:formalization}.}
\begin{proof}
Assume the quotient \(q\), ledger \(\lambda\), support relation
\(S\), weight assignment \(m\), and readout \(r\).  The first
equivalence is the quotient-descent criterion: a descended readout
\(\bar r:Q\to Z\) satisfies
\[
  \bar r(q(h))=r(h)\quad \forall h\in H,
\]
if and only if \(r\) is constant on every \(q\)-fiber.

There are two cases.  If the displayed constancy condition holds,
then each fiber has a single readout value, and the probability
record is deterministic on that fiber.  If it fails, then there are
\(h,h'\in H\) with
\[
  q(h)=q(h')\quad\text{and}\quad r(h)\ne r(h'),
\]
which is exactly unresolvedness.  The formula defining
\(\operatorname{Pr}_{q,r}\) assigns weights to the unresolved values
inside the fiber \(q^{-1}(a)\).  Such a weight assignment is a
probability record precisely when this function is invariant under
the allowed perturbations, i.e. when
\(\operatorname{Stable}(\operatorname{Pr}_{q,r})\) holds.
\end{proof}

The theorem does not determine numerical probability values and does
not apply to resolved fibers.  The measure-theoretic foundation of
probability is \citet{Kolmogorov1933Grundbegriffe}; the branch /
decoherent-histories interpretation that comes closest to a stable
unresolved-fiber reading appears in \citet{Griffiths1984ConsistentHistories}.

\subsection*{Layer instantiations}\label{layer-inst:F23}

\begin{description}
\item[\emph{Design-based survey sampling (social
statistics).}]
For a fixed finite population, the carrier is the set of
possible sample-draw realisations under the chosen
randomisation scheme; the source quotient collapses
draws sharing the same population-level summary; the
readout is the realised sample statistic.  The same
population-level summary supports many distinct sample
realisations.  The sampling design's inclusion-probability
measure is the ledger; the design-induced probability
law on draws is the weight assignment.  The
Horvitz-Thompson estimator and the design-based variance
estimator are the stable probability records derived
from this design measure
\citep{Cochran1977SamplingTechniques}.

\item[\emph{Chain-ladder loss-reserving in non-life
insurance (actuarial science).}]
The carrier is the space of individual
claim-development histories within an accident year;
the source quotient collapses histories sharing the
same observed cumulative-payment pattern through the
current development cell; the readout is the ultimate
loss, which is unresolved over the observed
development-pattern fibre.  Mack's distribution-free
assumptions on the cumulative payments constitute the
ledger and supply the first-and-second-moment structure
of incremental development factors as the weight
assignment.  The chain-ladder reserve estimator and
Mack's prediction-error formula are the stable
probability records used for reserve setting and
solvency capital \citep{Mack1993ChainLadder}.

\item[\emph{Wright-Fisher stochastic population genetics
(theoretical biology).}]
The carrier is the set of gamete-sampling realisations
within a generation; the source quotient collapses
realisations to the current-generation allele frequency;
the readout is the realised next-generation allele
count.  Two distinct gamete-sampling outcomes sharing
the same current frequency can produce different
next-generation counts, so the readout is unresolved
over the current-frequency fibre.  The ledger records
the effective population size and any
selection/mutation parameters; the weight assignment is
the Wright-Fisher binomial (or diffusion) transition
kernel they induce.  The resulting transition law is
the stable probability record that descends to the
population-level frequency dynamics
\citep{Wright1931Evolution}.

\item[\emph{Shannon noisy-channel coding (information
theory).}]
The carrier is the set of (input, output) trajectories
of a noisy memoryless channel; the source quotient
collapses trajectories sharing the same observed output
sequence; the readout is the transmitted input.  The
receiver cannot uniquely recover the input from the
output (unresolved fibre).  The channel input
distribution supplies the ledger; the channel law
supplies the weight assignment; the maximum mutual
information over input distributions, the channel
capacity, is the stable probability-theoretic quantity
bounding achievable code rates.  The noisy-channel
coding theorem is the stable-probability theorem on
this fibre \citep{Shannon1948Mathematical}.

\item[\emph{Weibull life-distribution model in
reliability engineering (engineering reliability).}]
The carrier is the set of individual component
life trajectories under specified operating conditions;
the source quotient collapses trajectories sharing the
same operating-condition class; the readout is the
realised time-to-failure of a particular component.
Individual times-to-failure are unresolved over the
operating-condition fibre.  The Weibull distribution
function \(F(t)=1-\exp(-(t/\eta)^{\beta})\) is the
standard ledger and supplies a stable two-parameter
hazard model as the weight assignment; fitted shape
\(\beta\) and scale \(\eta\) give the stable
probability record used for population-level
reliability prediction, warranty pricing, and burn-in
design \citep{Weibull1951Statistical}.
\end{description}

\subsection{Causality as Stable Intervention Transport}\label{sec:status-records-coherence:causality}

\begin{theorem}[Causality as Stable Intervention Transport]\label{thm:F25}
Let \(q:H\to Q\) be a context quotient, let
\(\alpha:A\to\bar A\) be an intervention abstraction, let
\(\Gamma:H\times A\to H'\) be an intervention transport, let
\(r:H'\to R\) be a target readout, and let
\(K:H\times A\to K_0\) be a kernel outcome.  Let
\(\operatorname{Adm}:Q\times\bar A\to\mathrm{Prop}\) and let
\(\operatorname{ResStable}\) and
\(\operatorname{ProbResStable}\) be residue-stability predicates.
Define
\[
  \mathcal O_{\Gamma}
  :=
  \{((h,a),(h',a')) :
    q(h)=q(h'),\ \alpha(a)=\alpha(a'),\
    r(\Gamma(h,a))\ne r(\Gamma(h',a'))\}.
\]
Write \(\operatorname{ITD}\), \(\operatorname{PITD}\), and
\(\operatorname{Caus}\) for the intervention-transport descent,
probabilistic intervention-transport descent, and causal-relation
predicates for this data.  Then
\begin{align*}
  \operatorname{ITD}(q,\alpha,\Gamma,r)
  &\iff \mathcal O_{\Gamma}=\emptyset\\
  &\iff
  \exists\,\bar\Gamma:Q\times\bar A\to R,\
    \bar\Gamma(q(h),\alpha(a))=r(\Gamma(h,a))\ \forall h,a,\\
  \operatorname{PITD}(q,\alpha,K)
  &\iff
  \begin{gathered}
    \exists\,\bar K:Q\times\bar A\to K_0,\
    \bar K(q(h),\alpha(a))=K(h,a)\ \forall h,a,\\
    \operatorname{ProbResStable}(q(h),\alpha(a),\alpha(a'))
  \end{gathered}\\
  \operatorname{Caus}(q,\alpha,\Gamma,r)
  &\iff
  \begin{gathered}
    \operatorname{Adm}\wedge
    \operatorname{ITD}(q,\alpha,\Gamma,r)\\
    {}\wedge\operatorname{ResStable}.
  \end{gathered}
\end{align*}
\end{theorem}
\footnote{\raggedright Verified in Lean as
\texttt{StatusRecordsCoherence.\allowbreak{}CausalityAsStable\allowbreak{}InterventionTransport.\allowbreak{}causality\_as\_stable\_\allowbreak{}intervention\_transport};
see \Cref{app:formalization}.}
\begin{proof}
Assume the context quotient \(q\), intervention abstraction
\(\alpha\), transport \(\Gamma\), target readout \(r\), and kernel
outcome \(K\).  If the intervention transport descends, then same
context class and same intervention class have the same target
readout, so \(\mathcal O_{\Gamma}=\emptyset\).  Conversely, if
\(\mathcal O_{\Gamma}=\emptyset\), define
\[
  \bar\Gamma(q(h),\alpha(a)):=r(\Gamma(h,a)).
\]
The obstruction is empty exactly when this definition is independent
of the representatives \(h\) and \(a\), which proves the first
factorization equivalence.

The probabilistic statement is identical with the target readout
replaced by the kernel outcome:
\[
  \bar K(q(h),\alpha(a)):=K(h,a).
\]
The extra hypothesis \(\operatorname{ProbResStable}\) says that the
resulting kernel is stable under admissible changes of intervention
class.  There are two cases.  If admissibility, transport descent,
and residue stability all hold, the relation depends only on
\((q(h),\alpha(a))\) and is causal.  If any of the three fails, the
claimed causal relation depends on undeclared context, undeclared
intervention, or unstable residue, and so it is not a stable causal
transport.
\end{proof}

The theorem does not identify causes with raw interventions and does
not apply to undeclared alternatives.  The intervention-counterfactual
hierarchy that this layer-agnostic formulation echoes is developed
in detail in \citet{Pearl2009Causality}.

\subsection*{Layer instantiations}\label{layer-inst:F25}

\begin{description}
\item[\emph{Randomised controlled trial / Fisher
randomisation (experimental design).}]
The carrier is the population of experimental units
(patients, field plots); the context quotient collapses
units to their pre-randomisation baseline stratum; the
intervention space is the set of treatment arms; the
abstraction collapses operational delivery details to
the trial-arm class; the transport is the trial's
delivery-and-follow-up process; the target readout is
the primary outcome.  Individual outcomes do not
deterministically descend through (baseline stratum,
arm); the F25 reading is its probabilistic form:
randomised assignment supplies admissibility, and the
stable residue is the arm-specific outcome distribution
or the within-stratum average treatment effect, secured
by pre-registration, blinding, and intention-to-treat
analysis \citep{Fisher1935DesignExperiments}.

\item[\emph{Mendelian randomisation
(instrumental-variable epidemiology).}]
A germline genetic variant associated with a modifiable
exposure is used as a natural instrumental intervention
to estimate the causal effect of the exposure on a
disease outcome.  The carrier is the study population;
the context quotient collapses subjects sharing the
same confounder profile; the intervention space ranges
over genotypes at the chosen instrument locus; the
abstraction collapses genotypes to their associated
exposure level; the transport is the lifelong exposure
process; the readout is the disease outcome.
Admissibility rests on the conjunction of relevance,
independence (incorporating population-stratification
control), and the exclusion restriction; the descended
residue is the local-average exposure-outcome causal
effect \citep{DaveySmith2003Mendelian}.

\item[\emph{CRISPR-Cas9 directed gene knockout
(molecular biology).}]
An engineered guide-RNA-directed nuclease introduces a
targeted double-strand break at a specified genomic
locus, producing loss-of-function alleles in an
otherwise isogenic background.  The carrier is the
population of treated cells or organisms; the context
quotient collapses cells to their genetic background and
cell type; the intervention space ranges over candidate
guide RNAs and delivery conditions; the abstraction
collapses guide RNAs to the target locus they edit; the
transport is the editing event plus downstream
biological development; the target readout is the
phenotypic measurement.  Within the declared
editing-efficiency and off-target-control regime, the
phenotypic readout descends through
(background-and-cell-type, target-locus)
\citep{Jinek2012CRISPR}.

\item[\emph{Particle-scattering cross-section
measurement (high-energy physics).}]
A scattering experiment fires a beam of particles of
defined type and energy at a target of defined
composition; the carrier is the ensemble of individual
beam-target interaction events; the context quotient
collapses events to the target-and-beam apparatus
configuration; the intervention space parameterises
beam-energy and target choices; the abstraction
collapses operational details to the beam-target class;
the transport is the scattering process; the target
readout of a single event is the realised scattering
angle and energy.  Individual events are stochastic; the
F25 reading is its probabilistic form, with the
differential cross-section \(d\sigma/d\Omega\) as the
stable kernel residue over repeated events.
Rutherford's 1911 alpha-particle scattering off thin
foils is the founding instance
\citep{Rutherford1911Scattering}.

\item[\emph{Regression-discontinuity natural experiment
(empirical econometrics).}]
A regression-discontinuity design uses a known
deterministic cutoff in an assignment variable (test
score for scholarship eligibility, age for retirement
benefit, vote share for incumbency) to identify the
local average treatment effect at the threshold.  The
carrier is the population of units near the threshold;
the context quotient collapses units sharing the same
local cutoff neighbourhood of the running variable; the
intervention space is the treatment-assignment
indicator; the abstraction collapses delivery-channel
details to the just-above / just-below assignment; the
transport is the policy delivery process; the target
readout is the outcome of interest.  Admissibility
holds in the limit at the threshold; residue stability
is established by continuity of potential outcomes and
no-manipulation of the running variable around the
cutoff \citep{ThistlethwaiteCampbell1960Regression}.
\end{description}

\subsection{Moduli / Contingency}\label{sec:status-records-coherence:moduli}

\begin{theorem}[Moduli / Contingency]\label{thm:F26}
Let \(\mathcal C\) be a closure space with lawfulness predicate
\(\operatorname{Lawful}:\mathcal C\to\mathrm{Prop}\), equivalence relation
\(\sim\) on lawful closures, moduli quotient
\(p:\mathcal C\to M\), readout \(r:\mathcal C\to V\), transport
\(T:\mathcal C\to\mathcal C'\), and selected closure
\(c_0\in\mathcal C\).  Assume \(p\) is surjective on lawful
closures.  Then
\begin{align*}
  \operatorname{ClosureObstructed}(\operatorname{Lawful})
  &\iff
  \neg\exists c\in\mathcal C,\ \operatorname{Lawful}(c),\\
  \operatorname{NontrivialModuli}(\operatorname{Lawful},\sim)
  &\iff
  \exists c,d\in\mathcal C,\,
    \operatorname{Lawful}(c)\wedge\operatorname{Lawful}(d)
    \wedge \neg(c\sim d),\\
  \operatorname{NecessaryReadout}(r)
  &\iff
  \exists\,\bar r:M\to V\ \text{with}\
    \bar r(p(c))=r(c)\quad
    \forall c\ \operatorname{Lawful}(c),\\
  \operatorname{ContingentReadout}(r)
  &\iff
  \exists c,d\in\mathcal C,\,
    \operatorname{Lawful}(c)\wedge\operatorname{Lawful}(d)
    \wedge p(c)\ne p(d)\wedge r(c)\ne r(d).
\end{align*}
\end{theorem}
\footnote{\raggedright Verified in Lean as
\texttt{StatusRecordsCoherence.ModuliContingency.\allowbreak{}moduli\_contingency};
see \Cref{app:formalization}.}
\begin{proof}
Assume the lawfulness predicate, equivalence relation, moduli quotient
\(p\), readout \(r\), and surjectivity of \(p\) on lawful closures.
The first equivalence is definitional: obstruction means that the
lawful locus is empty,
\[
  \operatorname{ClosureObstructed}(\operatorname{Lawful})
  \Longleftrightarrow
  \{c\in\mathcal C:\operatorname{Lawful}(c)\}=\emptyset.
\]
Nontrivial moduli means that the lawful locus has at least two
inequivalent representatives, which is the second displayed
condition.

For the readout, suppose \(r\) is necessary.  Then it is constant on
the fibers of \(p\) over lawful closures, so the formula
\[
  \bar r(p(c)):=r(c)
\]
is independent of the representative \(c\).  This gives the descended
readout \(\bar r:M\to V\).  Conversely, the existence of such
\(\bar r\) forces \(r\) to be stable on moduli classes.  If no such
factorization exists, then some lawful pair has different moduli and
different readout values, which is the contingent case.  Thus the
moduli/contingency distinction is the quotient-stability split.
\end{proof}

The theorem does not enumerate a carrier's moduli and does not
collapse stable parameters with contingent choices.

\subsection*{Layer instantiations}\label{layer-inst:F26}

\begin{description}
\item[\emph{Moduli space of elliptic curves (algebraic
geometry).}]
The carrier is the set of Weierstrass-form presentations
\(y^2=x^3+ax+b\) over the base field; the lawfulness
predicate restricts to non-singular cubics
\((4a^3+27b^2\ne 0)\); the equivalence is isomorphism of
elliptic curves; the moduli quotient sends a
presentation to its isomorphism class.  The j-invariant
\(j(a,b)\) is a necessary readout: distinct Weierstrass
presentations of isomorphic elliptic curves share the
same j-invariant, and the descended map
\(\bar\jmath:M\to\mathbb A^1\) is a bijection onto the
coarse moduli space.  The individual coefficients
\((a,b)\) are presentation-contingent
\citep{Mumford1965GIT}.

\item[\emph{Crystallographic space groups and Bravais
lattices (solid-state physics).}]
The carrier is the set of three-dimensional crystalline
arrangements with given chemical content; the lawfulness
predicate requires translational and rotational
periodicity; the equivalence is space-group isomorphism
up to choice of origin and basis; the moduli quotient
sends a crystal to its space-group label among the 230
enumerated space groups.  The space-group label is a
necessary readout forced by the lawful symmetry
structure.  The lattice metric (unit-cell lengths
\(a,b,c\), angles \(\alpha,\beta,\gamma\),
atomic-displacement parameters) is contingent: lawful
crystals with the same space group can have different
lattice metrics \citep{HahnITC2016}.

\item[\emph{Standard Model free parameters as
renormalisation moduli (particle physics).}]
The carrier is the space of Standard-Model Lagrangians
with the fixed gauge group
SU(3)\(\times\)SU(2)\(\times\)U(1), three-generation
fermion content, and Higgs doublet built into the
lawfulness predicate (renormalisability, gauge
invariance, anomaly cancellation, observed particle
content); the equivalence identifies Lagrangians giving
identical low-energy phenomenology under the
renormalisation group; the moduli quotient sends a
Lagrangian to its renormalised-parameter point.  The 19+
dimensionless parameters (three gauge couplings, nine
charged-fermion masses, four CKM-matrix entries,
neutrino-sector parameters, Higgs mass and
self-coupling, QCD vacuum angle) are the contingent
moduli coordinates; predicted phenomenology descends
through the renormalised-parameter point
\citep{Workman2022Review}.

\item[\emph{Instruction set architecture versus
micro-architecture (computer architecture).}]
The carrier is the space of computer-system
implementations; the lawfulness predicate requires
correct execution of the declared instruction set
architecture (ISA); the equivalence is observational
equivalence at the ISA boundary (same program, same
architectural-state update); the moduli quotient sends
an implementation to its ISA-conformant behaviour
class.  The ISA contract (instruction encodings,
register semantics, addressing modes, observable side
effects) is a necessary readout shared by all
conforming implementations.  The micro-architecture
(pipeline depth, cache hierarchy, branch-prediction
strategy, clock frequency) is contingent.  The
System/360 family, which fixed an ISA while shipping
many distinct micro-architectures, is the founding
instance \citep{AmdahlBlaauwBrooks1964System360}.

\item[\emph{Evolutionary contingency in macroevolution
(paleobiology).}]
The carrier is the space of possible biotic
trajectories on Earth; the lawfulness predicate
requires consistency with natural selection,
developmental constraints, and the physical
environment; the equivalence identifies trajectories
with identical sequences of macroevolutionary
bottleneck outcomes; the moduli quotient sends a
trajectory to its bottleneck-survival-pattern class.
Selection-driven adaptation, the existence of
developmental constraints, and thermodynamically
consistent matter/energy flow are necessary readouts.
The specific assemblage of surviving phyla on Earth
(e.g., the contingent survival of chordates after the
Cambrian extinction events) is a contingent readout
that varies across moduli classes; ``replaying the
tape of life'' would traverse a different moduli class
\citep{Gould1989WonderfulLife}.
\end{description}

\subsection{Conservation as Orbit Descent}\label{sec:status-records-coherence:conservation}

\begin{theorem}[Conservation as Orbit Descent]\label{thm:F27}
A quantity \(C:X\to V\) is conserved for a declared route or symmetry
structure if and only if it descends through the orbit quotient; that
is, \(C(x)=C(x')\) whenever \(x\) and \(x'\) lie in the same declared
orbit.
\end{theorem}
\footnote{\raggedright Verified in Lean as
\texttt{StatusRecordsCoherence.ConservationAsOrbitDescent.\allowbreak{}conservation\_as\_orbit\_descent};
see \Cref{app:formalization}.}
\begin{proof}
If \(C\) descends through the orbit quotient, it is constant on orbit
classes by definition.  Conversely, if \(C\) is constant on orbits,
define the quotient readout by \(\bar C([x])=C(x)\).  Orbit constancy
makes the definition independent of representative, and
\(C=\bar C\circ q_{\mathrm{orb}}\).  Thus conservation is orbit
descent.  Lawhood then requires that this descended readout be
declared at the current access layer.
\end{proof}

The theorem does not enumerate conserved quantities and does not
apply to undeclared routes or symmetries.  The classical
variational-symmetry instance of this orbit-descent reading is
Noether's theorem \citep{Noether1918Invariante}.

\subsection*{Layer instantiations}\label{layer-inst:F27}

\begin{description}
\item[\emph{Lavoisier conservation of mass in chemical
reactions (chemistry).}]
The carrier is the space of chemical configurations in
a closed reaction vessel; the declared route structure
is the set of elementary chemical reactions that
preserve atomic identity; the orbit quotient sends a
configuration to its multiset of atoms.  Total mass
\(C(x)=\sum_e n_e(x)\,m_e\) is constant across orbits
of chemical rearrangement and so descends through the
orbit quotient.  Lavoisier's stoichiometric balancing
of chemical equations is the founding expression of
this orbit-descent reading \citep{Lavoisier1789Traite}.

\item[\emph{Knot invariants under Reidemeister moves
(combinatorial knot theory).}]
The carrier is the set of planar knot diagrams; the
declared route structure is the Reidemeister moves (R1
twist, R2 poke, R3 slide) together with planar isotopy;
the orbit quotient sends a diagram to its
ambient-isotopy class.  Polynomial invariants such as
the Jones or Alexander polynomial are unchanged under
every Reidemeister move and so descend through the
orbit quotient to functions on knots.  Reidemeister's
theorem certifies that ambient isotopy of knots is
exactly generated by the three planar moves
\citep{Reidemeister1927Knotentheorie}.

\item[\emph{Subject reduction in operational semantics
(programming-language theory).}]
The carrier is the set of typable program expressions
under a fixed typing context \(\Gamma\) and type
system; the declared route structure is the small-step
reduction relation \(\to\) with its reachability
preorder.  Subject reduction (type preservation) is the
conservation statement: if
\(\Gamma\vdash e:\tau\) and \(e\to e'\) then
\(\Gamma\vdash e':\tau\).  The typing judgement
\(\Gamma\vdash\cdot:\tau\) is therefore constant along
each reduction trajectory and descends through the
reachability orbit.  Together with progress, it yields
the Wright-Felleisen formulation of type soundness
\citep{WrightFelleisen1994SyntacticTypeSoundness}.

\item[\emph{Meselson-Stahl semi-conservative {DNA}
replication (molecular biology).}]
The carrier is the population of DNA duplexes in a
dividing cell; the declared route structure is one
round of replication, in which each parental duplex
yields two daughter duplexes that each contain one
parental strand and one newly synthesised strand; the
orbit quotient sends a daughter duplex to its parental
template-strand identity.  The conserved quantity is
the template-strand sequence content carried by each
daughter molecule, which is constant across the
replication orbit (modulo replication-error
corrections).  The Meselson-Stahl density-gradient
experiment is the canonical demonstration,
distinguishing semi-conservative replication from
conservative and dispersive alternatives
\citep{MeselsonStahl1958Replication}.

\item[\emph{Martingale property under conditional
expectation (probability theory).}]
The carrier is the path space of a stochastic process
\((M_t)_{t\ge 0}\) adapted to a filtration
\((\mathcal F_t)\); the declared route structure is
filtration extension \(\mathcal F_s\subseteq\mathcal F_t\)
together with the conditional-expectation operator;
the orbit quotient sends a path to its
\(\mathcal F_s\)-equivalence class.  The martingale
property
\(\mathbb E[M_t\mid\mathcal F_s]=M_s\) almost surely
is the conservation statement: the conditional
expectation \(C=\mathbb E[M_t\mid\mathcal F_s]\) is
constant along the filtration orbit and descends to
the \(\mathcal F_s\)-quotient.  Optional stopping,
Doob's maximal inequalities, and martingale
convergence follow from this orbit-descent property
\citep{Doob1953Stochastic}.
\end{description}

\subsection{Composability}\label{sec:status-records-coherence:composability}

\begin{theorem}[Composability]\label{thm:F28}
Let \(X_A,X_B\) be component spaces with quotients
\(q_A:X_A\to Q_A\), \(q_B:X_B\to Q_B\), interface readouts
\(\beta_A:X_A\to I_A\), \(\beta_B:X_B\to I_B\), composition
\(U:X_A\times X_B\to X_{AB}\), composite quotient
\(q_{AB}:X_{AB}\to Q_{AB}\), admissibility predicate
\(\operatorname{Adm}:X_A\times X_B\to\mathrm{Prop}\), and residual
\(\delta:X_A\times X_B\to D\).  Define
\begin{align*}
  \mathcal O_U
  &:=
  \{((a,b),(a',b')) :
     q_A(a)=q_A(a'),\ q_B(b)=q_B(b'),\\
  &\hspace{7em}
     q_{AB}(U(a,b))\ne q_{AB}(U(a',b'))\},\\
  \mathcal O_{\delta}
  &:=
  \{((a,b),(a',b')) :
     q_A(a)=q_A(a'),\ q_B(b)=q_B(b'),\
     \delta(a,b)\ne\delta(a',b')\}.
\end{align*}
Write \(\operatorname{COD}\) and \(\operatorname{CRD}\) for
composite-outcome descent and cross-residual descent.  Then
\begin{align*}
  \operatorname{COD}
  &\iff
  \mathcal O_U=\emptyset
  \iff
  \exists\,\bar U:Q_A\times Q_B\to Q_{AB},\
    \bar U(q_Aa,q_Bb)=q_{AB}(U(a,b)),\\
  \operatorname{CRD}
  &\iff
  \mathcal O_{\delta}=\emptyset,\\
  \operatorname{Composable}
  &\iff
  \begin{gathered}
    \operatorname{Adm}\ \text{descends}\wedge
    \operatorname{COD}\\
    {}\wedge\operatorname{CRD}
    \wedge\operatorname{Assoc}\wedge\operatorname{Id}.
  \end{gathered}
\end{align*}
\end{theorem}
\footnote{\raggedright Verified in Lean as
\texttt{StatusRecordsCoherence.Composability.composability}; see
\Cref{app:formalization}.}
\begin{proof}
Assume the component quotients \(q_A,q_B\), interface readouts
\(\beta_A,\beta_B\), composition \(U\), composite quotient
\(q_{AB}\), admissibility predicate \(\operatorname{Adm}\), and
residual \(\delta\).  The composite outcome descends exactly when it
is independent of representatives of the \(Q_A\)- and \(Q_B\)-classes.
Thus, if \(\mathcal O_U=\emptyset\), the definition
\[
  \bar U(q_Aa,q_Bb):=q_{AB}(U(a,b))
\]
is well-defined; if \(\mathcal O_U\ne\emptyset\), a displayed pair in
\(\mathcal O_U\) is precisely a split-pair obstruction.

The same argument applied to \(\delta\) gives cross-residual descent.
There are two cases.  If admissibility, outcome descent, residual
descent, associativity, and identity all hold, the chained
\(\BirdInt\) audits compose and \Cref{thm:F20} forms the composite
record.  If any one of these conditions fails, either the composite is
not well-defined on quotient classes or the audit chain is invalid, so
the records do not compose.
\end{proof}

The theorem does not enumerate composable records and does not apply
to undeclared compositions.

\subsection*{Layer instantiations}\label{layer-inst:F28}

\begin{description}
\item[\emph{Universally composable cryptographic
protocols (theoretical computer science).}]
The carriers are implementations of protocol components;
the source quotients collapse implementations to their
UC ideal-functionality class; the interface readouts
record input-output behaviour visible to the
environment; the composition combines components by
message-passing; the composite quotient collapses the
composite to its ideal-functionality class.  The UC
composition theorem yields composite-outcome descent;
the residual records simulator output and descends
through the same component quotient.  Admissibility,
associativity, and identity correspond to syntactic
well-formedness of the protocol-composition operation
\citep{Canetti2001UC}.

\item[\emph{Retrosynthetic analysis and synthon-based
total synthesis (organic chemistry).}]
The carriers are organic compounds; the source quotients
collapse compounds to their synthon class (connectivity
and functional groups); the interface readouts are
reactive sites available for bond formation; the
composition is a declared coupling reaction within a
specified reaction family of declared functional-group
tolerance.  The residual records yield, regio- and
stereo-byproduct distribution, and
reaction-mass-efficiency residual.  Composite-outcome
descent holds
when the same target connectivity is obtained from
distinct yet retrosynthetically-equivalent precursors;
cross-residual descent holds when the residual
byproduct distribution is invariant under
representative-level choices within the same reaction
family \citep{Corey1989LogicSynthesis}.

\item[\emph{Intermodal maritime containerisation
(industrial logistics).}]
The ISO 668 standard fixes the external dimensions,
corner-fitting geometry, and rated load of intermodal
freight containers.  The carriers are containers loaded
by different shippers; the source quotients collapse
them to their ISO 668-conformant external-interface
class; the interface readouts are the ISO corner
fittings and the rated gross mass; the composition is
stacking on a vessel or chaining across
truck-train-vessel handoffs; the composite quotient
classifies the consolidated load by its
external-interface footprint.  The residual captures
internal load distribution, weight class, and
damage / customs / handling residuals.  Admissibility
restricts to rated compatible cargo (weight under
ratings, stacking-strength conformance, hazmat
compatibility) \citep{ISO668_2020}.

\item[\emph{{OSI} 7-layer protocol stack composition
(network engineering).}]
The OSI Basic Reference Model partitions network
communication into seven layers, each defined by
service primitives at Service Access Points (SAPs).
The carriers are layer-\(N\) entity implementations;
the source quotients collapse implementations to their
SAP-service-primitive contract; the interface readouts
record the service-primitive sequence delivered at the
SAP; the composition stacks a layer-\((N{+}1)\) entity
on top of a layer-\(N\) service by invoking layer-\(N\)
service primitives; the composite quotient classifies
the composite by its end-to-end service.  The residual
records PDU header overhead, fragmentation, and
retransmission residual.  Composite-outcome descent
holds because the layer-\((N{+}1)\) service depends
only on layer-\(N\) service primitives, not on the
specific implementation of layer \(N\)
\citep{ISOIEC7498OSI}.

\item[\emph{Golden Gate modular {DNA} cloning
(synthetic biology).}]
Golden Gate cloning uses Type IIS restriction enzymes
that cut DNA outside their recognition site, leaving
scarless four-base overhangs designed arbitrarily.
The carriers are donor plasmids carrying DNA parts
flanked by Type IIS recognition sequences; the source
quotients collapse plasmids to their part-class plus
overhang specification; the interface readouts are the
four-base overhang fingerprints; the composition is the
one-pot enzymatic digest-and-ligate reaction; the
composite quotient classifies the assembled multi-part
construct.  The residual records misligation products,
unligated linearised backbone, and off-target
re-religation events.  Composite-outcome descent holds
because overhang complementarity determines assembly
order independent of donor-plasmid backbone
\citep{Engler2008GoldenGate}.
\end{description}

\subsection{Presentation Invariance}\label{sec:status-records-coherence:presentation-invariance}

\begin{theorem}[Presentation Invariance]\label{thm:F29}
Let \(P\) be a presentation space with quotient
\(\pi:P\to\Pi\), let \(\sim\) be an equivalence relation compatible
with \(\pi\), let \(s:\Pi\to P\) be a section with
\(\pi(s(x))=x\), and let \(T:P\to V\) be a readout.  Let
\(\tau:\operatorname{Tuple}(P)\to\operatorname{Tuple}(\Pi)\) be the
tuple quotient, let \(R:\operatorname{Tuple}(P)\to\mathrm{Prop}\) be a
relation, and let a group \(G\) act on \(P\) and \(V\).  Then
\begin{align*}
  \operatorname{PresentationInvariant}(\pi,T)
  &\iff
  \forall p,p'\in P,\ \pi(p)=\pi(p')\Rightarrow T(p)=T(p')\\
  &\iff
  \exists\,\bar T:\Pi\to V\ \text{with}\ \bar T\circ\pi=T,\\
  \operatorname{EquivalentPresentationInvariant}(\sim,T)
  &\iff
  \forall p,p'\in P,\ p\sim p'\Rightarrow T(p)=T(p'),\\
  \operatorname{PresentationRelationInvariant}(\tau,R)
  &\iff
  \forall x,y,\ \tau(x)=\tau(y)\Rightarrow (R(x)\leftrightarrow R(y)),\\
  \operatorname{GroupEquivariant}(T)
  &\iff
  \forall g\in G,\forall p\in P,\ T(g\cdot p)=g\cdot T(p).
\end{align*}
\end{theorem}
\footnote{\raggedright Verified in Lean as
\texttt{StatusRecordsCoherence.PresentationInvariance.\allowbreak{}presentation\_invariance};
see \Cref{app:formalization}.}
\begin{proof}
Assume the quotient \(\pi\), compatible equivalence \(\sim\), section
\(s\), readout \(T\), tuple quotient \(\tau\), relation \(R\), and
group action of \(G\).  If \(T\) is presentation-invariant, it is
constant on \(\pi\)-fibers.  Conversely, if \(T\) is constant on
\(\pi\)-fibers, define
\[
  \bar T(x):=T(s(x)).
\]
Since \(\pi(s(x))=x\), this satisfies \(\bar T\circ\pi=T\), and
compatibility of \(\sim\) with \(\pi\) gives the equivalent
presentation-invariance condition.

The relation statement is the same descent criterion for the tuple
quotient \(\tau\).  The group-action statement is a separate
intertwining condition:
\[
  T(g\cdot p)=g\cdot T(p).
\]
When all four displayed conditions hold, readouts, relations, and
audit data transport through the presentation quotient.  Hence the
formed records are invariant under the presentation change.
\end{proof}

The theorem does not enumerate presentations and does not collapse
the presentation/object distinction.

\subsection*{Layer instantiations}\label{layer-inst:F29}

\begin{description}
\item[\emph{Coordinate-free formulation of general
relativity (mathematical physics).}]
The presentation space is the set of
coordinate-component representations of a geometric
tensor over a spacetime manifold; the quotient
collapses representations to the underlying
coordinate-free tensor field.  The diffeomorphism
group acts on representations by pushforward; the
chart components transform equivariantly
\(T(\phi\cdot p)=\phi_*T(p)\), while the tensor field
is the presentation-invariant readout.  General
covariance is the requirement that physically
meaningful readouts be presentation-invariant in this
sense \citep{MisnerThorneWheeler1973Gravitation}.

\item[\emph{Gr\"obner basis as canonical
polynomial-ideal presentation (computer algebra).}]
For a polynomial ideal \(I\subseteq k[x_1,\dots,x_n]\),
the presentation space is the set of generating sets
for \(I\); the quotient sends a generating set to the
ideal it generates.  Under a fixed monomial ordering,
the reduced Gr\"obner basis supplies a canonical
section: every ideal has a unique reduced Gr\"obner
basis.  Ideal membership, ideal equality, and ideal
intersection are presentation-invariant readouts
computed by normal-form reduction; Buchberger's
algorithm constructs the section explicitly
\citep{Buchberger1965ThesisGroebner}.

\item[\emph{Jordan canonical form and matrix similarity
(linear algebra).}]
For a finite-dimensional vector space \(V\) over an
algebraically closed field \(k\), the presentation
space is the set of matrix representatives of linear
endomorphisms; the quotient sends a matrix to its
similarity class.  Under a fixed eigenvalue and
block ordering, the Jordan canonical form supplies a
canonical section.  Trace, determinant, characteristic
polynomial, minimal polynomial, eigenvalue multiset
with algebraic and geometric multiplicities, rank, and
nullity are presentation-invariant readouts.  The
general linear group \(\mathrm{GL}_n(k)\) acts on
matrices by similarity; the listed invariants are the
\(\mathrm{GL}_n(k)\)-invariant readouts on the
presentation space \citep{HoffmanKunze1971LinearAlgebra}.

\item[\emph{De Bruijn indices for alpha-equivalence
(programming-language theory).}]
The presentation space is the set of lambda terms with
named bound variables; the quotient sends a term to its
alpha-equivalence class.  The de Bruijn-index encoding
gives a canonical nameless representative (each bound
variable replaced by the distance to its binder).  Two
terms have the same de Bruijn encoding iff they are
alpha-equivalent.  Beta-reduction, normal-form
computation, and operational-semantics evaluation are
presentation-invariant readouts; the de Bruijn section
solves the variable-capture problem by construction
\citep{deBruijn1972LambdaNotation}.

\item[\emph{Lewin transformational music theory and the
chromatic group (music theory).}]
The presentation space is the set of pitch-class
assignments to a musical motif; the chromatic group
\(G=\mathbb Z/12\mathbb Z\) acts by transposition, and
the dihedral group \(D_{12}\) extends this with
inversion.  Motif features transform equivariantly:
\(T(g\cdot p)=g\cdot T(p)\).  Pitch-class invariants
such as the interval-class content (a canonical
\(G\)-invariant readout) descend through the orbit
quotient of the \(D_{12}\)-action and are
presentation-invariant musical features
\citep{Lewin1987GMIT}.
\end{description}

\subsection{Explanation as Compression with Descent}\label{sec:status-records-coherence:explanation}

\begin{theorem}[Explanation as Compression with Descent]\label{thm:F30}
Let \(H\) be a history space with base readout \(b:H\to B\),
explanation readout \(e:H\to E\), target readout \(t:H\to T\),
continuation \(\kappa:H\to H'\), later explanation
\(e':H'\to E'\), and later target \(t':H'\to T'\).  Assume \(b\)
and \(e\) are surjective.  Define
\[
  \mathcal O_{\mathrm{src}}
  :=
  \{(h,h')\in H^2:b(h)=b(h')\ \text{and}\ e(h)\ne e(h')\},
\]
and
\[
  \mathcal O_{\mathrm{tgt}}
  :=
  \{(h,h')\in H^2:e(h)=e(h')\ \text{and}\ t(h)\ne t(h')\}.
\]
Then
\begin{align*}
  \operatorname{ExactExplanation}(b,e,t)
  &\iff
  \mathcal O_{\mathrm{src}}=\emptyset
  \wedge
  \mathcal O_{\mathrm{tgt}}=\emptyset\\
  &\iff
  \exists\,\alpha:B\to E,\ \exists\,\beta:E\to T,\
    \alpha\circ b=e\ \text{and}\ \beta\circ e=t,\\
  \operatorname{StrictCompression}(b,e)
  &\iff
  \exists\,\alpha:B\to E,\ \alpha\circ b=e
  \ \text{and}\ 
  \neg\exists\,\delta:E\to B,\ \delta\circ e=b.
\end{align*}
\end{theorem}
\footnote{\raggedright Verified in Lean as
\texttt{StatusRecordsCoherence.\allowbreak{}ExplanationAsCompression\allowbreak{}WithDescent.\allowbreak{}explanation\_as\_compression\_\allowbreak{}with\_descent};
see \Cref{app:formalization}.}
\begin{proof}
Assume the base readout \(b\), explanation readout \(e\), target
readout \(t\), and the surjectivity of \(b\) and \(e\).  If an exact
explanation exists, then \(e\) is determined by \(b\) and \(t\) is
determined by \(e\), so both obstruction sets are empty.  Conversely,
if \(\mathcal O_{\mathrm{src}}=\emptyset\), define
\[
  \alpha(b(h)):=e(h),
\]
which is independent of \(h\).  If
\(\mathcal O_{\mathrm{tgt}}=\emptyset\), define
\[
  \beta(e(h)):=t(h),
\]
again independently of representative.  These definitions give
\(\alpha\circ b=e\) and \(\beta\circ e=t\).

For strict compression, there are two cases.  If such an
\(\alpha\) exists but no \(\delta:E\to B\) satisfies
\(\delta\circ e=b\), then the explanation coarsens the base without
being equivalent to it.  If a \(\delta\) also exists, then \(b\) and
\(e\) are mutually recoverable and the compression is not strict.
\end{proof}

The theorem does not rank explanations and does not identify
explanation with the fact explained.  The Kolmogorov-complexity and
minimum-description-length traditions develop the
shortest-program / shortest-code reading of explanation in detail
\citep{LiVitanyi2008Kolmogorov}.

\subsection*{Layer instantiations}\label{layer-inst:F30}

\begin{description}
\item[\emph{Newtonian mechanics explaining
Kepler's laws (mathematical physics).}]
The history space is the set of solar-system
trajectories; the base records observed positions and
velocities.  The explanation is Newton's gravitational
dynamics: masses of the bodies, the inverse-square
gravitational law, and initial conditions.  The target
is the set of Kepler-type empirical regularities
(elliptical orbits, equal areas in equal times,
\(T^2\propto a^3\)) and the ephemeris of future
planetary positions.  The descended map identifies
masses and initial conditions from kinematics, and
the target map solves the two-body problem to derive
Kepler's laws; the compression is strict because
continuous trajectories are summarised by a finite
number of parameters \citep{Newton1687Principia}.

\item[\emph{Mendeleev periodic table explaining
chemical regularities (chemistry).}]
The history space is the set of chemical elements
with their observed properties (atomic weight,
valence, melting point, oxidation states, reaction
products).  Mendeleev's 1869 table is the
explanation: each element is assigned a
(period, group) position based on atomic weight and
chemical-similarity grouping (later refined by atomic
number).  The descended map sends a property profile
to its table position; the target map predicts the
existence and properties of then-unknown elements
(gallium, germanium, scandium were subsequently
discovered with properties anticipated from table
gaps).  The compression is strict: the table position
summarises chemical-family behaviour without
recovering all numerical property values
\citep{Mendeleev1869Periodic}.

\item[\emph{Principal component analysis as
dimensionality-reducing explanation (statistics).}]
The history space is observed records
\(x\in\mathbb R^p\); the base is the full record.
The explanation is the top-\(k\) principal-component
representation, computed as eigenvectors of the
empirical covariance matrix.  The target is the
low-rank reconstruction or rank-\(k\) orthogonal
projection visualisation, both exactly determined by
the principal-component representation through a
linear map.  The compression is strict whenever
\(k<p\) and the discarded components carry non-zero
variance; the explained-variance ratio measures the
compression fidelity \citep{Pearson1901PCA}.

\item[\emph{Boltzmann {H}-theorem and
statistical-mechanical irreversibility (statistical
physics).}]
The history space is the microscopic phase-space
trajectory set of a dilute gas; the base records the
full microstate.  The explanation is the one-particle
distribution function \(f(\mathbf r,\mathbf v,t)\) and
the H-functional
\(H[f]=\int f\log f\,d^3v\,d^3r\).  Under the
molecular-chaos hypothesis, \(H\) is non-increasing
along Boltzmann-equation evolution; identifying
macroscopic entropy with \(S=-k_B H\), the target is
the macroscopic thermodynamic readout.  The
compression is strict: distinct microstates with the
same one-particle distribution share the same
macroscopic readout \citep{Boltzmann1872Weitere}.

\item[\emph{Hart's \emph{The Concept of Law} as
jurisprudential summary (jurisprudence).}]
The history space is the recorded legal practice of a
jurisdiction (statutes, decided cases, official
enactments); the base records the full legal corpus.
Hart's explanation distinguishes primary rules of
obligation from secondary rules of recognition,
change, and adjudication, together with the
identified rule of recognition.  The target, in
Hart's construction, is the legal-validity verdict on
candidate rules: a candidate counts as valid law iff
it satisfies the rule of recognition.  The target map
is exact because the rule of recognition is itself
the criterion of validity; the compression is strict
because many distinct legal corpora share the same
rule-system structure \citep{Hart1961Concept}.
\end{description}

\subsection{Counterfactual Legitimacy}\label{sec:status-records-coherence:counterfactual}

\begin{theorem}[Counterfactual Legitimacy]\label{thm:F31}
Let \(q:H\to Q\) be a context quotient, let
\(\alpha:A\to\bar A\) be an intervention abstraction, let
\(\operatorname{Adm}:H\times A\to\mathrm{Prop}\) be admissibility, let
\(\Gamma:H\times A\to H'\) be a counterfactual route, let
\(q':H'\to R\) be a target quotient, and let
\(\tau:H'\to\mathrm{Bool}\) be target truth.  Define
\[
  \mathcal O_{\mathrm{adm}}
  :=
  \{((h,a),(h',a')) :
    q(h)=q(h'),\ \alpha(a)=\alpha(a'),\
    \operatorname{Adm}(h,a)\not\leftrightarrow
    \operatorname{Adm}(h',a')\},
\]
and
\[
  \mathcal O_{\mathrm{cf}}
  :=
  \left\{
  \begin{array}{l}
    ((h,a),(h',a')) :
    q(h)=q(h'),\ \alpha(a)=\alpha(a'),\\
    \operatorname{Adm}(h,a),\ \operatorname{Adm}(h',a'),\\
    \tau(\Gamma(h,a))\ne\tau(\Gamma(h',a'))
  \end{array}
  \right\}.
\]
Write \(\operatorname{CFL}\) for the counterfactual-legitimacy
predicate for this data.  Then
\begin{align*}
  \operatorname{CFL}(q,\alpha,\operatorname{Adm},\Gamma,q',\tau)
  &\iff
  \mathcal O_{\mathrm{adm}}=\emptyset
  \wedge
  \mathcal O_{\mathrm{cf}}=\emptyset\\
  &\iff
  \begin{gathered}
    \exists\,\bar\tau:Q\times\bar A\to\mathrm{Bool}
    \ \text{such that}\\
    \bar\tau(q(h),\alpha(a))=\tau(\Gamma(h,a))
    \ \text{for all admissible }(h,a).
  \end{gathered}
\end{align*}
\end{theorem}
\footnote{\raggedright Verified in Lean as
\texttt{StatusRecordsCoherence.CounterfactualLegitimacy.\allowbreak{}counterfactual\_legitimacy};
see \Cref{app:formalization}.}
\begin{proof}
Assume the context quotient \(q\), intervention abstraction
\(\alpha\), admissibility predicate \(\operatorname{Adm}\), route
\(\Gamma\), target quotient \(q'\), and target truth \(\tau\).  If a
counterfactual is legitimate, admissibility is constant on
\((q,\alpha)\)-classes, so
\(\mathcal O_{\mathrm{adm}}=\emptyset\).  Its target truth is also
constant on admissible \((q,\alpha)\)-classes, so
\(\mathcal O_{\mathrm{cf}}=\emptyset\).

Conversely, assume both obstruction sets are empty.  Define, for
admissible pairs,
\[
  \bar\tau(q(h),\alpha(a)):=\tau(\Gamma(h,a)).
\]
The emptiness of \(\mathcal O_{\mathrm{adm}}\) makes admissibility a
property of \((q(h),\alpha(a))\), and the emptiness of
\(\mathcal O_{\mathrm{cf}}\) makes \(\bar\tau\) independent of the
admissible representative.  There are two cases.  If both
obstructions are empty, the counterfactual descends to the declared
context/intervention orbit and is legitimate.  If either obstruction
is nonempty, the claim uses undeclared admissibility variation or an
undescended target truth, hence it is an overread.
\end{proof}

The theorem does not determine the truth value of counterfactuals and
does not apply to arbitrary undeclared possible worlds.  The
possible-worlds semantics in the background is the classical reference
\citep{Lewis1973Counterfactuals}.

\subsection*{Layer instantiations}\label{layer-inst:F31}

\begin{description}
\item[\emph{{MIL-STD-810} environmental-qualification
bench testing (engineering qualification).}]
US Department of Defense MIL-STD-810 prescribes a
catalogue of environmental qualification test methods
(thermal cycling, vibration, shock, humidity, salt fog,
etc.).  The context quotient collapses specimens to
their design baseline (model number, build
configuration); the intervention abstraction collapses
test profiles to the standardised method class;
admissibility requires the profile to lie within the
standard envelope; the counterfactual route is the
laboratory test execution under the prescribed
profile; the target truth records pass/fail according
to the acceptance criterion.  The legitimacy of a
qualification-test counterfactual --- ``this specimen
would pass under method X'' --- is the F31 statement:
admissibility and pass/fail descend through
(design-baseline, test-method) classes
\citep{MILSTD810H2019}.

\item[\emph{Fogel counterfactual economic history of US
growth without railroads (cliometrics).}]
The context quotient collapses economic trajectories
to their factor-endowment class (capital, labour,
natural resources); the intervention abstraction
collapses transportation mixes to (railroad available,
wagon and canal substitute) categories; the route is
the social-savings simulation that re-prices freight
using non-rail substitutes for each route; the target
truth records whether the counterfactual 1890 GNP
exceeds a specified threshold (Fogel's calculation
gives a deviation of less than 5\% of actual).  The
legitimacy of the counterfactual is the F31 statement:
the social-savings estimate descends through
factor-endowment and transportation-mix classes
\citep{Fogel1964RailroadsAmericanGrowth}.

\item[\emph{Legal ``but-for'' causation in tort law
(jurisprudence).}]
In common-law tort doctrine, but-for causation asks
whether the plaintiff's harm would have occurred but
for the defendant's act.  The context quotient
collapses factual histories to their legally cognisable
class; the intervention abstraction collapses
candidate acts to their legal characterisation;
admissibility encodes proximate-cause limits; the
route is the counterfactual reconstruction of ``what
would have happened'' without the defendant's act; the
target truth records whether the harm would still have
occurred.  The legitimacy of the but-for
counterfactual is the F31 statement: the verdict
descends through legally cognisable
(context, act) classes, and overreaches (e.g.,
multiple-sufficient-cause cases where but-for fails)
trigger
doctrinal corrections such as the substantial-factor
test \citep{HartHonore1985Causation}.

\item[\emph{Off-policy evaluation in reinforcement
learning (machine learning).}]
An agent estimates the expected return of a target
policy \(\pi_e\) using trajectories collected under a
different behaviour policy \(\pi_b\).  The context
quotient collapses environments to their (state-action
transition kernel, reward function) class; the
intervention abstraction collapses target policies to
their induced visitation; admissibility requires the
support condition \(\pi_b>0\) wherever \(\pi_e>0\);
the route is the importance-sampling re-weighting; the
target truth records whether the estimated return of
\(\pi_e\) exceeds a specified threshold.  The
Precup-Sutton-Singh result gives conditions under
which the off-policy estimator descends from
behaviour-policy trajectories to a legitimate
counterfactual verdict on the target policy's value
\citep{PrecupSuttonSingh2000Eligibility}.

\item[\emph{Climate-model counterfactual emission
scenarios (climate science).}]
In atmospheric general-circulation modelling, the
context quotient collapses climate trajectories to
their baseline-forcing class (pre-industrial CO\(_2\),
solar and orbital parameters); the intervention
abstraction collapses scenarios to an idealised
greenhouse-gas forcing class (e.g., doubled
atmospheric CO\(_2\) concentration); the route is the
numerical integration of the coupled atmosphere
model under the chosen forcing; the target truth
records whether a specified climate diagnostic (e.g.,
equilibrium global-mean surface temperature change)
exceeds a threshold.  Manabe \& Wetherald 1967 is the
founding instance, with its calculation of equilibrium
climate sensitivity to doubled CO\(_2\); legitimacy
is the F31 statement that the threshold-crossing
verdict descends through (baseline, idealised forcing)
classes \citep{ManabeWetherald1967Thermal}.
\end{description}

\subsection*{Status, Records, Coherence at a glance}\label{sec:status-records-coherence:summary}
\begin{table}[p]
  \centering
  \footnotesize
  \setlength{\tabcolsep}{4pt}
  \textbf{Status}
  \smallskip

  \begin{tabularx}{\textwidth}{@{}
      >{\raggedright\arraybackslash}p{0.24\textwidth}
      >{\raggedright\arraybackslash}X
      >{\raggedright\arraybackslash}p{0.16\textwidth}
      >{\raggedright\arraybackslash}p{0.10\textwidth}@{}}
    \toprule
    Paper-prose name & One-line claim & Status & Substrate cite \\
    \midrule
    Objectivity Theorem &
    Shared objects are common access quotients across declared
    presentations. &
    Theorem &
    \textemdash \\
    Lawhood-Descent Theorem &
    A candidate law is lawful exactly when it factors through the
    available access quotient. &
    Theorem &
    \textemdash \\
    Interface-Mediation Theorem &
    Interfaces mediate closures iff transported observables descend
    through the interface access quotient. &
    Theorem &
    \textemdash \\
    Probability-as-Stable-Fiber Theorem &
    Probability is the stable quotient record of an unresolved fiber. &
    Theorem &
    \textemdash \\
    \bottomrule
  \end{tabularx}

  \medskip
  \textbf{Records}
  \smallskip

  \begin{tabularx}{\textwidth}{@{}
      >{\raggedright\arraybackslash}p{0.24\textwidth}
      >{\raggedright\arraybackslash}X
      >{\raggedright\arraybackslash}p{0.16\textwidth}
      >{\raggedright\arraybackslash}p{0.10\textwidth}@{}}
    \toprule
    Paper-prose name & One-line claim & Status & Substrate cite \\
    \midrule
    Object-Persistence Theorem &
    An object persists under transport iff its identity predicate
    descends through the transport. &
    Theorem &
    \textemdash \\
    Fact--Record Theorem &
    A fact forms iff the claim is closure-of-record and its evidence
    descends through access. &
    Theorem &
    \textemdash \\
    Reflexive-Nonclosure Theorem &
    Same-layer self-audit cannot close without growing an audit layer. &
    Theorem &
    \textemdash \\
    Causality Theorem &
    Causality is stable target residue of lawful intervention
    transports. &
    Theorem &
    \textemdash \\
    \bottomrule
  \end{tabularx}

  \medskip
  \textbf{Coherence}
  \smallskip

  \begin{tabularx}{\textwidth}{@{}
      >{\raggedright\arraybackslash}p{0.24\textwidth}
      >{\raggedright\arraybackslash}X
      >{\raggedright\arraybackslash}p{0.16\textwidth}
      >{\raggedright\arraybackslash}p{0.10\textwidth}@{}}
    \toprule
    Paper-prose name & One-line claim & Status & Substrate cite \\
    \midrule
    Moduli--Contingency Theorem &
    Stable parameters descend through perturbation orbits; variation
    across lawful closures is contingency. &
    Theorem &
    \textemdash \\
    Conservation Theorem &
    Conserved quantities are exactly readouts that factor through
    route-orbit quotients. &
    Theorem &
    \textemdash \\
    Composability Theorem &
    Formed records compose exactly when their audit data chains. &
    Theorem &
    \textemdash \\
    Presentation-Invariance Theorem &
    Presentations yield the same records iff access and audit transport
    agree. &
    Theorem &
    \textemdash \\
    Explanation-as-Compression Theorem &
    Explanations are audited compressions that preserve source and
    target descent. &
    Theorem &
    \textemdash \\
    Counterfactual-Legitimacy Theorem &
    Counterfactuals are legitimate only for declared intervention
    routes with descended target truth. &
    Theorem &
    \textemdash \\
    \bottomrule
  \end{tabularx}

  \caption{Status, Records, Coherence at a glance: the fourteen laws of \Cref{sec:status-records-coherence}, grouped by sub-axis.}
  \label{tab:axis-status-records-coherence-summary}
\end{table}
  \section{Axis VI — Self-Reference and Limits}\label{sec:self-reference-limits}

The Self-Reference and Limits axis collects the laws that mark where
closure formation reaches a boundary.  Some boundaries are horizons:
finite data can be sufficient for a closure decision even when the
ambient carrier is larger.  Others concern information loss,
singular limits, locality, topology, and completeness: the apparatus
must say what is visible, what is glued globally, what survives a
limit, and what cannot be established from inside the same closure
without additional audit structure.

The eighteen laws proceed from local sufficiency to meta-coherence.
Horizon Sufficiency / Holographic Compression is the lead theorem:
finite horizon data can suffice for closure decisions.  It is followed by
Information Loss Normal Form, Scale Descent / Renormalization,
Singular Limit Completion, Complementarity as Non-Joint Access, Arrow
from Record Monotonicity, and Entropy as Fiber Volume.  The middle of
the chapter treats Irreducibility / No Short Closed Description,
Continuum Emergence, Criticality / Phase Boundary, Locality / Domain
of Dependence, Law--State Split, and Fine-Tuning as Small Selector
Region.  The final stretch gives Topology as Global Gluing Invariant,
Common-Source / Nonlocal Correlation Normal Form, Vacuum / Background
Selection, Unification as Common Refinement, and Closure Completeness
Boundary.  The last of these is the chapter's meta-coherence law and
is distinct from the Closure--Reality Boundary anchor of
\Cref{sec:anchor:closure-reality-boundary}.

All eighteen rows in this chapter are direct theorem rows.  After
\Cref{sec:apparatus}, the chapter is self-contained: it reuses the
same quotient, descent, interaction-status, and split-pair vocabulary
fixed earlier.  The at-a-glance table at the end groups the laws into
limits, emergence, and invariants.

\subsection{Horizon Sufficiency / Holographic Compression}\label{sec:self-reference-limits:horizon-sufficiency}

\begin{definition}[Horizon-bounded support]\label{def:horizon-bounded-support}
A target object is horizon-bounded when its declared support lies
inside a finite horizon record: a quotient of the apparatus boundary
through which every declared target readout for that object is meant
to be evaluated.
\end{definition}

\begin{theorem}[Horizon Sufficiency / Holographic Compression]\label{thm:F33}
Let \(H\) be a set of histories with an interior--exterior
decomposition \(H \cong I \times E\), and define
\[
  \pi : I \times E \longrightarrow B \times Q,
  \qquad
  \pi(i,e)=(\beta(i),\varepsilon(e)),
\]
where \(\beta:I\to B\) is the boundary readout and
\(\varepsilon:E\to Q\) is the exterior classifier. Assume \(\pi\) is
surjective, and let \(t:I\times E\to T\) be a declared target readout.
Then
\begin{align*}
  t \text{ descends through }\pi
  &\iff
  \forall (i,e),(i',e'),\
    \pi(i,e)=\pi(i',e')\Rightarrow t(i,e)=t(i',e')\\
  &\iff
  \exists\,\bar t:B\times Q\to T\ \text{ such that }\
    \bar t\circ\pi=t .
\end{align*}
For a family \((t_\sigma)_{\sigma\in\Sigma}\), horizon sufficiency is
the pointwise conjunction of these equivalent conditions for every
\(\sigma\in\Sigma\).
\end{theorem}
\footnote{\raggedright Verified in Lean as
\texttt{SelfReferenceLimits.\allowbreak{}HorizonSufficiency.\allowbreak{}horizon\_sufficiency};
see \Cref{app:formalization}.}
\begin{proof}
Assume the surjectivity of \(\pi\). If \(t=\bar t\circ\pi\), then for
any \((i,e),(i',e')\) with \(\pi(i,e)=\pi(i',e')\),
\[
  t(i,e)=\bar t(\pi(i,e))
        =\bar t(\pi(i',e'))
        =t(i',e').
\]
Conversely assume the displayed representative-independence condition.
For \((b,q)\in B\times Q\), choose \((i,e)\in\pi^{-1}(b,q)\) and set
\[
  \bar t(b,q):=t(i,e).
\]
The choice is independent of representative by the same condition, so
\(\bar t\circ\pi=t\). Equivalently, the obstruction set
\[
  \{((i,e),(i',e')):\pi(i,e)=\pi(i',e'),\ t(i,e)\ne t(i',e')\}
\]
is empty. Applying the argument to each \(t_\sigma\) gives the family
statement.
\end{proof}

The substrate-physics instance is the holographic principle, in which
the bulk information of a region is encoded on its boundary
\citep{Maldacena1997LargeN,Bousso2002Holographic}.  The
layer-agnostic statement here is the descent-through-boundary form
of that pattern; it does not assert the AdS/CFT correspondence or
any particular gravitational realization.

\subsection*{Layer instantiations}\label{layer-inst:F33}

\begin{description}
\item[\emph{Boundary-value problem for a linear elliptic
{PDE} (mathematical analysis).}]
For a fixed domain \(\Omega\subseteq\mathbb R^n\), the
interior parameterises admissible boundary data
\(g\) together with the prescribed boundary-condition
type; the boundary readout records \(g\) in a fixed
function space (e.g., \(H^{1/2}(\partial\Omega)\)).
The exterior parameterises source terms and coefficient
fields; the exterior classifier records the
source-and-coefficient data in their function-space
equivalence class.  The target readout is the entire
weak solution \(u\in H^1(\Omega)\).  By the Lax-Milgram
theorem, the target descends through the product
quotient: same boundary data plus same
source-coefficient class yield the same solution
\citep{Evans2010PartialDifferential}.

\item[\emph{Boundary representation in solid {CAD}
modelling (engineering).}]
A 3D rigid solid \(B\subset\mathbb R^3\) is represented
by its boundary surface together with adjacency and
orientation data: a B-rep records faces, edges, and
vertices with incidence relations.  The interior is the
open solid; the boundary readout records the B-rep with
face-orientation conventions; the exterior classifier
records the combinatorial topology (genus, handles,
through-holes).  Geometric and topological readouts ---
volume via the divergence theorem, mass distribution,
collision detection, boolean operations --- descend
through the product quotient
\citep{Mortenson2006GeometricModeling}.

\item[\emph{{TCP/IP} encapsulation (network protocols).}]
A network message is encapsulated as a sequence of
nested headers wrapping a payload.  The interior is the
full packet state (headers and payload together); the
boundary readout records the IP and TCP header fields
(source/destination addresses, ports, sequence numbers,
flags, checksum).  The exterior classifier records the
transport-protocol family (TCP, UDP, ICMP) and the
link-layer class.  Routing, flow-control, retransmission,
and demultiplexing decisions descend through the product
quotient: two packets sharing the same header fields and
the same protocol stack class receive the same routing
verdict, regardless of payload bytes
\citep{Postel1981RFC791,Postel1981RFC793}.

\item[\emph{Fisher sufficient statistic (mathematical
statistics).}]
For an exponential-family parametric model
\(p(x\mid\theta)\), a sufficient statistic
\(T(X_1,\dots,X_n)\) summarises a sample so that the
conditional distribution of the sample given \(T\) is
independent of \(\theta\).  The interior is the raw
sample space; the boundary readout is the sufficient
statistic; the exterior classifier is the ancillary
statistic.  By the Fisher-Neyman factorization theorem,
likelihood-based inference functionals --- likelihood
ratio, score, maximum-likelihood estimate, and
likelihood-equivalent functionals --- descend through
the product quotient \citep{Fisher1922Mathematical}.

\item[\emph{Cauchy's integral formula in complex analysis
(complex analysis).}]
For a function \(f\) holomorphic on a simply connected
domain \(D\subseteq\mathbb C\) and a point \(z_0\) in
the interior of a positively oriented simple closed
contour \(\gamma\subset D\),
\[
  f(z_0)=\frac{1}{2\pi i}\oint_\gamma \frac{f(z)}{z-z_0}\,
  dz.
\]
The interior parameterises the choice of contour
\(\gamma\) together with the holomorphic-function trace
\(f|_\gamma\) on it; the boundary readout records that
trace; the exterior classifier records the holomorphy
class of \(f\).  The target value \(f(z_0)\) at any
interior point descends exactly through the product
quotient: it is given by the Cauchy kernel integral
applied to the boundary trace
\citep{Ahlfors1979ComplexAnalysis}.
\end{description}

\subsection{Information Loss Normal Form}\label{sec:self-reference-limits:information-loss}

\begin{theorem}[Information Loss Normal Form]\label{thm:F34}
Let \(H_0,H_1\) be source and later carriers with transport
\(\tau:H_0\to H_1\). Let \(\sigma:H_0\to I\) be source information,
\(\rho:H_1\to R\) a recovery quotient, \(q_0:H_0\to Q_0\) a source
quotient, \(q_1:H_1\to Q_1\) a later quotient, and
\(s:H_0\to S\) a declared signature. Assume that
\(\rho\circ\tau\), \(q_0\), and \(q_1\circ\tau\) are surjective onto
their images. Define
\[
  \mathcal O_\rho
  =\{(h,h'):\rho(\tau h)=\rho(\tau h')\ \text{and}\
      \sigma(h)\ne\sigma(h')\}.
\]
Then
\begin{align*}
  \operatorname{Recoverable}(\sigma,\tau,\rho)
  &\iff \mathcal O_\rho=\emptyset\\
  &\iff
  \exists\,\bar\sigma:R\to I\ \text{ such that }\
    \bar\sigma\circ\rho\circ\tau=\sigma .
\end{align*}
Information loss is legitimate exactly when every declared later
identification made by \(\rho\circ\tau\) is invisible to \(\sigma\).
\end{theorem}
\footnote{\raggedright Verified in Lean as
\texttt{SelfReferenceLimits.\allowbreak{}InformationLoss.\allowbreak{}information\_loss\_normal\_form};
see \Cref{app:formalization}.}
\begin{proof}
Assume the named maps and the surjectivity of \(\rho\circ\tau\). If
\(\sigma=\bar\sigma\circ\rho\circ\tau\), then
\[
  \rho(\tau h)=\rho(\tau h')
  \quad\Longrightarrow\quad
  \sigma(h)=\bar\sigma(\rho(\tau h))
           =\bar\sigma(\rho(\tau h'))=\sigma(h'),
\]
so \(\mathcal O_\rho=\emptyset\). Conversely assume
\(\mathcal O_\rho=\emptyset\). For \(r\in R\) in the image of
\(\rho\circ\tau\), choose \(h\) with \(\rho(\tau h)=r\) and define
\[
  \bar\sigma(r):=\sigma(h).
\]
If \(h'\) is another choice, then \((h,h')\notin\mathcal O_\rho\), so
\(\sigma(h)=\sigma(h')\); hence \(\bar\sigma\) is well-defined and
\(\bar\sigma\circ\rho\circ\tau=\sigma\). If
\(\mathcal O_\rho\ne\emptyset\), some pair \(h,h'\) has the same
recovered class but different source information, which is precisely
loss across an access class.
\end{proof}

The famous physical instance is the black-hole information paradox
posed in \citet{Hawking1976BHThermodynamics}; the layer-agnostic
statement above records the descent obstruction abstracted from
that substrate.

\subsection*{Layer instantiations}\label{layer-inst:F34}

\begin{description}
\item[\emph{Huffman entropy coding (information theory).}]
The source carrier is the set of source-symbol sequences
over an alphabet \(\Sigma\); the source information is
the original symbol sequence.  The transport encodes
into a prefix-free binary codeword sequence; the
recovery quotient is the prefix decoder.  Distinct
source sequences produce distinct codeword streams; the
obstruction set is empty and the source is fully
recoverable.  Compression achieves the entropy bound to
within one bit per symbol
\citep{Huffman1952Method}.

\item[\emph{Reed-Solomon forward error correction
(coding theory).}]
For a Reed-Solomon \((n,k)\) code over a finite field
\(\mathbb F_q\), the source information is a
k-symbol message; the transport evaluates the
message polynomial at \(n\) distinct field points and
transmits through a channel injecting a bounded error
vector of Hamming weight at most
\(t=\lfloor(n-k)/2\rfloor\).  The later carrier is the
set of received words within the correction ball.  The
recovery quotient is the Berlekamp-Massey decoding
algorithm.  Within bounded-error capacity, the
obstruction set is empty and the source message is
recoverable; Reed-Solomon codes underlie CD/DVD, QR
codes, and deep-space communication
\citep{ReedSolomon1960Codes}.

\item[\emph{Cryptographic one-way hash function
(cryptography).}]
For a cryptographic hash function
\(H:\{0,1\}^*\to\{0,1\}^n\) (e.g., SHA-256), the
source information is the original bitstring.  The
transport maps bitstrings to fixed-length digests;
the recovery quotient is the identity on digests.
Many distinct source bitstrings collide to the same
digest, so the obstruction set is non-empty.  The
cryptographic hash is the canonical illegitimate-loss
instance, exploited to make digital signatures,
password verifiers, and Merkle trees efficient and
secure \citep{NIST2015FIPS180}.

\item[\emph{Database write-ahead logging and {ARIES}
recovery (database systems).}]
The source information is the committed-state effects
of database transactions.  The transport writes log
records (before-image, after-image, LSN, transaction
ID) to stable storage before any modified page is
written.  The recovery quotient is the ARIES
three-pass algorithm (analysis, redo, undo) applied
to the on-disk log and data pages after a crash.
Within the failure model, the obstruction set is
empty and the committed-state effects are recoverable
\citep{Mohan1992ARIES}.

\item[\emph{Magnetic-resonance image reconstruction
(medical imaging).}]
The source information is the tissue proton-density
map at the declared spatial bandlimit.  The transport
samples k-space (the Fourier dual of physical space)
under controlled magnetic-field gradient sequences;
the later carrier is the set of sampled k-space
datasets.  The recovery quotient is the inverse
Fourier transform, in clinical practice combined with
parallel-imaging unfolding and partial-Fourier
filling.  When k-space is sampled at or above the
Nyquist rate over the field of view, the obstruction
set is empty and the proton-density map is
recoverable up to the declared resolution; below the
Nyquist rate, aliasing makes the loss illegitimate
\citep{Lauterbur1973MRI}.
\end{description}

\subsection{Scale Descent / Renormalization}\label{sec:self-reference-limits:scale-descent}

\begin{theorem}[Scale Descent / Renormalization]\label{thm:F35}
Let \(H\) be a carrier with consecutive scale quotients
\(q_n:H\to Q_n\) and \(q_{n+1}:H\to Q_{n+1}\). Let
\(c:Q_n\to Q_{n+1}\) be a surjective coarsening with
\[
  c\circ q_n=q_{n+1},
\]
and let \(t:H\to T\) be a target law. Put
\[
  \mathcal O_n(t)=
  \{(h,h'):q_n(h)=q_n(h')\ \text{and}\ t(h)\ne t(h')\}.
\]
Then
\begin{align*}
  t \text{ descends to } Q_n
  &\iff \mathcal O_n(t)=\emptyset\\
  &\iff
  \exists\,\bar t_n:Q_n\to T\ \text{ such that }\
    \bar t_n\circ q_n=t .
\end{align*}
The renormalized law at scale \(Q_{n+1}\) exists exactly when
\(\bar t_n\) is constant on the fibers of \(c\).
\end{theorem}
\footnote{\raggedright Verified in Lean as
\texttt{SelfReferenceLimits.\allowbreak{}ScaleDescentRenormalization.\allowbreak{}scale\_descent\_renormalization};
see \Cref{app:formalization}.}
\begin{proof}
Assume \(c\circ q_n=q_{n+1}\) and the surjectivity of \(c\). If
\(t=\bar t_n\circ q_n\), then
\[
  q_n(h)=q_n(h') \Longrightarrow
  t(h)=\bar t_n(q_n(h))=\bar t_n(q_n(h'))=t(h'),
\]
so \(\mathcal O_n(t)=\emptyset\). Conversely, if
\(\mathcal O_n(t)=\emptyset\), define \(\bar t_n(q_n(h)):=t(h)\);
the empty obstruction is exactly the representative-independence
condition. For the next scale, there are two cases:
\[
  \begin{cases}
  \bar t_n(x)=\bar t_n(x')\ \text{whenever }c(x)=c(x'),
    &\text{then }\bar t_{n+1}(c(x)):=\bar t_n(x)\text{ is defined},\\
  \exists x,x'\ c(x)=c(x')\ \text{and}\ \bar t_n(x)\ne\bar t_n(x'),
    &\text{then exact renormalization fails.}
  \end{cases}
\]
Thus scale descent is quotient descent iterated along the tower.
\end{proof}

The Wilsonian renormalization group of statistical and quantum field
theory \citep{Wilson1975RG} is the standard substrate realization in
which scale-by-scale quotient descent is the central calculational
device.

\subsection*{Layer instantiations}\label{layer-inst:F35}

\begin{description}
\item[\emph{Hydrodynamic limit via Chapman-Enskog
expansion (kinetic theory).}]
For a rarefied gas, the carrier is the space of
phase-space distribution functions \(f(x,v,t)\).  The
scale quotient at the kinetic scale records the
one-particle distribution at Knudsen number
\(\mathrm{Kn}\sim O(1)\); the coarsening sends \(f\)
to its hydrodynamic moments (density, momentum
density, energy density) in the limit
\(\mathrm{Kn}\to 0\).  The target law at the kinetic
scale is the Boltzmann equation; the renormalised law
at the hydrodynamic scale is the Navier-Stokes (or
Euler) system, derived by the Chapman-Enskog
asymptotic expansion.  Scale descent succeeds in
\(\mathrm{Kn}\ll 1\) regimes and fails near shocks
and boundary layers \citep{Cercignani1988Boltzmann}.

\item[\emph{Multigrid methods on a nested mesh
hierarchy (numerical analysis).}]
For a linear elliptic PDE discretised on a nested
sequence of meshes
\(V_0\subset V_1\subset V_2\subset\cdots\) with mesh
sizes \(h_0>h_1>h_2>\cdots\), the carrier is the
finest-grid function space \(V_N\).  The scale quotient
is the restriction \(R_n:V_N\to V_n\); the coarsening
is the restriction operator between successive
coarser levels.  The discrete Galerkin variational
form \(A_n u_n=f_n\) at each level is the target law;
Galerkin coarsening \(A_{n-1}=R_n A_n P_n\) (with
prolongation \(P_n\)) makes the coarse operator the
exact descended renormalised law.  The multigrid
V-cycle exploits this descent for asymptotically
optimal solution costs \citep{Hackbusch1985Multigrid}.

\item[\emph{Burt-Adelson Laplacian image pyramid
(computer vision and signal processing).}]
For a digital image, the carrier is the space of
images at all dyadic resolutions.  The scale quotient
sends an image to its Gaussian-pyramid level at
resolution \(2^{-n}\) (Gaussian filter followed by
downsampling); the coarsening sends level \(n\) to
level \(n+1\) by repeating the filter-and-downsample
step.  The band-pass Laplacian layer
\(L_n=G_n-\mathrm{up}(G_{n+1})\) is the target law at
level \(n\) and descends through the scale quotient.
The Laplacian pyramid is a complete representation:
the original image is recovered exactly from the
bottom level and the Laplacian residues
\citep{BurtAdelson1983LaplacianPyramid}.

\item[\emph{Coarse-grained molecular dynamics
(computational physics and materials).}]
For a biomolecular system, the carrier is the
all-atom phase-space trajectory.  The scale quotient
groups atoms into coarse-grained beads (e.g.,
MARTINI mapping of heavy atoms to lipid beads); the
coarsening iterates the grouping to even coarser
representations.  The target law at each level is a
structural-statistical observable matched between
the all-atom and coarse-grained models (radial
distribution functions, free-energy landscapes,
thermodynamic averages); the renormalised law is the
effective bead Hamiltonian or potential of mean force
chosen to reproduce those observables.  Scale descent
succeeds in the equilibrium structural regime and
becomes path-dependent for dynamical observables
sensitive to integrated-out fast atomic motions
\citep{Voth2008CoarseGraining}.

\item[\emph{Lucas critique on policy-invariant micro
parameters (macroeconomics).}]
The carrier is the set of individual agent decision
histories under a declared policy regime.  The scale
quotient collapses individual decisions to
macroeconomic aggregates (output, inflation,
employment); the coarsening sends aggregates under
one policy regime to aggregates under another via
macro-econometric extrapolation.  The Lucas critique
observes that reduced-form macro laws do not descend
through this coarsening when the policy change
itself alters agents' decision rules: the obstruction
set is non-empty.  Legitimate scale descent requires
identifying policy-invariant deep-structural micro
parameters and aggregating through them
\citep{Lucas1976Critique}.
\end{description}

\subsection{Singular Limit Completion}\label{sec:self-reference-limits:singular-limit}

\begin{theorem}[Singular Limit Completion]\label{thm:F36}
Let \(P\) be a set of limiting processes, let \(\ell:P\to\widehat Q\)
be a surjective limit quotient, let \(j:Q\to\widehat Q\) be a
completion map, let \(u:P\to\widehat T\) be the target-limit readout,
let \(\sigma:P\to\widehat S\) be a declared signature-limit readout,
let \(u^{\mathrm{ren}}:P\to\widehat T^{\mathrm{ren}}\) be a declared
renormalized target-limit readout, and let \(p_0\in P\) be a selected
process.  Define the singular target obstruction
\[
  \mathcal O_{\ell}(u)
  =\{(p_1,p_2)\in P^2:
      \ell(p_1)=\ell(p_2)\ \text{and}\ u(p_1)\ne u(p_2)\}
\]
and the corresponding family obstruction
\(\mathcal O_{\ell}(\sigma)\).  Write
\[
  \rho_{u,\ell}:P\to\widehat Q\times\widehat T,
  \qquad
  \rho_{u,\ell}(p):=(\ell(p),u(p)),
\]
for the joint \((\ell,u)\) refinement, and let
\(\operatorname{Current}_{j}(\ell,p_0)\) abbreviate
\(\exists c\in Q,\ j(c)=\ell(p_0)\).  Then
\begin{align*}
  \operatorname{Completed}(u,\ell)
  &\iff \mathcal O_{\ell}(u)=\emptyset,\\
  \operatorname{Completed}(u,\ell)
  &\iff
  \exists\,\bar u:\widehat Q\to\widehat T,\ \bar u\circ\ell=u,\\
  \operatorname{Completed}(\sigma,\ell)
  &\iff \mathcal O_{\ell}(\sigma)=\emptyset,\\
  \operatorname{PathDependent}(u,\ell)
  &\iff \neg\operatorname{Completed}(u,\ell),\\
  \operatorname{Removable}(j,\ell,\sigma,p_0)
  &\iff
  \neg\operatorname{Current}_{j}(\ell,p_0)
  \wedge
  \operatorname{Completed}(\sigma,\ell),\\
  &\operatorname{Refines}(\rho_{u,\ell},\ell),\\
  &\operatorname{Completed}(u,\rho_{u,\ell}),\\
  \operatorname{Completed}(u^{\mathrm{ren}},\ell)
  &\iff
  \operatorname{Completed}(u^{\mathrm{ren}},\ell).
\end{align*}
Equivalently, completion is descent of \(u\) through \(\ell\); the
canonical path-memory repair \(\rho_{u,\ell}\) is a refinement of
\(\ell\) that completes \(u\); a renormalized target inherits the
same completion criterion as its raw target.
\end{theorem}
\footnote{\raggedright Verified in Lean as
\texttt{SelfReferenceLimits.\allowbreak{}SingularLimitCompletion.\allowbreak{}singular\_limit\_completion};
see \Cref{app:formalization}.}
\begin{proof}
Assume the surjectivity of \(\ell\) and the displayed maps.  We prove
the eight conjuncts.

\emph{(1) Completed iff singular obstruction empty.}
If \(u=\bar u\circ\ell\) for some
\(\bar u:\widehat Q\to\widehat T\), then
\[
  \ell(p_1)=\ell(p_2)
  \Longrightarrow
  u(p_1)=\bar u(\ell(p_1))=\bar u(\ell(p_2))=u(p_2),
\]
so \(\mathcal O_{\ell}(u)=\emptyset\).  Conversely, if
\(\mathcal O_{\ell}(u)=\emptyset\), choose
\(p\in\ell^{-1}(\hat q)\), possible by surjectivity of \(\ell\), and
set
\[
  \bar u(\hat q):=u(p).
\]
The definition is independent of the chosen process because any two
choices have the same \(\ell\)-class and hence the same \(u\)-value.
Thus \(\bar u\circ\ell=u\).

\emph{(2) Completed iff factorization.}
The factorization form is the same descent statement written as an
existential:
\[
  \operatorname{Completed}(u,\ell)
  \iff
  \exists\,\bar u:\widehat Q\to\widehat T,\ \bar u\circ\ell=u.
\]
This is exactly the construction and verification from (1).

\emph{(3) Target family completed iff family obstruction empty.}
Apply the same descent argument to the signature-limit readout
\(\sigma:P\to\widehat S\).  The obstruction is
\[
  \mathcal O_{\ell}(\sigma)
  =
  \{(p_1,p_2)\in P^2:
      \ell(p_1)=\ell(p_2)\ \text{and}\ \sigma(p_1)\ne\sigma(p_2)\},
\]
so \(\operatorname{Completed}(\sigma,\ell)\) holds exactly when this
set is empty.

\emph{(4) Path-dependent iff not completed.}
Path dependence is witnessed by a split pair:
\[
  \operatorname{PathDependent}(u,\ell)
  \iff
  \exists(p_1,p_2)\in\mathcal O_{\ell}(u).
\]
By (1), \(\mathcal O_{\ell}(u)\ne\emptyset\) is equivalent to
\(\neg\operatorname{Completed}(u,\ell)\).

\emph{(5) Removable iff boundary-completed plus family-completed.}
For the selected process \(p_0\), current realization means
\[
  \operatorname{Current}_{j}(\ell,p_0)
  \iff
  \exists c\in Q,\ j(c)=\ell(p_0).
\]
Boundary-completedness is the negation of this current-realization
condition.  Therefore removable-by-completion is exactly
\[
  \operatorname{Removable}(j,\ell,\sigma,p_0)
  \iff
  \neg\operatorname{Current}_{j}(\ell,p_0)
  \wedge
  \operatorname{Completed}(\sigma,\ell).
\]

\emph{(6) The path repair refines \(\ell\).}
The canonical path-memory repair is
\[
  \rho_{u,\ell}(p)=(\ell(p),u(p)).
\]
Projection to the first coordinate recovers the original limit:
\[
  \operatorname{pr}_{\widehat Q}\circ\rho_{u,\ell}=\ell.
\]
Thus each \(\rho_{u,\ell}\)-fiber lies inside an \(\ell\)-fiber, so
\(\rho_{u,\ell}\) is a source refinement of \(\ell\).

\emph{(7) After the path repair, the target is completed.}
Projection to the second coordinate recovers the target-limit readout:
\[
  \operatorname{pr}_{\widehat T}\circ\rho_{u,\ell}=u.
\]
Hence \(u\) factors through the repaired quotient by the explicit
factor \(\operatorname{pr}_{\widehat T}\), and therefore
\[
  \operatorname{Completed}(u,\rho_{u,\ell}).
\]

\emph{(8) Renormalized completed iff completed.}
The renormalized-target completion predicate is the ordinary target
completion predicate applied to the declared renormalized readout
\(u^{\mathrm{ren}}\).  Thus its expansion is the reflexive
equivalence
\[
  \operatorname{Completed}(u^{\mathrm{ren}},\ell)
  \iff
  \operatorname{Completed}(u^{\mathrm{ren}},\ell),
\]
recording that renormalization is a relabeling of the declared target.
\end{proof}

\subsection*{Layer instantiations}\label{layer-inst:F36}

\begin{description}
\item[\emph{Riemann removable singularity in complex
analysis (complex analysis).}]
For a function \(f\) holomorphic on a punctured disk,
the set of limiting processes is the family of
approach paths to the puncture; the limit quotient
collapses paths sharing the same limit point; the
target readout is the path-limit value of \(f\).  By
Riemann's removable-singularity theorem, if \(f\) is
bounded on the punctured disk then the limit is
independent of approach path and the obstruction set
is empty.  Poles and essential singularities
instantiate the path-dependent case
\citep{Ahlfors1979ComplexAnalysis}.

\item[\emph{L'H\^opital's rule for indeterminate
forms (real analysis).}]
For functions \(f,g\) with \(f(a)=g(a)=0\), the
literal substitution \(\lim_{x\to a}f(x)/g(x)\)
yields the indeterminate form \(0/0\), which by
itself does not descend through the limit quotient.
L'H\^opital's rule provides the completion: if \(f\)
and \(g\) are differentiable with \(g'(a)\ne 0\),
the limit equals \(f'(a)/g'(a)\).  The path-memory
repair records the order of vanishing of \(f\) and
\(g\) at \(a\), completing the limit through the
derivative-ratio readout
\citep{Cauchy1821CoursAnalyse}.

\item[\emph{Prandtl boundary-layer matched asymptotic
expansion (fluid dynamics).}]
For the incompressible Navier-Stokes equations with
small viscosity \(\nu\), the set of limiting
processes is the family of solutions parameterised by
\(\nu\).  The limit quotient sends each solution to
its formal \(\nu\to 0\) limit; the target readout
records the velocity field.  The naive Euler limit
fails to satisfy the no-slip boundary condition;
Prandtl's boundary-layer theory introduces an inner
expansion in the rescaled wall-normal coordinate
\(y/\sqrt{\nu}\) and matches it to the outer Euler
solution.  The matched-asymptotic structure is the
canonical path-memory repair, carrying the
boundary-layer thickness as an additional quotient
parameter that completes the singular limit
\citep{Prandtl1904BoundaryLayer}.

\item[\emph{Critical-point scaling at second-order
phase transitions (statistical mechanics).}]
For a magnetic system near its critical temperature
\(T_c\), the set of limiting processes is the family
of approach paths to \((T,h)=(T_c,0)\) in the
temperature-field plane.  Order parameter,
susceptibility, and correlation length diverge or
vanish with universal critical exponents but with
prefactors that depend on the approach direction.
The Widom scaling hypothesis is the path-memory
repair: it records the universal scaling functions
parameterised by \(h/|t|^{\beta\delta}\) (with
\(t=(T-T_c)/T_c\)), completing the critical-point
readout by labelling each approach direction
\citep{Stanley1971PhaseTransitions}.

\item[\emph{Riemann integral as limit of Riemann sums
(classical analysis).}]
For a bounded \(f:[a,b]\to\mathbb R\), the set of
limiting processes is the family of tagged
partitions with mesh tending to zero; the target
readout is the Riemann sum.  The function is
Riemann-integrable precisely when the obstruction
set is empty, i.e., when the limit of Riemann sums
is independent of the choice of partition
refinement and tag selection.  The completed map is
the Riemann integral
\(\int_a^b f\,dx\); non-Riemann-integrable functions
(e.g., the Dirichlet function) instantiate the
path-dependent case
\citep{Riemann1854Habilitationsschrift}.
\end{description}

\subsection{Complementarity as Non-Joint Access}\label{sec:self-reference-limits:complementarity}

The classical quantum-mechanical instance of two access quotients
that do not admit a joint refinement is the
position--momentum complementarity formalized by
\citet{Bohr1928Quantum}; the layer-agnostic statement below
abstracts the non-joint-access shape from that substrate.

\begin{theorem}[Complementarity as Non-Joint Access]\label{thm:F37}
Let \(H\) be a carrier with two access quotients
\(q_A:H\to Q_A\) and \(q_B:H\to Q_B\). A joint access is a quotient
\(j:H\to J\) together with maps \(a:J\to Q_A\) and \(b:J\to Q_B\)
satisfying
\[
  a\circ j=q_A,\qquad b\circ j=q_B .
\]
Assume \(q_A\) and \(q_B\) are separately admissible. Then
\begin{align*}
  q_A \perp_{\mathrm{acc}} q_B
  \iff
  \neg\exists (J: \mathrm{Type})(j:H\to J)(a:J\to Q_A)(b:J\to Q_B),\
  a\circ j=q_A\ \wedge\ b\circ j=q_B .
\end{align*}
Thus complementarity is precisely non-existence of an admissible joint
quotient carrying both accesses.
\end{theorem}
\footnote{\raggedright Verified in Lean as
\texttt{SelfReferenceLimits.\allowbreak{}Complementarity.\allowbreak{}complementarity\_as\_non\_joint\_access};
see \Cref{app:formalization}.}
\begin{proof}
Assume the separate admissibility of \(q_A\) and \(q_B\). There are two
cases. If a tuple \((J,j,a,b)\) satisfying
\[
  a\circ j=q_A,\qquad b\circ j=q_B
\]
exists, then a single quotient \(j\) carries both accesses, so the pair
is not complementary. If no such tuple exists, any proposed joint
readout would have to add a distinction not present in either access
quotient. By the quotient criterion for carrying both maps, that added
distinction is exactly the missing joint quotient. Therefore the
separately admissible accesses are complementary precisely in the
non-existence case.
\end{proof}

\subsection*{Layer instantiations}\label{layer-inst:F37}

\begin{description}
\item[\emph{Time-frequency Gabor uncertainty (signal
processing).}]
For a square-integrable signal \(s\in L^2(\mathbb R)\),
the access quotient \(q_A\) records its time-marginal
\(|s(t)|^2\) and \(q_B\) records its
frequency-marginal \(|\hat s(f)|^2\).  Within the
admissibility class of Cohen-class bilinear
time-frequency representations satisfying the correct
marginals, the covariance properties of bilinear
distributions, and non-negativity (probability-density
interpretation), no joint quotient exists: the
Wigner-Ville distribution satisfies marginals and
covariance but is not non-negative, and no member of
Cohen's class is jointly compatible.  The access
quotients are F37-complementary
\citep{Cohen1995TimeFrequency}.

\item[\emph{Joint diagonalisation of non-commuting
Hermitians (linear algebra).}]
For a finite-dimensional Hilbert space \(H\) and two
Hermitian operators \(A,B\) on \(H\), the access
quotient \(q_A\) records the spectral labelling of
\(A\) in a basis and \(q_B\) records that of \(B\).
A joint admissible quotient is a single orthonormal
basis simultaneously diagonalising both operators.
By the spectral theorem, such a basis exists iff the
commutator \([A,B]=AB-BA\) vanishes.  When
\([A,B]\ne 0\), no admissible joint quotient exists;
the access quotients are F37-complementary within
the orthonormal-basis class
\citep{Halmos1974FiniteDimensional}.

\item[\emph{Intuitionistic vs classical logic
(mathematical logic).}]
For a propositional language \(L\) with atomic
propositions, the access quotient \(q_A\) sends a
formula \(\varphi\) to its intuitionistic-validity
class under the Brouwer-Heyting-Kolmogorov (BHK)
interpretation; \(q_B\) sends \(\varphi\) to its
classical-proof class requiring the law of excluded
middle (LEM) and proof by contradiction.  Within the
admissibility class of BHK-respecting proof systems,
no joint quotient validates LEM for undecidable
propositions \(P\): classical proof of
\(P\vee\neg P\) is admissible classically but lacks
a constructive witness.  The two proof-system
quotients are F37-complementary, formalised by
Heyting's intuitionistic logic
\citep{Heyting1930Logik}.

\item[\emph{{CAP} theorem in distributed systems
(distributed systems).}]
For a replicated distributed data store, \(q_A\)
records linearisable-consistency observations and
\(q_B\) records always-available reads and writes.
The admissibility constraint is partition tolerance:
the implementation must continue under arbitrary
network partitions.  Gilbert-Lynch's formalisation of
Brewer's CAP conjecture shows that, under network
partitions, no admissible implementation jointly
realises both \(q_A\) and \(q_B\); the access
quotients are F37-complementary within the
partition-tolerant implementation class
\citep{GilbertLynch2002CAP}.

\item[\emph{Russell's paradox under naive
comprehension (set-theoretic foundations).}]
For a candidate set \(R\), the access quotient
\(q_A\) imposes the self-membership predicate
\(x\in R\iff x\notin x\); the access quotient
\(q_B\) requires \(R\) to be a set in the admissible
class.  Substituting \(R\) for \(x\) yields
\(R\in R\iff R\notin R\), so no joint admissible
quotient exists under unrestricted comprehension.
The standard remedies (ZFC's separation and
foundation axioms, Russell-Whitehead type
stratification) refine admissibility to block this
joint quotient \citep{Russell1902Letter}.
\end{description}

\subsection{Arrow from Record Monotonicity}\label{sec:self-reference-limits:arrow}

The classical reading of an arrow of time as monotone growth of an
audit record of the universe is due to \citet{Eddington1928Nature};
the formulation below records the layer-agnostic structural form
without committing to a specific physical realization.

\begin{theorem}[Arrow from Record Monotonicity]\label{thm:F38}
Let \(H_s,H_t\) be source and target history sets with continuation
\(\kappa:H_s\to H_t\). Let \(r_s:H_s\to R_s\) and
\(r_t:H_t\to R_t\) be surjective record quotients, and let
\(s_s:R_s\to S\), \(s_t:R_t\to S\) be status readouts into a preorder
\((S,\le_S)\). Then
\begin{align*}
  r_t\circ\kappa \text{ descends through } r_s
  &\iff
  \exists\,\bar\kappa:R_s\to R_t\ \text{ such that }\
    \bar\kappa\circ r_s=r_t\circ\kappa,\\
  \operatorname{Arrow}(\kappa)
  &\iff
  \exists\,\bar\kappa:R_s\to R_t,\
  \bar\kappa\circ r_s=r_t\circ\kappa
  \ \wedge\
  \forall x\in R_s,\ s_s(x)\le_S s_t(\bar\kappa x).
\end{align*}
The arrow is the monotone direction of the descended record transport.
\end{theorem}
\footnote{\raggedright Verified in Lean as
\texttt{SelfReferenceLimits.\allowbreak{}ArrowFromRecordMonotonicity.\allowbreak{}arrow\_from\_record\_monotonicity};
see \Cref{app:formalization}.}
\begin{proof}
Assume \(r_s\) and \(r_t\) are surjective. If
\(\bar\kappa\circ r_s=r_t\circ\kappa\), then for
\(h,h'\in H_s\) with \(r_s(h)=r_s(h')\),
\[
  r_t(\kappa h)=\bar\kappa(r_s h)
              =\bar\kappa(r_s h')
              =r_t(\kappa h'),
\]
so continuation is record-invariant. Conversely, if this
record-invariance holds, define \(\bar\kappa(r_s h):=r_t(\kappa h)\);
surjectivity and representative independence make \(\bar\kappa\)
well-defined. The arrow condition then splits into the two displayed
requirements:
\[
  \text{descent of }r_t\circ\kappa\quad\text{and}\quad
  s_s(x)\le_S s_t(\bar\kappa x)\ \text{ for all }x\in R_s .
\]
These are necessary and sufficient for monotone record growth.
\end{proof}

\subsection*{Layer instantiations}\label{layer-inst:F38}

\begin{description}
\item[\emph{Git commit {DAG} and version-control history
(software engineering).}]
The source history set is a project-state at a given
commit; the target is the project-state after a child
commit; \(\kappa\) records the commit operation.  The
record quotient sends a state to its commit hash, with
status readout equal to the DAG-ancestry generation
(the longest ancestor chain length).  The hash of a
new commit is the cryptographic digest of its content
together with parent-commit references, so the
descended map carries each parent hash to the child
hash.  The status readout is strictly monotone: each
child commit has generation greater than each of its
parents \citep{ChaconStraub2014ProGit}.

\item[\emph{Steno's law of superposition (geology).}]
The source history set is the configuration of
sediment-deposition processes; \(\kappa\) is the
geological-time evolution.  The record quotient sends
each deposition event to its stratum class (named bed,
geological epoch); the status readout is the
stratigraphic depth.  Steno's law of superposition
asserts that, in undisturbed sequences, older strata
lie below younger strata.  Folding, faulting, and
overturning are F38-obstructions that require
unfolding to recover the arrow \citep{Steno1669Prodromus}.

\item[\emph{Dendrochronological tree-ring sequence
(paleoclimate and archaeology).}]
The source history set is the sequence of annual
tree-ring formation events; \(\kappa\) records one
annual growth cycle.  The record quotient sends a
tree-ring to its ring-pattern signature; the status
readout is the calendrical year identified by master
tree-ring chronology cross-dating.  Each successive
ring is added outside the previous ones, producing a
strictly monotone chronological sequence.  Douglass's
master chronology cross-dates archaeological wood to
its year of growth via pattern matching against the
reference sequence \citep{Douglass1919ClimaticCycles}.

\item[\emph{Libby radiocarbon dating (archaeology and
nuclear chemistry).}]
The source history set is the moment of organism death
(after which exchange of \(^{14}\)C with the atmosphere
ceases); \(\kappa\) records radioactive decay over
time.  The record quotient sends a sample to its
measured \(^{14}\)C\(/^{12}\)C ratio; the status readout
is the calibrated radiocarbon age inferred from the
exponential decay law with half-life \(\approx 5730\)
years.  The ratio decreases monotonically, providing a
chronological arrow up to about 50{,}000 years before
present, calibrated against tree-ring and varve
chronologies \citep{Libby1955RadiocarbonDating}.

\item[\emph{Brugmann historical-comparative
phonological reconstruction (historical linguistics).}]
The source history set is the lexicon of a
reconstructed proto-language; the target is the
lexicon of an attested daughter language; \(\kappa\)
records the regular sound-change rules carrying
proto-forms forward to daughter reflexes.  The record
quotient sends a proto-form to its phonological class;
the status readout is the chronological layer in the
family tree, increasing from the proto-language toward
each attested daughter.  Brugmann's regularity
principle asserts that sound changes are exceptionless
within a defined phonological environment, providing a
monotone forward arrow along the family tree
\citep{Brugmann1886Grundriss}.
\end{description}

\subsection{Entropy as Fiber Volume}\label{sec:self-reference-limits:entropy}

The fiber-volume reading of entropy combines two classical substrate
sources: the statistical-mechanical interpretation
\(S=k\log W\) developed in \citet{Boltzmann1877Beziehung} and the
communication-theoretic entropy of \citet{Shannon1948Mathematical}.
The layer-agnostic statement below records the structural form
common to both.

\begin{theorem}[Entropy as Fiber Volume]\label{thm:F39}
Let \(q:H\to Q\) be a surjective access quotient with target readout
\(t:H\to T\).  Let \(q_F:H\to Q_F\) be a finer quotient and
\(q_C:H\to Q_C\) a coarsening, related by a factorization
\(c:Q_F\to Q_C\) with
\[
  c\circ q_F=q_C.
\]
Let \((\mathcal E,\le_E)\) be a preorder of entropy values, and let
\[
  E:\{\,q':H\to Q'\,\}\longrightarrow \mathcal E
\]
be an entropy functional assigning a fiber-volume value to every
declared quotient.  Suppose the apparatus declares a mixing
functional \(m\) and a combiner
\(\odot:\mathcal E\times\mathcal E\to\mathcal E\) satisfying the
chain-rule identity
\[
  E(q_C)=E(q_F)\odot m(q_F,q_C)
\]
and the combine-monotonicity condition
\[
  \forall a,b\in\mathcal E,\quad a\le_E a\odot b.
\]
Define the raw split-pair and target obstructions
\begin{align*}
  \mathcal O_q^{\mathrm{raw}}
  &:=
  \{(h,h')\in H^2:q(h)=q(h')\ \text{and}\ h\ne h'\},\\
  \mathcal O_q(t)
  &:=
  \{(h,h')\in H^2:q(h)=q(h')\ \text{and}\ t(h)\ne t(h')\}.
\end{align*}
Then
\begin{align*}
  E(q)=0
  &\iff
  \mathcal O_q^{\mathrm{raw}}=\emptyset,\\
  \operatorname{TargetEntropyZero}(q,t)
  &\iff
  t\ \text{descends through}\ q,\\
  \operatorname{TargetUnresolved}(q,t)
  &\iff
  \neg(t\ \text{descends through}\ q),\\
  \forall b\in Q_C,\ \forall h\in H,\quad
  q_C(h)=b
  &\iff
  \exists a\in Q_F,\ c(a)=b\ \wedge\ q_F(h)=a,\\
  E(q_F)
  &\le_E
  E(q_C).
\end{align*}
\end{theorem}
\footnote{\raggedright Verified in Lean as
\texttt{SelfReferenceLimits.\allowbreak{}EntropyAsFiberVolume.\allowbreak{}entropy\_as\_fiber\_volume};
see \Cref{app:formalization}.}
\begin{proof}
Assume \(q\) is surjective, \(c\circ q_F=q_C\), the chain-rule
identity \(E(q_C)=E(q_F)\odot m(q_F,q_C)\), and combine-monotonicity
\(\forall a,b,\ a\le_E a\odot b\).  We prove the five conjuncts.

\emph{(1) Raw entropy zero iff no raw split pair.}
If \(E(q)=0\), every \(q\)-fiber has trivial unresolved volume, so
any \(h,h'\) in the same \(q\)-fiber must coincide.  Hence
\(\mathcal O_q^{\mathrm{raw}}=\emptyset\).  Conversely, if no
\(q\)-fiber contains distinct elements, every fiber is a singleton;
the declared fiber volume of a singleton is the unit of the entropy
ledger, and its logarithmic value is \(0\).

\emph{(2) Target entropy zero iff target descends.}
Target entropy is measured by the target obstruction
\(\mathcal O_q(t)\).  Its emptiness is exactly the quotient-descent
condition for \(t\) through \(q\):
\[
  \mathcal O_q(t)=\emptyset
  \iff
  \exists\,\bar t:Q\to T,\ \bar t\circ q=t.
\]
This is \(\operatorname{TargetEntropyZero}(q,t)\) iff target descent.

\emph{(3) Target unresolved iff target does not descend.}
The target-unresolved predicate is witnessed by a split pair
\((h,h')\in\mathcal O_q(t)\).  By (2), nonemptiness of
\(\mathcal O_q(t)\) is the negation of target descent.

\emph{(4) Coarse fibers decompose into fine fibers.}
Let \(b\in Q_C\) and \(h\in H\).  If \(q_C(h)=b\), set
\(a:=q_F(h)\).  The factorization \(c\circ q_F=q_C\) gives
\[
  c(a)=c(q_F(h))=q_C(h)=b,
\]
and by construction \(q_F(h)=a\).  Conversely, if there exists
\(a\in Q_F\) with \(c(a)=b\) and \(q_F(h)=a\), then the same
factorization gives
\[
  q_C(h)=c(q_F(h))=c(a)=b.
\]

\emph{(5) Entropy is monotone under coarsening.}
Apply the chain-rule identity to the factorized coarsening pair
\((q_F,q_C)\):
\[
  E(q_C)=E(q_F)\odot m(q_F,q_C).
\]
Combine-monotonicity with \(a:=E(q_F)\) and \(b:=m(q_F,q_C)\) gives
\[
  E(q_F)\le_E E(q_F)\odot m(q_F,q_C)=E(q_C).
\]
Thus \(E(q_F)\le_E E(q_C)\).
\end{proof}

\subsection*{Layer instantiations}\label{layer-inst:F39}

\begin{description}
\item[\emph{Simpson concentration index (ecology).}]
For a community of \(N\) individuals classified into
species, the access quotient sends a pair of randomly
selected individuals to its species-pair class.  The
Simpson concentration index
\(\lambda=\sum_i p_i^2\) (with \(p_i\) the abundance of
species \(i\)) is the fiber-volume entropy on the
same-species fiber.  Taxonomic coarsening (genus,
family) combines species classes; by the
Cauchy-Schwarz inequality
\(\sum_g p_g^2\ge\sum_i p_i^2\), so the index is
monotone non-decreasing under coarsening
\citep{Simpson1949Diversity}.

\item[\emph{Herfindahl-Hirschman market-concentration
index (industrial economics).}]
For a market with firms having shares \(s_i\)
(\(\sum_i s_i=1\)), the access quotient sends each
customer transaction to the firm serving it.  The
Herfindahl-Hirschman index
\(\mathrm{HHI}=\sum_i s_i^2\) is the probability that
two random transactions are served by the same firm.
Under firm grouping (mergers, parent-company
aggregation), Cauchy-Schwarz forces the sum of squared
shares to non-decrease, so HHI is monotone under
coarsening.  The US Department of Justice merger
guidelines use HHI thresholds to identify market
concentration concerns \citep{Hirschman1964HHI}.

\item[\emph{Burnside orbit-counting in finite group
theory (group theory).}]
For a finite group \(G\) acting on a finite set \(X\),
the access quotient sends each element to its
\(G\)-orbit.  The fiber over an orbit \(O\) has
cardinality \(|O|\); the fiber-volume entropy is
\(\log(|X|/|X/G|)\), with \(|X/G|\) given by
Burnside's lemma as the average number of fixed
points of group elements.  Coarsening to a supergroup
\(G'\supseteq G\) merges orbits into larger fibers,
non-decreasing the average orbit size
\citep{Burnside1897Theory}.

\item[\emph{Logarithm of the multinomial coefficient
(combinatorics).}]
For \(N\) distinguishable items distributed into \(k\)
labelled cells, the fine quotient sends an assignment
to its cell-content profile
\((n_1,\dots,n_k)\); the fine fiber has cardinality
\(\binom{N}{n_1,\dots,n_k}=N!/\prod_i n_i!\).
Coarsening merges two cells \(i,j\) into a single
labelled cell of size \(n_{\{i,j\}}=n_i+n_j\); the
coarse fiber is the disjoint union of fine fibers
with \(n_i'+n_j'=n_{\{i,j\}}\), so its cardinality is
\(\sum_{n_i'+n_j'=n_{\{i,j\}}}
\binom{N}{n_i',n_j',\dots}\), at least any single
summand.  The combinatorial fiber-volume entropy is
thus monotone non-decreasing under cell coarsening
\citep{Knuth1997TAOCP}.

\item[\emph{Kingman coalescent effective population
size (population genetics).}]
For a haploid neutral genealogical model, the access
quotient sends a genealogy of \(n\) sampled lineages
to its lineage-merger pattern.  The Kingman coalescent
gives pairwise coalescence rate \(1/N_e\), where
\(N_e\) is the effective population size; the
coalescent fiber over a coalesced event has time
length proportional to \(N_e\).  Demographically
coarsening to a model with larger effective
population size grows the coalescent fiber volume,
instantiating the F39 inequality on a model-by-model
basis \citep{Kingman1982Coalescent}.
\end{description}

\subsection{Irreducibility / No Short Closed Description}\label{sec:self-reference-limits:irreducibility}

The notion of an object's shortest closed description is the central
object of Kolmogorov complexity and algorithmic information theory;
the standard reference is \citet{LiVitanyi2008Kolmogorov}.

\begin{theorem}[Irreducibility / No Short Closed Description]\label{thm:F42}
Let \(H\) be a set of objects with base map \(b:H\to B\), compression
map \(e:H\to E\), target readout \(t:H\to T\), and route updates
\(r_\gamma:H\to H\) for \(\gamma\in\Gamma\). Let
\(\operatorname{cost}:E\to C\) take values in a strict preorder
\((C,<)\). A compression \(e\) is closed for \(h\in H\) when \(t\) and
all \(e\circ r_\gamma\) descend through the \(e\)-class of \(h\). Then
\begin{align*}
  \operatorname{ClosedCompression}(e,h)
  &\iff
  \mathcal O_t(e,h)=\emptyset
  \ \wedge\
  \forall\gamma\in\Gamma,\ \mathcal O_\gamma(e,h)=\emptyset,\\
  \operatorname{Irreducible}(h)
  &\iff
  \neg\exists e'\!,\
  \operatorname{ClosedCompression}(e',h)
  \ \wedge\
  \operatorname{cost}(e')<\operatorname{cost}(e(h)).
\end{align*}
Irreducibility is therefore absence of a shorter closed description.
\end{theorem}
\footnote{\raggedright Verified in Lean as
\texttt{SelfReferenceLimits.\allowbreak{}Irreducibility.\allowbreak{}irreducibility\_no\_short\_closed\_description};
see \Cref{app:formalization}.}
\begin{proof}
Assume the named maps and the strict cost preorder. A compression is
closed exactly when no target or route update distinguishes two
representatives of the same compressed description. In symbols,
\[
  \mathcal O_t(e,h)=
  \{(x,x'):e(x)=e(x')=e(h),\ t(x)\ne t(x')\},
\]
and similarly \(\mathcal O_\gamma(e,h)\) is formed using
\(e(r_\gamma x)\). Hence closedness is precisely the simultaneous
emptiness of these obstruction sets. There are two cases. If a shorter
closed \(e'\) exists, then \(e'\) carries all declared readouts with
smaller cost, so \(h\) is reducible. If no such \(e'\) exists, every
shorter candidate fails a target or route obstruction, and the audit
description \(e(h)\) is minimal among closed descriptions.
\end{proof}

\subsection*{Layer instantiations}\label{layer-inst:F42}

\begin{description}
\item[\emph{Abel-Ruffini-Galois unsolvability of the
quintic (algebra).}]
For a degree-\(n\ge 5\) polynomial over the rationals,
the admissible compression class restricts \(e\) to
nested-radical expressions; routes are arithmetic
operations and root-extractions within the radical
class.  A generic quintic with Galois group \(S_5\),
which contains the simple unsolvable \(A_5\), has no
root expression in this class at all.  Within the
radical-compression class, no closed description
exists; F42-irreducibility holds unconditionally
relative to the class \citep{Abel1824Quintic}.

\item[\emph{Quine-McCluskey minimum sum-of-products
form (digital-logic design).}]
For a Boolean function \(f:\{0,1\}^n\to\{0,1\}\), the
admissible compression class is sum-of-products
expressions; routes are Boolean-algebra equivalences
(De Morgan, absorption, distribution).  Quine-McCluskey
yields a minimum-literal expression via prime
implicants and an NP-hard essential-cover step; the
function is F42-irreducible at the minimum-cover
expression \citep{Quine1952Boolean}.

\item[\emph{Information-theoretic
\(\Omega(n\log n)\) sorting lower bound (analysis of
algorithms).}]
For comparison-sorting on \(n\) distinct keys, the
admissible compression class is comparison-based
decision trees; routes are two-key comparisons; the
sorted permutation is the target.  A correct decision
tree must have at least \(n!\) leaves, so its
worst-case depth is \(\Omega(\log n!) = \Omega(n\log
n)\).  No asymptotically shallower comparison-based
decision tree exists; the F42 irreducibility is the
unconditional \(\Omega(n\log n)\) lower bound
\citep{Knuth1998TAOCPv3}.

\item[\emph{Minimal genome for self-replicating cell
(synthetic biology).}]
Within the JCVI-syn3.0 Mycoplasma chassis design
space and declared laboratory growth conditions, the
admissible compression is the genome's gene
complement; routes are single-gene deletion /
insertion edits.  The Hutchison-Gibson-Venter team
reached a 531-kilobase genome encoding 473 genes;
further single-gene deletions in the tested design
space disrupted either metabolism or replication.
Relative to that scoped design space and growth
medium, the minimal viable genome is
F42-irreducible at JCVI-syn3.0
\citep{Hutchison2016MinimalGenome}.

\item[\emph{Crystallographic irreducible Brillouin
zone (solid-state physics).}]
For a periodic crystalline solid with space group
\(G\), the admissible compression is the
restriction of crystal momentum \(\mathbf k\) to the
irreducible Brillouin zone (IBZ), the fundamental
domain of the reciprocal-lattice point-group action;
routes are point-group operations.  The
Bouckaert-Smoluchowski-Wigner analysis shows the
IBZ is the minimal subset of the full Brillouin
zone that, together with symmetry operations,
generates all distinct electronic states.  The IBZ
is F42-irreducible by construction
\citep{BouckaertSmoluchowskiWigner1936Theory}.
\end{description}

\subsection{Continuum Emergence}\label{sec:self-reference-limits:continuum-emergence}

\begin{theorem}[Continuum Emergence]\label{thm:F43}
Let \((I,\le)\) be a directed index set, let \(A_i\) be a discrete
approximation at index \(i\), and let
\(\phi_{ij}:A_j\to A_i\) be comparison maps for \(i\le j\). Let
\(p_i:L\to A_i\) be projections from a candidate limit \(L\), let
\(\xi_i\in R_i\) be residuals, and let
\(\ell:L\to Q_\infty\) be a surjective limit quotient. Then
\begin{align*}
  \operatorname{StableLimitQuotient}(L)
  \iff\
  &\operatorname{TowerCoherent}(\phi_{ij})
  \ \wedge\
  \operatorname{Cofinal}(I)\\
  &\wedge\
  \operatorname{Vanishes}(\xi_i)
  \ \wedge\
  \exists\,\bar\ell:Q_\infty\to L_{\mathrm{rec}}\ \text{ compatible with all }p_i .
\end{align*}
Continuum emergence is this stable quotient-completion of the
coherent approximation tower.
\end{theorem}
\footnote{\raggedright Verified in Lean as
\texttt{SelfReferenceLimits.\allowbreak{}ContinuumEmergence.\allowbreak{}continuum\_emergence};
see \Cref{app:formalization}.}
\begin{proof}
Assume the directed tower \((A_i,\phi_{ij})\), residuals \((\xi_i)\),
and surjective limit quotient \(\ell\). If a stable limit quotient is
formed, its projections satisfy
\[
  \phi_{ij}\circ p_j=p_i\qquad (i\le j),
\]
so the tower is coherent; every declared finite test factors through
some \(A_i\), giving cofinality; and the residuals vanish because the
limit has no remaining declared finite obstruction. Conversely assume
the four displayed conditions. Coherence and cofinality give a
well-defined inverse-system endpoint, vanishing residuals remove the
singular obstruction, and Singular Limit Completion applied to
\(\ell\) gives the quotient-level map \(\bar\ell\). The endpoint is
therefore stable, and the continuum has emerged as the completion of
the discrete tower.
\end{proof}

\subsection*{Layer instantiations}\label{layer-inst:F43}

\begin{description}
\item[\emph{Dedekind cut construction of the reals
(analysis).}]
The discrete approximations are rational intervals
indexed by precision; comparison maps are
inclusions of finer intervals into coarser ones.
A Dedekind cut is a non-empty downward-closed
subset \(A\subset\mathbb Q\) without greatest
element.  The candidate limit is the set of cuts;
projections send a cut to its rational interval
approximant; the surjective limit quotient sends
a cut to its real value.  Tower coherence is
nested-approximation consistency, cofinality is
density of \(\mathbb Q\) in \(\mathbb R\), and
vanishing residuals correspond to absence of
gaps in the real line \citep{Dedekind1872Stetigkeit}.

\item[\emph{p-adic completion in number theory
(algebraic number theory).}]
For a prime \(p\), the discrete approximations are
the finite quotient rings \(\mathbb Z/p^i\mathbb Z\)
with reduction maps
\(\phi_{ij}:\mathbb Z/p^j\to\mathbb Z/p^i\).  The
candidate limit is the p-adic integers
\(\mathbb Z_p=\varprojlim\mathbb Z/p^i\mathbb Z\),
with projections to each finite quotient.  Tower
coherence is congruence compatibility, cofinality
is directedness of p-power indices, and vanishing
residuals correspond to convergence in the p-adic
metric.  The non-archimedean continuum emerges as
the p-adic completion of \(\mathbb Q\),
distinct from the archimedean reals
\citep{Hensel1908Theorie}.

\item[\emph{Thermodynamic limit in statistical
mechanics (statistical physics).}]
For a classical-spin lattice model on a finite
volume \(\Lambda\subset\mathbb Z^d\), the discrete
approximations are finite-volume Gibbs measures
at inverse temperature \(\beta\) with restriction
maps for nested volumes.  The candidate limit is
the set of Dobrushin-Lanford-Ruelle (DLR) states
on \(\{+1,-1\}^{\mathbb Z^d}\), with projections
to finite-volume marginals.  Tower coherence is
the DLR consistency condition, cofinality is
exhaustion by finite volumes, and vanishing
residuals correspond to absence of unstable
boundary contributions.  The infinite-volume Gibbs
state emerges as the F43 stable limit quotient
\citep{Ruelle1969StatisticalMechanics}.

\item[\emph{Krull absolute Galois group as inverse
limit (field theory).}]
For a perfect field \(K\) with separable closure
\(K^{\mathrm{sep}}\), the discrete approximations
are finite Galois groups
\(\mathrm{Gal}(L_i/K)\) for finite Galois
extensions \(L_i\), with restriction maps for
\(L_i\subseteq L_j\).  The candidate limit is the
absolute Galois group
\(\mathrm{Gal}(K^{\mathrm{sep}}/K)
=\varprojlim_i\mathrm{Gal}(L_i/K)\), with
projections by restriction.  Tower coherence is
compatibility of restrictions, cofinality is
directedness of finite Galois extensions, and
vanishing residuals correspond to profinite
separation by the finite quotients.  Krull's
infinite Galois theory establishes the absolute
Galois group with its profinite (Krull) topology
as the F43 stable limit quotient
\citep{Krull1928Galois}.

\item[\emph{Born-Huang continuum elasticity from
atomic lattice (solid-state physics).}]
For a crystalline solid, the discrete
approximations are atomic-lattice models at
finer-and-finer lattice spacings; comparison maps
are inter-lattice coarse-graining operations.
The candidate limit is the continuum elasticity
theory governing displacement fields on
\(\mathbb R^3\); projections restrict continuum
solutions to lattice points.  Tower coherence is
compatibility of inter-lattice strain
representations, cofinality is the existence of
sufficiently fine lattices for any continuum
scale, and vanishing residuals correspond to the
long-wavelength approximation being valid.  The
Born-Huang derivation of elastic constants from
interatomic potentials is the canonical
solid-state F43 continuum emergence
\citep{BornHuang1954Dynamical}.
\end{description}

\subsection{Criticality / Phase Boundary}\label{sec:self-reference-limits:criticality}

The substrate instance is the theory of phase transitions and
critical phenomena; the standard textbook treatment is
\citet{Stanley1971PhaseTransitions}.  The layer-agnostic statement
below records the descent-and-margin shape that the substrate
realizes.

\begin{theorem}[Criticality / Phase Boundary]\label{thm:F44}
Let \(\mathcal C\) be a control space with quotient
\(c:\mathcal C\to C\), phase readout \(\varphi:\mathcal C\to\Phi\),
formedness predicate \(F:\mathcal C\to\mathrm{Prop}\), and
neighborhood relation \(N(x,U,y)\). Let
\(S,M,U,R:\mathcal C\to\mathrm{Prop}\) denote, respectively,
status-change, margin-vanishing, non-uniform response, and
residuals-statused. For a readout \(r:\mathcal C\to V\), define
\[
  \mathcal O_c(r)=
  \{(x,y):c(x)=c(y)\ \text{and}\ r(x)\ne r(y)\}.
\]
Then, for \(x,\ell,r_0\in\mathcal C\),
\begin{align*}
  r \text{ descends through }c
  &\iff \mathcal O_c(r)=\emptyset,\\
  \operatorname{PhaseTransition}(\ell,r_0)
  &\iff \varphi(\ell)\ne\varphi(r_0),\\
  \operatorname{PhaseBoundary}(x)
  &\iff F(x)\ \wedge\
    \neg\exists U,\ \forall y,\ N(x,U,y)\Rightarrow \varphi(y)=\varphi(x),\\
  \operatorname{Critical}(x)
  &\iff \operatorname{PhaseBoundary}(x)\wedge (M(x)\vee U(x))\wedge R(x).
\end{align*}
\end{theorem}
\footnote{\raggedright Verified in Lean as
\texttt{SelfReferenceLimits.\allowbreak{}Criticality.\allowbreak{}criticality\_phase\_boundary};
see \Cref{app:formalization}.}
\begin{proof}
Assume the quotient \(c\), phase readout \(\varphi\), neighborhood
relation \(N\), and predicates \(F,S,M,U,R\). The descent clause is the
ordinary quotient criterion:
\[
  r=\bar r\circ c
  \quad\Longleftrightarrow\quad
  \forall x,y,\ c(x)=c(y)\Rightarrow r(x)=r(y)
  \quad\Longleftrightarrow\quad
  \mathcal O_c(r)=\emptyset .
\]
The phase-transition clause is definitional: the two controls occupy
different phases exactly when their phase values differ. For boundary
and criticality there are two cases:
\[
  \begin{cases}
  \exists U\,\forall y,\ N(x,U,y)\Rightarrow\varphi(y)=\varphi(x),
    &\text{\(x\) is locally phase-constant,}\\
  \neg\exists U\,\forall y,\ N(x,U,y)\Rightarrow\varphi(y)=\varphi(x),
    &\text{\(x\) lies on a phase boundary.}
  \end{cases}
\]
Criticality is the boundary case in which either the separating margin
vanishes or response is non-uniform, and the residual is still typed by
\(R(x)\).
\end{proof}

\subsection*{Layer instantiations}\label{layer-inst:F44}

\begin{description}
\item[\emph{Allee effect threshold (population
ecology).}]
For a species population \(\dot N=N\,r(N)\), the
control space is population size; the phase readout
records ``extinction'' vs ``persistence''.  The Allee
effect makes the per-capita growth \(r(N)\) negative
below an Allee threshold \(N_A\), producing a phase
boundary at \(N=N_A\): the marginal growth rate
vanishes, response to environmental perturbations is
non-uniform, and the residual demographic-noise
status respects the declared deterministic-model
envelope \citep{Allee1931Aggregations}.

\item[\emph{Thom cusp catastrophe (catastrophe
theory).}]
The cusp catastrophe is the family of potentials
\(V_a(x)=x^4/4+a_1 x^2/2+a_2 x\) parameterised by
\((a_1,a_2)\in\mathbb R^2\); the phase readout
records the number of stable equilibria.  The cusp
manifold \(27 a_2^2+8 a_1^3=0\) is the phase
boundary: two stable equilibria merge with a saddle,
the energy margin vanishes, the response is
non-uniform, and the residual smooth-perturbation
status is classified by the elementary catastrophe
normal form \citep{Thom1972Stabilite}.

\item[\emph{{NBER} business-cycle peak/trough
(macroeconomics).}]
The control space is macroeconomic time series of
GDP, employment, industrial production, and personal
income; the phase readout records expansion vs
recession.  NBER-defined peaks and troughs are
phase boundaries: leading-indicator marginal growth
vanishes, sectoral responses are non-uniform, and
the residual status is handled by the NBER Business
Cycle Dating Committee's revision protocol over
noisy indicators \citep{BurnsMitchell1946Measuring}.

\item[\emph{Reynolds laminar-turbulent transition
(fluid dynamics).}]
For incompressible smooth-pipe flow, the control
space is the Reynolds number
\(\mathrm{Re}=\rho v D/\mu\); the phase readout
records laminar vs turbulent.  Reynolds's pipe-flow
experiments identify a critical
\(\mathrm{Re}_c\approx 2300\) as the phase
boundary at which laminar flow loses empirical
robustness to finite-amplitude perturbations under
the declared controlled-inflow protocol.  The
empirical robustness margin vanishes, the
intermittent-puff response is non-uniform, and the
residual status is described by transition
intermittency statistics
\citep{Reynolds1883Turbulent}.

\item[\emph{Hopf bifurcation onset of oscillation
(nonlinear dynamics).}]
For a parameter-dependent ODE
\(\dot{\mathbf x}=f(\mathbf x;\mu)\), the control
space is \(\mu\); the phase readout records ``stable
fixed point'' vs ``limit cycle''.  A Hopf bifurcation
at \(\mu=\mu_H\) occurs when a complex-conjugate
eigenvalue pair of the Jacobian crosses the
imaginary axis: the real part vanishes, oscillation
amplitude grows as \(\sqrt{\mu-\mu_H}\), and the
residual perturbation is statused by smooth
normal-form coefficients \citep{Hopf1942Abzweigung}.
\end{description}

\subsection{Locality / Domain of Dependence}\label{sec:self-reference-limits:locality}

The classical relativistic instance of a local quotient with a
domain-of-dependence structure is the causal structure of
Lorentzian spacetime; the standard reference for the
domain-of-dependence formalism is \citet[Chapter~8]{Wald1984GR}.

\begin{theorem}[Locality / Domain of Dependence]\label{thm:F45}
Let \(H\) be a history set with local quotient \(\ell:H\to L\),
target readout \(t:H\to T\), signature \(\sigma:H\to\Sigma\), update
map \(u:H\to H'\), and target quotient \(q':H'\to Q'\). Assume
\(\ell\) is surjective. Define
\[
  \mathcal O_\ell(t)=
  \{(h,h'):\ell(h)=\ell(h')\ \text{and}\ t(h)\ne t(h')\}.
\]
Then
\begin{align*}
  \operatorname{LocalToDomain}(\ell,t)
  &\iff \mathcal O_\ell(t)=\emptyset,\\
  \operatorname{TargetFamilyLocal}(\ell,\sigma)
  &\iff
  \forall a,b\in H,\ \ell(a)=\ell(b)\Rightarrow \sigma(a)=\sigma(b),\\
  \operatorname{DynamicLocality}(\ell,u,q')
  &\iff
  \forall a,b\in H,\ \ell(a)=\ell(b)\Rightarrow q'(u(a))=q'(u(b)),\\
  \operatorname{MinimalDomain}(\ell,\sigma)
  &\iff \operatorname{TargetFamilyLocal}(\ell,\sigma)\ \wedge\\
  &\qquad
  \neg\exists \ell',\ \ell'\!<\ell\ \wedge\
  \operatorname{TargetFamilyLocal}(\ell',\sigma).
\end{align*}
\end{theorem}
\footnote{\raggedright Verified in Lean as
\texttt{SelfReferenceLimits.\allowbreak{}Locality.\allowbreak{}locality\_domain\_of\_dependence};
see \Cref{app:formalization}.}
\begin{proof}
Assume the maps \(\ell,t,\sigma,u,q'\) and the surjectivity of
\(\ell\). If \(t=\bar t\circ\ell\), then
\[
  \ell(a)=\ell(b)\Rightarrow
  t(a)=\bar t(\ell(a))=\bar t(\ell(b))=t(b),
\]
so \(\mathcal O_\ell(t)=\emptyset\). Conversely, if
\(\mathcal O_\ell(t)=\emptyset\), define \(\bar t(\ell(h)):=t(h)\);
surjectivity and empty obstruction make this independent of the
representative. The family and dynamic clauses repeat the same
argument with \(t\) replaced by \(\sigma\) and \(q'\circ u\):
\[
  \ell(a)=\ell(b)\Rightarrow \sigma(a)=\sigma(b),
  \qquad
  \ell(a)=\ell(b)\Rightarrow q'(u(a))=q'(u(b)).
\]
Minimality is then the case split between \(\ell\) itself carrying the
family and some strictly smaller \(\ell'\) carrying it.
\end{proof}

\subsection*{Layer instantiations}\label{layer-inst:F45}

\begin{description}
\item[\emph{Markov property in stochastic processes
(probability theory).}]
For a stochastic process \((X_t)\) on a state space
\(S\), the history set is the path space; the local
quotient sends a path to its present state \(X_t\).
The target readout is the conditional future
distribution \(\mathcal L(X_{t+\Delta}\mid\cdot)\).
The Markov property states that
\(\mathcal L(X_{t+\Delta}\mid\mathcal F_t)
=\mathcal L(X_{t+\Delta}\mid X_t)\), i.e., the
future distribution descends through the
present-state quotient.  Within the declared
process and signature, the minimal domain is the
present state; lumpable chains may admit coarser
sufficient quotients in that admissibility class
\citep{Markov1906Investigation}.

\item[\emph{Sheaf locality on a topological space
(algebraic geometry).}]
For a sheaf \(\mathcal F\) on a topological space
\(X\), the history set consists of pairs
\((U,s)\) of an open subset and a section; the
local quotient sends \((U,s)\) to the family of
germs \(\{\mathcal F_x\}_{x\in U}\).  The target
readout is the section.  The sheaf condition
states that two sections agreeing on all germs
coincide, and any compatible family of germs glues
to a unique section: the section descends through
the germ-family quotient.  The minimal domain is
the germ / restriction cover required by the
sheaf's declared gluing data
\citep{Grothendieck1957Tohoku}.

\item[\emph{Cellular automaton local-update rule
(theoretical computer science).}]
For a cellular automaton on \(\mathbb Z^d\) with
finite state space, the history set is global
configurations; the local quotient sends a
configuration to its restriction to a radius-\(r\)
neighbourhood of a cell.  The local update rule
computes the next state from the radius-\(r\)
neighbourhood; the next-state readout descends
through the neighbourhood quotient.  Within the
declared rule's actual dependency radius, the
minimal domain is the radius-\(r\) neighbourhood
\citep{vonNeumann1966SelfReproducing}.

\item[\emph{Lamport happens-before relation
(distributed systems).}]
For a distributed system of asynchronous processes
exchanging messages, the history set is ordered
event pairs; the local quotient sends a pair
\((e_1,e_2)\) to its happens-before classification
under Lamport's relation \(\to\), the smallest
transitive relation satisfying: events within the
same process are ordered by occurrence, and every
message-send precedes its matching message-receive.
A pair is classified as \(e_1\to e_2\), \(e_2\to
e_1\), or \(e_1\parallel e_2\) (concurrent).  The
causal-ordering verdict descends through the
classifier; the minimal domain is the happens-before
classifier itself \citep{Lamport1978HappensBefore}.

\item[\emph{Voronoi cell as nearest-neighbour
domain (computational geometry).}]
For a finite set of seed points \(P\subset
\mathbb R^d\), the history set is the ambient
space; the local quotient sends each point to the
index of its nearest seed.  The local domain is the
Voronoi cell
\(V_i=\{x:\|x-p_i\|\le\|x-p_j\|\,\forall j\}\), and
the nearest-neighbour classification descends
through this quotient.  On interior points the
verdict is robust to small perturbations; on
Voronoi boundary edges and vertices the verdict is
degenerate under the declared tie-handling rule.
The minimal domain is the Voronoi-cell partition
\citep{Voronoi1908Nouvelles}.
\end{description}

\subsection{Law--State Split}\label{sec:self-reference-limits:law-state-split}

\begin{theorem}[Law--State Split]\label{thm:F46}
Let \(\mathcal X\) be a closure space with moduli quotient
\(\mu:\mathcal X\to M\), readout \(r:\mathcal X\to V\), and descended
readout \(\bar r:M\to V\) satisfying \(\bar r\circ\mu=r\). Let
\(\tau:\mathcal X\to\mathcal X_2\) be a transport and
\(\mu_2:\mathcal X_2\to M_2\) a target moduli quotient. Then
\begin{align*}
  r \text{ respects }\mu
  &\iff
  \forall x,y\in\mathcal X,\ \mu(x)=\mu(y)\Rightarrow r(x)=r(y),\\
  \operatorname{LawReadout}(r)
  &\iff
  \forall m,m'\in M,\ \bar r(m)=\bar r(m'),\\
  \operatorname{StateDependent}(r)
  &\iff \neg\operatorname{LawReadout}(r),\\
  \operatorname{LawReadout}(r)\vee\operatorname{StateDependent}(r),
  &\qquad
  \neg(\operatorname{LawReadout}(r)\wedge\operatorname{StateDependent}(r)).
\end{align*}
\end{theorem}
\footnote{\raggedright Verified in Lean as
\texttt{SelfReferenceLimits.\allowbreak{}LawStateSplit.\allowbreak{}law\_state\_split};
see \Cref{app:formalization}.}
\begin{proof}
Assume \(\mu\) is surjective and \(\bar r\circ\mu=r\). The first
equivalence follows by substitution:
\[
  \mu(x)=\mu(y)\Rightarrow
  r(x)=\bar r(\mu(x))=\bar r(\mu(y))=r(y).
\]
If \(\bar r\) is constant on \(M\), then every perturbation class has
the same value and \(r\) is a law-readout. If not, there exist
\(m,m'\in M\) with
\[
  \bar r(m)\ne\bar r(m'),
\]
so the value depends on the selected moduli class and the readout is
state-dependent. These two cases are complements by the displayed
definition, which gives both the disjunction and the impossibility of
their conjunction.
\end{proof}

\subsection*{Layer instantiations}\label{layer-inst:F46}

\begin{description}
\item[\emph{Clapeyron ideal gas law (thermodynamics).}]
For an ideal-gas system, the closure space is the
set of equilibrium thermodynamic states; the moduli
quotient sends a state to its
\((P,V,T,n)\) coordinates.  The readout
\(PV/(nT)\) equals the universal gas constant
\(R=8.314\;\mathrm{J\,mol^{-1}\,K^{-1}}\) on
every ideal-gas state, so the descended readout is
constant on the moduli space and is a law-readout
in the F46 sense.  The state-dependent readouts
(individual \(P\), \(V\), \(T\), \(n\)) vary across
moduli \citep{Clapeyron1834Memoire}.

\item[\emph{Hardy-Weinberg principle (population
genetics).}]
For a randomly mating diploid population with two
alleles at frequencies \(p,q\) under the
Hardy-Weinberg assumptions, the closure space is
the set of equilibrium populations and the moduli
quotient collapses populations to their
allele-frequency pair \((p,q)\).  The residual readout
\((f_{AA}-p^2,\,f_{Aa}-2pq,\,f_{aa}-q^2)\) equals
\((0,0,0)\) on every equilibrium population: the
descended readout is constant and is a law-readout.
The individual genotype frequencies
\(f_{AA},f_{Aa},f_{aa}\) vary across moduli and are
state-dependent \citep{Hardy1908Mendelian}.

\item[\emph{Pythagorean theorem (Euclidean
geometry).}]
For right triangles in Euclidean \(\mathbb R^2\),
the moduli quotient sends a triangle to its leg
lengths and hypotenuse \((a,b,c)\).  The Pythagorean
readout \(c^2-a^2-b^2\) equals \(0\) on every right
triangle; the descended readout is constant and is
a law-readout.  Triangle area \(ab/2\), interior
angles, and perimeter \(a+b+c\) are state-dependent
\citep{Euclid300BCElements}.

\item[\emph{Coulomb's law (electromagnetism).}]
For a system of point charges in vacuum, the moduli
quotient sends configurations to charge values and
pairwise distances.  The readout
\(F r_{12}^2/(q_1 q_2)\) equals the Coulomb constant
\(k_e=8.99\times 10^9\;\mathrm{N\,m^2\,C^{-2}}\)
on every vacuum point-charge configuration; the
descended readout is constant and is a law-readout.
Individual force magnitudes, charge values, and
distances are state-dependent
\citep{Coulomb1785Memoires}.

\item[\emph{Bayes' rule for probability updating
(probability theory).}]
For a probabilistic inference apparatus, the moduli
quotient sends an inference tuple to its prior /
likelihood / data coordinates.  The Bayes residual
readout \(P(H\mid D)\,P(D)-P(D\mid H)\,P(H)\)
equals zero on every valid tuple; the descended
readout is constant and is a law-readout.  The
specific posterior, prior, and data values are
state-dependent \citep{Bayes1763Essay}.
\end{description}

\subsection{Fine-Tuning as Small Selector Region}\label{sec:self-reference-limits:fine-tuning}

The canonical physical instance of a small target-compatible region
inside a moduli space is Weinberg's cosmological-constant analysis
\citep{Weinberg1989CosmologicalConstant}.  The layer-agnostic
statement below abstracts the smallness gate from that substrate
discussion.

\begin{theorem}[Fine-Tuning as Small Selector Region]\label{thm:F47}
Let \(M\) be a moduli space, \(r:M\to V\) a readout, and
\(T:V\to\mathrm{Prop}\) a target-region predicate.  Let
\(\mathcal M\) be the carrier type of measure-record values, let
\(\mu:\mathrm{Pow}(M)\to\mathcal M\) be a measure record on subsets
of \(M\), let \(\mu_{\mathrm{tot}}\in\mathcal M\) be the total
measure, let \(\theta\) be a threshold, and let
\(\odot:\Theta\times\mathcal M\to\mathcal M\) be the declared
threshold scaling.  Let \(\operatorname{Pos}:\mathcal M\to
\mathrm{Prop}\) and \({\le_{\mathcal M}}\subseteq
\mathcal M\times\mathcal M\) be the declared positivity and order
predicates on the measure record.  Let \(\mathrm{sel}\in M\) be a
selected modulus, let \(\mathsf{selDec}\) be a selector-declared
predicate, and let
\(\mathsf{modDec},\mathsf{tgtDec},\mathsf{meaDec},\mathsf{thrDec},
\mathsf{presStat}\) be the five required apparatus-declaration
predicates.  Define the selector region
\[
  \Sigma(r,T):=\{m\in M:T(r(m))\}.
\]
Then
\begin{align*}
  m\in\Sigma(r,T)
  &\iff
  T(r(m)),\\
  \operatorname{Small}(\Sigma(r,T))
  &\iff
  \operatorname{Pos}(\mu(\Sigma(r,T)))
  \wedge
  \mu(\Sigma(r,T))\le_{\mathcal M}\theta\odot\mu_{\mathrm{tot}},\\
  \operatorname{Realized}(\Sigma(r,T),\mathrm{sel},\mathsf{selDec})
  &\iff
  \mathsf{selDec}
  \wedge
  \mathrm{sel}\in\Sigma(r,T)
  \wedge
  \operatorname{Small}(\Sigma(r,T)),\\
  \operatorname{TargetImpossible}(\Sigma(r,T))
  &\Longrightarrow
  \neg\operatorname{Realized}(\Sigma(r,T),\mathrm{sel},\mathsf{selDec}),\\
  \operatorname{WellTyped}
  &\iff
  \mathsf{modDec}\wedge\mathsf{tgtDec}\wedge\mathsf{meaDec}
  \wedge\mathsf{thrDec}\wedge\mathsf{presStat}.
\end{align*}
\end{theorem}
\footnote{\raggedright Verified in Lean as
\texttt{SelfReferenceLimits.\allowbreak{}FineTuning.\allowbreak{}fine\_tuning\_small\_selector\_region};
see \Cref{app:formalization}.}
\begin{proof}
Assume the moduli space \(M\), readout \(r\), target predicate \(T\),
measure record \(\mu\), threshold \(\theta\) with scaling \(\odot\),
declared positivity \(\operatorname{Pos}\) and order
\(\le_{\mathcal M}\), selected modulus \(\mathrm{sel}\), and the
declaration predicates.  We prove the five conjuncts.

\emph{(1) Compatibility iff target preimage.}
The compatible region is the preimage of the target predicate under
the moduli readout:
\[
  \Sigma(r,T)=r^{-1}(T)=\{m\in M:T(r(m))\}.
\]
Hence \(m\in\Sigma(r,T)\) iff \(T(r(m))\), by the definition of
preimage.

\emph{(2) Smallness iff positive measure within the threshold gate.}
Smallness of the selector region is the conjunction of two apparatus
tests on the measure record \(\mu(\Sigma(r,T))\): positivity under
the declared \(\operatorname{Pos}\) predicate, and the order gate
\[
  \mu(\Sigma(r,T))\le_{\mathcal M}\theta\odot\mu_{\mathrm{tot}}.
\]
The displayed iff is the unfolding of this defining conjunction.

\emph{(3) Realized fine-tuning iff selected, compatible, and small.}
The realized-fine-tuning predicate is the conjunction of a selector
declaration \(\mathsf{selDec}\), the selected modulus belonging to
the compatible region, and smallness of that compatible region:
\[
  \mathsf{selDec}
  \wedge
  \mathrm{sel}\in\Sigma(r,T)
  \wedge
  \operatorname{Small}(\Sigma(r,T)).
\]

\emph{(4) Target-impossible blocks realized fine-tuning.}
If \(\operatorname{TargetImpossible}(\Sigma(r,T))\), then by
definition every \(m\in M\) satisfies
\[
  \neg(m\in\Sigma(r,T)).
\]
In particular \(\mathrm{sel}\notin\Sigma(r,T)\), contradicting the
second conjunct of \(\operatorname{Realized}\).  Hence
\[
  \neg\operatorname{Realized}(\Sigma(r,T),\mathrm{sel},\mathsf{selDec}).
\]

\emph{(5) Well-typedness iff the five declarations hold.}
Well-typedness is the conjunction of the five declaration predicates:
\[
  \operatorname{WellTyped}
  \iff
  \mathsf{modDec}\wedge\mathsf{tgtDec}\wedge\mathsf{meaDec}
  \wedge\mathsf{thrDec}\wedge\mathsf{presStat}.
\]
\end{proof}

\subsection*{Layer instantiations}\label{layer-inst:F47}

\begin{description}
\item[\emph{Circumstellar habitable zone (astrobiology).}]
The moduli space is the set of orbital and
atmospheric parameter values for a planet (semi-major
axis, eccentricity, atmospheric composition, greenhouse
mass column); the readout is the equilibrium surface
temperature; the target predicate restricts to the
liquid-water range.  The selector region is the
circumstellar habitable zone; the measure on
(orbital, atmospheric) parameter space shows it forms
a small fraction of the total.  Earth realises the
fine-tuning predicate within astrobiological criteria
\citep{Kasting1993Habitable}.

\item[\emph{Pharmaceutical therapeutic window
(clinical pharmacology).}]
The moduli space is the set of dose-and-timing
protocols; the readout is the steady-state plasma
concentration; the target predicate restricts to
concentrations between the minimum effective
concentration (MEC) and the maximum tolerated
concentration (MTC).  The therapeutic-window selector
region is small for drugs with narrow therapeutic
indices (warfarin, lithium, digoxin); realized
fine-tuning is the clinically attained dose protocol
placing plasma concentration within the window
\citep{GoodmanGilman2018Pharmacology}.

\item[\emph{Mechanical tolerance stack-up gap
(manufacturing engineering).}]
The moduli space is the set of component dimension
tuples within their tolerance bands; the readout is
the resulting critical assembly dimension; the
target predicate restricts to assemblies meeting the
critical-dimension specification.  The selector
region shrinks as the assembly chain lengthens;
realized fine-tuning is the actually manufactured
component-dimension tuple landing inside the
allowable region \citep{Bjorke1989Tolerance}.

\item[\emph{Hutchinson ecological niche (ecology).}]
The moduli space is the multi-dimensional
environmental parameter space (temperature,
humidity, pH, salinity, prey abundance); the readout
is the species' population growth rate; the target
predicate restricts to non-negative growth.  Within
the declared environmental envelope for the species,
the niche occupies a small fraction of total moduli
volume; realized fine-tuning is the species' actual
occupancy of the niche
\citep{Hutchinson1957Concluding}.

\item[\emph{Nyquist closed-loop stability region
(control theory).}]
The moduli space is the set of (plant,
controller-parameter) combinations; the readout is
the closed-loop pole location in the complex plane;
the target predicate restricts to all poles having
negative real part.  The Nyquist criterion
characterises the stability selector region: the
closed-loop system is stable iff the number of
counter-clockwise encirclements of \(-1+0i\) equals
the number of open-loop right-half-plane poles.  For
systems with large gain, delay, or open-loop
instability, the stability region can be small,
instantiating F47 fine-tuning of controller
parameters \citep{Nyquist1932Regeneration}.
\end{description}

\subsection{Topology as Global Gluing Invariant}\label{sec:self-reference-limits:topology}

The classical mathematical instance is sheaf theory, in which
topological information emerges as the consistent global gluing of
locally defined data; the standard reference is
\citet{MacLaneMoerdijk1992Sheaves}.

\begin{theorem}[Topology as Global Gluing Invariant]\label{thm:F48}
Let \(G\) be a set of global states with gluing quotient
\(\gamma:G\to\Gamma\), readout \(r:G\to V\), transport
\(\tau:G\to G'\), and target gluing quotient \(\gamma':G'\to\Gamma'\).
Assume \(\gamma\) is surjective, and define
\[
  \mathcal O_\gamma(r)=
  \{(g,g'):\gamma(g)=\gamma(g')\ \text{and}\ r(g)\ne r(g')\}.
\]
Then
\begin{align*}
  \operatorname{TopologicalReadout}(r)
  &\iff \mathcal O_\gamma(r)=\emptyset,\\
  \operatorname{CompleteInvariant}(r)
  &\iff \operatorname{TopologicalReadout}(r)\wedge
    \forall g,g',\ r(g)=r(g')\Rightarrow \gamma(g)=\gamma(g'),\\
  \operatorname{CoarseInvariant}(r)
  &\iff \operatorname{TopologicalReadout}(r)\wedge
    \exists g,g',\gamma(g)\ne\gamma(g')\wedge r(g)=r(g'),\\
  \operatorname{TransportDescends}(\tau)
  &\iff
    \exists\,\bar\tau:\Gamma\to\Gamma'\ \text{ with }\
    \bar\tau\circ\gamma=\gamma'\circ\tau .
\end{align*}
\end{theorem}
\footnote{\raggedright Verified in Lean as
\texttt{SelfReferenceLimits.\allowbreak{}Topology.\allowbreak{}topology\_global\_gluing\_invariant};
see \Cref{app:formalization}.}
\begin{proof}
Assume the gluing quotient \(\gamma\) is surjective. The first
equivalence is the quotient criterion:
\[
  r=\bar r\circ\gamma
  \quad\Longleftrightarrow\quad
  \forall g,g',\ \gamma(g)=\gamma(g')\Rightarrow r(g)=r(g')
  \quad\Longleftrightarrow\quad
  \mathcal O_\gamma(r)=\emptyset .
\]
If \(r\) also separates \(\Gamma\)-classes, it is complete; if two
distinct gluing classes have the same \(r\)-value, it is coarse. These
are the two cases
\[
  \forall g,g',\ r(g)=r(g')\Rightarrow\gamma(g)=\gamma(g')
  \quad\text{or}\quad
  \exists g,g',\gamma(g)\ne\gamma(g')\wedge r(g)=r(g').
\]
The transport clause is the same descent criterion applied to
\(\gamma'\circ\tau\).
\end{proof}

\subsection*{Layer instantiations}\label{layer-inst:F48}

\begin{description}
\item[\emph{Euler characteristic (algebraic
topology).}]
For a finite cell complex, the gluing quotient
collapses cell-complex presentations to their
underlying topological space.  The Euler
characteristic
\(\chi=V-E+F\) (alternating sum of cell counts in
higher dimensions) is invariant under
cell-decomposition refinement and homotopy, hence a
TopologicalReadout.  It is a coarse invariant: the
torus and the Klein bottle both have \(\chi=0\) but
the torus is orientable and the Klein bottle is
not.  Transport descends under homotopy
\citep{Euler1750Polyhedron}.

\item[\emph{Genus of a closed orientable surface
(topology).}]
For closed orientable 2-manifolds, the gluing
quotient collapses polygon-with-edge-identification
presentations to homeomorphism classes.  The genus
\(g(\Sigma)\) — the number of handles, equivalently
\((2-\chi)/2\) — is a TopologicalReadout, and the
classification theorem of closed surfaces makes
genus a CompleteInvariant for closed orientable
surfaces \citep{Riemann1857Theorie}.

\item[\emph{Fundamental group (algebraic topology).}]
For a topological space with basepoint, the gluing
quotient sends a space to its based homotopy type.
The fundamental group \(\pi_1(X,x_0)\) — homotopy
classes of based loops — is homotopy-invariant,
hence a TopologicalReadout.  It is coarse: the
2-sphere and 3-sphere both have trivial \(\pi_1\)
but are not homotopy-equivalent.  Continuous
basepoint-preserving maps induce \(\pi_1\)-compatible
group homomorphisms \citep{Poincare1895Analysis}.

\item[\emph{{DNA} linking number of closed circular
{DNA} (molecular biology).}]
For closed circular double-stranded DNA, the
gluing quotient collapses spatial configurations
related by smooth ambient isotopy in
\(\mathbb R^3\) (no strand breaks).  The linking
number \(\mathrm{Lk}\in\mathbb Z\) — the signed
Gauss linking integral counting strand wrappings
— is invariant under any smooth deformation that
does not cleave a strand, hence a
TopologicalReadout.  The Călugăreanu-White formula
\(\mathrm{Lk}=\mathrm{Tw}+\mathrm{Wr}\) decomposes
\(\mathrm{Lk}\) into non-topological twist and
writhe; only the sum is topological.
Topoisomerase enzymes implement the transport that
changes \(\mathrm{Lk}\) by strand cleavage and
religation \citep{BatesMaxwell2005DNATopology}.

\item[\emph{Reeb graph from {M}orse-function level
sets (geometric topology).}]
For a smooth manifold with a Morse function \(f\),
the gluing quotient collapses configurations
continuously deformable along the Morse flow.  The
Reeb graph quotients the manifold by the
``connected component of \(f^{-1}(c)\)'' relation,
encoding topology changes of level sets at
critical values.  Within the declared
function-isotopy class, the Reeb graph is a
TopologicalReadout, and a coarse invariant
(distinct Morse functions can yield isomorphic
Reeb graphs) \citep{Reeb1946Points}.
\end{description}

\subsection{Common-Source / Nonlocal Correlation Normal Form}\label{sec:self-reference-limits:common-source}

The famous physical instance of correlations that admit no local
common-source explanation is the Einstein--Podolsky--Rosen / Bell
situation \citep{Bell1964EPR}; the layer-agnostic statement below
records the descent-to-common-source obstruction abstracted from
that substrate.

\begin{theorem}[Common-Source / Nonlocal Correlation Normal Form]\label{thm:F49}
Let \(H\) be a history space with source quotient \(\sigma:H\to S\),
shared interface \(\iota:H\to I\), access quotients
\(q_A:H\to L_A\), \(q_B:H\to L_B\), and readouts
\(r_A:H\to A\), \(r_B:H\to B\). Let
\(\kappa:H\to K\) be the correlation readout. Then
\begin{align*}
  \operatorname{CommonSource}(\sigma)
  &\iff
  \exists\,\bar r_A:S\to A,\bar r_B:S\to B,\
    \bar r_A\circ\sigma=r_A\ \wedge\ \bar r_B\circ\sigma=r_B,\\
  \operatorname{InterfaceFactorization}(\iota)
  &\iff
  \exists\,\hat r_A:I\to A,\hat r_B:I\to B,\
    \hat r_A\circ\iota=r_A\ \wedge\ \hat r_B\circ\iota=r_B,\\
  \operatorname{LocallyExplainable}(\kappa)
  &\iff \operatorname{CommonSource}(\sigma)
    \vee \operatorname{InterfaceFactorization}(\iota).
\end{align*}
If neither disjunct holds, the correlation is a statused nonlocal
residual.
\end{theorem}
\footnote{\raggedright Verified in Lean as
\texttt{SelfReferenceLimits.\allowbreak{}CommonSourceNonlocal.\allowbreak{}common\_source\_nonlocal\_correlation};
see \Cref{app:formalization}.}
\begin{proof}
Assume the source quotient \(\sigma\), interface \(\iota\), accesses
\(q_A,q_B\), and readouts \(r_A,r_B,\kappa\). If
\(\bar r_A\circ\sigma=r_A\) and \(\bar r_B\circ\sigma=r_B\), then both
local readouts depend only on the shared source class:
\[
  \sigma(h)=\sigma(h')\Rightarrow
  r_A(h)=r_A(h')\ \wedge\ r_B(h)=r_B(h').
\]
This is a common-source explanation. The interface case is identical
with \(\iota\) replacing \(\sigma\). Conversely, any explanation by a
common source or interface gives exactly such descended readouts. Thus
there are two local-explanation cases,
\[
  \operatorname{CommonSource}(\sigma)
  \quad\text{or}\quad
  \operatorname{InterfaceFactorization}(\iota),
\]
and failure of both leaves the correlation as a nonlocal residual.
\end{proof}

\subsection*{Layer instantiations}\label{layer-inst:F49}

\begin{description}
\item[\emph{{MISTRA} twin studies (behavioral
genetics).}]
For the Minnesota Study of Twins Reared Apart, the
source quotient sends a twin pair to its
shared-genotype class; the readouts are conditional
trait
distributions
\(P(\mathrm{trait}\mid\sigma)\) under a declared
quantitative-genetic model.  The CommonSource
hypothesis holds when conditional distributions
descend through the shared-genotype quotient; the
interface captures shared rearing environment
(absent for the reared-apart sample).  The Falconer
formula \(h^2=2(r_{MZ}-r_{DZ})\) operationalises the
F49 LocallyExplainable decomposition
\citep{Bouchard1990Sources}.

\item[\emph{Doolittle lateral gene transfer
(microbial phylogenetics).}]
For prokaryotic phylogenomics, the source quotient
sends a gene-distribution configuration to its
vertical-inheritance most-recent-common-ancestor
class; the readouts are per-gene tree topologies
and presence patterns.  Vertical inheritance gives
CommonSource correlations across co-inherited
genes; Doolittle observed that horizontal gene
transfer events mediated by phages, plasmids, and
conjugation break the common-source pattern,
constituting an interface channel.  Locally
Explainable correlations decompose into vertical
(common source) and lateral (interface)
contributions \citep{Doolittle1999Phylogenetic}.

\item[\emph{Plagiarism detection via shared-source
fingerprinting (digital forensics).}]
For pairs of documents, the source quotient sends
a pair to its shared-source text together with
declared edit / paraphrase / template
transformations; each document is fingerprinted via
n-gram or hashed-shingle representations; the
correlation readout is an overlap metric (Jaccard
similarity, MinHash).  CommonSource holds when the
overlap descends through the
shared-source-plus-declared-transformations
quotient; the interface
captures legitimate citation, common literature,
and formulaic templates not classified as
plagiarism \citep{Heintze1996Scalable}.

\item[\emph{Mosteller-Wallace stylometric authorship
attribution (statistical humanities).}]
For the disputed Federalist Papers, the source
quotient assigns a paper to a candidate author; the
readouts are author-conditioned function-word
frequency distributions
\(P(\mathrm{wordfreq}\mid\sigma)\) estimated from
each candidate's known corpus.  The Bayesian
posterior \(P(\sigma\mid\mathrm{wordfreq})\)
operationalises the CommonSource attribution; joint
authorship or formulaic shared templates
instantiate the interface, and residual anomalies
index unmodeled influences
\citep{MostellerWallace1964Federalist}.

\item[\emph{Common-cause failure analysis in
fault trees (safety engineering).}]
For a safety-critical system, the source quotient
sends a failure event to its underlying
common-cause basic event (loss of off-site power,
fire,
earthquake).  Under nominal independence, redundant
subsystem failure readouts are uncorrelated;
observed correlation indicates a common-cause
failure mode mediated by the source quotient.  Fault-tree analysis
identifies the basic events that constitute the
source, codifying methodology for nuclear-plant
reliability and risk assessment
\citep{Vesely1981FaultTree}.
\end{description}

\subsection{Vacuum / Background Selection}\label{sec:self-reference-limits:vacuum}

The most discussed contemporary substrate instance is the string
landscape of vacua surveyed in \citet{Susskind2003Landscape};
the layer-agnostic statement below records the selector-singleton
shape without committing to that physical realization.

\begin{theorem}[Vacuum / Background Selection]\label{thm:F50}
Let \(\mathcal X\) be a closure space with moduli quotient
\(\mu:\mathcal X\to M\), raw background \(b_0:\mathcal X\to B\),
descended background \(b:M\to B\), target \(t:M\to T\), and vacuum
criterion \(V\subseteq M\). Assume \(\mu\) is surjective. Then
\begin{align*}
  b_0 \text{ descends through }\mu
  &\iff
  \exists\,b:M\to B\ \text{ with } b\circ\mu=b_0,\\
  \operatorname{NecessaryBackground}(b)
  &\iff \forall m,m'\in M,\ b(m)=b(m'),\\
  \operatorname{BackgroundModuli}(b)
  &\iff \exists m,m'\in M,\ b(m)\ne b(m'),\\
  \operatorname{Vacuum}(m_\ast)
  &\iff V=\{m_\ast\}\ \wedge\ t(m_\ast)\ \text{is formed}.
\end{align*}
The necessary/moduli alternatives are exhaustive and disjoint.
\end{theorem}
\footnote{\raggedright Verified in Lean as
\texttt{SelfReferenceLimits.\allowbreak{}VacuumBackgroundSelection.\allowbreak{}vacuum\_background\_selection};
see \Cref{app:formalization}.}
\begin{proof}
Assume \(\mu\) is surjective. The descent clause follows from the
same representative-independence calculation:
\[
  \mu(x)=\mu(y)\Rightarrow b_0(x)=b_0(y)
  \quad\Longleftrightarrow\quad
  \exists b,\ b\circ\mu=b_0 .
\]
For the descended background \(b\), either \(b\) is constant on \(M\)
or there exist \(m,m'\) with \(b(m)\ne b(m')\):
\[
  \bigl(\forall m,m', b(m)=b(m')\bigr)
  \quad\vee\quad
  \bigl(\exists m,m', b(m)\ne b(m')\bigr).
\]
These alternatives are mutually exclusive. A vacuum is the singleton
case \(V=\{m_\ast\}\) together with formedness of the selected target;
totality and single-valuedness are exactly this singleton condition.
\end{proof}

\subsection*{Layer instantiations}\label{layer-inst:F50}

\begin{description}
\item[\emph{Waddington epigenetic landscape
(developmental biology).}]
For a developing embryo, the moduli quotient
collapses stem-cell trajectory configurations
sharing the same cell-type endpoint; the
descended background records the cell-fate class.
Waddington's epigenetic landscape depicts cell-fate
decision as a ball rolling down a branching
landscape: each branching point is a vacuum-style
selection of one cell-fate from a set of accessible
attractors \citep{Waddington1957Strategy}.

\item[\emph{Ostwald rule of stages: polymorph
selection (crystallography).}]
For a supersaturated solution, the moduli quotient
collapses crystallisation trajectories sharing the
same crystalline polymorph; the descended
background records the polymorph type.  Ostwald's
rule of stages predicts kinetic selection of a
metastable polymorph close in free energy to the
supersaturated state.  The vacuum predicate
identifies the ultimately stable polymorph as the
singleton target-formed modulus
\citep{Ostwald1897Studien}.

\item[\emph{Zermelo axiomatisation (foundations
of mathematics).}]
The closure space is the set of formal theories
purporting to be foundations of mathematics; the
moduli quotient sends a theory to its axiom-system
class (Zermelo 1908 set theory,
Fraenkel-Skolem-extended ZF, ZFC, NBG, intuitionistic
ZF,
Martin-Löf type theory).  Zermelo's 1908
axiomatisation selected a specific moduli value —
Zermelo set theory with separation and choice —
as a singleton background; subsequent extensions
(Fraenkel Replacement, Foundation, large cardinals)
refine that selection within the same F50
background-moduli framework
\citep{Zermelo1908Untersuchungen}.

\item[\emph{Hund's rules for atomic ground state
(atomic physics).}]
For a multi-electron atom, the moduli quotient
sends Pauli-allowed configurations to their
\((S,L,J)\) term-symbol class; the descended
background is the ground-state term symbol.
Hund's rules select the ground state: (i) maximize
total \(S\); (ii) for given \(S\), maximize \(L\);
(iii) for given \((S,L)\), minimize \(J\) if the
subshell is less than half-filled, otherwise
maximize \(J\).  Within the LS-coupling regime, the
vacuum predicate is the Hund-selected ground term
symbol \citep{Hund1925Linienspektren}.

\item[\emph{Pittendrigh circadian phase entrainment
(chronobiology).}]
For a circadian oscillator under a periodic
zeitgeber, the moduli quotient sends oscillator
trajectories sharing the same entrained phase to a
class; the descended background records the
entrained phase.  Pittendrigh's phase-response-curve
analysis explains how zeitgebers select a definite
entrained phase from the continuous phase moduli
space; the vacuum predicate is the biologically
realised entrained phase
\citep{Pittendrigh1960Circadian}.
\end{description}

\subsection{Unification as Common Refinement}\label{sec:self-reference-limits:unification}

The standard particle-physics instance is the electroweak unification
introduced in \citet{Weinberg1967Leptons}.  The layer-agnostic
statement below records the common-refinement shape that any such
unification realizes.

\begin{theorem}[Unification as Common Refinement]\label{thm:F51}
Let \(\mathcal C_U,\mathcal C_A,\mathcal C_B\) be parent and child
closures with quotients \(q_U:\mathcal C_U\to Q_U\),
\(q_A:\mathcal C_A\to Q_A\), and \(q_B:\mathcal C_B\to Q_B\). Let
\(\iota_A:\mathcal C_U\to\mathcal C_A\) and
\(\iota_B:\mathcal C_U\to\mathcal C_B\) be interpretations, and let
\(\pi_A:Q_U\to Q_A\), \(\pi_B:Q_U\to Q_B\) be quotient projections.
For claim sets \(K_A,K_B,K_U\), lifts
\(\lambda_A:K_A\to K_U\), \(\lambda_B:K_B\to K_U\), and status maps
\(s_A,s_B,s_U\), unification is equivalent to
\begin{align*}
  \pi_A\circ q_U=q_A\circ\iota_A
  &\quad\wedge\quad
  \pi_B\circ q_U=q_B\circ\iota_B,\\
  \pi_A\circ R_U^A=R_A\circ\pi_A
  &\quad\wedge\quad
  \pi_B\circ R_U^B=R_B\circ\pi_B,\\
  s_U(\lambda_A k)=\alpha_A(s_A k)
  &\quad\wedge\quad
  s_U(\lambda_B \ell)=\alpha_B(s_B \ell)
  \quad\text{for all }k\in K_A,\ell\in K_B .
\end{align*}
\end{theorem}
\footnote{\raggedright Verified in Lean as
\texttt{SelfReferenceLimits.\allowbreak{}Unification.\allowbreak{}unification\_as\_common\_refinement};
see \Cref{app:formalization}.}
\begin{proof}
Assume the parent and child quotients, interpretations, projections,
claim lifts, and status maps. If a unified parent exists, projecting
objects to the two children gives
\[
  \pi_A(q_U x)=q_A(\iota_A x),
  \qquad
  \pi_B(q_U x)=q_B(\iota_B x),
\]
so \(Q_U\) is a common refinement. Compatibility of routes gives the
two commutative squares in the theorem, and compatibility of claim
lifts gives the displayed status equations. Conversely, if these
equations hold, \(Q_U\) carries both child object quotients, both route
systems commute through the projections, and lifted claims receive
compatible statuses. Objectivity gives the common quotient, and
composability gives the parent audit ledger; hence the two theories
unify.
\end{proof}

\subsection*{Layer instantiations}\label{layer-inst:F51}

\begin{description}
\item[\emph{Maxwell unification of electromagnetism
(classical physics).}]
The parent closure is Maxwell's electromagnetic
field theory; the children are pre-Maxwellian
electrostatics (Coulomb-Gauss) and magnetostatics
(Amp\`ere-Biot-Savart).  Setting time derivatives to
zero in Maxwell's equations yields each child as a
specialisation.  The unified parent predicts
electromagnetic waves at the speed of light,
identifying light as electromagnetic — a
prediction not derivable from either child alone
\citep{Maxwell1865Dynamical}.

\item[\emph{Cartesian analytic geometry
(mathematics).}]
The parent closure is Cartesian analytic geometry;
the children are Euclidean geometry (synthetic
proofs) and algebra (symbolic equations).  Cartesian
coordinates supply interpretations in both
directions: every coordinate-representable
algebraic curve maps to a geometric object and
vice versa.  Euclidean theorems become algebraic
identities on coordinates, and algebraic equations
acquire geometric interpretation as curves and
surfaces \citep{Descartes1637Geometrie}.

\item[\emph{Banach functional analysis (functional
analysis).}]
The parent closure is Banach's framework of complete
normed vector spaces with bounded linear operators;
the children are finite-dimensional linear algebra
and classical function-space analysis.  Both embed
as concrete Banach spaces, and theorems
(Hahn-Banach, open mapping, closed graph, uniform
boundedness) supply unified results across both
children within the framework's declared scope
\citep{Banach1932Theorie}.

\item[\emph{L\'evi-Strauss structural anthropology
(anthropology).}]
The parent closure is L\'evi-Strauss's structural
analysis of myth via mythemes (binary oppositions,
mediating figures); the children are
ethnographically distinct mythological corpora.
Each corpus is interpreted in terms of its
structural elements, and the parent identifies
mythematic isomorphisms across geographically
isolated cultures \citep{LeviStrauss1958Anthropologie}.

\item[\emph{Huxley modern synthesis (evolutionary
biology).}]
The parent closure is the modern synthesis (Huxley
1942), unifying Mendelian genetics and Darwinian
natural selection together with cytology,
paleontology, and biogeography.  Population
genetics supplies the common refinement:
allele-frequency dynamics under Mendelian
inheritance extend, with selection coefficients, to a
unified
predictive framework recovering both Mendelian
ratios and Darwinian adaptive change as special
cases \citep{Huxley1942Evolution}.
\end{description}

\subsection{Closure Completeness Boundary}\label{sec:self-reference-limits:closure-completeness}

The proof-theoretic prototype of a closure-completeness boundary is
G\"odel's incompleteness theorem \citep{Godel1931Unentscheidbare},
and the modal-logic systematization of provability obstructions is
collected in \citet{Boolos1993LogicProvability}.  The layer-agnostic
meta-coherence statement below records the apparatus-side scope
discipline that complements those classical results.

\begin{definition}[Catalog boundary]\label{def:catalog-boundary}
The catalog boundary is the audited frontier between \FIV\ laws
covered by the layer-agnostic language, laws reducible to covered
normal forms, and claims explicitly marked out of scope because they
require substrate-specific data or stronger apparatus.
\end{definition}

\begin{theorem}[Closure Completeness Boundary]\label{thm:F52}
Let \(R\) be a residual type with predicates
\(\operatorname{hasStatus}:R\to\mathrm{Prop}\) and
\(\operatorname{adequate}:R\to\mathrm{Prop}\). Let \(D\) be a direction
type with predicates \(\operatorname{inScope}:D\to\mathrm{Prop}\),
\(\operatorname{claimed}:D\to\mathrm{Prop}\), and
\(\operatorname{namedOutside}:D\to\mathrm{Prop}\), and let
\[
  v\in\{\operatorname{accepted},\operatorname{open},\operatorname{refused}\}
\]
be a meta-audit verdict. Then
\begin{align*}
  \operatorname{ResidualCoverage}
  &\iff
  \forall r\in R,\ \operatorname{hasStatus}(r)\wedge\operatorname{adequate}(r),\\
  \operatorname{MetaCovered}(v)
  &\iff v=\operatorname{accepted},\\
  \operatorname{ScopedClosureComplete}
  &\iff
  \operatorname{formed}\wedge\operatorname{scopeDeclared}\wedge
  \operatorname{ResidualCoverage}\\
  &\qquad\wedge\
  \forall d\in D,\ \operatorname{claimed}(d)\Rightarrow
    \bigl(\operatorname{inScope}(d)\vee\operatorname{namedOutside}(d)\bigr),\\
  \operatorname{MetaComplete}
  &\iff \operatorname{ScopedClosureComplete}\wedge\operatorname{MetaCovered}(v).
\end{align*}
\end{theorem}
\footnote{\raggedright Verified in Lean as
\texttt{SelfReferenceLimits.\allowbreak{}ClosureCompletenessBoundary.\allowbreak{}closure\_completeness\_boundary};
see \Cref{app:formalization}.}
\begin{proof}
Assume formedness, declared scope, residual predicates, direction
predicates, and meta-audit verdict \(v\). Residual coverage means two
obstruction sets are empty:
\[
  \{r:\neg\operatorname{hasStatus}(r)\}=\emptyset,
  \qquad
  \{r:\operatorname{hasStatus}(r)\wedge\neg\operatorname{adequate}(r)\}
  =\emptyset .
\]
This is equivalent to the displayed universal statement over \(R\).
For directions, there are three cases:
\[
  \operatorname{inScope}(d),\qquad
  \operatorname{namedOutside}(d),\qquad
  \neg\operatorname{inScope}(d)\wedge\neg\operatorname{namedOutside}(d).
\]
Scoped completeness excludes the third case for every claimed
direction. Meta-coverage is exactly the accepted verdict. Combining
residual coverage, direction coverage, formedness, and accepted
meta-audit gives meta-completeness; omitting any conjunct gives the
corresponding boundary failure.
\end{proof}

\subsection*{Layer instantiations}\label{layer-inst:F52}

\begin{description}
\item[\emph{{CONSORT} statement for randomised-trial
reporting (clinical research).}]
The residuals are the methodological items required
for adequate trial reporting; directions are
declarable claims and outcomes.  CONSORT 2010
requires 25 explicit reporting items, with primary
and secondary outcomes \emph{inScope} and
exploratory analyses \emph{namedOutside}.  Residual
coverage holds when every CONSORT item has status
and is adequate; the meta-audit is journal-editor
acceptance \citep{Schulz2010CONSORT}.

\item[\emph{{PMBOK} project-scope management
(project management).}]
The residuals are deliverables and acceptance
criteria; directions are project activities and
stakeholder requests.  The PMBOK scope statement
declares deliverables (\emph{inScope}), exclusions
(\emph{namedOutside}), assumptions, constraints,
and acceptance criteria; scope creep is a
stakeholder request that is neither.  The
meta-audit is project-sponsor sign-off
\citep{PMI2017PMBOK}.

\item[\emph{International Building Code
jurisdictional scope (construction codes).}]
The residuals are code-mandated safety,
structural, fire-protection, and accessibility
requirements; directions are declarable
construction features.  The IBC declares scope by
occupancy classification, construction type, and
jurisdictional adoption, with historic, agricultural,
and accessory-structure exemptions
\emph{namedOutside}.  The meta-audit is the
building inspector's certificate of occupancy
\citep{ICC2021IBC}.

\item[\emph{Patent-claim scope under {M}arkman
hearings (intellectual-property law).}]
The residuals are patent-claim terms requiring
construction; directions are declarable claim
scopes and accused infringing embodiments.  Under
\emph{Markman v. Westview Instruments}, claim
construction is a matter of law: each disputed term
is construed (\emph{inScope}) or held indefinite
under \S112 (\emph{namedOutside} the valid claim).
The meta-audit is the Federal Circuit affirmance of
district-court claim construction
\citep{Markman1996Westview}.

\item[\emph{{FDA} drug-label indication scope
(regulatory affairs).}]
The residuals are safety and efficacy items required
on the FDA-approved label; directions are medical
uses.  21 CFR 201.57 prescribes the format and
content of the ``Indications and Usage'' section,
with FDA-approved indications \emph{inScope} and
off-label uses \emph{namedOutside}.  The meta-audit
is FDA approval of the (Supplemental) New Drug
Application \citep{FDA21CFR20157}.
\end{description}

\subsection*{Self-Reference and Limits at a glance}\label{sec:self-reference-limits:summary}
\begin{table}[p]
  \centering
  \footnotesize
  \setlength{\tabcolsep}{4pt}
  \textbf{Limits}
  \smallskip

  \begin{tabularx}{\textwidth}{@{}
      >{\raggedright\arraybackslash}p{0.24\textwidth}
      >{\raggedright\arraybackslash}X
      >{\raggedright\arraybackslash}p{0.16\textwidth}
      >{\raggedright\arraybackslash}p{0.10\textwidth}@{}}
    \toprule
    Paper-prose name & One-line claim & Status & Substrate cite \\
    \midrule
    Horizon-Sufficiency Theorem &
    Finite horizon data suffice for closure decisions about
    horizon-bounded objects. &
    Theorem &
    \textemdash \\
    Information-Loss Theorem &
    Legitimate compression loses information only along access fibers. &
    Theorem &
    \textemdash \\
    Singular-Limit Theorem &
    Singular completions are legitimate exactly when completion data
    descends and audits consistently. &
    Theorem &
    \textemdash \\
    Locality Theorem &
    A quantity depends only on the declared domain whose transports
    reach the target. &
    Theorem &
    \textemdash \\
    \bottomrule
  \end{tabularx}

  \medskip
  \textbf{Emergence}
  \smallskip

  \begin{tabularx}{\textwidth}{@{}
      >{\raggedright\arraybackslash}p{0.24\textwidth}
      >{\raggedright\arraybackslash}X
      >{\raggedright\arraybackslash}p{0.16\textwidth}
      >{\raggedright\arraybackslash}p{0.10\textwidth}@{}}
    \toprule
    Paper-prose name & One-line claim & Status & Substrate cite \\
    \midrule
    Scale-Descent Theorem &
    Fine-scale laws descend through a scale chain exactly when they
    factor through each coarser quotient. &
    Theorem &
    \textemdash \\
    Continuum-Emergence Theorem &
    Continuum limits form from cofinal refinement chains with
    convergent audit data. &
    Theorem &
    \textemdash \\
    Fine-Tuning Theorem &
    Fine-tuning is smallness of the selector region under a declared
    parameter measure. &
    Theorem &
    \textemdash \\
    Vacuum Theorem &
    A vacuum is the unique total background selection over the declared
    closure orbit. &
    Theorem &
    \textemdash \\
    Unification Theorem &
    Theories unify exactly through a common refinement with composable
    audit ledgers. &
    Theorem &
    \textemdash \\
    \bottomrule
  \end{tabularx}

  \medskip
  \textbf{Invariants}
  \smallskip

  \begin{tabularx}{\textwidth}{@{}
      >{\raggedright\arraybackslash}p{0.24\textwidth}
      >{\raggedright\arraybackslash}X
      >{\raggedright\arraybackslash}p{0.16\textwidth}
      >{\raggedright\arraybackslash}p{0.10\textwidth}@{}}
    \toprule
    Paper-prose name & One-line claim & Status & Substrate cite \\
    \midrule
    Complementarity Theorem &
    Complementarity is the failure of a declared joint refinement in
    the apparatus quotient lattice. &
    Theorem &
    \textemdash \\
    Arrow Theorem &
    The arrow is the canonical direction of monotone record growth. &
    Theorem &
    \textemdash \\
    Entropy Theorem &
    Entropy is the quotient-invariant log-volume of an access fiber. &
    Theorem &
    \textemdash \\
    Irreducibility Theorem &
    Irreducibility is the audit-ledger lower bound on closed
    description length. &
    Theorem &
    \textemdash \\
    Criticality Theorem &
    Criticality is the locus where the apparatus quotient lattice
    changes regime. &
    Theorem &
    \textemdash \\
    Law--State Split Theorem &
    Laws are perturbation-stable predicates; states are
    perturbation-contingent data. &
    Theorem &
    \textemdash \\
    Topology Theorem &
    Topology is the invariant content of the declared global gluing
    quotient. &
    Theorem &
    \textemdash \\
    Common-Source Theorem &
    Nonlocal correlations are explained by common sources exactly when
    they factor through a common refinement. &
    Theorem &
    \textemdash \\
    Closure Completeness Theorem &
    The catalog is complete for its declared scope when residuals and
    exclusions are audited. &
    Theorem &
    \textemdash \\
    \bottomrule
  \end{tabularx}

  \caption{Self-Reference and Limits at a glance: the eighteen laws of \Cref{sec:self-reference-limits}, grouped by sub-axis.}
  \label{tab:axis-self-reference-limits-summary}
\end{table}
  \section{Closure--Reality Boundary anchor}\label{sec:anchor:closure-reality-boundary}

The final row of the catalog is not another internal normal form.  It
marks the boundary between formal closure-of-record and bare physical
realisability.  A closure-of-record is a carrier whose closure data,
residuals, statuses, and nonclaims are audited inside the apparatus.
Bare physical realisability only says that some physical or
operational system instantiates a model in an ordinary sense.  The
Closure--Reality Boundary asks whether these two notions can be
identified.  They cannot be identified without further structure.

This boundary gets its own short chapter because its positive
replacement is anchored in the apparatus of reflective closure in
\PaperAOR\ \citep{Tsiokos2026AOR}.  At the layer-agnostic level,
\FIV\ records the normal form: unrestricted closure-equals-reality
fails in both directions, and the correct replacement is a reflective
operational notion in which realisability is already typed, bounded,
recorded, and audited.  The substrate paper supplies the
carrier-specific reflective apparatus; this chapter states the
boundary theorem that the catalog imports.

Let \(C\) be a carrier equipped with a declared scope
\(\mathfrak S\), recording the targets, routes, instruments,
interfaces, residuals, budgets, and nonclaims under discussion.  The
naive closure-equals-reality principle would identify bare physical
realisability of \(C\) with \(C\) being formally closed over
\(\mathfrak S\).  The direct direction says that physical
realisability implies closure-of-record.  The converse says that
closure-of-record implies physical realisability.  The two theorems
below isolate why neither unrestricted direction is valid.

\begin{definition}[Reflective closure operator]\label{def:reflective-closure}
The AOR apparatus assigns to a carrier \(C\) over a scope
\(\mathfrak S\) a reflective closure \(\ClSB(C)\), a unit
\[
  \etaSB(C):C\longrightarrow \ClSB(C),
\]
and a residual \(\XiSB(C)\).  The object \(\ClSB(C)\) is the audited
closure completion of \(C\); the unit compares the raw carrier with
that completion; and \(\XiSB(C)\) records what must be added,
repaired, budgeted, weakened, or excluded to pass from \(C\) to the
closed carrier over \(\mathfrak S\).
\end{definition}

\begin{theorem}[Direct CER no-go]\label{thm:CER-direct-no-go}
Let \(\mathfrak S\) be a declared scope, let \(\mathcal C\) be a
class of carriers, let
\(\operatorname{BareReal}_{\mathfrak S}:\mathcal C\to\mathrm{Prop}\)
be bare physical realisability, and let
\(\mathrm{Cl}_{\mathfrak S}:\mathcal C\to\mathcal C\) be the formal
closure operator over \(\mathfrak S\).  Under the reduct-obstruction
hypothesis there exists a carrier \(C\in\mathcal C\) such that
\[
  \operatorname{BareReal}_{\mathfrak S}(C)
  \quad\text{and}\quad
  C\ne \mathrm{Cl}_{\mathfrak S}(C).
\]
Equivalently, the direct CER implication
\[
  \operatorname{BareReal}_{\mathfrak S}(C)
  \Longrightarrow
  C=\mathrm{Cl}_{\mathfrak S}(C)
\]
fails in general.
\end{theorem}
\footnote{The Lean encoding records the direct no-go as a derived
consequence of an opaque reduct-obstruction assumption; see
\Cref{app:formalization}.}
\begin{proof}
Assume the scope \(\mathfrak S\), the predicate
\(\operatorname{BareReal}_{\mathfrak S}\), the closure operator
\(\mathrm{Cl}_{\mathfrak S}\), and the reduct-obstruction
hypothesis.  The hypothesis supplies a rich carrier \(C^+\) and a
reduct \(C\) obtained by forgetting a closure-relevant datum \(d\)
while retaining a target claim whose audit uses \(d\).  Thus
\[
  \operatorname{BareReal}_{\mathfrak S}(C^+)
  \quad\text{and}\quad
  C=\operatorname{Reduct}(C^+,d).
\]
The reduct is still bare-realised, since it is a coarser abstraction
of the same realised system:
\[
  \operatorname{BareReal}_{\mathfrak S}(C).
\]
But \(C\) is not a closure-of-record.  The retained claim has a
native residual whose required record or status was precisely the
forgotten datum \(d\), so closure adds that missing datum:
\[
  C\ne \mathrm{Cl}_{\mathfrak S}(C).
\]
This \(C\) is the displayed witness to failure of the direct CER
implication.
\end{proof}

\subsection*{Layer instantiations}\label{layer-inst:CER-direct}

\begin{description}
\item[\emph{Common Rule informed-consent record
(medical research ethics).}]
The carrier \(C\) is the clinical procedure as
performed; the scope \(\mathfrak S\) imposes the
Common Rule audit (protocol, IRB approval,
informed consent, adverse-event reporting).  The
closure-relevant datum \(d\) is the signed
consent form with its IRB version history.  The
procedure can be bare-realised (physically
performed) while \(d\) is missing, so the reduct
satisfies \(\operatorname{BareReal}_{\mathfrak S}\)
but is not a closure-of-record: the retained
claim ``subject enrolled under IRB-approved
protocol'' has a residual whose required record
is the missing consent \citep{DHHS2018CommonRule}.

\item[\emph{Preclinical reproducibility record
(research methodology).}]
The carrier \(C\) is the reported experimental
result; \(\mathfrak S\) imposes the
reproducibility audit (lab notebook, raw data,
randomisation log, blinded scoring).  The
closure-relevant datum \(d\) is the
replication-enabling record.  Begley and Ellis
report that the bench experiments behind
``landmark'' preclinical findings were
bare-realised yet the audit record required to
close the published claim was inadequate, so the
reduct satisfies \(\operatorname{BareReal}
_{\mathfrak S}\) but the reproducibility
residual has no satisfying record
\citep{Begley2012Drug}.

\item[\emph{Software bill of materials (software
supply chain).}]
The carrier \(C\) is the compiled binary in
production; \(\mathfrak S\) imposes the
supply-chain audit (component inventory, licence
terms, vulnerability mapping, build provenance).
The closure-relevant datum \(d\) is the SBOM
listing every linked component with version,
hash, and origin per the NTIA minimum elements.
The binary runs without the SBOM, so the reduct
satisfies \(\operatorname{BareReal}_{\mathfrak S}\)
yet the retained claim ``free of
known-vulnerable components'' has a residual
whose audit datum (the SBOM) was never produced
\citep{NTIA2021SBOM}.

\item[\emph{Photometric flux without calibration
log (observational astronomy).}]
The carrier \(C\) is the recorded photon count on
a CCD; \(\mathfrak S\) imposes the photometric
audit (flat field, dark current, atmospheric
extinction, zero-point).  The closure-relevant
datum \(d\) is the per-exposure calibration log
binding ADU counts to absolute flux, as in the
SDSS \"ubercal solution.  The raw integration is
bare-realised even when \(d\) is absent, so the
reduct satisfies \(\operatorname{BareReal}
_{\mathfrak S}\) yet the retained claim ``object
has AB magnitude \(m\)'' has a residual whose
audit datum is the missing calibration
\citep{Padmanabhan2008Improved}.

\item[\emph{Chain of title for real property
(real-property law).}]
The carrier \(C\) is the parcel together with
the occupant's actual possession; \(\mathfrak S\)
imposes the marketable-title audit (recorded
deeds, mortgages, easements, liens).  The
closure-relevant datum \(d\) is the unbroken
chain of recorded conveyances from a sovereign
source.  Adverse occupancy bare-realises the
parcel even when an unrecorded prior deed or
undischarged mortgage breaks the chain, so the
reduct satisfies \(\operatorname{BareReal}
_{\mathfrak S}\) but the retained claim
``marketable title'' has a residual whose
required record is the missing deed
\citep{Palomar2003Patton}.
\end{description}

\begin{theorem}[Converse CER no-go]\label{thm:CER-converse-no-go}
Let \(\mathfrak S\), \(\mathcal C\),
\(\operatorname{BareReal}_{\mathfrak S}\), and
\(\mathrm{Cl}_{\mathfrak S}\) be as above.  Under the
closure-instantiation-gap hypothesis there exists a carrier
\(C\in\mathcal C\) such that
\[
  C=\mathrm{Cl}_{\mathfrak S}(C)
  \quad\text{and}\quad
  \neg\operatorname{BareReal}_{\mathfrak S}(C).
\]
Equivalently, the converse CER implication
\[
  C=\mathrm{Cl}_{\mathfrak S}(C)
  \Longrightarrow
  \operatorname{BareReal}_{\mathfrak S}(C)
\]
fails in general.
\end{theorem}
\footnote{The Lean encoding records the converse no-go as a derived
consequence of an opaque closure-instantiation-gap assumption; see
\Cref{app:formalization}.}
\begin{proof}
Assume the scope \(\mathfrak S\), the predicate
\(\operatorname{BareReal}_{\mathfrak S}\), the closure operator
\(\mathrm{Cl}_{\mathfrak S}\), and the closure-instantiation-gap
hypothesis.  The hypothesis supplies a formal carrier \(C\) whose
records, residual statuses, bridges, and nonclaims are internally
closed:
\[
  C=\mathrm{Cl}_{\mathfrak S}(C).
\]
It also supplies the missing-instantiation witness:
\[
  \neg\operatorname{BareReal}_{\mathfrak S}(C).
\]
The point is that formal closure only checks the declared audit data
inside \(\mathfrak S\).  It does not by itself construct an external
physical system realising \(C\).  Therefore the displayed carrier
\(C\) is closed over \(\mathfrak S\) but not bare-realised, which is
the desired witness to failure of the converse CER implication.
\end{proof}

\subsection*{Layer instantiations}\label{layer-inst:CER-converse}

\begin{description}
\item[\emph{Construction documents for a
never-constructed building (architectural
design).}]
The carrier \(C\) is the building considered as a
closure-of-record: sealed construction documents,
stamped structural calculations, life-safety
analysis, code-review approvals, and issued
permit.  Under \(\mathfrak S\) the
architectural-design audit is closed at the
construction-document set.  When the owner
declines to break
ground or funding lapses, no physical building
instantiates \(C\), so \(C=\mathrm{Cl}_{\mathfrak
S}(C)\) yet \(\neg\operatorname{BareReal}
_{\mathfrak S}(C)\)
\citep{AllenIano2019Fundamentals}.

\item[\emph{Ratified RFC without deployed
implementation (Internet standards).}]
The carrier \(C\) is the standards-track RFC
considered as a closure-of-record: working-group
consensus, IETF Last Call, IESG approval, RFC
Editor publication, IANA registry allocation per
RFC 2026.  Under \(\mathfrak S\) the standards
process is closed at publication.  When no
production implementation deploys the protocol
(orphan or stillborn RFCs), no Internet host
realises \(C\), so \(C=\mathrm{Cl}_{\mathfrak S}
(C)\) yet \(\neg\operatorname{BareReal}_{\mathfrak
S}(C)\) \citep{Bradner1996RFC2026}.

\item[\emph{Statute fallen into desuetude
(statutory law).}]
The carrier \(C\) is the enacted statute
considered as a closure-of-record: bill passed by
both chambers, presented and signed (or enacted
over veto), codified in the official compilation.
Under \(\mathfrak S\) the legislative procedure is
closed at codification.  When prosecutors decline
to charge, courts decline to apply, and citizens
disregard the statute over sustained time, no
enforcement instance realises \(C\), so
\(C=\mathrm{Cl}_{\mathfrak S}(C)\) yet \(\neg
\operatorname{BareReal}_{\mathfrak S}(C)\); Calabresi
treats desuetude as the recognised judicial
response \citep{Calabresi1982Common}.

\item[\emph{Superconducting Super Collider: closed
design, cancelled construction (megaproject).}]
The carrier \(C\) is the proposed 20\,TeV
proton--proton collider considered as a
closure-of-record: Conceptual Design Report,
Site-Specific
Conceptual Design Report, technical design
reports, DOE design reviews, environmental impact
statement, congressional authorisations.  Under
\(\mathfrak S\) the megaproject design-review
audit is closed at the review-completion
boundary.  Congressional cancellation in 1993
terminated construction before operations, so no
running collider instantiates \(C\).  Hence
\(C=\mathrm{Cl}_{\mathfrak S}(C)\) yet
\(\neg\operatorname{BareReal}_{\mathfrak S}(C)\)
\citep{Riordan2015Tunnel}.

\item[\emph{Standard {ML} signature without a
structure (programming-language type systems).}]
The carrier \(C\) is a Standard ML signature
considered as a closure-of-record: abstract type
declarations, value specifications with types,
exception declarations, all type-checked against
the SML definition.  Under \(\mathfrak S\) the
type-system audit is closed at signature
elaboration.  Without a matching structure binding
no runtime value inhabits the abstract types and
no function body computes the declared values, so
\(C=\mathrm{Cl}_{\mathfrak S}(C)\) yet \(\neg
\operatorname{BareReal}_{\mathfrak S}(C)\)
\citep{MacQueen1984Modules}.
\end{description}

\begin{proposition}[Closure--Reality Boundary resolution]\label{prop:F53}
The AOR apparatus supplies the positive replacement for the failed
closure-equals-reality principle.  Via \(\ClSB\), \(\etaSB\), and
\(\XiSB\), bare carriers are reflected into audited operational
closures, and the strengthened equivalence
\[
  C \text{ closed over } \mathfrak S
  \quad\Longleftrightarrow\quad
  C \text{ AOR-realisable over } \mathfrak S
\]
holds at the AOR scope.
\end{proposition}
\footnote{\raggedright Anchored to the substrate Lean at
\texttt{ClosureRealityBoundary.\allowbreak{}closure\_reality\_boundary}
(in
\texttt{SixBirdsMetaMath.FoundationsIV.\allowbreak{}ClosureRealityBoundary});
the substrate proof appears in \citep{Tsiokos2026AOR}.  See
\Cref{app:formalization} for the anchor-disclosure discipline.}
\begin{proof}
The direct and converse no-go theorems show that the failed CER
predicate used the wrong right-hand side: bare realisability is too
weak to remember closure records, and formal closure alone is too
weak to produce an external instantiation.  AOR strengthens
realisability by requiring finite audited operational structure over
\(\mathfrak S\): typed interfaces, bounded access, finite records,
route and holonomy data, residual statuses, nonclaim registration,
and meta-audited coverage.

Given a raw carrier \(C\), the operator \(\ClSB\) adds or records
exactly the data required by this AOR scope.  The unit \(\etaSB(C)\)
compares \(C\) to the reflected carrier, and the residual \(\XiSB(C)\)
records the remaining gap.  If \(C\) is already closed over
\(\mathfrak S\), the residual vanishes in the audited sense and the
carrier is AOR-realisable.  Conversely, if \(C\) is AOR-realisable,
then the operational realisability certificate already includes the
records, statuses, and nonclaims required for closure over
\(\mathfrak S\), so \(C\) is closed.  The reflective adjunction in the
AOR substrate gives the universal property: maps compatible with the
audited operational structure factor through \(\ClSB(C)\).  Thus AOR
realisability is the positive replacement for bare realisability at
the closure--reality boundary.
\end{proof}

\subsection*{Layer instantiations}\label{layer-inst:F53}

\begin{description}
\item[\emph{Good Manufacturing Practice
(pharmaceutical regulatory affairs).}]
The raw carrier \(C\) is a manufactured drug
substance; the AOR scope \(\mathfrak S\) imposes
the GMP audit (batch records, validated processes,
raw-material traceability, quality-unit release,
deviation investigations, change control).  The
reflective closure \(\mathrm{Cl}^{\mathrm{SB}}\)
adds the GMP documentation; \(\eta^{\mathrm{SB}}\)
compares the raw production data to the audited
substance; \(\Xi^{\mathrm{SB}}\) catalogs
non-conformances.  At the AOR scope, GMP-closed
equates with GMP-realisable under FDA-inspectable
manufacturing, strengthening the failed bare
``we made the drug'' \citep{ICH2016Q7}.

\item[\emph{{GAAP}-conformant audited financial
statements (financial reporting).}]
The raw carrier \(C\) is the firm's financial
position; the AOR scope \(\mathfrak S\) imposes
the GAAP/auditor audit (ledgers, transaction
evidence, internal-control assessment, workpapers,
attestation report).  The reflective closure
\(\mathrm{Cl}^{\mathrm{SB}}\) adds the
GAAP-conformant statements with auditor opinion
and ICFR attestation; \(\eta^{\mathrm{SB}}\)
records missing transaction data;
\(\Xi^{\mathrm{SB}}\) catalogs control
deficiencies.  At the AOR scope, GAAP-closed
equates with audit-realisable under
PCAOB-inspectable workpapers, strengthening the
failed bare ``we have this revenue''
\citep{FASB2009ASC}.

\item[\emph{{ACM} Artifact Review and Badging
(computational reproducibility).}]
The raw carrier \(C\) is a reported computational
result; the AOR scope \(\mathfrak S\) imposes the
artefact-review audit (executable artefact,
archived data, reproduction instructions,
environment specification, reviewer reproduction
log).  The reflective closure \(\mathrm{Cl}
^{\mathrm{SB}}\) adds the reviewed package with
Available / Functional / Reusable / Results
Reproduced / Results Replicated badges;
\(\eta^{\mathrm{SB}}\) records missing reproduction
data; \(\Xi^{\mathrm{SB}}\) catalogs
reviewer-detected gaps.  At the AOR scope,
artefact-closed
equates with reproduction-realisable under
independent review, strengthening the failed bare
``we ran the code'' \citep{ACM2020Artifact}.

\item[\emph{{CompCert} formally verified compiler
(programming-language verification).}]
The raw carrier \(C\) is the compilation function
from C source to assembly; the AOR scope
\(\mathfrak S\) imposes the verification audit (Coq
formalisation of source and target semantics,
mechanically checked semantic-preservation
theorem, small trusted base).  The reflective
closure \(\mathrm{Cl}^{\mathrm{SB}}\) adds the
end-to-end Coq refinement proof;
\(\eta^{\mathrm{SB}}\) records unverified
preservation; \(\Xi^{\mathrm{SB}}\) catalogs
unverified passes and trust assumptions.  At the
AOR scope, verification-closed equates with
semantic-preservation-realisable under
machine-checked refinement, strengthening the
failed bare ``the compiler emits assembly that
runs'' \citep{Leroy2009Formal}.

\item[\emph{End-to-end verifiable elections
(election integrity).}]
The raw carrier \(C\) is the reported vote tally;
the AOR scope \(\mathfrak S\) imposes the
end-to-end-verifiability audit (encrypted ballot
receipts, public bulletin board, zero-knowledge
proofs of correct shuffle and decryption,
individual voter verification).  The reflective
closure \(\mathrm{Cl}^{\mathrm{SB}}\) adds the
E2E-verifiable record per the Helios protocol;
\(\eta^{\mathrm{SB}}\) records missing
verification data; \(\Xi^{\mathrm{SB}}\) catalogs
disputed ballots or proof failures.  At the AOR
scope, verifiability-closed equates with
E2E-realisable under public cryptographic audit
while preserving ballot secrecy, strengthening
the failed bare ``we counted the votes''
\citep{Adida2008Helios}.
\end{description}

This anchor closes the catalog by identifying where formal closure
must be replaced by anchored reflective operational realisability.
The no-go theorems establish the need to strengthen the right-hand
side from bare reality to AOR-realisability; the AOR resolution states
the strengthened equivalence.  The chapter does not re-prove the AOR
apparatus, does not derive ad hoc raw counterexamples, and does not
extend the equivalence to non-AOR reflective frameworks.  The
broader philosophical setting in which this boundary lives is the
structural-realism debate launched by \citet{Worrall1989StructuralRealism};
the present chapter does not take a stand in that debate, only
records the formal closure / realisability gap that any such stand
must confront.
 \section{Substrate instances at a glance}\label{sec:substrate-instances}

This chapter introduces no new claim. It is a substrate-instance map
for readers arriving from a substrate paper or following a citation
from an axis chapter: the table and short paragraphs below say which
layer-agnostic \FIV\ row is instantiated, anchored, or contextually
supported by which substrate paper.

The No-Needles Stability theorem (\Cref{thm:F6} in
\Cref{sec:stability}) is the layer-agnostic stability statement
instantiated by \PaperNeedleKiller\ \citep{Tsiokos2026NeedleKiller} and
\PaperNS\ \citep{Tsiokos2026NavierStokes}. The former supplies the
layer-dissolving membrane endpoint; the latter supplies the
Navier--Stokes carrier context. The formalization appendix records this
row as an Anchor.

The Strict Extension lemma (\Cref{lem:F8} in \Cref{sec:stability}) is a
direct \FIV\ theorem whose Schur-calculus intuition is supplied by
\PaperAR\ \citep{Tsiokos2026XiCompanion}. The substrate citation is
descriptive context for residual composition along refinement chains,
not a delegation. The formalization appendix records this row as a
Theorem.

The Hiddenness Normal Form theorem (\Cref{thm:F13a} in
\Cref{sec:access}) is the layer-agnostic quotient/currentizer statement
associated with \PaperPvNP\ \citep{Tsiokos2026PvNP}. The substrate paper
supplies the saturated-SAT carrier expected to discharge the hiddenness
source condition. The formalization appendix records this row as a
Partial wrapper.

The CSL Formed-Closure Law (\Cref{thm:F13b} in \Cref{sec:access}) is
the layer-agnostic formed-closure statement associated with \PaperCSL\
\citep{Tsiokos2026CSL}. The source context is diagnosis-grade CSL
formed-closure, so the body does not treat it as an unconditional direct
theorem. The formalization appendix records this row as an Obligation
wrapper.

The Duality Fixity theorem (\Cref{thm:F14} in \Cref{sec:symmetry}) is
the layer-agnostic self-dual trace-confinement statement associated with
\PaperSDTC\ \citep{Tsiokos2026SDTC} and \PaperRH\
\citep{Tsiokos2026RH}. Either substrate source supplies the
trace-confinement bridge, matching the disjunctive source condition
discussed in \Cref{thm:F14}. The formalization appendix records this
row as a Partial wrapper.

The Spontaneous Symmetry Breaking / Branch-Selection Formation theorem
(\Cref{thm:F15b} in \Cref{sec:symmetry}) is the layer-agnostic
future-source branch-selection statement parallel to the CSL diagnosis
pattern in \PaperCSL\ \citep{Tsiokos2026CSL}. It records the structural
form of a formed branch-selection record once a future branch-selection
source is supplied. The formalization appendix records this row as an
Obligation wrapper.

The Closure--Reality Boundary resolution (\Cref{prop:F53} in
\Cref{sec:anchor:closure-reality-boundary}) is the layer-agnostic
boundary statement anchored by \PaperAOR\ \citep{Tsiokos2026AOR}. The
substrate paper supplies the reflective closure operator \(\ClSB\), unit
\(\etaSB\), residual \(\XiSB\), and adjunction underlying the AOR
positive replacement for bare closure-equals-reality. The formalization
appendix records this row as an Anchor.

\begin{table}[t]
  \centering
  \footnotesize
  \begin{tabularx}{\textwidth}{@{}
      >{\raggedright\arraybackslash}p{0.18\textwidth}
      >{\raggedright\arraybackslash}p{0.18\textwidth}
      >{\raggedright\arraybackslash}X
      >{\raggedright\arraybackslash}p{0.22\textwidth}@{}}
    \toprule
    F-row & Substrate paper(s) & Substrate result & F-IV treatment status \\
    \midrule
    F6 No-Needles &
    \citet{Tsiokos2026NeedleKiller, Tsiokos2026NavierStokes} &
    Layer-dissolving membrane theorem plus Navier--Stokes
    instance. &
    \statAnchor \\
    F13a Hiddenness &
    \citet{Tsiokos2026PvNP} &
    Saturated-SAT closure construction. &
    \statPartial \\
    F13b CSL Formed-Closure &
    \citet{Tsiokos2026CSL} &
    CSL diagnosis-grade formed-closure. &
    \statObligation \\
    F14 Duality Fixity &
    \citet{Tsiokos2026SDTC, Tsiokos2026RH} &
    SDTC apparatus theorem plus RH trace
    specialization. &
    \statPartial\ (disjunction guard) \\
    F15b Spontaneous SSB &
    \citet{Tsiokos2026CSL} (parallel) &
    Parallel to CSL diagnosis; future-source. &
    \statObligation \\
    F53 Closure--Reality Boundary &
    \citet{Tsiokos2026AOR} &
    AOR reflective-closure resolution. &
    \statAnchor \\
    \bottomrule
  \end{tabularx}
  \caption{Substrate-instance map for \FIV\ anchor, partial, and obligation rows.}
  \label{tab:substrate-instances}
\end{table}
 
Readers arriving from a substrate paper can use
\Cref{tab:substrate-instances} to locate the corresponding
layer-agnostic statement and then jump to the displayed theorem,
lemma, or proposition named above. The table is navigational: the axis
chapters contain the mathematical statements and proofs, while
\Cref{app:formalization} carries the formalization ledger.
 \section{Sketches and future work}\label{sec:sketches}

This chapter records pointers, not new laws. The preceding chapters
state the \NumLaws\ laws that belong to the \FIV\ catalog. Here we
record structural candidates that appeared in the atlas work but were
not promoted to catalog rows. The reason is not that the candidates
are uninteresting; it is that their recognition source has not yet
been identified. A recognition source is the substrate, apparatus
layer, or carrier discipline that makes the named structure precise
enough to support the normal-form template of \Cref{sec:apparatus}.
Without such a source, a phrase may be suggestive but still fail to
determine the setup, theorem, proof spine, cases, interaction status, and
nonclaims required of a \FIV\ law.

The catalog is therefore deliberately non-exhaustive. Several structural
candidates documented in the atlas await a recognition source before
they can be cast in normal form; this section lists them with that
caveat. This is a scope discipline rather than a promise that every
sketch will later become layer-agnostic. Some sketches may turn out to
belong to a substrate paper only, because the named structure is
carrier-specific. Others may become \FIV-style laws once a future paper
or a new apparatus layer supplies the missing recognition source.

The clearest deferred item is dimension as stable rank of access. The
sketch asks whether dimension should be read, at the layer-agnostic
level, as a presentation-invariant rank of the access quotient: roughly,
the stable number of independent directions that survive the declared
observation discipline. The structural appeal of this reading is that it
would explain why dimension is neither raw carrier data nor merely a
coordinate convention. Yet it is not a catalog law here. To state it as a
normal form one must know which rank notion is being recognized, which
class of quotient presentations preserves it, and which interaction
judgment certifies that the rank is observed rather than imposed. Those
data are not part of the present \FIV\ apparatus. Until they are supplied
by a substrate or a new layer of general theory, the item remains a
sketch.

A second deferred item is smoothness as coherent linearization. The
sketch asks whether smooth structure can be characterized as the
existence of compatible local linearizations across the apparatus's
declared overlaps. In familiar settings, smoothness is tied to tangent
data, coordinate changes, differentiability, or analytic estimates. The
\FIV\ temptation is to abstract this as a gluing law: local
linearization data should globalize exactly when the local linear
approximations descend through the relevant access and overlap
quotients. But that sentence still hides the recognition source. The
catalog would need a substrate-independent account of what counts as
linearization, which overlaps preserve it, and how the associated
coherence ledger is audited. Without that apparatus, the item is not a
law; it is a pointer to a possible future bridge between the local-global
obstruction pattern of \Cref{sec:transport:local-global} and a
smoothness-specific substrate.

A third deferred item is stability of nothing, or vacuum non-emptiness.
The sketch starts from the observation that a ``nothing'' state at a
layer is rarely absolute absence. It is more naturally a terminal,
initial, null, or background class relative to a quotient and an audit
discipline. The catalog already contains a direct theorem on
Vacuum / Background Selection in
\Cref{sec:self-reference-limits:vacuum}: a vacuum exists when the
declared background selector is total over the audit-bounded
perturbations. The deferred sketch is stronger and more delicate. It
would ask when the very stability of the null or empty class forces a
non-empty residual structure. That question needs a recognition source
for the null class and for the residual that is supposed to survive it.
Absent that source, the safe statement remains the catalog law on
background selection; the stronger ``nothing is structurally non-empty''
reading is kept out of scope.

The atlas notes also contain other candidates around boundary
compression, topology, continuum behavior, critical loci, uncertainty,
and inaccessible determining state. Many of these were absorbed into
the catalog after their normal-form shape became clear: Horizon
Sufficiency, Information Loss, Continuum Emergence, Criticality,
Topology as Global Gluing Invariant, and Hiddenness all appear as named
rows in the preceding chapters. The remaining atlas-only formulations
are not repeated as shadow theorems here. They function as a record of
search directions: places where the same quotient, descent, gluing,
memory, and audit vocabulary appeared to be doing real work, but where
the paper did not yet have enough typed substrate to state a theorem
without overclaiming.

A note on the eight-axis grouping. ``Axis'' is used as an
organizing label for laws that address the same kind of
obstruction; the eight axes are not a structured object but an
editorial scheme to make the catalog navigable. The axis-topology
schematic mirrors that role: it shows how the axes and drafting
chapters relate, not a theorem about an object called an axis. The
deferred sketches listed above come from the same planning record
that produced the axis grouping, and they are not a shadow
inventory of laws.

This distinction is important for readers who arrive with a favored
carrier theory in mind. If a substrate paper supplies a recognition
source for dimension, smoothness, or vacuum non-emptiness, then a later
project can test whether the candidate is genuinely layer-agnostic. The
test is the same one used throughout this paper: identify the setup,
state the quotient or closure normal form, give the proof spine, record
the cases and interaction status, and write the nonclaims. If the
candidate passes that test without depending on special carrier
features, it can be promoted to the catalog. If the candidate depends on
special analytic, geometric, arithmetic, or physical structure, then it
belongs in the corresponding substrate paper instead.

The future-work posture is therefore conservative. \FIV\ does not claim
that the \NumLaws-row catalog is complete under future developments. The
Closure Completeness Boundary of
\Cref{sec:self-reference-limits:closure-completeness} is a
meta-coherence statement about the declared \FIV\ language: within the
present layer-agnostic scope, a proposed direction is either in the
catalog, reducible to catalog laws, or marked out of scope. It is not an
absolute completeness theorem for all future apparatus. New recognition
sources may move a deferred sketch from the third class into the first
or second class.

Deferral is not an inadequacy claim. It prevents the paper from
using evocative vocabulary before the apparatus can support it.
Dimension, smoothness, and nothingness are substantial mathematical
notions, but their layer-agnostic forms must be earned. Until the
recognition source is explicit, the appropriate treatment is to
record the candidate, explain why it was not counted, and leave the
normal-form template untouched.
 \section{Nonclaims and caveats}\label{sec:nonclaims}

This chapter is the consolidated nonclaims landing for \FIV. The
per-axis chapters state the local caveats where they matter; here the
scope fence and the prepared objection mitigations are collected in one
place.

\begin{enumerate}[label=\textbf{N\arabic*.}, leftmargin=*]
\item \textbf{\FIV\ does not re-prove substrate theorems.}
The paper catalogs layer-agnostic structural laws; it does not re-derive
the carrier-specific theorems of the substrate papers. Anchor, partial,
and obligation rows identify their provenance in the body, while
\Cref{app:formalization} records the formal bookkeeping.

\item \textbf{\FIV\ is not a foundational reset.}
The apparatus inherited from \FoundI, \FoundII, and \FoundIII\ is used
as a black box after the review in \Cref{sec:apparatus}. Closure
formation, admissibility, quotient packaging, and the \(\BirdInt\)
judgment are inputs here, not new foundations being rederived.

\item \textbf{\FIV\ is not a unification claim.}
The axes are an organizing device for recurring obstruction types. The
paper does not assert that the axes form a categorical, fibrational,
monadic, or otherwise privileged structure.

\item \textbf{\FIV\ is not a complete catalog.}
The \NumLaws-row inventory is the closed scope of this paper, not an
exhaustiveness theorem. Deferred candidates appear in
\Cref{sec:sketches} only when their recognition source is not yet
identified.

\item \textbf{The tagged trust-base axioms are obligations or
pending substrate dependencies.}
The catalog's four trust-base axioms are recorded as obligations or
substrate-pending entries with opaque source-predicate guards;
\Cref{app:formalization} gives the per-axiom register.

\item \textbf{\FIV\ does not mechanize the self-dual trace
confinement, Riemann Hypothesis, P versus NP, or hiddenness
substrate proofs.}
Those substrate proofs \citep{Tsiokos2026SDTC, Tsiokos2026RH,
Tsiokos2026PvNP, Tsiokos2026CSL} do not have Lean mechanization in
this repository.

\item \textbf{The catalog discipline is the contribution.}
The contribution is the uniform normal-form treatment of \NumLaws\
layer-agnostic laws, together with explicit proof, provenance, and
nonclaim discipline. The substrate proofs themselves live in the
cited substrate papers.

\item \textbf{The eight-axis scheme is not minimal or canonical.}
The 8-axis scheme is an editorial choice, not a theorem of the
apparatus. Alternative groupings are legitimate if they preserve the
same statements, disclosures, and exclusions.
\end{enumerate}

The same fence answers the predictable reviewer objections.

\begin{enumerate}[label=\textbf{O\arabic*.}, leftmargin=*]
\item \textbf{Why call this ``Foundations IV'' if it is a catalog?}
``Foundations IV'' names the series position, not a reset. The
contribution is the catalog discipline: a uniform normal-form
template applied to \NumLaws\ layer-agnostic laws, each carrying
its proof or its substrate disclosure.

\item \textbf{Are the anchor rows merely restating substrate theorems?}
No. The layer-agnostic statement quantifies over closures of a class,
whereas the substrate paper specializes to one carrier. Recording the
layer-agnostic statement is therefore a distinct contribution.

\item \textbf{Why are some dependencies recorded as guarded assumptions
rather than derived here?}
The relevant substrate paper has not been mechanized in this repository,
or the source is explicitly future-facing. The guarded entry records
that dependency and discharges only when the corresponding source
certificate lands.

\item \textbf{How do the four trust-base axioms avoid soundness
pitfalls?}
Each axiom's hypothesis is an opaque source predicate without a public
constructor, except Duality Fixity where the guard is a definitional
disjunction of two such predicates. Thus the axiom alone cannot
manufacture an instance of its hypothesis.

\item \textbf{Why a single monograph rather than per-axis papers?}
The page budget supports a single monograph, and the single format
avoids repeating the apparatus in each axis. The packaging decision can
be revisited only if the page budget or review process forces that
choice.

\item \textbf{Why these \NumLaws\ laws and not others?}
These laws are the recurring structural obstructions and closures that
surfaced across the broader Six Birds project. The catalog is not
exhaustive; \Cref{sec:sketches} records candidates that remain outside
the inventory until their recognition sources are identified.

\item \textbf{Why are the Status/Records/Coherence and
Self-reference/Limits chapters so dense?}
Those axes cover the broadest range of record, coherence, information,
limit, and self-reference phenomena. Short per-law subsections and
summary tables keep the dense material navigable.

\item \textbf{Do the apparatus restatements lose nuance from the
foundations papers?}
Yes; the review in \Cref{sec:apparatus} is intentionally compact.
It is enough to read this paper, and the cited foundations papers
carry the operational nuance.
\end{enumerate}

These caveats are part of the normal-form discipline. A catalog of
structural laws is useful only if its hypotheses, provenance, proof
status, and exclusions are visible alongside its theorems.
 

\appendix
\section{Lean formalization}\label{app:formalization}

\subsection{Wording discipline}\label{subsec:lean-wording-discipline}

This appendix records the wording discipline used throughout the body.
Its purpose is to prevent a paper-side theorem statement, a Lean
theorem, a record projection, a metadata entry, and a guarded
trust-base wrapper from being described with the same verb. Body prose
uses the paper-prose theorem names; Lean declaration names and trust-base
tags are confined to this appendix.

\begin{center}
\small
\begin{tabularx}{\textwidth}{@{}
    >{\raggedright\arraybackslash}p{0.27\textwidth}
    >{\raggedright\arraybackslash}X
    >{\raggedright\arraybackslash}X@{}}
  \toprule
  Lean coverage mode & Permitted body-prose wording & Forbidden wording \\
  \midrule
  \texttt{direct\_proof} &
  ``Lean proves \ldots'', ``a Lean derivation establishes \ldots'',
  ``verified in Lean'', or ``checked in Lean''. &
  None, provided the statement is the same statement being claimed in
  the prose. \\
  \texttt{projection\_over\_record} &
  ``tracked by the formalization harness'', ``encoded as a
  proof-carrying record field'', or ``checked as a conditional schema
  over explicit hypotheses''. &
  ``Lean proves'', ``verified in Lean'', ``mechanized in Lean'', or
  ``a Lean derivation establishes \ldots''. \\
  \texttt{definition\_metadata} &
  ``recorded in the traceability metadata'', or no Lean mention. &
  ``formalized as a Lean definition'', ``the Lean definition of
  \ldots'', or ``Lean carries the definition of \ldots''. \\
  \texttt{local\_definition} &
  ``realized in Lean'' when the local definition is faithful; ``partially
  realized in Lean'' when the registry marks only a partial
  realization. &
  Unqualified ``realized in Lean'' for partial realizations or entries
  that still need a real definition. \\
  \bottomrule
\end{tabularx}
\end{center}

\begin{center}
\small
\begin{tabularx}{\textwidth}{@{}
    >{\raggedright\arraybackslash}p{0.27\textwidth}
    >{\raggedright\arraybackslash}X
    >{\raggedright\arraybackslash}X@{}}
  \toprule
  Lean coverage mode & Permitted body-prose wording & Forbidden wording \\
  \midrule
  \texttt{theorem} (mechanize-now rows) &
  ``Lean proves \ldots'' or equivalent direct-proof wording. &
  None beyond the ordinary requirement that the prose statement match
  the Lean theorem. \\
  \texttt{anchor\_substrate} &
  ``The layer-agnostic statement is recorded here; the substrate
  instance was proved in the cited substrate paper. The Lean realization
  in this paper delegates to inherited substrate Lean.'' &
  ``Lean proves \ldots'' without the delegation qualifier, or
  ``\FIV\ proves \ldots'' when the proof is actually inherited from a
  substrate. \\
  \texttt{partial} &
  ``The Lean realization in this paper is a partial theorem wrapper
  applying a same-module substrate-pending axiom guarded by an opaque
  source predicate; the axiom discharges when the substrate
  mechanization lands.'' &
  ``Lean proves \ldots'' without the axiom-applied qualifier, or
  ``\FIV\ derives \ldots'' when the result depends on the guarded
  pending substrate. \\
  \texttt{obligation} &
  ``The Lean realization in this paper is a theorem wrapper applying a
  same-module obligation axiom guarded by an opaque source predicate; no
  instantiation of the axiom can derive False from the axiom alone.'' &
  ``Lean proves \ldots'' without the axiom-applied qualifier, ``\FIV\
  derives \ldots'', or silently dropping the opacity-guard disclosure. \\
  \bottomrule
\end{tabularx}
\end{center}

The semantic-alignment field overlays these modes. A \texttt{faithful}
row uses the base wording without qualification. A
\texttt{narrowed\_surrogate} row must say that Lean covers the
typed or finite version while the broader statement is carried by the
paper proof. A \texttt{weakened\_genericized} row must say that Lean
covers the genericized class while the substrate-specific statement is
kept in the cited paper. A \texttt{needs\_strengthening} row must name
the upgrade debt rather than implying closure. A
\texttt{projection\_packaged} row uses the
\texttt{projection\_over\_record} wording. An
\texttt{anchored\_to\_substrate} row uses the
\texttt{anchor\_substrate} wording. A \texttt{not\_applicable} row
receives no Lean mention. In the \FIV\ statements of record, all 46
direct theorem rows are faithful; F6 and F53 are anchored to substrate;
and F13a, F13b, F14, and F15b are explicitly marked as requiring
strengthening through their guarded wrapper discipline.

\subsection{Opacity-guard pattern}\label{subsec:opacity-guard-pattern}

This subsection is the appendix-level restatement of the opacity-guard
pattern introduced in \Cref{sec:apparatus}. Body sections that mention
partial or obligation rows refer to the apparatus chapter for the
reader-facing explanation and to this subsection for the formal
per-axiom shape.

A trust-base axiom in this paper is always paired with an opaque
source-predicate hypothesis: either a Lean predicate of type
\texttt{Prop} declared \texttt{opaque} and lacking a public constructor,
or a definitional combination, typically a disjunction, of such opaque
predicates. The axiom asserts that, given an instance of the source
predicate, the conclusion holds. Because the underlying opaque
predicates have no public constructors, no automated tactic and no user
proof can construct an instance of the guard predicate, or its negation,
without a separate axiom or opaque definition that produces an instance
of one of the underlying opaque predicates. As a consequence, no
instantiation of the trust-base axiom can derive \texttt{False} from the
axiom alone; deriving \texttt{False} would require additionally asserting
an instance of one of the underlying opaque source predicates, and that
assertion would itself be a trust-base commitment recorded in the trust
base.

The point of the pattern is to separate the paper-side structural
statement from the substrate-side certificate that is not available in
this Lean tree. The guarded axiom has the shape
\[
  \texttt{SourcePredicate}(\mathcal S) \to
  \texttt{Conclusion}(\mathcal S),
\]
where \(\mathcal S\) is substrate-shaped data. The source predicate is
not a transparent definition such as \(\top\), a decidable case split,
or a proposition whose constructor is exported by the module. The
wrapper is therefore parametric in a source certificate rather than an
unconditional theorem about arbitrary data; the remaining obligation is
located in the future or external substrate mechanization that produces
an instance of the guard.

\begin{center}
\footnotesize
\begin{tabularx}{\textwidth}{@{}
    >{\raggedright\arraybackslash}p{0.28\textwidth}
    >{\raggedright\arraybackslash}p{0.32\textwidth}
    >{\raggedright\arraybackslash}X@{}}
  \toprule
  Guard name & Underlying Lean shape & Per-guard description \\
  \midrule
  \texttt{Paper7SaturatedSAT}\\\texttt{Instance} &
  Single \texttt{opaque} predicate. &
  F13a Hiddenness guard over a saturated-SAT instance shape. \\
  \texttt{Paper8CSLDiagnosis}\\\texttt{Source} &
  Single \texttt{opaque} predicate. &
  F13b CSL Formed-Closure guard over a CSL diagnosis source. \\
  \texttt{DualityFixity}\\\texttt{SubstrateAnchor} &
  Definitional disjunction
  \(\texttt{Paper2SDTCApparatus} \lor
    \texttt{Paper6RHTraceSpecialization}\), with both disjuncts
  \texttt{opaque}. &
  F14 Duality Fixity guard. The disjunction caveat is essential: the
  guard is not itself a single opaque declaration, but the soundness
  posture is preserved because neither underlying disjunct has public
  constructors. \\
  \texttt{FutureBranchSelection}\\\texttt{FormationSource} &
  Single \texttt{opaque} predicate. &
  F15b Spontaneous SSB guard over a future branch-selection source. \\
  \bottomrule
\end{tabularx}
\end{center}

Thus the four guarded rows have the same soundness posture even though
F14 uses a definitional disjunction. In all four cases, the formal
wrapper cannot be instantiated to prove \texttt{False} from the axiom
alone. Any proof that supplies a guard instance is a separate
commitment, and the review burden moves to that commitment rather than
to the guarded implication itself.

\subsection{Narrowing and parametric disclosure for the axiom-applied wrappers}
\label{subsec:narrowing-disclosure}

The body footnotes for F13a, F13b, F14, and F15b point here for the
precise narrowing disclosure. In all four cases the Lean theorem proves
the layer-agnostic claim parametrically in a source predicate; it applies
to instances supplied by the relevant substrate or future source and
does not manufacture such an instance internally.

\textbf{F13a, Hiddenness Normal Form.}
The Lean theorem
\path|Access.HiddennessNormalForm.hiddenness_normal_form| proves
the layer-agnostic hiddenness statement under an instance of the opaque
predicate \texttt{Paper7SaturatedSATInstance}. The substrate
certification of saturated-SAT closure is not derived in this paper's
Lean tree; it enters as the hypothesis. The narrowing is the parametric
dependency on this opaque hypothesis: the body theorem holds for every
instance that future saturated-SAT substrate Lean
\citep{Tsiokos2026PvNP} will provide.

\textbf{F13b, CSL Formed-Closure Law.}
The Lean theorem
\path|Access.CSLFormedClosure.csl_formed_closure| proves the
layer-agnostic formed-closure statement under an instance of the opaque
predicate \texttt{Paper8CSLDiagnosisSource}. The CSL diagnosis source
\citep{Tsiokos2026CSL} is not reconstructed here; it is the external
certificate that supplies the guard. The wrapper therefore records the formed-closure law
for every certified CSL diagnosis source and not for arbitrary
uncertified carriers.

\textbf{F14, Duality Fixity.}
The Lean theorem \path|Symmetry.DualityFixity.duality_fixity|
proves the layer-agnostic duality-fixity statement under an instance of
the definitional disjunction
\(\texttt{Paper2SDTCApparatus} \lor
  \texttt{Paper6RHTraceSpecialization}\).
Either disjunct suffices; this is recorded by the
\texttt{DualityFixitySubstrateAnchor} definition. The narrowing is the
disjunction dependence: the theorem holds for every instance of either
opaque source predicate supplied by future self-dual trace confinement
\citep{Tsiokos2026SDTC} or Riemann Hypothesis \citep{Tsiokos2026RH}
substrate Lean, without asserting either source unconditionally.

\textbf{F15b, Spontaneous Symmetry Breaking / Branch-Selection Formation.}
The Lean theorem
\path|Symmetry.SpontaneousSSB.spontaneous_ssb_formation| proves the
layer-agnostic branch-selection formation statement under an instance
of the opaque predicate
\texttt{FutureBranchSelectionFormationSource}. The future-source
character means that no completed substrate paper is cited as the
discharge point; the source formulation remains future work. The formal
theorem tracks the layer-agnostic conclusion for any future certified
branch-selection source while keeping the obligation open.

\subsection{Tagged-axiom register}\label{subsec:tagged-axiom-register}

\begin{landscape}
\begin{table}[p]
  \centering
  \footnotesize
  \begin{tabularx}{\linewidth}{@{}
      >{\raggedright\arraybackslash}p{0.08\linewidth}
      >{\raggedright\arraybackslash}p{0.18\linewidth}
      >{\raggedright\arraybackslash}p{0.14\linewidth}
      >{\raggedright\arraybackslash}p{0.24\linewidth}
      >{\raggedright\arraybackslash}X
      >{\raggedright\arraybackslash}p{0.20\linewidth}@{}}
    \toprule
    Row & Axiom name & Tag & Opacity-guard predicate &
    Conclusion shape & Expected discharge \\
    \midrule
    F13b &
    \texttt{csl\_formed\_closure} &
    \texttt{\# obligation} &
    \texttt{Paper8CSLDiagnosis}\newline\texttt{Source} (single opaque) &
    diagnosis-source \(\Rightarrow\) formed-closure law &
    When \PaperCSL\ \citep{Tsiokos2026CSL} ships its substrate Lean \\
    F15b &
    \texttt{spontaneous\_ssb\_}\newline\texttt{formation} &
    \texttt{\# obligation} &
    \texttt{FutureBranchSelection}\newline\texttt{FormationSource}
    (single opaque) &
    future branch-selection source \(\Rightarrow\) formed selection record &
    Future-source; parallel to F13b \\
    F13a &
    \texttt{hiddenness\_pending} &
    \texttt{\# substrate\_}\newline\texttt{pending} &
    \texttt{Paper7SaturatedSAT}\newline\texttt{Instance} (single opaque) &
    saturated-SAT source \(\Rightarrow\) hiddenness normal form &
    When \PaperPvNP\ \citep{Tsiokos2026PvNP} ships its substrate Lean \\
    F14 &
    \texttt{duality\_fixity\_}\newline\texttt{pending} &
    \texttt{\# substrate\_}\newline\texttt{pending} &
    \texttt{DualityFixity}\newline\texttt{SubstrateAnchor}
    (definitional disjunction over
    \texttt{Paper2SDTCApparatus} \(\lor\)
    \texttt{Paper6RHTrace}\newline\texttt{Specialization}, both opaque) &
    anchor \(\Rightarrow\) duality-fixity normal form &
    When \PaperSDTC\ \citep{Tsiokos2026SDTC} or \PaperRH\
    \citep{Tsiokos2026RH} ships its substrate Lean \\
    \bottomrule
  \end{tabularx}
  \caption{\FIV\ trust-base axioms: tagged-axiom register.}
  \label{tab:tagged-axiom-register}
\end{table}
\end{landscape}
 
\subsection{Per-row formalization coverage}\label{subsec:formalization-coverage}

The per-row coverage table below records the paper label, paper-prose
name, Lean coverage status, semantic alignment, and Lean declaration for
each of the 52 catalog rows. The table is presented in landscape to fit
the full namespaced declaration path. The status column uses the same
vocabulary as the wording-discipline table
(\Cref{subsec:lean-wording-discipline}); the alignment column uses the
semantic-alignment vocabulary defined there.

\begin{landscape}
\begin{table}[p]
  \centering
  \scriptsize
  \renewcommand{\arraystretch}{0.90}
  \begin{tabularx}{\linewidth}{@{}
      >{\raggedright\arraybackslash}p{0.055\linewidth}
      >{\raggedright\arraybackslash}p{0.25\linewidth}
      >{\raggedright\arraybackslash}p{0.10\linewidth}
      >{\raggedright\arraybackslash}p{0.13\linewidth}
      >{\raggedright\arraybackslash}X@{}}
    \toprule
    F\# & Paper-prose short name & Status & Alignment & Lean decl (prefix dropped) \\
    \midrule
    \multicolumn{5}{@{}l}{\textbf{Transport}} \\
    F2 & Descent--Repair Normal Form & \statTheorem & faithful & \path|Transport.DescentRepair.descent_repair_normal_form| \\
    F3 & Holonomy-Memory Repair Normal Form & \statTheorem & faithful & \path|Transport.HolonomyMemoryRepair.holonomy_memory_repair_normal_form| \\
    F4 & Local--Global Obstruction Normal Form & \statTheorem & faithful & \path|Transport.LocalGlobalObstruction.local_global_obstruction_normal_form| \\
    F5 & Recombination Comparison (Branch-Merge Transport) & \statTheorem & faithful & \path|Transport.RecombinationComparison.recombination_comparison| \\
    \addlinespace
    \multicolumn{5}{@{}l}{\textbf{Stability}} \\
    F6 & No-Needles Stability & \statAnchor & anchored & \path|Stability.NoNeedles.no_needles_stability| \\
    F7 & Sufficiency Closure & \statTheorem & faithful & \path|Stability.SufficiencyClosure.sufficiency_closure| \\
    F8 & Strict Extension & \statTheorem & faithful & \path|Stability.StrictExtension.strict_extension| \\
    F9 & No Unstatused Residual & \statTheorem & faithful & \path|Stability.NoUnstatusedResidual.no_unstatused_residual| \\
    \addlinespace
    \multicolumn{5}{@{}l}{\textbf{Access}} \\
    F10 & Quotientality (Layer-Sameness via Quotient) & \statTheorem & faithful & \path|Access.Quotientality.quotientality| \\
    F11 & Adequacy / No-Overread & \statTheorem & faithful & \path|Access.Adequacy.adequacy_no_overread| \\
    F12 & No-Free-Distinction & \statTheorem & faithful & \path|Access.NoFreeDistinction.no_free_distinction| \\
    F13a & Hiddenness -- Quotient/Currentizer Normal Form & \statPartial & needs-strengthening & \path|Access.HiddennessNormalForm.hiddenness_normal_form| \\
    F13b & CSL Formed-Closure Law & \statObligation & needs-strengthening & \path|Access.CSLFormedClosure.csl_formed_closure| \\
    F24 & Layer Multiplicity & \statTheorem & faithful & \path|Access.LayerMultiplicity.layer_multiplicity| \\
    \bottomrule
  \end{tabularx}
  \caption{Per-row Lean realization map for the 52 \FIV\ catalog rows (Transport, Stability, and Access).}
  \label{tab:formalization-coverage}
\end{table}

\begin{table}[p]
  \centering
  \scriptsize
  \renewcommand{\arraystretch}{0.90}
  \begin{tabularx}{\linewidth}{@{}
      >{\raggedright\arraybackslash}p{0.055\linewidth}
      >{\raggedright\arraybackslash}p{0.25\linewidth}
      >{\raggedright\arraybackslash}p{0.10\linewidth}
      >{\raggedright\arraybackslash}p{0.13\linewidth}
      >{\raggedright\arraybackslash}X@{}}
    \toprule
    F\# & Paper-prose short name & Status & Alignment & Lean decl (prefix dropped) \\
    \midrule
    \multicolumn{5}{@{}l}{\textbf{Symmetry}} \\
    F14 & Duality Fixity (Self-Dual Trace Confinement) & \statPartial & needs-strengthening & \path|Symmetry.DualityFixity.duality_fixity| \\
    F15a & Selection Obstruction / Repair Normal Form & \statTheorem & faithful & \path|Symmetry.SelectionObstruction.selection_obstruction| \\
    F15b & Spontaneous Symmetry Breaking / Branch-Selection Formation & \statObligation & needs-strengthening & \path|Symmetry.SpontaneousSSB.spontaneous_ssb_formation| \\
    F40 & Anomaly / Symmetry Obstruction & \statTheorem & faithful & \path|Symmetry.Anomaly.anomaly_symmetry_obstruction| \\
    F41 & Duality Equivalence & \statTheorem & faithful & \path|Symmetry.DualityEquivalence.duality_equivalence| \\
    \addlinespace
    \multicolumn{5}{@{}l}{\textbf{Status / Records / Coherence}} \\
    F17 & Objectivity as Common Quotient & \statTheorem & faithful & \path|StatusRecordsCoherence.ObjectivityAsCommonQuotient.objectivity_as_common_quotient| \\
    F18 & Lawhood Descent & \statTheorem & faithful & \path|StatusRecordsCoherence.LawhoodDescent.lawhood_descent| \\
    F19 & Object Persistence & \statTheorem & faithful & \path|StatusRecordsCoherence.ObjectPersistence.object_persistence| \\
    F20 & Fact / Record Formation & \statTheorem & faithful & \path|StatusRecordsCoherence.FactRecordFormation.fact_record_formation| \\
    F21 & Reflexive Nonclosure & \statTheorem & faithful & \path|StatusRecordsCoherence.ReflexiveNonclosure.reflexive_nonclosure| \\
    F22 & Interface Mediation & \statTheorem & faithful & \path|StatusRecordsCoherence.InterfaceMediation.interface_mediation| \\
    F23 & Probability as Stable Unresolved Fiber & \statTheorem & faithful & \path|StatusRecordsCoherence.ProbabilityAsStableUnresolvedFiber.probability_as_stable_unresolved_fiber| \\
    F25 & Causality as Stable Intervention Transport & \statTheorem & faithful & \path|StatusRecordsCoherence.CausalityAsStableInterventionTransport.causality_as_stable_intervention_transport| \\
    F26 & Moduli / Contingency & \statTheorem & faithful & \path|StatusRecordsCoherence.ModuliContingency.moduli_contingency| \\
    F27 & Conservation as Orbit Descent & \statTheorem & faithful & \path|StatusRecordsCoherence.ConservationAsOrbitDescent.conservation_as_orbit_descent| \\
    F28 & Composability & \statTheorem & faithful & \path|StatusRecordsCoherence.Composability.composability| \\
    F29 & Presentation Invariance & \statTheorem & faithful & \path|StatusRecordsCoherence.PresentationInvariance.presentation_invariance| \\
    F30 & Explanation as Compression with Descent & \statTheorem & faithful & \path|StatusRecordsCoherence.ExplanationAsCompressionWithDescent.explanation_as_compression_with_descent| \\
    F31 & Counterfactual Legitimacy & \statTheorem & faithful & \path|StatusRecordsCoherence.CounterfactualLegitimacy.counterfactual_legitimacy| \\
    \bottomrule
  \end{tabularx}
  \caption{Per-row Lean realization map for the 52 \FIV\ catalog rows (Symmetry and Status/Records/Coherence).}
\end{table}

\begin{table}[p]
  \centering
  \scriptsize
  \renewcommand{\arraystretch}{0.90}
  \begin{tabularx}{\linewidth}{@{}
      >{\raggedright\arraybackslash}p{0.055\linewidth}
      >{\raggedright\arraybackslash}p{0.25\linewidth}
      >{\raggedright\arraybackslash}p{0.10\linewidth}
      >{\raggedright\arraybackslash}p{0.13\linewidth}
      >{\raggedright\arraybackslash}X@{}}
    \toprule
    F\# & Paper-prose short name & Status & Alignment & Lean decl (prefix dropped) \\
    \midrule
    \multicolumn{5}{@{}l}{\textbf{Self-reference / Limits}} \\
    F33 & Horizon Sufficiency / Holographic Compression & \statTheorem & faithful & \path|SelfReferenceLimits.HorizonSufficiency.horizon_sufficiency| \\
    F34 & Information Loss Normal Form & \statTheorem & faithful & \path|SelfReferenceLimits.InformationLoss.information_loss_normal_form| \\
    F35 & Scale Descent / Renormalization & \statTheorem & faithful & \path|SelfReferenceLimits.ScaleDescentRenormalization.scale_descent_renormalization| \\
    F36 & Singular Limit Completion & \statTheorem & faithful & \path|SelfReferenceLimits.SingularLimitCompletion.singular_limit_completion| \\
    F37 & Complementarity as Non-Joint Access & \statTheorem & faithful & \path|SelfReferenceLimits.Complementarity.complementarity_as_non_joint_access| \\
    F38 & Arrow from Record Monotonicity & \statTheorem & faithful & \path|SelfReferenceLimits.ArrowFromRecordMonotonicity.arrow_from_record_monotonicity| \\
    F39 & Entropy as Fiber Volume & \statTheorem & faithful & \path|SelfReferenceLimits.EntropyAsFiberVolume.entropy_as_fiber_volume| \\
    F42 & Irreducibility / No Short Closed Description & \statTheorem & faithful & \path|SelfReferenceLimits.Irreducibility.irreducibility_no_short_closed_description| \\
    F43 & Continuum Emergence & \statTheorem & faithful & \path|SelfReferenceLimits.ContinuumEmergence.continuum_emergence| \\
    F44 & Criticality / Phase Boundary & \statTheorem & faithful & \path|SelfReferenceLimits.Criticality.criticality_phase_boundary| \\
    F45 & Locality / Domain of Dependence & \statTheorem & faithful & \path|SelfReferenceLimits.Locality.locality_domain_of_dependence| \\
    F46 & Law--State Split & \statTheorem & faithful & \path|SelfReferenceLimits.LawStateSplit.law_state_split| \\
    F47 & Fine-Tuning as Small Selector Region & \statTheorem & faithful & \path|SelfReferenceLimits.FineTuning.fine_tuning_small_selector_region| \\
    F48 & Topology as Global Gluing Invariant & \statTheorem & faithful & \path|SelfReferenceLimits.Topology.topology_global_gluing_invariant| \\
    F49 & Common-Source / Nonlocal Correlation Normal Form & \statTheorem & faithful & \path|SelfReferenceLimits.CommonSourceNonlocal.common_source_nonlocal_correlation| \\
    F50 & Vacuum / Background Selection & \statTheorem & faithful & \path|SelfReferenceLimits.VacuumBackgroundSelection.vacuum_background_selection| \\
    F51 & Unification as Common Refinement & \statTheorem & faithful & \path|SelfReferenceLimits.Unification.unification_as_common_refinement| \\
    F52 & Closure Completeness Boundary & \statTheorem & faithful & \path|SelfReferenceLimits.ClosureCompletenessBoundary.closure_completeness_boundary| \\
    \addlinespace
    \multicolumn{5}{@{}l}{\textbf{Closure--Reality Boundary}} \\
    F53 & Closure--Reality Boundary & \statAnchor & anchored & \path|ClosureRealityBoundary.closure_reality_boundary| \\
    \bottomrule
  \end{tabularx}
  \caption{Per-row Lean realization map for the 52 \FIV\ catalog rows (Self-reference/Limits and Closure--Reality Boundary).}
\end{table}
\end{landscape}
 
\subsection{Coverage summary by axis}\label{subsec:lean-coverage-by-axis}

\begin{table}[!ht]
  \centering
  \footnotesize
  \begin{tabularx}{\textwidth}{@{}
      >{\raggedright\arraybackslash}X
      >{\raggedleft\arraybackslash}p{0.11\textwidth}
      >{\raggedleft\arraybackslash}p{0.10\textwidth}
      >{\raggedleft\arraybackslash}p{0.10\textwidth}
      >{\raggedleft\arraybackslash}p{0.13\textwidth}
      >{\raggedleft\arraybackslash}p{0.08\textwidth}@{}}
    \toprule
    Axis & Theorem & Anchor & Partial & Obligation & Total \\
    \midrule
    Transport & 4 & 0 & 0 & 0 & 4 \\
    Stability & 3 & 1 & 0 & 0 & 4 \\
    Access & 4 & 0 & 1 & 1 & 6 \\
    Symmetry & 3 & 0 & 1 & 1 & 5 \\
    Status/Records/Coherence & 14 & 0 & 0 & 0 & 14 \\
    Self-reference/Limits & 18 & 0 & 0 & 0 & 18 \\
    Closure--Reality Boundary & 0 & 1 & 0 & 0 & 1 \\
    \midrule
    \textbf{Total} & \textbf{46} & \textbf{2} & \textbf{2} & \textbf{2} & \textbf{52} \\
    \bottomrule
  \end{tabularx}
  \caption{\FIV\ Lean coverage by axis and status.}
  \label{tab:lean-coverage-by-axis}
\end{table}
 
\subsection{Reviewer-cycle log}\label{subsec:reviewer-cycle-log}

The mechanization pass went through one axis review followed by three
fix cycles. The axis-level flow review on 2026-05-22 caught significant
issues on F2, F10, and F53. The review also applied the opacity lessons
learned from the sibling meta-math axis, where earlier obligation
audits had shown that unguarded trust-base assumptions are too easy to
instantiate unsafely.

The F2 fix cycle, commit \texttt{be4a629b}
(\emph{fix: F2 target repair to generated equivalence}), replaced the
prior terminal-\texttt{Unit} fallback for target coarsening. The
canonical target repair quotient is now built from the equivalence
generated by identifying \(rF(x)\) and \(rF(x')\) whenever
\(q(x)=q(x')\), matching the source statement's minimal-repair
universal property.

The F10 fix cycle, commit \texttt{767c28da}
(\emph{fix: F10 canonical qI from ObservationSignature}), changed the
access quotient from an assumed quotient to a constructed one. The
canonical access quotient \(q_I\) is now obtained from a lawful
\texttt{ObservationSignature}, and the universal-property formulation
covers observable factorization together with the coarsest-quotient
property.

The F53 fix cycle, commit \texttt{a5ebd266}
(\emph{fix: F53 derived CER no-go + status hygiene}), changed the
closure--reality no-go statements from raw counterexample hypotheses to
derived consequences of opacity-guarded structural assumptions:
reduct-stability on the direct side and closure-instantiation gap on
the converse side. This brought the row into alignment with its
anchor-substrate status and with the AOR substrate delegation.

\subsection{Validator-chain instructions}\label{subsec:validator-chain-instructions}

The full validator chain can be rerun from the repository root.

\begin{enumerate}
\item Run \texttt{lake build}. This checks the \FIV\ Lean tree as part
of the repository build (164 jobs at the close of mechanization) and
verifies that the Lean files compile with no \texttt{sorry} and no
\texttt{admit} outside the four tagged trust-base axioms recorded in
\texttt{lean/manifests/trust\_base.txt}.
\item Run
\texttt{python3 scripts/check\_statements\_of\_record\_foundations\_iv.py}.
This cross-checks the statements of record against the inventory,
manifest, trust base, source files, and bibliography selections.
\item Run \texttt{make paper-preflight-foundations\_iv}. This builds
\texttt{paper/foundations\_iv/main.tex}, runs
\texttt{scripts/paper\_lint.sh}, checks the build log for layout and
reference problems, and reruns the \FIV\ statements-of-record validator.
\item During early drafting only, if the body is still in a skeleton
state, use
\texttt{PREP\_SKELETON\_MODE=1 make paper-preflight-foundations\_iv}.
That escape hatch suppresses the documented empty-skeleton
\texttt{Overfull \textbackslash vbox} warning and does not mask other
overfull boxes, underfull boxes, undefined references, citation
failures, or \texttt{Cref-Cref} artifacts.
\end{enumerate}

At paper close the strict gate is the required one:
\texttt{make paper-preflight-foundations\_iv} must pass without
\texttt{PREP\_SKELETON\_MODE}. The close-time pass must have no
undefined references or citations, no lint-detected layout warnings, and
a clean statements-of-record validator result.
 

\clearpage
\bibliographystyle{plainnat}
\bibliography{references}

\end{document}
